
\documentclass[10pt]{article} % For LaTeX2e
\usepackage{tmlr}
% If accepted, instead use the following line for the camera-ready submission:
%\usepackage[accepted]{tmlr}
% To de-anonymize and remove mentions to TMLR (for example for posting to preprint servers), instead use the following:
%\usepackage[preprint]{tmlr}

% Optional math commands from https://github.com/goodfeli/dlbook_notation.
\input{math_commands.tex}

\usepackage{hyperref}
\usepackage{url}
\usepackage{booktabs}
\usepackage{multirow}


\usepackage{amsthm}
\usepackage{graphicx}
\usepackage{amssymb}
\usepackage{subcaption}
\usepackage{mathtools}
\newtheorem{theorem}{Theorem}
\newtheorem{lemma}{Lemma}
\newtheorem{corollary}{Corollary}
\newtheorem{definition}{Definition}  

\newcommand\JW[1]{{\color{purple} [BT: #1]}}
\newcommand\todo[1]{{\color{red} [TODO: #1]}}
\newcommand\LY[1]{{\color{green} [Chang: #1]}}
\newcommand{\minisection}[1]{\vspace{0.04in}\noindent{\bf #1}}

\usepackage[table]{xcolor}
% \usepackage[dvipsnames]{xcolor}
\definecolor{LightPurple}{rgb}{0.88,0.88,1}

\title{Which Flatness Matters in LoRA-Based Continual Learning?}

% \\Journal Submissions

% Authors must not appear in the submitted version. They should be hidden
% as long as the tmlr package is used without the [accepted] or [preprint] options.
% Non-anonymous submissions will be rejected without review.

\author{\name Kyunghyun Cho \email kyunghyun.cho@nyu.edu \\
      \addr Department of Computer Science\\
      University of New York
      \AND
      \name Raia Hadsell \email raia@google.com \\
      \addr DeepMind
      \AND
      \name Hugo Larochelle \email hugolarochelle@google.com\\
      \addr Mila, Universit\'e de Montr\'eal \\
      Google Research\\
      CIFAR Fellow}

% The \author macro works with any number of authors. Use \AND 
% to separate the names and addresses of multiple authors.

\newcommand{\fix}{\marginpar{FIX}}
\newcommand{\new}{\marginpar{NEW}}

\def\month{MM}  % Insert correct month for camera-ready version
\def\year{YYYY} % Insert correct year for camera-ready version
\def\openreview{\url{https://openreview.net/forum?id=XXXX}} % Insert correct link to OpenReview for camera-ready version


\begin{document}


\maketitle

\begin{abstract}

% \begin{abstract}

% Flat minima in the lossland have been associated with improved generalization and with reduced catastrophic forgetting in continual learning. In LoRA-Based Continual Learning,  update directions Constrict in the lora subspace and LoRA-based continual learning methods have explicitly exploited this structure  to alleviate forgetting and interferenc. Traditional viewpoints  suggest that full-parameter perturbation is needed to achieve global flatness. However unclear whether continual generalization is governed by full-space flatness or by subspace flatness.

% a frozen pretrained backbone is adapted through task-conditioned low-rank update direction
% update directions in the adapter subspace despite a frozen backbone, lots of work  
% Yet in LoRA-based continual learning, a frozen pretrained backbone is adapted through task-conditioned low-rank update directions. 
% However,  It is therefore unclear whether continual generalization is governed by full-space flatness or by subspace flatness.

Flat minima in the loss landscape have long been associated with improved generalization and with reduced catastrophic forgetting in continual learning. In LoRA-based continual learning, adaptation is confined to task-conditioned subspace and recent methods exploit this subspace structure to reduce forgetting and interference. However, existing flatness-based method are inherited from full-parameter training and measure perturbations and curvature over the whole model. This creates a central ambiguity in LoRA-based continual learning, whether continual generalization should be understood through full-space flatness or through flatness restricted to the admissible subspace. We develop a sequential hierarchical PAC-Bayesian analysis for LoRA-based continual learning and derive a bound with an explicit hyperposterior drift term. Under a frozen backbone, we show that both the KL complexity and the sharpness-dependent empirical term reduce exactly to quantities supported on the admissible LoRA update subspace. This identifies adapter-subspace flatness as the relevant notion of flatness in this regime.
Experiments on vision and language benchmarks support this conclusion. Sharpness-aware optimization restricted to LoRA updates consistently improves continual performance.
% \end{abstract}

% \begin{abstract}
% Flat minima have long been associated with improved generalization and with reduced catastrophic forgetting in continual learning. However in LoRA-based continual learning, this notion becomes nontrivial: a frozen pretrained backbone is adapted only through task-conditioned low-rank update directions, so continual adaptation is confined to an admissible update subspace rather than the full parameter space. As a result, it is unclear whether continual generalization should be understood through flatness of the whole model, or through a local notion of flatness restricted to the admissible LoRA update geometry.
% We answer this question through a sequential hierarchical PAC-Bayesian analysis for LoRA-based continual learning. We derive a generalization bound with an explicit hyperposterior drift term. Under a frozen backbone, we show that both the KL complexity and the sharpness-dependent empirical term reduce to quantities supported on the admissible LoRA update subspace. This identifies task-conditioned local flatness along admissible LoRA directions, rather than arbitrary full-space flatness, as the relevant notion of flatness in this regime.
% Experiments on vision and language continual learning benchmarks support this conclusion. 

% Sharpness-aware optimization restricted to LoRA updates consistently improves continual performance, whereas perturbations that extend into frozen backbone parameters provide no systematic benefit and can harm stability. These results provide a principled foundation for subspace-restricted flatness in LoRA-based continual learning.
\end{abstract}


% \begin{abstract}
% LoRA based continual learning adapts a frozen pretrained backbone through low-rank task updates, so continual adaptation is confined to a task-conditioned admissible update subspace rather than the full parameter space. This raises a basic question: which notion of flatness matters for continual generalization and forgetting in this regime? 

% Flatness is widely associated with generalization and robustness, but in LoRA-based continual learning it is no longer obvious which notion of flatness matters. Because a frozen pretrained backbone is adapted only through task-conditioned low-rank update directions, continual adaptation is confined to an admissible update subspace rather than the full parameter space.


% Existing flatness-based arguments are mostly inherited from full-parameter training, but in LoRA based continual learning, full-space flatness is not automatically the object that governs future adaptation.
% We develop a sequential hierarchical PAC-Bayesian analysis for LoRA based continual learning and derive a bound with an explicit hyperposterior drift term. Under a frozen backbone, we show that both the KL complexity and the sharpness-dependent empirical term reduce to quantities supported on the admissible LoRA update subspace. This identifies task-conditioned local flatness along admissible LoRA directions as the relevant notion of flatness in this regime.
% Experiments on vision and language continual learning benchmarks support this conclusion. Our results provide a principled basis for subspace-restricted flatness in LoRA based continual learning.

%  Sharpness-aware optimization restricted to LoRA updates consistently improves continual performance, while perturbations that extend into frozen backbone parameters provide no systematic benefit and can harm stability. 
% \end{abstract}

% LoRA based continual learning adapts a frozen pretrained backbone through low rank task updates, so future adaptation is confined to a task dependent admissible update subspace rather than the full parameter space. This raises a fundamental question: in such a constrained regime, which notion of flatness actually matters for continual generalization and forgetting? Existing flatness based methods are largely inherited from full parameter training, where perturbations and curvature are defined over the whole model. In LoRA based continual learning, however, this full space notion is not automatically the one that governs future adaptation.
% We answer this question through a sequential hierarchical PAC Bayesian analysis for LoRA based continual learning. We derive a generalization bound with an explicit hyperposterior drift term that captures the sequential evolution of task priors, and show that under a frozen backbone, both the KL complexity and the sharpness dependent empirical term reduce to quantities supported on the admissible LoRA update subspace. The bound therefore identifies task conditioned local flatness along admissible LoRA directions, rather than arbitrary full space flatness, as the relevant notion of flatness in this regime.
% Experiments on vision and language continual learning benchmarks support this prediction. Sharpness aware optimization restricted to LoRA updates consistently improves continual performance, whereas perturbations that extend into frozen backbone parameters provide no systematic gains and can degrade stability. These results provide a principled foundation for studying flatness in LoRA based continual learning and offer guidance for geometry aware design in parameter efficient continual adaptation.



% Low-Rank Adaptation (LoRA) enables parameter-efficient continual learning by adapting a frozen pre-trained backbone through low-rank task updates. In this regime, continual adaptation is confined to a task-dependent admissible update subspace, while most flatness-based arguments are inherited from full-parameter training and are therefore defined in the full parameter space. This creates a fundamental mismatch: it is unclear whether continual generalization and forgetting are governed by full-space flatness, or by a local notion of flatness restricted to the LoRA update geometry.  
% We address this question through a sequential hierarchical PAC-Bayesian analysis for LoRA-based continual learning. We derive a generalization bound with an explicit hyperposterior drift term that captures the sequential evolution of task priors, and show that under a frozen backbone, both the KL complexity and the sharpness-dependent empirical term reduce to quantities supported on the admissible LoRA update subspace. This identifies task-conditioned, solution-local flatness along admissible LoRA directions as the relevant notion of flatness in this setting. 







% In this regime, continual adaptation is confined to a task-dependent admissible update subspace, while most flatness-based arguments are inherited from full-parameter training and are therefore defined in the ambient parameter space. This creates a fundamental mismatch: it is unclear whether continual generalization and forgetting are governed by full-space flatness, or by a local notion of flatness restricted to the LoRA update geometry. 

% We address this question through a sequential hierarchical PAC-Bayesian analysis for LoRA-based continual learning. We derive a generalization bound with an explicit hyperposterior drift term that captures the sequential evolution of task priors, and show that under a frozen backbone, both the KL complexity and the sharpness-dependent empirical term reduce to quantities supported on the admissible LoRA update subspace. This identifies task-conditioned, solution-local flatness along admissible LoRA directions as the relevant notion of flatness in this setting. 

% Guided by this result, we compare perturbation scopes and sharpness types in LoRA-based continual learning. Across vision and language benchmarks, subspace-restricted sharpness-aware optimization consistently improves continual performance, whereas extending perturbations into frozen backbone parameters provides no systematic benefit and can harm stability. Our results offer a principled foundation for studying flatness in LoRA-based continual learning and provide guidance for geometry-aware design in parameter-efficient continual adaptation.






% --------------
% Parameter efficient continual learning (PECL) adapts large pre trained models through adapters while keeping the backbone fixed to stable previous kenowledge and learn new task, but there are still  trainable adapter are shared across all tasks,

% Parameter efficient continual learning (PECL) adapts large pretrained models through lightweight adapters while freezing the backbone so that previously acquired knowledge remains stable.
% Parameter-efficient continual learning (PECL) adapts large pre-trained models to sequential tasks by tuning lightweight adapters while freezing the backbone to preserve prior knowledge.
% From a loss-landscape perspective, flat minima enhance generalization, yet in PECL, task interference concentrates in the new trained shared adapter subspace despite a frozen backbone. 
% Traditional viewpoints (such as Flat-LoRA) suggest that full-parameter perturbation is needed to achieve global flatness, but this contradicts the original intention of PECL to "freeze the backbone and fine-tune efficiently".
% % To overcome forgetting, traditional viewpoints (such as Flat-LoRA) suggest that full-parameter perturbation is needed to achieve global flatness, but this contradicts the original intention of PECL to "freeze the backbone and fine-tune efficiently".
% This raises a central question: in this regime, should one pursue flat minima in the full parameter space, or is it sufficient to control flatness only within the adapter subspace? 
% We answer this question through a PAC Bayesian analysis and derive a sharpness aware generalization bound tailored to PECL. 
% % Once the backbone is frozen, the bound shows that both the complexity and flatness terms are entirely determined by the trainable LoRA manifold, so that adapter subspace sharpness becomes the relevant notion of flatness. 
% Once the backbone is frozen, the bound shows that both the complexity term and the flatness term are determined by the trainable LoRA manifold, which identifies adapter subspace sharpness as the relevant notion of flatness.
% This provides a principled justification for flatness oriented optimisation confined to the adapter subspace and offers a conceptual basis for geometry aware designs of LoRA based continual learning methods. Experiments onvision benchmarks and language benchmark support this conclusion, confirming that flatness promotion restricted to the adapter subspace is sufficient and preferable in practice.

% \end{abstract}

\section{Introduction}




% % Continual learning with parameter-efficient adaptation (PECL) has become a practical way to reuse large pre-trained models on evolving task sequences \todo{(CITES)}.  Among these methods, Low-Rank Adaptation (LoRA) \todo{(CITE)} freezes the backbone and trains compact adapters, offering substantial memory and compute savings while preserving upstream capacities. \todo{short symmary on main works in this direction could also be paragraph.} However, in sequential settings, partially shared adapter subspaces still suffer from task interference, which challenges the trade-off between stability and plasticity.



% % Continual learning (CL) has become increasingly relevant for adapting large-scale pre-trained models to evolving tasks in dynamic environments~\cite{PARISI201954,10444954,a2923225ea7d94e6c86faff97e23c69be,10.1145/3735633}. 
% % While full fine-tuning remains the most expressive adaptation strategy, its application in sequential settings is computationally demanding and highly susceptible to catastrophic forgetting, since newly acquired knowledge may interfere with previously learned representations~\cite{howard2018universallanguagemodelfinetuning,doi:10.1073/pnas.1611835114,pmlr-v234-zhai24a,luo2025empiricalstudycatastrophicforgetting}.
% % Parameter efficient continual learning methods (PECL)~\cite{pmlr-v97-houlsby19a,dingParameterefficientFinetuningLargescale2023,He_2025_CVPR,NEURIPS2024_b0bc711f}, such as Low-Rank Adaptation (LoRA)~\cite{hu2022lora}, mitigate these issues by introducing lightweight task-specific modules while freezing the pre-trained backbone.
% % In these approaches, most task-specific adaptation is funneled through a small set of adapters that are reused or combined across tasks, so that interference is concentrated in trainable adapters that are shared across all tasks.
% % % because trainable parameters are shared across all tasks
% % % This concentration implies that, beyond parameter efficiency, a central challenge in PECL is to understand and control the local geometry of this shared adapter subspace.
% % This concentration implies that, beyond parameter efficiency, a central challenge in PECL is to understand and control the local geometry of this shared adapter subspace, because the shape of the loss landscape in this region governs how new gradients interact with existing adapters, mediates the stability–plasticity trade-off.







% % ----------------

% % \section{Introduction}

% Parameter-efficient continual learning (PECL) has emerged as a practical approach for adapting large pretrained models to sequential tasks with limited computational and storage overhead~\cite{pmlr-v97-houlsby19a,Wang_2022_CVPR,dingParameterefficientFinetuningLargescale2023,xu2023parameterefficientfinetuningmethodspretrained}. Among PECL methods, Low-Rank Adaptation (LoRA) is particularly attractive because it enables continual adaptation through lightweight low-rank updates while keeping the pretrained backbone fixed~\cite{hu2022lora}. This design naturally turns continual adaptation into a constrained geometric problem: future tasks do not modify the model freely in the full parameter space, but only through task-dependent low-rank update directions induced by the trainable LoRA parameters. Recent LoRA-based continual learning methods have explicitly exploited this structure through adapter modularity, orthogonal constraints, and task-specific subspace allocation to alleviate forgetting and interference~\cite{Gao_2023_ICCV,NEURIPS2020_70d85f35,wang2023orthogonal,wu2025sdlorascalabledecoupledlowrank}. However, these advances are primarily algorithmic. They do not yet provide a principled understanding of how such admissible update geometry shapes continual generalization, forgetting, and the notion of solution stability that should matter in this regime.

% Another line of work studies continual learning through the geometry of the loss landscape. Flat minima have long been associated with improved generalization and robustness~\cite{FlatMinima1997,wu2017understandinggeneralizationdeeplearning,chaudhariEntropySGDBiasingGradient2019,NEURIPS2020_518a38cc,NEURIPS2021_995f5e03,JMLRAnEmpiricalInvestigatio,yangRevisitingFlatnessAwareOptimization2025}, and sharpness-aware optimization methods such as SAM and its variants operationalize this principle by favoring solutions that remain stable under local perturbations~\cite{foretSharpnessAwareMinimizationEfficiently2021,pmlr-v151-bisla22a,Zhang_GAM,NEURIPS2024_C-Flat}. Prior work has further suggested that flatter solutions can mitigate catastrophic forgetting in continual learning~\cite{NEURIPS2020_518a38cc,JMLRAnEmpiricalInvestigatio,shiOvercomingCatastrophicForgetting2021,liRevisitingCatastrophicForgetting2024}. However, most existing flatness-based arguments are inherited from full-parameter training, where perturbations and curvature are defined over the whole model~\cite{li2024flatlora}. In LoRA-based continual learning, by contrast, future adaptation is confined to admissible low-rank update directions. A model may therefore appear flat in many directions that future tasks can never realize, yet remain highly sensitive along the directions through which continual adaptation actually occurs. The key issue is therefore not whether flatness is beneficial, but which notion of flatness remains relevant once continual adaptation is restricted to admissible LoRA directions.





% % Low-Rank Adaptation (LoRA) has become a widely used mechanism for parameter-efficient continual learning because it enables sequential task adaptation through lightweight low-rank updates while keeping the pretrained backbone fixed. This design makes continual adaptation scalable and storage-efficient, but it also fundamentally changes the geometry of learning. In LoRA-based continual learning, future tasks do not update the model freely in the full parameter space. Instead, adaptation is confined to a task-dependent admissible update subspace induced by the trainable LoRA parameters. As a result, both cross-task interference and future correction are mediated through a restricted set of directions, rather than through the full model. 



















% % Flat minima in the loss landscape are widely believed to promote generalization and robustness, since flatter solutions are less sensitive to parameter perturbations and data distribution shifts~\citep{FlatMinima1997,wu2017understandinggeneralizationdeeplearning,chaudhariEntropySGDBiasingGradient2019,NEURIPS2020_518a38cc,NEURIPS2021_995f5e03,JMLRAnEmpiricalInvestigatio,yangRevisitingFlatnessAwareOptimization2025}. Sharpness-aware optimizers such as SAM~\citep{foretSharpnessAwareMinimizationEfficiently2021}, RWP~\citep{pmlr-v151-bisla22a}, GAM~\citep{zhangGradientNormAware2023} and C-FLAT\cite{NEURIPS2024_C-Flat} instantiate this idea in practice by explicitly penalizing increases of the loss under small parameter perturbations.
% % While such methods have proven effective in full fine-tuning regimes, their behavior in the restricted optimization landscape of PECL remains under-explored. 
% % LoRA-SAM~\citep{NEURIPS2024_4eb2c0ad} provides evidence that limiting sharpness-aware updates to the LoRA coordinates already yields substantial implicit regularization, whereas Flat-LoRA~\citep{li2025flatloralowrankadaptationflat} argues that flatness restricted to the subspace may be insufficient and therefore advocates perturbing all model parameters.
% % In the PECL regime that we focus on, however, the core efficiency principle is to strictly freeze the pre-trained backbone and confine optimization to a trainable adapter subspace $W_\Delta$. 
% % % However, such a strategy fundamentally conflicts with the core efficiency principle of PECL, which relies on freezing the pre-trained backbone to isolate upstream knowledge and optimization is strictly confined to a adapter subspace $W_\Delta$.
% % This raises a central question: in PECL, is it necessary to pursue global flatness by perturbing all model parameters, or is controlling flatness solely within the trainable adapter manifold sufficient to govern stability and plasticity?


% % This tension highlights a critical gap in the literature: it is currently unknown whether the pursuit of global flatness is truly indispensable, or if controlling the local geometry solely within the adapter subspace is sufficient to govern the stability-plasticity balance of the overall model.






% % This latter position effectively posits that generalization and robustness are governed primarily by flatness in the full model space, which is conceptually misaligned with practical PECL designs that deliberately restrict learning to adapter parameters.


% % Flat minima in the loss landscape are widely believed to promote generalization and robustness, since flatter solutions are less sensitive to parameter perturbations and distribution shifts~\citep{FlatMinima1997,wu2017understandinggeneralizationdeeplearning,chaudhariEntropySGDBiasingGradient2019,NEURIPS2020_518a38cc,NEURIPS2021_995f5e03,JMLRAnEmpiricalInvestigatio,yangRevisitingFlatnessAwareOptimization2025}.
% % Sharpness-aware optimizers such as SAM, RWP, GAM, and MC-FLAT instantiate this idea in practice by explicitly penalizing worst case or expected increases of the loss under small parameter perturbations~\citep{foretSharpnessAwareMinimizationEfficiently2021,pmlr-v151-bisla22a,zhangGradientNormAware2023,bahri2022sharpnessawareminimizationimproveslanguage}.
% % In LoRA settings, LoRA-SAM~\citep{NEURIPS2024_4eb2c0ad} provides evidence that limiting sharpness-aware updates to the LoRA coordinates can already yield substantial implicit regularization. By contrast, Flat-LoRA~\citep{li2025flatloralowrankadaptationflat} contends that solutions that are flat in the LoRA space may remain sharp in the full parameter space and therefore advocates perturbing all parameters rather than only the adapters.
% % This latter position effectively posits that generalization and robustness are governed primarily by flatness in the full model space, which is conceptually misaligned with practical PECL designs that deliberately restrict learning to adapter parameters.
% % In PECL this misalignment becomes more pronounced, since the backbone and previously learned adapters are frozen while optimization is strictly confined to a LoRA subspace $W_\Delta$.
% % Despite extensive work on sharpness-aware training, there is still no theoretical account that explains how the geometry of this local adapter subspace shapes the predictive behaviour of the global model under continual updates. In particular, it remains unclear whether controlling flatness only within $W_\Delta$ is sufficient to improve generalization and resistance to forgetting, or whether some form of flatness in the full parameter space is fundamentally required.









% % Several studies associate flat minima in the loss landscape with improved generalization and robustness, since flatter solutions are more stable under parameter perturbations~\citep{wu2017understandinggeneralizationdeeplearning,chaudhariEntropySGDBiasingGradient2019,NEURIPS2020_518a38cc,NEURIPS2021_995f5e03,JMLRAnEmpiricalInvestigatio,yangRevisitingFlatnessAwareOptimization2025}.
% % Prior work has explored flat minima in CL through sharpness-aware optimization~\cite{pmlr-v151-bisla22a,foretSharpnessAwareMinimizationEfficiently2021,zhangGradientNormAware2023}, mode connectivity~\cite{izmailov2019averagingweightsleadswider,mirzadeh2020linearmodeconnectivitymultitask}, and improved representation learning~\cite{hu2022doesselfsupervisedpretrainingperform,ramasesh2022effect}.
% % While flatness-aware training has been effective in full-model regimes, its implications for an adapter subspace remain unclear.
% % In particular, it is unknown whether reducing sharpness only within the LoRA subspace suffices to improve the stability-plasticity balance, and a theoretical connection between a subspace-restricted flatness constraint and generalization guarantees is currently lacking in PECL.

% % {Motivating Observation: LoRA Subspace Flatness Leads to a Flat Global Landscape}


% % \begin{figure}[htbp]
% %     \centering
% %     \begin{minipage}{0.495\textwidth}
% %         \centering
% %         \includegraphics[width=0.9\linewidth]{figures_TRML/AAA_curve_Imagenet-R_summary_inc10_rank16.pdf}
% %         \subcaption{Average Accuracy (ACC)  on ImageNet-R}
% %         \label{fig:curves_imageR_1_AAA}
% %     \end{minipage}%
% %     \begin{minipage}{0.495\textwidth}
% %         \centering
% %         \includegraphics[width=0.9\linewidth]{figures_TRML/NME_curve_Imagenet-R_summary_inc10_rank16.pdf}
% %         \subcaption{Nearest Mean of Exemplars (NME) on ImageNet-R}
% %         \label{fig:curves_imageR_2_NME}
% %     \end{minipage}
% %     \caption{
% %         AAA and NME summary curves on ImageNet-R ($T=20, r=16$). Across various CL configurations (SeqFT, SeqLoRA, IncLoRA, and OLoRA ), SAM consistently outperforms SGD, achieving higher accuracy in both AAA and NME throughout the incremental learning sequence. The performance gap typically widens on later tasks, indicating that SAM learns more stable representations with reduced forgetting.
% %     }
% %     \label{fig:combined_curves_imageR}
% % \end{figure}




% % \begin{figure}[htbp]
% % \centering
% % \begin{minipage}{0.495\linewidth}
% % \centering
% % \includegraphics[width=0.9\linewidth]{figures_TRML/AAA_curve_Imagenet-R_summary_inc10_rank16.pdf}
% % \caption{$\mathrm{AAA}$ Curve}
% % \label{fig: 1_AAA}
% % \end{minipage}
% % \begin{minipage}{0.495\linewidth}
% % \centering
% % \includegraphics[width=0.9\linewidth]{figures_TRML/NME_curve_Imagenet-R_summary_inc10_rank16.pdf}
% % \caption{$\mathrm{NME}$ Curve}
% % \label{fig: 2_NME}
% % \end{minipage}
% % \end{figure}



% % \begin{figure}[htbp]
% % \centering
% % \begin{minipage}{0.9\linewidth}
% % \centering
% % \includegraphics[width=0.9\linewidth]{figures_TRML/2D_lossLandscape_task0_lora_opt_sam_vs_sgd_3d_compare2.pdf}
% % \caption{SGD}
% % \label{fig: EFM_seqFT}
% % \end{minipage}

% % \end{figure}




% % In PECL, optimization is confined to a low-dimensional LoRA adapter subspace, with the pre-trained backbone and previously learnt adapters remaining frozen. This design naturally raises a central question: if local flatness  within  adapter subspace sufficient to improve the performance of the global model in CL?


% % LoRA-SAM~\citep{NEURIPS2024_4eb2c0ad} explores the implicit regularisation of SAM for scale-invariant problems and finds that SAM promotes balance, whereas Flat-LoRA~\citep{li2025flatloralowrankadaptationflat} argues that solutions that are flat in the LoRA space may remain sharp in the full parameter space and therefore recommends that all parameters be perturbed.  
% % However, this central tension between local and global flatness remains unresolved in CL. Moreover, Flat-LoRA's proposed "global perturbation" approach fundamentally conflicts with the core principle of PECL, i.e. isolating pre-trained knowledge and avoiding global parameter fluctuations. Overall, existing studies have not yet clarified how the geometry of the local adapter subspace influences the predictive behavior of the global model.


% % Our empirical observations shed new light on this issue. As shown in Figure~\ref{fig:3d_loss_landscape}, when only the low-rank LoRA subspace is optimized using a flatness-aware method, the loss landscape of the entire model around the final solution already becomes wider and smoother than that achieved by SGD. This suggests that enhancing flatness in the adapter subspace can indeed lead to a flatter global landscape, which in turn motivates the first research question of this paper: Does increasing flatness in the LoRA subspace consistently improve continual learning performance?


% % This raises the fundamental, unresolved question of the relationship between the flatness of the local LoRA subspace and the robust performance of the global model in the PECL.



% % This raises a key question: Is enforcing flatness only within the adapter subspace helps improve model performance?  


% % Moreover, our empirical studies reveal a consistent pattern (Figure \ref{fig:3d_loss_landscape}): enforcing flatness only within the low-rank LoRA subspace do improves whole model's flatness. This raises the fundamental, unresolved question of the relationship between the flatness of the local LoRA subspace and the robust performance of the global model in the PECL. Existing research has not clarified how the geometry of the local subspace affects the predictive behaviour of the global model, nor has it explained why optimising only local flatness can improve global CL performance. The connection between the local and global remains ambiguous.

% % Our empirical observations further refine this debate. As illustrated in Figure~\ref{fig:3d_loss_landscape}, when only the low rank LoRA subspace is optimised with a flatness aware method, the loss landscape of the entire model around the final solution already becomes wider and smoother than under SGD. This suggests that enhancing flatness in the adapter subspace can induce a flatter global landscape and motivates the first question studied in this paper: does increasing LoRA subspace flatness consistently translate into improved continual learning performance?





% % Our empirical studies reveal a consistent pattern (Figure \ref{fig:combined_curves_imageR}): enforcing flatness only within the low-rank LoRA subspace improves continual-learning outcomes, raising final accuracy and reducing forgetting. However, this observation challenges the common intuition that local (subspace) geometry is less important than global loss-landscape properties.





% % However, the proposed 'global perturbation' solution

% % However, these discussions are not situated within the context of continual learning, leaving conclusions in this specific setting largely unexplored. Nevertheless,

% % , as neither approach has been rigorously evaluated

% % LoRA-SAM ~\citep{NEURIPS2024_4eb2c0ad} explores the implicit regularization of SAM for scale-invariant problems and finds that SAM promotes a better stability-plasticity balance, whereas Flat-LoRA (Li et al., 2025a) argues that solutions that are flat in the LoRA space may remain sharp in the full parameter space and therefore advocates global perturbations.


\begin{figure}[htbp]
\centering
% \begin{minipage}{0.8\linewidth}
\centering
\includegraphics[width=0.8\linewidth]{figures_TRML/2D_lossLandscape_task0_lora_opt_sam_vs_sgd_3d_compare_withxyz.pdf}
\caption{Loss landscape visualization for a ViT model adapted with LoRA, comparing SGD and adapter-only SAM. Although perturbation is restricted to the LoRA subspace, full-space probes around the final solution show that adapter-only SAM leads to a flatter local landscape than SGD.}
\label{fig:3d_loss_landscape}
% \end{minipage}
\end{figure}



% \includegraphics[width=0.72\linewidth]{figures_TRML/ACC_dataset_pair_Robustness_to_Corruptions_on_ImageNet-C_Robustness_to_Perturbation_on_ImageNet-P2.pdf}

% \begin{subfigure}[t]{0.48\linewidth}
%     \centering
%     \includegraphics[width=\linewidth,height=4.2cm,keepaspectratio,trim=1.2cm 0.8cm 1.2cm 0.6cm,clip]{figures_TRML/2D_lossLandscape_task0_lora_opt_sam_vs_sgd_3d_compare_withxyz.pdf}
%     \caption{...}
% \end{subfigure}
% \hfill
% \begin{subfigure}[t]{0.48\linewidth}
%     \centering
%     \includegraphics[width=\linewidth,height=4.2cm,keepaspectratio]{figures_TRML/ACC_dataset_pair_Robustness_to_Corruptions_on_ImageNet-C_Robustness_to_Perturbation_on_ImageNet-P2.pdf}
%     \caption{...}
% \end{subfigure}


% \usepackage{subcaption}

% \begin{figure}[htbp]
% \centering

% \begin{subfigure}[t]{0.48\linewidth}
%     \centering
%     \includegraphics[height=4.6cm]{figures_TRML/2D_lossLandscape_task0_lora_opt_sam_vs_sgd_3d_compare_withxyz.pdf}
%     \caption{Loss landscape visualization ...}
%     \label{fig:3d_loss_landscape}
% \end{subfigure}
% \hfill
% \begin{subfigure}[t]{0.48\linewidth}
%     \centering
%     \includegraphics[height=4.6cm]{figures_TRML/ACC_dataset_pair_Robustness_to_Corruptions_on_ImageNet-C_Robustness_to_Perturbation_on_ImageNet-P2.pdf}
%     \caption{Robustness of perturbation types under CL ...}
%     \label{fig:robustness_imagenet_cp}
% \end{subfigure}

% \caption{Overall comparison.}
% \label{fig:main_compare}
% \end{figure}



% We ground our investigation in a motivating empirical observation that calls into question the need for global parameter perturbations in PECL. As visualized in Figure \ref{fig:3d_loss_landscape}, applying sharpness-aware minimization restricted to the low-rank adapter subspace is already sufficient to induce a loss landscape that appears wider and flatter when probed in the full parameter space, compared with an SGD baseline. This suggests that the geometry of the frozen backbone and that of the trainable adapters are not dynamically independent. Stabilizing the local curvature within the trainable subspace can implicitly regularize the effective geometry of the whole model. Motivated by this phenomenon, we hypothesize that subspace flatness, understood as geometric control confined to the trainable manifold, may already suffice to substantially improve robustness and achieve a better balance between stability and plasticity in PECL.
% % can already be sufficient to improve robustness, and a better balance between stability and plasticity in PECL.


% To investigate this hypothesis, we answer this question through a sequential hierarchical PAC-Bayesian analysis for LoRA-based continual learning. We derive a generalization bound with an explicit hyperposterior drift term that captures the sequential evolution of task priors across tasks. Under a frozen backbone, we show that both the KL complexity and the sharpness-dependent empirical term reduce to quantities supported on the admissible LoRA update subspace. This identifies task-conditioned local flatness along admissible LoRA directions, rather than arbitrary full-space flatness, as the relevant notion of flatness in this regime. Guided by this result, we further study perturbation scope and sharpness type in LoRA-based continual learning, and show empirically that sharpness-aware optimization restricted to LoRA updates consistently improves continual performance, whereas perturbations that extend into frozen backbone parameters provide no systematic benefit and can harm stability.



% we derive a PAC-Bayesian generalization bound specialized to PECL with a frozen backbone and Gaussian posteriors supported on the adapter subspace. Under this modelling assumption, the complexity and flatness contributions to the bound can be expressed entirely in terms of quantities restricted to update subspace rather than the full parameter space. Beyond this bound, we clarify the mechanism linking weight-space flatness to feature-space robustness. We show that the LoRA is the implement to the abstract update subspace, and that optimizing flatness in factor coordinates is equivalent to suppressing dominant eigenvalues of the local Fisher information matrix. Taken together, these results provide a principled account of how focusing flatness regularization on the trainable adapter manifold can shape the global loss landscape, stabilize feature representations and improve performance in PECL.


% To investigate this hypothesis, we establish a theoretical framework that unifies probabilistic generalization bounds with geometric stability. We first derive a Subspace-aware PAC-Bayesian bound specifically tailored for PECL and show that, once the backbone is frozen, the complexity and sharpness terms that govern generalization are determined entirely by the adapter subspace rather than by the full parameter space.
% Beyond theory bound, we elucidate the underlying mechanism by linking this weight-space flatness to feature-space robustness. We demonstrate that the LoRA factor space is geometrically isomorphic to the abstract update subspace, implying that flatness optimization in factor coordinates is equivalent to suppressing the dominant eigenvalues of the local Fisher Information Matrix. In combination, these results provide a principled account of how focusing flatness regularization on the trainable adapter manifold can shape the global loss landscape, stabilize feature representations, and achieve better peformance in PECL.




% , which theoretically confirms that the generalization error is upper-bounded strictly by the expected sharpness within the adapter subspace, rendering global perturbations theoretically redundant.

% Beyond sufficiency, we elucidate the underlying mechanism by linking this weight-space flatness to feature-space robustness. We demonstrate that the LoRA factor space is geometrically isomorphic to the abstract update subspace, implying that flatness optimization in factor coordinates is equivalent to suppressing the dominant eigenvalues of the local Fisher Information Matrix. This spectral suppression limits the sensitivity of feature representations to parameter updates, thereby directly mitigating catastrophic forgetting.


% Our extensive empirical evaluation on the distribution drift benchmark ImageNet-R,  the corruption benchmarks ImageNet-C, and the perturbation benchmarks ImageNet-P using SeqLoRA, IncLoRA and OLoRA adapters corroborates these theoretical findings. 

% Extensive empirical evaluation across the distribution drift benchmark, corruption benchmarks, and perturbation benchmarks corroborates these theoretical findings under three distinct subspace sharing strategies: fully shared, partially fused , and constraint-augmented. 
% Across the three subspace sharing strategies, we demonstrate that flatness-aware optimization confined to the adapter subspace consistently improves continual learning metrics and feature stability compared to standard baselines. Furthermore, by directly comparing local versus global perturbation strategies, we reveal that subspace-restricted optimization not only outperforms global approaches but also mitigates the instability associated with perturbing frozen parameters. These results collectively confirm that controlling the geometry of the adapter subspace is a sufficient and effective strategy for robust PECL, validating the practical implications of our PAC-Bayesian framework.

% We then conduct an extensive empirical study on vision benchmarks. We evaluate ViT-B/16 adapters on ImageNet-R as a distribution-shift benchmark and on ImageNet-C and Tiny-ImageNet-C/P as corruption and perturbation benchmarks. Across three adapter-subspace sharing strategies, corresponding to fully shared (SeqLoRA), partially fused (IncLoRA) and constraint-augmented (OLoRA) architectures, flatness-aware optimization confined to the adapter subspace consistently improves continual learning metrics and feature stability compared to SGD-based baselines. Furthermore, by directly comparing local versus global perturbation strategies within the same PECL models, we observe that subspace-restricted optimization can match or outperform global approaches on these vision benchmarks, while often mitigating the instability associated with perturbing frozen backbone weights.



% Parameter-efficient continual learning has emerged as a practical approach for adapting large pretrained models to sequential tasks with limited computational and storage overhead~\cite{pmlr-v97-houlsby19a,Wang_2022_CVPR,dingParameterefficientFinetuningLargescale2023,xu2023parameterefficientfinetuningmethodspretrained}. Low-Rank Adaptation (LoRA) is particularly attractive because it enables continual adaptation through lightweight low-rank updates while keeping the pretrained backbone fixed~\cite{hu2022lora}. 

LoRA based continual learning has emerged as a practical approach for adapting large pretrained models to sequential tasks with limited computational and storage overhead, since it enables continual adaptation through lightweight low rank updates while keeping the pretrained backbone fixed~\cite{hu2022lora}. This design turns continual adaptation into a constrained geometric problem: future tasks do not modify the model freely in the full parameter space, but only through task-conditioned low-rank update directions induced by the trainable LoRA parameters. Recent LoRA based continual learning methods have explicitly exploited this structure through orthogonal constraints, and task-specific subspace allocation to alleviate forgetting and interference~\cite{Gao_2023_ICCV,NEURIPS2020_70d85f35,wang2023orthogonal,wu2025sdlorascalabledecoupledlowrank}. 

% However, these advances are primarily algorithmic. They do not yet provide a principled understanding of how such admissible update geometry shapes continual generalization and forgetting in this regime.

Another line of work studies continual learning through the geometry of the loss landscape. Flat minima have long been associated with improved generalization and robustness~\cite{FlatMinima1997,wu2017understandinggeneralizationdeeplearning,chaudhariEntropySGDBiasingGradient2019,NEURIPS2020_518a38cc,NEURIPS2021_995f5e03,JMLRAnEmpiricalInvestigatio,yangRevisitingFlatnessAwareOptimization2025}, and sharpness-aware optimization methods such as SAM operationalize this principle by favoring solutions that remain stable under local perturbations~\cite{foretSharpnessAwareMinimizationEfficiently2021,pmlr-v151-bisla22a,Zhang_GAM,NEURIPS2024_C-Flat}. Prior work has further suggested that flatter solutions can mitigate catastrophic forgetting in continual learning~\cite{NEURIPS2020_518a38cc,JMLRAnEmpiricalInvestigatio,shiOvercomingCatastrophicForgetting2021,liRevisitingCatastrophicForgetting2024}. However, most existing flatness-based arguments are inherited from full-parameter training, where perturbations and curvature are defined over the whole model~\cite{li2024flatlora}. In LoRA based continual learning, by contrast, flatness is still defined relative to the task-conditioned loss, while only the admissible low-rank update directions are available to fit the current task data. It therefore remains unclear whether continual generalization should be understood through full-space flatness, or through task-conditioned local flatness along admissible LoRA directions.


Surprisingly, as visualized in Figure~\ref{fig:3d_loss_landscape}, applying sharpness-aware minimization only within the LoRA update subspace already yields a solution whose surrounding loss landscape appears flatter than that obtained by SGD when probed in the full parameter space. Thus, even though flatness control is imposed only on the admissible LoRA directions, its effect is already visible under full-space. More importantly, restricting flatness optimization to LoRA updates also yields substantial gains in continual learning performance. Taken together, these observations create a direct tension with flatness notions inherited from full-parameter training and claims that effective flatness control in LoRA requires perturbing the whole model~\cite{li2024flatlora}.

% We sharpen this ambiguity through a motivating empirical observation. As visualized in Figure~\ref{fig:3d_loss_landscape}, applying sharpness-aware minimization only within the LoRA update subspace already yields a solution whose surrounding loss landscape appears flatter than that obtained by SGD when probed in the full parameter space. This suggests that geometric control within the trainable LoRA subspace can propagate to measurable differences in full-model sensitivity. More importantly, restricting flatness optimization to LoRA updates also produces substantial gains in continual learning performance. Taken together, these observations create a direct tension with flatness notions inherited from full-parameter training, as well as with claims that effective flatness control in LoRA requires perturbing the whole model~\cite{li2024flatlora}. They suggest that adapter-restricted flatness control is already sufficient to induce both measurable changes in full-model sensitivity and significant performance gains, yet they do not explain why this happens or which notion of flatness is actually being controlled.

To answer this question, we develop a sequential hierarchical PAC-Bayesian analysis for LoRA based continual learning. We derive a generalization bound with an explicit hyperposterior drift term that captures the sequential evolution of task priors across tasks. Under a frozen backbone, we show that both the KL complexity and the sharpness-dependent empirical term reduce to quantities supported on the admissible LoRA update subspace. This identifies task-conditioned local flatness along admissible LoRA directions, rather than arbitrary full-space flatness, as the relevant notion of flatness in this regime. Guided by this result, we further derive empirical predictions about perturbation scope and sharpness type, and verify them experimentally: sharpness-aware optimization restricted to LoRA updates consistently improves continual performance, whereas perturbations that extend into frozen backbone parameters provide no systematic benefit and can harm stability.












% Parameter-efficient continual learning (PECL) has emerged as a practical approach for adapting large pretrained models to sequential tasks with limited computational and storage overhead~\cite{pmlr-v97-houlsby19a,Wang_2022_CVPR,dingParameterefficientFinetuningLargescale2023,xu2023parameterefficientfinetuningmethodspretrained}.% ~\cite{pmlr-v97-houlsby19a,Wang_2022_CVPR,dingParameterefficientFinetuningLargescale2023,xu2023parameterefficientfinetuningmethodspretrained}
% Among PECL methods, Low-Rank Adaptation (LoRA) is particularly attractive because it enables continual adaptation through lightweight low-rank updates while keeping the pretrained backbone fixed~\cite{hu2022lora}. This design naturally turns continual adaptation into a constrained geometric problem: future tasks do not modify the model freely in the full parameter space, but only through task-dependent low-rank update directions induced by the trainable LoRA parameters. Recent LoRA-based continual learning methods have explicitly exploited this structure through adapter modularity, orthogonal constraints, and task-specific subspace allocation to alleviate forgetting and interference~\cite{Gao_2023_ICCV,NEURIPS2020_70d85f35,wang2023orthogonal,wu2025sdlorascalabledecoupledlowrank}. However, these advances are primarily algorithmic. They do not yet provide a principled understanding of how such admissible update geometry shapes continual generalization, forgetting, and the notion of solution stability that should matter in this regime.


% Another line of work studies continual learning through the geometry of the loss landscape. Flat minima have long been associated with improved generalization and robustness~\cite{FlatMinima1997,wu2017understandinggeneralizationdeeplearning,chaudhariEntropySGDBiasingGradient2019,NEURIPS2020_518a38cc,NEURIPS2021_995f5e03,JMLRAnEmpiricalInvestigatio,yangRevisitingFlatnessAwareOptimization2025}, and sharpness-aware optimization methods such as SAM and its variants operationalize this principle by favoring solutions that remain stable under local perturbations~\cite{foretSharpnessAwareMinimizationEfficiently2021,pmlr-v151-bisla22a,Zhang_GAM,NEURIPS2024_C-Flat}. Prior work has further suggested that flatter solutions can mitigate catastrophic forgetting in continual learning~\cite{NEURIPS2020_518a38cc,JMLRAnEmpiricalInvestigatio,shiOvercomingCatastrophicForgetting2021,liRevisitingCatastrophicForgetting2024}. However, most existing flatness-based arguments are inherited from full-parameter training, where perturbations and curvature are defined over the whole model~\cite{li2024flatlora}. In LoRA-based continual learning, by contrast, adaptation is confined to task-conditioned admissible low-rank update directions. This means that flatness measured over the whole model is not automatically the flatness notion that governs continual adaptation, because flatness is defined relative to the task-conditioned loss, while under frozen-backbone LoRA adaptation, only the admissible low-rank update directions are available to fit the current task data. It therefore remains unclear whether continual generalization should be understood through full-space flatness, or through task-conditioned local flatness along admissible LoRA directions.

% We sharpen this ambiguity through two empirical observations. First, as visualized in Figure~\ref{fig:3d_loss_landscape}, applying sharpness-aware minimization only within the low-rank LoRA update subspace already yields a solution whose surrounding loss landscape appears flatter than that obtained by SGD when probed in the full parameter space. Second, this is not merely a local geometric artifact: restricting flatness optimization to LoRA updates also produces substantial gains in continual learning performance, with improvements becoming more pronounced in later tasks in Figure~\ref{fig:robustness_imagenet_cp}. Taken together, these observations create a direct tension with flatness notions inherited from full-parameter training, as well as with claims that effective flatness control in LoRA requires perturbing the whole model~\cite{li2024flatlora}. They suggest that adapter-restricted flatness control is already sufficient to induce both measurable changes in full-model sensitivity and significant performance gains, yet they do not explain why this happens or which notion of flatness is actually being controlled. This is precisely the question we address.


% At the same time, such an observation does not resolve which notion of flatness is theoretically relevant: it leaves open whether full-space flatness is itself the governing object for continual generalization, or whether it is only a byproduct of controlling the admissible LoRA subspace geometry.

% To resolve this question, we develop a sequential hierarchical PAC-Bayesian analysis for LoRA-based continual learning. We derive a generalization bound with an explicit hyperposterior drift term that captures the sequential evolution of task priors across tasks. Under a frozen backbone, we show that both the KL complexity and the sharpness-dependent empirical term reduce to quantities supported on the admissible LoRA update subspace. This identifies task-conditioned local flatness along admissible LoRA directions, rather than arbitrary full-space flatness, as the relevant notion of flatness in this regime. The analysis further yields concrete empirical predictions about perturbation scope and sharpness type, which we test experimentally: sharpness-aware optimization restricted to LoRA updates consistently improves continual performance, whereas perturbations that extend into frozen backbone parameters provide no systematic benefit and can harm stability.

% Another line of work studies continual learning through the geometry of the loss landscape. Flat minima have long been associated with improved generalization and robustness~\cite{FlatMinima1997,wu2017understandinggeneralizationdeeplearning,chaudhariEntropySGDBiasingGradient2019,NEURIPS2020_518a38cc,NEURIPS2021_995f5e03,JMLRAnEmpiricalInvestigatio,yangRevisitingFlatnessAwareOptimization2025}, and sharpness-aware optimization methods such as SAM and its variants operationalize this principle by favoring solutions that remain stable under local perturbations~\cite{foretSharpnessAwareMinimizationEfficiently2021,pmlr-v151-bisla22a,Zhang_GAM,NEURIPS2024_C-Flat}. Prior work has further suggested that flatter solutions can mitigate catastrophic forgetting in continual learning~\cite{NEURIPS2020_518a38cc,JMLRAnEmpiricalInvestigatio,shiOvercomingCatastrophicForgetting2021,liRevisitingCatastrophicForgetting2024}. However, most existing flatness-based arguments are inherited from full-parameter training, where perturbations and curvature are defined over the whole model~\cite{li2024flatlora}. In LoRA-based continual learning, by contrast, adaptation is confined to task-conditioned admissible low-rank update directions. This means that flatness measured over the whole model is not automatically the flatness notion that governs continual adaptation, because flatness is defined relative to the task-conditioned loss, while under frozen-backbone LoRA adaptation, only the admissible low-rank update directions are available to fit the current task data. The central question is thus not whether flatness is beneficial in general, but whether continual generalization is governed by full-space flatness or by task-conditioned local flatness along admissible LoRA directions.

% We ground this question in a motivating empirical observation. As visualized in Figure~\ref{fig:3d_loss_landscape}, applying sharpness-aware minimization only within the low-rank adapter subspace already yields a solution whose surrounding loss landscape appears wider and flatter than that obtained by SGD when probed in the full parameter space. This observation suggests that geometric control within the trainable LoRA subspace can translate into measurable differences in the effective sensitivity of the full model. At the same time, it does not explain when such a transfer should occur, why it should matter for continual generalization, or which notion of flatness should be regarded as theoretically relevant in this setting.

% To answer this question, we develop a sequential hierarchical PAC-Bayesian analysis for LoRA-based continual learning. We derive a generalization bound with an explicit hyperposterior drift term that captures the sequential evolution of task priors across tasks. Under a frozen backbone, we show that both the KL complexity and the sharpness-dependent empirical term reduce to quantities supported on the admissible LoRA update subspace. This identifies task-conditioned local flatness along admissible LoRA directions, rather than arbitrary full-space flatness, as the relevant notion of flatness in this regime. Guided by this result, we further study perturbation scope and sharpness type in LoRA-based continual learning, and show empirically that sharpness-aware optimization restricted to LoRA updates consistently improves continual performance, whereas perturbations that extend into frozen backbone parameters provide no systematic benefit and can harm stability.



Our main contributions are summarized as follows:
\begin{itemize}
    \item \textbf{Which flatness matters in LoRA-based continual learning.}
    We formulate this question in a sequential PECL setting and show that, under frozen-backbone LoRA adaptation, continual learning takes place in a task-conditioned admissible update geometry rather than in the full parameter space. This motivates a distinction between arbitrary full-space flatness and task-conditioned local flatness along admissible LoRA directions.
    \item \textbf{Sequential hierarchical PAC-Bayesian analysis for LoRA-based PECL.}
    We derive a sequential hierarchical PAC-Bayesian bound with an explicit hyperposterior drift term that captures the sequential evolution of task priors across tasks. Under frozen-backbone LoRA adaptation with adapter-supported posteriors, both the KL complexity and the sharpness-dependent empirical term reduce to quantities supported on the admissible LoRA update subspace $\mathcal{W}_\Delta$. This identifies adapter-subspace flatness as the bound-relevant notion of flatness in PECL.
    \item \textbf{Perturbation design under admissible LoRA geometry.}
    Guided by the theory, we study perturbation design in LoRA-based continual learning along two axes: perturbation scope and sharpness type. Across multiple vision and language benchmarks, as well as different LoRA organizations, we show that sharpness-aware optimization restricted to LoRA updates consistently improves continual performance. Curvature diagnostics further show that lower subspace sharpness is associated with higher continual accuracy and reduced forgetting across tasks.
\end{itemize}



% Our main contributions are summarized as follows: 
% \begin{itemize}
%     \item Theoretical Framework: 
%     % We formulate a novel 
%     % subspace-aware PAC-Bayesian 
%     % bound tailored for PECL, theoretically proving that under the PECL regime with frozen backbone, the generalization gap is governed entirely by expected sharpness and KL divergence within the adapter subspace, which makes subspace-restricted flatness control both sufficient and preferable.
%     % generalization error is governed by the expected sharpness within the adapter subspace, thereby establishing the sufficiency of local flatness control.
%     We derive a PAC-Bayesian generalization bound specialized to PECL with a frozen backbone and adapter-subspace posteriors. In this regime, both the sharpness and the KL complexity that govern the bound can be written in terms of the trainable adapter subspace $\mathcal{W}_\Delta$. This identifies adapter-subspace flatness as the relevant notion of flatness for PECL.
    
%     \item  Geometric Mechanism: 
%     % We unveil the mechanism linking weight-space flatness to feature-space robustness, demonstrating that LoRA updates are equivalent to the abstract subspace and that minimizing factor-space curvature suppresses the Fisher Information Matrix to stabilize learned representations.
%     We clarify the mechanism linking weight-space flatness to feature-space robustness in LoRA-based PECL. We show that LoRA updates can be identified with an abstract update subspace and that minimizing curvature in factor space corresponds to suppressing dominant eigenvalues of the Fisher information matrix, which stabilizes learned representations.
    
%     \item Empirical Validation: 
%     We provide empirical evidence across multiple benchmarks and adapter architectures that subspace-restricted flatness optimization enhances both stability and plasticity in PECL. Our experiments further illustrate that full-parameter perturbations do not consistently improve performance in this regime and can introduce additional instability,
%     % We provide comprehensive empirical evidence across multiple benchmarks and adapter architectures, confirming that subspace-restricted flatness optimization enhances both stability and plasticity, while clarifying the limits of global perturbations in the PECL regime.
% \end{itemize}



% * Theoretical Framework: We formulate a novel subspace-aware PAC-Bayesian bound tailored for PECL, theoretically proving that generalization error is governed by the expected sharpness within the adapter subspace, thereby establishing the sufficiency of local flatness control. * Geometric Mechanism: We unveil the mechanism linking weight-space flatness to feature-space robustness, demonstrating that LoRA updates are geometrically equivalent to the abstract subspace and that minimizing factor-space curvature suppresses the restricted Fisher Information Matrix to stabilize learned representations. * Empirical Validation: We provide comprehensive empirical evidence across multiple benchmarks and adapter architectures, confirming that subspace-restricted flatness optimization enhances both stability and plasticity, while clarifying the limits of global perturbations in the PECL regime


% We therefore ask whether and how LoRA subspace flatness affects PECL. To make this question operational, we decompose it into three parts: (i) sufficiency, (ii) mechanism, and (iii) governing factors.
% \begin{itemize}
%   \item \textbf{Q1 (Sufficiency).} Under matched compute and memory budgets, does increasing LoRA subspace flatness consistently improve continual learning performance across NLP and vision, as measured by current task accuracy, forgetting (stability), and forward transfer (plasticity)? How does it shift the balance between stability and plasticity as quantified by these metrics?
%  \item \textbf{Q2 (Mechanism).} How does flatness in parameter space affect the loss? Specifically, is its effect mediated by the feature space, such that increased LoRA subspace flatness shapes the learned representations and yields smoother and more stable features that account for the observed performance gains?
% \item \textbf{Q3 (Governing factors)} How do LoRA architectures and flatness strategies modulate the magnitude and limits of the effect?
% \end{itemize}
% Taken together, Q1-Q3 define our central question: whether adapter-subspace flatness is sufficient to govern continual learning performance, through which mechanism it operates, and how its effect is modulated by adapter architectures and flatness strategies. 


% {%
% % {%
% %
% % \setlength{\intextsep}{4pt}
% % \setlength{\textfloatsep}{4pt}
% % \setlength{\abovecaptionskip}{6pt}
% % \setlength{\belowcaptionskip}{pt}
% % \captionsetup{skip=10pt}
% \begin{figure}[ht]
% \centering
% \includegraphics[width=0.6\linewidth]{figures_TRML/ACC_dataset_pair_Robustness_to_Corruptions_on_ImageNet-C_Robustness_to_Perturbation_on_ImageNet-P2.pdf}
% % \includegraphics[width=0.9\linewidth]{figures_TRML/hebing_robutness2.pdf}
% \vspace{-6pt}
% \caption{\textbf{Robustness of perturbation types under CL} Classification average accuracy $\operatorname{Acc}_t^{\mathrm{cls}}$ over tasks evaluated under distribution shift on ImageNet-C  and ImageNet-P.
% Across tasks and LoRA organizations, both $\mathrm{AS}^{(0)}$ and $\mathrm{AS}^{(1)}$ consistently outperform SGD, and the gap widens in later tasks.}
% \label{fig:robustness_imagenet_cp}
% \vspace{-10pt}
% \end{figure}
% }%

\section{Related Work}

\minisection{LoRA-Based CL.}
PECL has emerged as a practical solution for incrementally adapting large pre-trained models
% , with LoRA becoming a representative approach due to its efficiency and compatibility with frozen backbones
~\cite{pmlr-v97-houlsby19a,Wang_2022_CVPR,dingParameterefficientFinetuningLargescale2023,xu2023parameterefficientfinetuningmethodspretrained}. By introducing trainable low-rank matrices, LoRA enables continual adaptation with minimal overhead, and has demonstrated strong performance~\cite{hu2022lora}.
However, LoRA-based methods remain vulnerable to forgetting and task interference
% , particularly in long and diverse task sequence
s~\cite{Gao_2023_ICCV,10.1145/3730436.3730456,yang2024parametercollisionhinderingcontinual}. 
 % Gated variants further select or combine adapters with task-dependent weights~\citep{Gao_2023_ICCV,10.1145/3730436.3730456,zhao-etal-2024-sapt,cheng2025exploiting,wu2025sdlorascalabledecoupledlowrank}.
Recent variants such as SDLoRA~\cite{wu2025sdlorascalabledecoupledlowrank} and OLoRA~\cite{NEURIPS2020_70d85f35,wang2023orthogonal} aim to alleviate these issues via adapter modularity and orthogonal constraints. 
Although such designs yield empirical improvements, they lack a principled understanding of how adapter structure affects forgetting and generalization.
% In particular, the interaction between adapter subspace design and solution stability remains insufficiently explored.

\minisection{Flat Minima.}
Prior work has explored flat minima in CL through sharpness-aware optimization~\cite{pmlr-v151-bisla22a,Zhang_GAM,NEURIPS2024_C-Flat,yangRevisitingFlatnessAwareOptimization2025,chen2026sharpnessflatnessdecompositionframework}, mode connectivity~\cite{izmailov2019averagingweightsleadswider,mirzadeh2020linearmodeconnectivitymultitask}, and improved representation learning~\cite{hu2022doesselfsupervisedpretrainingperform,ramasesh2022effect}.
These methods proposed for full model training do not account for continual training dynamics in the low-rank subspace~\cite{yangRevisitingFlatnessAwareOptimization2025}. Whether flatness in the low-rank subspace alone can ensure generalization in CL remains unclear, motivating further investigation.


\minisection{PAC-Bayesian Theory.}
The PAC-Bayesian framework provides a principled foundation for analyzing generalization by balancing empirical risk and posterior complexity~\cite{mcallesterPACBayesianModelAveraging1999,mcallesterPACBayesianStochasticModel2003,pentinaPACBayesianBoundLifelong2014,germainPACBayesianTheoryMeets2016,lotfiPACBayesCompressionBounds2022}.
Online PAC-Bayes~\citep{NEURIPS2022_onlineBayes,NEURIPS2024_xinrui} allows time-varying, data-dependent priors, but it is formulated at the level of individual examples rather than tasks and does not explicitly model task-level inheritance in continual learning.
Hierarchical PAC-Bayes bounds for meta learning~\citep{pentinaPACBayesianBoundLifelong2014,JMLR:v24:22-1254} introduce hyperposteriors over task priors, but they assume a single stationary hyperposterior and therefore do not model the sequential, history-dependent evolution of the inductive bias that underlies forgetting in continual learning. Paper \cite{friedman2026pac} is a concurrent PAC-Bayes analysis for sequential CL, deriving a cumulative-loss bound. Our contribution differs in two respects. First, we introduce a hierarchical sequential PAC-Bayes framework, separating the hyperposterior level from task-specific priors. Second, we specialize the analysis to PECL, where both the KL complexity and the sharpness-smoothed empirical term reduce to the task-permissible LoRA subspace .
% In continual learning, the inductive bias must adapt to the observed task history, and this non-stationarity is particularly pronounced in LoRA-based PECL, where the update scope and the initialization of incremnental adapters could change with training.


\section{Preliminaries}

\subsection{LoRA-Based CL}
\label{sec:pecl-setup}

Let $\mathcal X$ be the input space, $\mathcal Y$ the label space, and $\mathcal Z = \mathcal X \times \mathcal Y$ the data space.
A model with parameters $W \in \mathbb R^d$ induces a predictor $f_W : \mathcal X \to \mathcal Y$.
We consider a bounded loss $\ell : \mathcal H \times \mathcal Z \to [0,1]$ and write $\ell(W,z)$ for $\ell(f_W,z)$. 
Given a distribution $\mathcal D$ over $\mathcal Z$ and a training set $S=\{z_j\}_{j=1}^m$ drawn i.i.d.\ from $\mathcal D$, the population and empirical risks are
$
L_{\mathcal D}(W)
=
\mathbb E_{z\sim\mathcal D}\ell(W,z),
\hat L_{S}(W)
=
\frac{1}{m}\sum_{j=1}^{m}\ell(W,z_{j}).
$



% LoRA-based PECL adapts a frozen backbone by learning low-rank updates across tasks. A key design axis is how these updates are organized in parameter space, since the resulting sharing pattern determines the degree of knowledge transfer and cross-task interference. We focus on \emph{shared-subspace} PECL, where tasks interact through a common adapter mechanism, and summarize three representative organizations.

% \textbf{Fully shared.}
% SeqLoRA reuses a single LoRA pair $(A,B)$ for all tasks and updates it sequentially on the fixed backbone $W_0$. This maximizes parameter sharing but also couples tasks most strongly, making interference accumulate within the same low-rank directions.

% \textbf{Incremental shared.}
% Incremental designs expand the adapter set over time by introducing a new pair $(A_t,B_t)$ at task $t$ while keeping $W_0$ and earlier adapters fixed, and then composing multiple adapters at inference.

% \textbf{Constraint incremental shared.}
% Constraint-based methods regularize the geometry of the adapter space to reduce conflicts based incremental designs. OLoRA~\citep{wang2023orthogonal} encourages orthogonality between task-specific directions to reduce parameter overlap, while gradient-based approaches align or project updates to avoid previously occupied directions~\citep{yang2024parametercollisionhinderingcontinual}.

In continual learning with \(T\) tasks indexed by \(t=1,\dots,T\), each task provides a sample \(S_t =\{z_{tj}\}_{j=1}^{m_t}\) drawn i.i.d. from a task distribution \(\mathcal{D}_t\). LoRA-based PECL adapts a frozen backbone by learning task-wise low-rank updates. A key design axis is how these updates are organized in parameter space. We focus on methods that share a common parameter space and distinguish three categories.
% LoRA-based PECL adapts a frozen backbone by learning low-rank updates across tasks and A key design axis is how these updates are organized in parameter space.  We focus on methods with a shared parameter space and
% distinguish three categories.

% In the context of CL, we focus on three representative LoRA-
% based training strategies:
% \begin{itemize}
% \vspace{-1em}
% \item \textbf{Fully shared (SeqLoRA)} : A single LoRA pair $(A, B)$ is trained sequentially and shared across tasks while keeping the backbone weights $W_0$ frozen. The final model after $N$ tasks is given by $W_0 + AB$.
% \vspace{-0.5em}
% \item \textbf{Incremental shared(IncLoRA)} : A new LoRA pair $(A_t, B_t)$ is incrementally added for each task $t$, and only the current pair is updated while all previous modules and $W_0$ are fixed. The resulting model is $W_0 + \sum_{t=1}^N A_t B_t$.
% \vspace{-0.25em}
% \item \textbf{Constraint shared(OLoRA)} : Extending IncLoRA, OLoRA~\cite{wang2023orthogonal} introduces an orthogonality constraint across tasks to encourage adapter subspace disentanglement. Specifically, during task $t$, a penalty $\mathcal{L}_{\text{orth}}^{(t)} = \sum_{i=1}^{t-1} \|A^{(i)\top} A^{(t)}\|_F^2$ is applied, where $A^{(i)}$ denotes the LoRA-A matrix from task $i$.
% \end{itemize}

% For completeness, some methods allocate fully disjoint adapters per task and activate them conditionally by task identification at inference. Such isolated designs are not the focus of our analysis.

% \subsection{PAC-Bayesian Bound}

% The PAC-Bayesian viewpoint~\cite{mcallesterPACBayesianModelAveraging1999,mcallesterPACBayesianStochasticModel2003} place probability measures over the hypothesis space $\mathcal{H}$ rather than focusing on a single function: starting from a prior $P\in\mathcal M(\mathcal H)$, the learner observes data $S$ and outputs a posterior $Q(\cdot\mid P,S)\in\mathcal M(\mathcal H)$. The corresponding posterior-averaged population and empirical risks  are $
% L_{\mathcal D}(Q)
% =
% \mathbb E_{W \sim Q(\cdot \mid P,S)}\mathbb E_{z\sim\mathcal D}\ell(W,z)
% $ and
% $
% \hat L_{S}(Q_t)
% =
% \mathbb E_{W \sim Q(\cdot \mid P,S)}\frac{1}{m}\sum_{j=1}^{m_t}\ell(W,z_{j}).
% $
% provides a generalization guarantee for a single task under a fixed data distribution and a prior that is fixed before observing the corresponding training set. Continual learning instead proceeds over a task sequence and the learner updates its hypothesis after observing task data. This sequential setting implies that the inductive bias should adapt with the task history, so the task prior at time $t$ must reflect information from tasks $1,\dots,t-1$, which cannot be captured by treating $\{P_t\}$ as externally specified constants.






% A classical PAC--Bayes bound~\cite{mcallesterPACBayesianStochasticModel2003} for a single task under a fixed distribution is stated as follows.

% \begin{theorem}[PAC-Bayes Generalization Bound for single task]
% \label{theorem:classical}
% For any loss $\ell: \mathcal{F} \times \mathcal{X} \times \mathcal{Y} \to [0,1]$, with probability at least $1-\delta$ over the choice of training set $S$:
% % {%
% % % \setlength{\abovedisplayskip}{2pt}
% % \setlength{\belowdisplayskip}{3pt}
% % % \setlength{\abovedisplayshortskip}{3pt}
% % \setlength{\belowdisplayshortskip}{3pt}
% %     \[\mathbb{E}_{f \sim Q}[L_{\mathcal{D}}^{\ell}(f)] \leq \mathbb{E}_{f \sim Q}[\hat{L}_S^{\ell}(f)] + \sqrt{\frac{KL(Q\|P) + \log \frac{m}{\delta}}{2(m-1)}}\]
% % }
% % \end{theorem}

\subsection{PAC-Bayesian Bound and Sharpness}
% Flat minima are often associated with improved generalization and robustness in standard supervised learning.

The PAC-Bayesian viewpoint~\cite{mcallesterPACBayesianModelAveraging1999,mcallesterPACBayesianStochasticModel2003} place probability measures over the hypothesis space $\mathcal{H}$ rather than focusing on a single function: starting from a prior $P\in\mathcal M(\mathcal H)$, the learner observes data $S$ and outputs a posterior $Q(\cdot\mid P,S)\in\mathcal M(\mathcal H)$. The corresponding posterior-averaged population and empirical risks  are $
L_{\mathcal D}(Q)
=
\mathbb E_{W \sim Q(\cdot \mid P,S)}\mathbb E_{z\sim\mathcal D}\ell(W,z)
$ and
$
\hat L_{S}(Q_t)
=
\mathbb E_{W \sim Q(\cdot \mid P,S)}\frac{1}{m}\sum_{j=1}^{m_t}\ell(W,z_{j}).
$

Motivated by PAC-Bayes analysis, SAM~\citep{foretSharpnessAwareMinimizationEfficiently2021} seeks to reduce the empirical contribution in generalization bounds by controlling the loss elevation in a local neighborhood. In particular, SAM can be interpreted as minimizing an \textbf{adversarial zeroth-order sharpness}~\citep{AWP2020,foretSharpnessAwareMinimizationEfficiently2021} on the training set $S$ within a radius $\rho$:
{%
\[\mathrm{AS}^{(0)}_S(W,\rho) =\max_{ \| \epsilon \|\leq \rho } \Big[ \hat{L}_{S}(W+{\epsilon})-\hat{L}_{S}(W)\ \Big].\]
}
% \vspace{-5mm}
A complementary formulation considers randomized perturbations and minimizes the expected loss elevation under Gaussian noise, referred to as \textbf{random zeroth-order sharpness}~\citep{liRevisitingRandomWeight2024,jiang2019fantasticgeneralizationmeasures}:
{%
\[
\mathrm{RS}_S^{(0)}({W}, \Sigma)
 =
 \mathbb{E}_{\varepsilon \sim \mathcal{N}(0,\Sigma)}
 \big[
 \hat{L}_{S}({W}+\varepsilon) - \hat{L}_{S}({W})
 \big],
 \]}Typical choices for $\mathcal{N}(0,\Sigma)$ include isotropic Gaussians and variants such as filter-wise Gaussians~\citep{pmlr-v151-bisla22a}. Moreover, SAM can be interpreted as optimizing the mean loss under a fixed perturbation variance: for a suitable correspondence between radius $\rho$ and covariance $\Sigma$, 
 % $\mathrm{AS}^{(0)}_S(W,\rho)$ provides an upper bound on $\mathrm{RS}_S^{(0)}({W}, \Sigma)$~\cite{mollenhoffSAMOptimalRelaxation2022}.
 the adversarial zeroth-order sharpness provides an upper bound on the random zeroth-order sharpness~\cite{mollenhoffSAMOptimalRelaxation2022}.

% Beyond loss elevation, sharpness can be measured through the neighborhood gradient norm.
The \textbf{adversarial first-order sharpness} proposed in GAM~\citep{Zhang_GAM} is defined as the maximum gradient norm within a radius $\rho$ on the training set:
{%
\[\mathrm{AS}^{(1)}_S(W,\rho) = \max_{ \| \epsilon \|\leq \rho } \| \nabla_{W} \hat{L}_S(W+\epsilon)\|.\]}Compared with zeroth-order sharpness which measures  loss increase, the first-order criterion emphasizes the sharpness of the gradient itself, and is closely tied to dominant curvature directions. For any $W$ and fixed $\rho$, the adversarial first-order sharpness dominates the adversarial zeroth-order sharpness, hence minimizing $\mathrm{AS}^{(1)}_S$ also reduces $\mathrm{AS}^{(0)}_S$, making it a stronger flatness measure~\citep{Zhang_GAM}.












% \section{Preliminaries}

% \subsection{Parameter efficient Continual Learning}
% Let $\mathcal{X}$ be the input space, $\mathcal{Y}$ the label space, $\mathcal{Z} = \mathcal{X} \times \mathcal{Y}$ the data space and $\mathcal{H}$ denote the hypothesis space.
% Let $W \in \mathbb{R}^d$ denote the model parameters, which induce a predictor
% $f_W \in \mathcal{H}$, with $f_W : \mathcal{X} \to \mathcal{Y}$.
% We fix a bounded loss \(\ell:\mathcal{H}\times\mathcal{Z}\to[0,1]\) and, for brevity, write \(\ell(W,z)\) for \(\ell(f_W,z)\). 
% In continual learning with \(T\) tasks indexed by \(t=1,\dots,T\), each task provides a sample \(S_t =\{z_{tj}\}_{j=1}^{m_t}\) drawn i.i.d. from a task distribution \(\mathcal{D}_t\). Define the empirical risk and expected risk for a model \(W\) as:
% $$
% \hat{L}_{S_t}(W) = \frac{1}{m_t} \sum_{j=1}^{m_t} \ell(W, z_{tj}), \quad
% L_{\mathcal{D}_t}(W) = \mathbb{E}_{z \sim \mathcal{D}_t} \ell(W, z).
% $$

% %=\overline{\mathrm{ACC}}_T


% % For each task \(t \in \{1, \dots, T\}\), let \(S_t = \{z_{tj}\}_{j=1}^{m_t}\) be i.i.d. draws from \(\mathcal{D}_t\), and let \(\ell(w, z) \in [0,1]\) be the loss function. Define the empirical risk and expected risk for a model \(W\) as:
% % \[
% % \hat{L}_{S_t}(W) = \frac{1}{m_t} \sum_{j=1}^{m_t} \ell(W, z_{tj}), \quad
% % L_{\mathcal{D}_t}(W) = \mathbb{E}_{z \sim \mathcal{D}_t} \ell(W, z).
% % \]


% % We consider a bounded loss $\ell : \mathcal{H} \times \mathcal{Z} \to \mathbb{R}^+$,
% % uniformly bounded by a constant $K>0$.
% % For notational convenience we write $\ell(W,z)$ to mean $\ell(f_W,z)$.


% % To analyze the role of sharpness in CL, we adopt a hierarchical PAC-Bayesian framework.
% % For each task \(t \in \{1, \dots, T\}\), let \(S_t = \{z_{tj}\}_{j=1}^{m_t}\) be i.i.d. draws from \(\mathcal{D}_t\), and let \(\ell(w, z) \in [0,1]\) be the loss function. Define the empirical risk and expected risk for a model \(W\) as:
% % \[
% % \hat{L}_{S_t}(W) = \frac{1}{m_t} \sum_{j=1}^{m_t} \ell(W, z_{tj}), \quad
% % L_{\mathcal{D}_t}(W) = \mathbb{E}_{z \sim \mathcal{D}_t} \ell(W, z).
% % \]


% In the context of PECL, we categorize adapter based methods by how task adapters interact on a frozen backbone $W_0$, contrasting shared adapter subspaces with task specific ones.
% % In the context of PECL, we categorize adapter-based methods by whether task adapters share a parameter space on a frozen backbone $W_0$. 
% Approaches with disjoint parameters that train separately and activate one adapter at inference do not induce cross-task sharing and are outside the scope of this paper\cite{Yu_2024_CVPR}.%\todo{CITE}
% We focus on methods with a shared parameter space and distinguish three categories.

% \begin{itemize}
%     \item  \textbf{Fully shared subspace}. A single low-rank subspace is reused across tasks. In \textbf{SeqLoRA}, one LoRA pair $(A,B)$ is updated sequentially, yielding $W = W_0 + AB .$ Training and inference are both restricted to this common subspace.
%     \item  \textbf{Partially fused subspaces}. \textbf{IncLoRA} incrementally add a new pair $(A_t,B_t)$ at task $t$ while keeping $W_0$ and earlier adapters fixed:  $W_0 + \sum_{t=1}^N A_t B_t$.  Gated fusion activates several adapters for input $x$, $W = W_0 + \sum_{j \in \mathcal{A}(x)} \pi_j(x) A_j B_j ,$ where $\mathcal{A}(x)$ is the active adapter set and the weights satisfy $\pi_j(x)\ge 0$~\citep{zhao-etal-2024-sapt,cheng2025exploiting}.
 
%     \item  \textbf{Constraint-augmented sharing}. Explicit regularization shapes inter-task geometry. \textbf{OLoRA}~\citep{wang2023orthogonal} enforces orthogonality among adapter directions via $\sum_{i<t}|A^{(i)\top}A^{(t)}|_F^2$, promoting disentangled subspaces. Gradient-level designs, exemplified by {InfLoRA}~\citep{Liang_2024_CVPR}, reduce cross-task interference through alignment penalties or projection.
% \end{itemize}


% \subsection{Flat minima}
% \label{sec: flat minima}


% Flat minima are widely associated with improved generalization and robustness because they reduce the sensitivity of the empirical risk to small perturbations in parameter space~\citep{FlatMinima1997,wu2017understandinggeneralizationdeeplearning,NEURIPS2020_518a38cc,bahri2022sharpnessawareminimizationimproveslanguage,NEURIPS2021_995f5e03,JMLRAnEmpiricalInvestigatio,yangRevisitingFlatnessAwareOptimization2025}.
% We formalize flatness through sensitivity to parameter perturbations and introduce representative  sharpness definitions widely used in the literature.
% The zeroth order sharpness measures the maximal loss increase within a radius $\rho$~\citep{AWP2020,foretSharpnessAwareMinimizationEfficiently2021}: 
% $\mathrm{Sh}^{(0)}_S(W,\rho) :=\max_{{\epsilon}\in B(W,\rho)} \Big[\ \hat{L}_{S}(W+{\epsilon})-\hat{L}_{S}(W)\ \Big].$
% A %distributional 
% stochastic variant averages the loss increase under random perturbations drawn from a prescribed distribution~\citep{liRevisitingRandomWeight2024,jiang2019fantasticgeneralizationmeasures}:
% $\mathrm{E\text{-}Sh}_t(\hat{W}_t, \Sigma_t)
%  :=
%  \mathbb{E}_{\varepsilon \sim \mathcal{N}(0,\Sigma_t)}
%  \big[
%  \hat{L}_{S_t}(\hat{W}_t+\varepsilon) - \hat{L}_{S_t}(\hat{W}_t)
%  \big],$
% Typical choices for $\mathcal{N}(0,\Sigma_t)$ include an isotropic Gaussian and a filter-wise Gaussian~\citep{pmlr-v151-bisla22a}. The first order sharpness controls the maximal gradient norm in a neighborhood~\citep{Zhang_GAM}, which is consistent with the zeroth order definition:
% $\mathrm{Sh}^{(1)}_S(W,\rho) := \rho\cdot \max_{W'\in B(W,\rho)} \big\|\nabla_{W} \hat{L}_S(W')\big\|_2,$





% A second-order Taylor expansion of $ \hat{L}_S$ around $W$ gives
% $\hat{L}_S(W+\epsilon) =  \hat{L}_S(W) + g(W)^{\top}{\epsilon} + \tfrac12\,{\epsilon}^{\top}H(W)\,{\epsilon} + R_3(W,{\epsilon}),$
% % \begin{equation}
% %     \hat{L}_S(W+\epsilon) =  \hat{L}_S(W) + g(W)^{\top}{\epsilon} + \tfrac12\,{\epsilon}^{\top}H(W)\,{\epsilon} + R_3(W,{\epsilon}),
% % \end{equation}
% with $g(W)=\nabla \hat L_S(W)$, $H(W)=\nabla^2 \hat L_S(W)$, and remainder $R_3(W,{\epsilon})=O(|{\epsilon}|^3)$.
% Near a flat local minima, $g(\hat{W})\approx \mathbf 0$, so for small $\|{\epsilon}\|$,
% \begin{equation}
% \label{eq: approx}
%     \hat{L}_S(\hat{W}+{\epsilon})-\hat{L}_S(\hat{W})
% \approx \tfrac12\,{\epsilon}^{\top}H(\hat{W})\,{\epsilon}.
% \end{equation}
% % $
% % \hat{L}_S(\hat{W}+{\epsilon})-\hat{L}_S(\hat{W})
% % \approx \tfrac12\,{\epsilon}^{\top}H(\hat{W})\,{\epsilon}.
% % $
% % This approximation yields concrete correspondences between the sensitivity-based notions above and spectral properties of $H(\hat{W})$.
% Thus all sharpness notions introduced above reduce to different ways of aggregating the quadratic form $\epsilon^{\top}H(\hat W)\epsilon$: Sh$(0)$ approximates its maximal value over a small norm ball, Sh$(1)$ approximates its value along a prescribed perturbation direction, and $\mathrm{E\text{-}Sh}$ approximates its expectation under the chosen perturbation distribution. These identifications justify the use of spectral proxies such as $\lambda_{\max}\big(H(\hat{W})\big)$ and $\mathrm{tr}\big(H(\hat{W})\big)$, as well as their Fisher surrogates $\lambda_{\max}\big(F(\hat{W})\big)$ and $\mathrm{tr}\big(F(\hat{W})\big)$, where the Fisher coincides with the expected Hessian under standard regularity conditions~\citep{dinhSharpMinimaCan2017,caoTradeoffFlatnessOptimization2024,tsuzukuNormalizedFlatMinima2020,miMakeSharpnessAwareMinimization2022}. 
% % In practice, the empirical Fisher is also a common estimator~\citep{dinhSharpMinimaCan2017,caoTradeoffFlatnessOptimization2024,tsuzukuNormalizedFlatMinima2020,miMakeSharpnessAwareMinimization2022}.
% % Hence, zeroth order sharpness aligns with the dominant curvature, distributional sharpness aligns with the aggregate curvature, and first order sharpness controls the local growth through gradient magnitude, which is governed by the Hessian near a minimum. These links justify the use of spectral proxies such as $\lambda_{\max}\big(H(\hat{W})\big)$ and $\mathrm{tr}\big(H(\hat{W})\big)$, as well as Fisher-based $\lambda_{\max}\big(F(\hat{W})\big)$ and $\mathrm{tr}\big(F(\hat{W})\big)$, where the Fisher equals the expected Hessian under standard regularity conditions. In practice, the empirical Fisher is also a common estimator~\citep{dinhSharpMinimaCan2017,caoTradeoffFlatnessOptimization2024,tsuzukuNormalizedFlatMinima2020,miMakeSharpnessAwareMinimization2022}.
% To visualize local sharpness, one can plot one-dimensional loss slices along random or adversarial directions, two-dimensional planes spanned by normalized directions, or three-dimensional surface plots~\citep{goodfellowQualitativelyCharacterizingNeural2015,izmailovAveragingWeightsLeads2019,chaSWADDomainGeneralization2021,liVisualizingLossLandscape2018}.


% \subsection{PAC-Bayesian Generalization Bound}


% % Let $P$ be a prior and $Q$ a posterior over hypothesis space $P,Q \sim \mathcal{M}(\mathcal{H})$

% % The flatness metric above quantifies the local sensitivity of empirical loss to parameter perturbations. 
% To formalize generalization guarantees, the PAC settings place probability measures on the hypothesis space rather than focusing on a single function. The object of learning on a task has been a hypothesis $h \in \mathcal{H}$ over which one presumes a \textit{prior distribution} $P\sim \mathcal{M}(\mathcal{H})$ that is updated into a \textit{posterior} $Q\sim \mathcal{M}(\mathcal{H})$ based on observed data. For any distributions \(Q\) and \(P\) with absolute continuity \(Q \ll P\), define the KL divergence as \(\mathrm{KL}(Q \| P) = \int \log\left( \frac{dQ}{dP} \right) dQ\). With high probability over $S$, PAC-Bayesian bounds control the generalization risk  by the empirical risk plus a complexity term that scales with $\mathrm{KL}(Q\Vert P)$ and the sample size~\citep{shawe1997pac,mcallesterPACBayesianStochasticModel2003}.  Algorithmically, SAM~\citep{foretSharpnessAwareMinimizationEfficiently2021} minimizes the worst-case empirical loss within a small parameter neighborhood, which serves as a practical surrogate that reduces sharpness and aligns with the PAC-Bayesian objective.
% %$\max_{\|\varepsilon\|\le\rho}\hat L_{S}(W+\varepsilon)$


% % PAC-Bayesian bound~\citep{pentinaPACBayesianBoundLifelong2014,JMLR:v24:22-1254}  have been employed for multi-task knowledge transfer within a hierarchical Bayesian framework that learns a However, their analysis assumes that all task priors $P$ are independent draws from a {fixed} hyperprior, which is suitable for multi-task transfer learning but inadequate for CL. In CL, the hyper-distribution must evolve sequentially to  reflect the collection of historical information, and their framework does not support such dynamic adaptation. 
% % Recent advances in online PAC-Bayesian learning ~\citep{NEURIPS2022_onlineBayes,NEURIPS2024_xinrui} introduce time-varying, data-dependent sequence of priors \( P_t \) (each \( \mathcal{F}_{t-1}\)-measurable filtration) and dynamically updated posteriors \(Q_t\) for sequential data, offering flexibility for sequential data at the point level.  However, these methods focus on data-point sequences rather than task-level progression, and their theoretical foundations remain misaligned with structured continual learning frameworks like PECL.

% % % Algorithmically, the NsCE framework decomposes online PAC-Bayesian bound into empirical risk, throughput, and task-divergence terms~\citep{NEURIPS2024_xinrui}.

% % Despite these contributions, current theoretical frameworks exhibit two major limitations in the context of PECL. First, they are geometry-agnostic: they do not incorporate the local flatness of the loss landscape into the divergence terms that govern generalization. Second, they overlook the adapter subspace central to PECL, failing to link curvature within this subspace to PAC-Bayesian divergence. To bridge these gaps, we propose a novel PAC-Bayesian formulation that is both geometry-aware and subspace-aware, specifically tailored to the PECL framework.

% Building on classical PAC-Bayesian theory, hierarchical bounds for lifelong and meta-learning~\citep{pentinaPACBayesianBoundLifelong2014,JMLR:v24:22-1254} learn a hyper-posterior over task-specific priors and thus induce a two-level complexity term, but they usually assume that all priors are independent draws from a fixed hyper-prior and therefore do not model the sequential dependence and information accumulation that characterise continual learning. In parallel, online PAC-Bayesian work~\citep{NEURIPS2022_onlineBayes,NEURIPS2024_xinrui}introduces time-varying, data-dependent priors with guarantees that depend on cumulative data, yet it is formulated at the level of individual examples rather than tasks and does not consider task-conditioned hyper-posteriors.

% % Our analysis combines these two ideas in a task-level CL setting by using a time-varying hyper-posterior over priors that is measurable with respect to past tasks, which yields a CL-specific bound whose complexity decomposes into task-wise posterior divergences and a hyper-posterior drift term across tasks, and which we later specialise to parameter-efficient continual learning with LoRA adapters.

% % Our analysis combines these two lines of work and adapts them to a structured continual learning setting. 
% Compared to existing hierarchical PAC-Bayesian bounds for meta-learning, our Theorem provides a new task-level CL bound with a time-varying hyper-posterior adapted to the task filtration.
% At the task level, we use a time-varying hyper-posterior over priors, which inherits the filtration-based measurability of online PAC-Bayes while retaining the hierarchical structure. This yields a CL-specific PAC-Bayesian bound in which the complexity decomposes into a sum of task-wise posterior divergences and a hyper-posterior drift term across tasks. 

% % In the next section we instantiate this framework and then specialize it to parameter-efficient continual learning with LoRA adapters.



% \section{Revisiting Sharpness for PECL}
% \label{sec:pac-pecl}

\section{PECL as constrained adaptation geometry}


\section{PAC-Bayesian Bound for Continual Learning}
\label{sec:single_to_cl}

We model transferable inductive bias by treating the task prior as a random object and introducing a hyperposterior $\mathcal P\in\mathcal M(\mathcal M(\mathcal H))$ over task-level priors. Our derivation does build on PAC-Bayesian tools, but the technical challenge is setting-specific. Prior work\cite{pentinaPACBayesianBoundLifelong2014}, although framed as lifelong learning, assumes that tasks are sampled i.i.d. from a fixed environment distribution. As a result, its bound is invariant to task order and does not model trajectory-dependent interference, making it closer to a multi-task formulation than to sequential continual learning (Fig 2). In fact, it\cite{pentinaPACBayesianBoundLifelong2014} explicitly identifies relaxing the i.i.d. task assumption as future work, which is precisely the regime we address. Concretely, common CL mechanisms make the prior sequence nonstationary, including curvature-based regularization that updates the effective prior through Fisher or Hessian summaries~\citep{EWC2017,SI2017} and scope-allocation strategies that modify which parameters remain trainable over time~\citep{progressive2016,HAT2018,Supermasks2020,pmlr-v97-houlsby19a}, which naturally introduces a notion of drift between priors.

{%
\setlength{\intextsep}{8pt}        % [h] 浮动体与正文的上下间距（最关键）
\setlength{\textfloatsep}{6pt}     % [t]/[b] 时浮动体与正文间距（备用）
\setlength{\abovecaptionskip}{2pt} % 图与caption之间（上方）间距
\setlength{\belowcaptionskip}{0pt} % caption下方到正文的间距
\captionsetup{skip=4pt}            % 你已在用：图与caption的距离（更细粒度）

\begin{figure}[t]
\centering
\includegraphics[width=0.6\linewidth]{figures_TRML/multitask5.pdf}
\caption{
Time-varying Hyperposterior in Continual Learning. \textbf{(Left)} In multi-task learning, task priors are drawn from a single time-invariant hyperdistribution, so there is no sequential drift. \textbf{(Right)} In continual learning, the hyperposterior is updated over time, inducing a history-dependent cumulative drift over tasks.
% Time-varying Hyperposterior in Continual Learning. \textbf{Left:} Multi-task uses a single hyperdistribution drift term to sample task priors.
% \textbf{Right:} Continual learning updates the hyperposterior over time, yielding history-dependent cumulate hyperdistribution drift.
}
\vspace{-10pt}
\label{fig:showPAC3}
\end{figure}
}% 结束局部分组

Let $\mathcal F_{t-1}=\sigma(S_1,\dots,S_{t-1})$ denote the information available before observing task $t$. We model continual learning as a hyperposterior process $\{\mathcal P_t\}_{t=0}^T$ where $\mathcal P_{t-1}$ is $\mathcal F_{t-1}$ measurable, which yields a pathwise bound with an explicit hyperposterior drift penalty that quantifies the information budget of sequential belief updates. In continual learning, the prior used at step 
 is produced by a history-dependent. This creates two challenges: the prior is no longer a fixed object, but a sequential process adapted to the task history and the generalization analysis must accumulate pathwise over the task sequence.
 
We define the hierarchical test and empirical risks
$
\mathcal L_t
=
\mathbb E_{P_t \sim \mathcal{P}_{t-1}}
\,
\mathbb E_{W_t \sim Q_t(\cdot \mid P_t,S_t)}
\big[
L_{\mathcal D_t}(W_t)
\big]
$ and
$
\hat{\mathcal L}_t
=
\mathbb E_{P_t \sim \mathcal{P}_{t-1}}
\,
\mathbb E_{W_t \sim Q_t(\cdot \mid P_t,S_t)}
\big[
\hat L_{S_t}(W_t)
\big].
$ We propose the following PAC-Bayes bound for continual learning:

% \vspace{0.5em}
\begin{theorem}[PAC Bayes with time varying hyperposterior]
\label{thm:pathwise_pac}
For any $\delta\in(0,1]$, with probability at least $1-\delta$ over $S_{1:T}$,
{%
% \setlength{\abovedisplayskip}{2pt}
% \setlength{\belowdisplayskip}{3pt}
% \setlength{\abovedisplayshortskip}{3pt}
% \setlength{\belowdisplayshortskip}{3pt}
\begin{equation*}
% \resizebox{\linewidth}{!}{$
\begin{aligned} 
&\quad\frac{1}{T} \sum_{t=1}^T \mathcal L_t
\le  
\frac{1}{T} \sum_{t=1}^T \hat{\mathcal L}_t  + 
\sqrt{
\frac{
\displaystyle
\sum_{t=1}^T
\mathbb E_{P_t\sim\mathcal P_{t-1}}
\mathrm{KL}\bigl(Q_t \,\Vert\, P_t\bigr)
+
\sum_{t=1}^T
\mathrm{KL}\bigl(\mathcal P_t\Vert\mathcal P_{t-1}\bigr)
+
\log\frac{1}{\delta}}
{2\displaystyle\sum_{t=1}^T (m_t - 1)}
}.
\end{aligned}
% $
% }
\end{equation*}
}
\end{theorem}

Theorem~\ref{thm:pathwise_pac} bounds the trajectory-averaged hierarchical test risk by the corresponding trajectory-averaged empirical risk plus a complexity penalty, aligning with standard continual learning evaluation over a task sequence.
The complexity decomposes into a within task adaptation term $\mathbb E_{P_t\sim\mathcal P_{t-1}}\mathrm{KL}(Q_t\Vert P_t)$ and a hyperdistribution drift term $\mathrm{KL}(\mathcal P_t\Vert\mathcal P_{t-1})$. The former quantifies how much the learner must deviate from the task prior to fit $S_t$, while the latter quantifies how the task prior generation mechanism evolves as new tasks are incorporated.
This drift is intrinsic to continual learning because the prior used at step $t$ is typically produced by a history-dependent rule rather than fixed externally,
for example by specifying the permissible update scope~\citep{HAT2018,Supermasks2020,hu2022lora}, or by shaping the prior geometry using curvature summaries as regularization~\citep{EWC2017,SI2017}.
Theorem~\ref{thm:pathwise_pac} also contains the classical exchangeable-task hierarchical PAC-Bayes bound~\citep{pentinaPACBayesianBoundLifelong2014} as a special case: if the hyperposterior is not a sequence change then the result reduces to a single shared hyperdistribution over task priors. We refer to Appendix~\ref{proof:theorem all} for the detailed proof of Theorem~\ref{thm:pathwise_pac}.



\section{Which flatness matters? to LoRA}
\label{sec:kl-decomp-pecl}

% In PECL, we represent the task-$t$ parameters in an affine form
% \[
% W_t = W^{\mathrm{froz}}_t + \Delta W_t ,
% \]
% where $W^{\mathrm{froz}}_t\in\mathbb R^d$ collects all parameters that are frozen while learning task $t$
% (with $d=\dim(\mathrm{vec}(W))$ under the chosen vectorization), and $\Delta W_t$ denotes the trainable
% increment induced by the current-task LoRA coordinates.
% Accordingly, the feasible learning dynamics are restricted to a task-dependent \emph{increment-direction subspace}:
% \[
% \Delta W_t \in \mathcal W_{\Delta,t}\subseteq \mathbb R^d .
% \]
% Here $\mathcal W_{\Delta,t}$ denotes the task-admissible linear space of \emph{increment directions} in a
% neighborhood of the learned solution, which provides the geometric support for the PAC--Bayesian
% distributions and is instantiated from the LoRA implementation via the local Jacobian construction
% in Appendix~\ref{app:lora_subspace_bridge}.
\begin{figure}[t]
\centering
% \begin{minipage}{1\linewidth}
\centering
\includegraphics[width=0.6\linewidth]{figures_TRML/KLshiyi1.pdf}
% {figures_TRML/KLshiyitu1.pdf}
% \caption{Illustration of the KL divergence decomposition and coordinate equivalence. 
%     (1) \textbf{KL Decomposition}: In PECL, the overall model complexity, measured by $\mathrm{KL}(Q^{\mathrm{full}} \| P^{\mathrm{full}})$, is solely determined by the distribution of the updated parameters in the adapter subspace $\mathcal{W}_\Delta$, formalized as $\mathrm{KL}(Q^{\Delta} \| P^{\Delta})$, while the frozen backbone contributes no divergence. 
%     (2) \textbf{Two equivalent views}: The optimization in the low-rank factor coordinates $(A, B)$ is theoretically isomorphic to operating directly in the update subspace $\mathcal{W}_\Delta$. This equivalence preserves geometric properties such as flatness and curvature, which allows sharpness-aware optimization to be confined to the factorized LoRA parameters.}
% \vspace{-4pt}
%$T_t(\Delta)=W_t^{\mathrm{froz}}+\Delta$
%: Full-Model KL vs Adapter-Space KL.
\caption{\textbf{KL Decomposition}
 The posterior/prior defined on the adapter update subspaces $(Q_t^\Delta,P_t^\Delta)$ induce full-model distributions $(Q_t,P_t)$ with $\mathrm{KL}(Q_t\Vert P_t)=\mathrm{KL}(Q_t^\Delta\Vert P_t^\Delta)$, so the complexity term depends only on task permissible update subspace.}
\vspace{-10pt}
% \caption{ Moving from Full-Finetuning to LoRA: Full-Model KL vs Adapter-Space KL.
% $\mathrm{KL}(Q^{\mathrm{full}}\|P^{\mathrm{full}})=\mathrm{KL}(Q^{\Delta}\|P^{\Delta})$, so the complexity depends only on the adapter subspace $\mathcal W_{\Delta}$.
% LoRA factors $(A,B)$ parameterize the same update subspace, allowing sharpness control to be applied in factor space.
% }
\label{fig: KL_decompose}
% \end{minipage}
\end{figure}



In PECL, the task-$t$ parameter vector can be written as
$
W_t = W^{\mathrm{froz}}_t + \Delta W_t,
$
where $W^{\mathrm{froz}}_t$ refers to
% $W^{\mathrm{froz}}_t\in\mathbb R^d$ collects
%from previous task while
all frozen parameters when learning task $t$, and $\Delta W_t$ are the trainable LoRA parameters for the current-task. The permissible update dynamics are thus restricted to a task-dependent incremental subspace
$
\Delta W_t \in \mathcal W_{\Delta,t}
$
where $\mathcal W_{\Delta,t}$ denotes the task-specific space of possible directions around the learned
solution (formalized in Appendix \ref{app:lora_subspace_bridge}). Consequently, the hypothesis class at task $t$ is the affine
set $W^{\mathrm{froz}}_t + \mathcal W_{\Delta,t}$ and all task-dependent randomness is supported
on $\mathcal W_{\Delta,t}$ and then pushed forward to the affine class by translation. 

% in the PAC--Bayes analysis 

% \in \mathcal W_{\Delta,t}
We formalize the LoRA adapter posterior on a Gaussian distribution in the permissible space $\mathcal W_{\Delta,t}$ as
$
Q_t^\Delta = \mathcal N(\widehat{\Delta W}_t,\Sigma_t)$ where
$\widehat{\Delta W}_t$ refers to the learned weight update 
and $\Sigma_t$ is a covariance around $\widehat{\Delta W}_t$ on $\mathcal W_{\Delta,t}$ .
Equivalently, sampling $W_t\sim Q_t$ on $W^{\mathrm{froz}}_t+\mathcal W_{\Delta,t}$ amounts to 
% sampling $ \Delta W_t \sim Q_t^{\Delta}$ \text{ on } $W_{\Delta,t}$ and 
setting $W_t=W^{\mathrm{froz}}_t+\Delta W_t$ and sampling $ \Delta W_t \sim Q_t^{\Delta}$ where $\Delta W_t = \widehat{\Delta W}_t + \varepsilon$
and
$
\varepsilon \sim \mathcal N(0,\Sigma_t).$

% \text{ on }  \mathcal W_{\Delta,t}.$

% and setting $W=W_{\mathrm{frozen}}+\Delta W$.


We propose the following lemma to formalize the resulting translation-invariance of the KL divergence:
% , and its proof is included in Appendix~\ref{proof:Lemma1}.
\begin{lemma}
\label{lem:kl-decomp-pecl}
Assume $Q_t^\Delta \ll P_t^\Delta$. Then $Q_t \ll P_t$ and
\[
\mathrm{KL}\bigl(Q_t \,\Vert\, P_t\bigr)
=
\mathrm{KL}\bigl(Q_t^\Delta \,\Vert\, P_t^\Delta\bigr).
\]
\end{lemma}
Lemma~\ref{lem:kl-decomp-pecl} shows that, under a frozen backbone, the  KL divergence in the full parameter space is exactly the KL divergence on the task-permissible update subspace. We refer to Appendix~\ref{proof:Lemma1} for the proof.

 
Now, we define the adapter subspace sharpness as the expected loss elevation under perturbations supported on the adapter update subspace $\mathcal W_{\Delta,t}$ for task $t$. The following lemma~\ref{lem:subspace-loss-pecl} isolates an additive loss elevation caused by perturbations supported on $\mathcal W_{\Delta,t}$.

\begin{lemma}[Task adapter subspace sharpness]
\label{lem:subspace-loss-pecl}
For each task $t$,
{%
\setlength{\abovedisplayskip}{1pt}
\setlength{\belowdisplayskip}{2pt}
\setlength{\abovedisplayshortskip}{1pt}
\setlength{\belowdisplayshortskip}{2pt}
\begin{equation*}
\label{eq:subspace-loss-pecl}
\begin{aligned}
&\mathbb E_{W_t\sim Q_t}\bigl[\hat L_{S_t}(W_t)\bigr] \\
&= \mathbb{E}_{\varepsilon\sim\mathcal N(0,\Sigma_t)}
    \big[\hat{L}_{S_t}(W^{\mathrm{froz}}_t+\widehat{\Delta W}_t+\varepsilon)\big] \\
& =\hat L_{S_t}(W^{\mathrm{froz}}_t + \widehat{\Delta W}_t)
+
\mathrm{TS}_t^{\Delta}(\widehat{\Delta W}_t,\Sigma_t).
\end{aligned}
\end{equation*}
}where the task adapter subspace sharpness is $
\mathrm{TS}_t^{\Delta}(\widehat{\Delta W}_t,\Sigma_t)
= \mathbb{E}_{\varepsilon\sim\mathcal N(0,\Sigma_t)}
   \big[\hat{L}_{S_t}(W^{\mathrm{froz}}_t+\widehat{\Delta W}_t+\varepsilon)
   - \hat{L}_{S_t}(W^{\mathrm{froz}}_t+\widehat{\Delta W}_t)\big] \text{\,and\, } 
     \varepsilon\sim\mathcal N(0,\Sigma_t) \in  \mathcal W_{\Delta,t}$.
     % \begin{equation*}
     % \begin{aligned}
     % &\mathrm{RS}_t^{\Delta}(\hat{\Delta W}_t,\Sigma_t):= \\
     % &\qquad \mathbb{E}_{\varepsilon\sim\mathcal N(0,\Sigma_t)}
     %    \big[\hat{L}_{S_t}(W_\mathrm {frozen }+\hat{\Delta W}_t+\varepsilon) \\
     % & \qquad- \hat{L}_{S_t}(W_\mathrm {frozen }+\hat{\Delta W}_t)\big] \\
     % &\qquad \varepsilon\sim\mathcal N(0,\Sigma_t) \in \mathcal W_\Delta
     % \end{aligned}
     % \end{equation*}
     \end{lemma}

We refer to Appendix~\ref{proof:lemma2} for the proof.



\section{PAC-Bayesian Bound for PECL}

In PECL, Lemma~\ref{lem:kl-decomp-pecl} and Lemma~\ref{lem:subspace-loss-pecl} reformulates the complexity and empirical terms entirely in the task adpter subspace $\mathcal W_{\Delta,t}$ through the task KL divergence $\mathrm{KL}\bigl(Q_t^\Delta \,\Vert\, P_t^\Delta\bigr)$ and the adapter subspace sharpness $\mathrm{TS}_t^{\Delta}(\widehat{\Delta W}_t,\Sigma_t)$ respectively. 
% See Fig.~\ref{fig: showPAC3} and Fig.~\ref{fig: KL_decompose} for an overview of the subspace reduction and the hierarchical construction.
Substituting these reductions into the bound from Theorem~\ref{thm:pathwise_pac}, we obtain the PAC-Bayesian bound for PECL as follows:
\begin{theorem}
\label{thm:pecl-bound}
Consider a PECL model with frozen backbone $W^{\mathrm{froz}}_t$ and Gaussian posteriors 
$Q_t^\Delta = \mathcal N(\widehat{\Delta W}_t,\Sigma_t)$ supported on $\mathcal W_{\Delta,t}$. 
Let $P_t^\Delta$ be task priors sampled from a hyper posterior $\mathcal P_{t-1}$ over adapter 
distributions. For any $\delta\in(0,1]$, with probability at least $1-\delta$ over $S_{1:T}$,
{%
\setlength{\abovedisplayskip}{2pt}
\setlength{\belowdisplayskip}{3pt}
\setlength{\abovedisplayshortskip}{3pt}
\setlength{\belowdisplayshortskip}{3pt}
\begin{equation*}
\label{eq:pecl-main-bound}
\resizebox{\linewidth}{!}{$
\begin{aligned}
&\frac{1}{T}\sum_{t=1}^T
\mathbb E_{P_t^\Delta\sim\mathcal P_{t-1}}
\mathbb E_{W_t\sim Q_t}
L_{\mathcal D_t}(W_t) \\
&\le
\frac{1}{T}\sum_{t=1}^T
\mathbb E_{P_t^\Delta\sim\mathcal P_{t-1}}
\bigl[
\hat L_{S_t}(W^{\mathrm{froz}}_t + \widehat{\Delta W}_t)
+
\mathrm{TS}_t^{\Delta}(\widehat{\Delta W}_t,\Sigma_t)
\bigr] \\
&+
\sqrt{
\frac{
\displaystyle
\sum_{t=1}^T
\mathbb E_{P_t^\Delta\sim\mathcal P_{t-1}}
\mathrm{KL}\bigl(Q_t^\Delta \,\Vert\, P_t^\Delta\bigr)
+
\sum_{t=1}^T
\mathrm{KL}\bigl(\mathcal P_t\Vert\mathcal P_{t-1}\bigr)
+
\log\frac{1}{\delta}}
{2\displaystyle\sum_{t=1}^T (m_t - 1)}
}.
\end{aligned}
$}
\end{equation*}
}
\end{theorem}
Theorem~\ref{thm:pecl-bound} specializes the Theorem~\ref{thm:pathwise_pac} to PECL by exploiting the fact that all task dependent updates are confined to the permissible task subspace $\mathcal W_{\Delta,t}$. 
% This restriction is not merely technical.
Under a frozen backbone $\mathcal W_{\Delta,t}$ is exactly the set of directions along which the model can be modified by future tasks.
Therefore, in PECL the relevant notion of flatness for generalization is the effective geometry induced on $\mathcal W_{\Delta,t}$, rather than a global property of the full parameter space. We provide the proof in Appendix~\ref{proof:theorem}.

% Since subsequent adaptation is restricted to $\mathcal W_{\Delta,t}$, this term quantifies the sensitivity of the task $t$ solution to admissible variations, and thus how robust the learned solution is with respect to the only perturbations that future tasks can effectively realize.
% The complexity penalty is expressed through the adapter level divergence $\mathrm{KL}(Q_t^\Delta\Vert P_t^\Delta)$ together with the hyperposterior drift, yielding a structurally aligned form of complexity control compared with full finetuning.

\subsection{Perturbation Scope and Type}

Having established the PAC-Bayesian bound for PECL, we now explore where and how to apply perturbations to obtain a flatter minima and improve the performance in PECL.

% \vspace{-1pt}
%the backbone $W^{\mathrm{froz}}_t$ is frozen and 
\minisection{Scope - Why permissible perturbations should lie in $\mathcal W_{\Delta,t}$?}
In PECL, learning is confined to the LoRA update subspace $\mathcal W_{\Delta,t}$. Consequently, only directions in $\mathcal W_{\Delta,t}$ are reachable by future task updates and can affect learning along the training trajectory.
Theorem \ref{thm:pecl-bound} makes this restriction explicit: both the task-wise KL term and the sharpness-dependent empirical term depend only on distributions supported on $\mathcal W_{\Delta,t}$.
Therefore, perturbing the parameter outside $\mathcal W_{\Delta,t}$ imposes impact along unreachable directions, which cannot be compensated by permissible LoRA updates under PECL constraint and can introduce unnecessary loss elevation and unstable optimization behavior.

\begin{figure}[tbp]
    \centering
\includegraphics[width=\linewidth]{figures_TRML/lossland_curv1d_combine_loss_title2.pdf}
    % \vspace{-6pt}
    \caption{
    % {How (type)}: 
    We visualize the loss landscape of the models by a 1D probe along the leading eigen-direction of the full curvature.
\textbf{Top}: The probe is restricted to the trainable LoRA parameters, using the leading eigen-direction of the restricted curvature.
\textbf{Bottom}: The probe is applied in the full parameter space, using the leading eigen-direction of the full curvature.
%     We visualize the loss of the same trained solutions by a 1D probe along the leading eigen-direction under two evaluation scopes.
% Top: The probe is restricted to the trainable LoRA parameters and follows the leading eigen-direction of the restricted curvature $C_{\Delta,t}$.
% Bottom: The probe is applied in the full parameter space and follows the leading eigen-direction of the full curvature $C_t$.
% Models trained with different flatness objectives show similar profiles in the LoRA-restricted probe but clearly different sensitivity under full-parameter perturbations.
     % Models trained with different flatness objectives yet different sensitivity under full-parameter perturbations. 
    % Different perturbation types shape the geometry of updates within the subspace, leading to different full-parameter loss landscapes. 
    }
    \label{fig:scope_type2}
    \vspace{-10pt}
\end{figure}

\minisection{Type - How different perturbations shape flatness?}
Let $\Pi_{\Delta,t}$ be the projector onto $\mathcal W_{\Delta,t}$ and let $C_t=C(\hat W_t;S_t)$ be a task-dependent curvature operator of $\hat L_{S_t}$ at $\widehat W_t$.
In PECL, task-relevant parameter updates and task-permissible sharpness are restricted to $\mathcal W_{\Delta,t}$, hence only the restricted curvature $C_{\Delta,t}=\Pi_{\Delta,t}C_t\Pi_{\Delta,t}$ and perturbations supported on $\mathcal W_{\Delta,t}$ can contribute.
Under a local second-order approximation and $\varepsilon\sim\mathcal N(0,\Sigma_t)$ supported on $\mathcal W_{\Delta,t}$,
{%
% \setlength{\abovedisplayskip}{4pt}
% \setlength{\belowdisplayskip}{3pt}
% \setlength{\abovedisplayshortskip}{4pt}
% \setlength{\belowdisplayshortskip}{3pt}
\[
\mathrm{RS}_t^{\Delta}(\widehat{\Delta W}_t,\Sigma_t)
% =
% \mathbb E_{\varepsilon}\big[\hat L_{S_t}(\hat W_t+\varepsilon)-\hat L_{S_t}(\hat W_t)\big]
\approx
\frac12\operatorname{tr}\!\big(C_{\Delta,t}\Sigma_t\big),
\]
}where $\mathcal W_{\Delta,t}$ fixes the permissible directions, $C_{\Delta,t}$ specifies the restricted curvature on them, $\Sigma_t$ allocates perturbation variance within them, and $\operatorname{tr}(C_{\Delta,t}\Sigma_t)$ measures curvature–noise alignment inside the task permissible geometry.

{%
% 1) table* 与上下正文的距离（页顶/页底浮动体）
\setlength{\textfloatsep}{2pt}      % 可在 4pt--10pt 之间调
\setlength{\floatsep}{2pt}          % 多个浮动体之间（备用）
\setlength{\dbltextfloatsep}{2pt}   % 某些模板对 table* 用这个（双栏更关键）
\begin{table*}[ht]
\centering
\caption{
Class incremental performance and weight space flatness metrics on ImageNet-R ($T=20$, $R=16$).
We report standard class-incremental metrics together with sharpness and empirical curvature diagnostics (Hessian and Fisher) {evaluated on the trainable LoRA parameters only}, namely restricted to the adapter-induced subspace rather than the full model parameter space.
% The table reports standard class incremental metrics  together with the sharpness and empirical Hessian and Fisher statistics restricted in the weight space.
}
\label{table:q1_ci_with_weight_flatness_reorg}
\setlength{\tabcolsep}{2.5pt}       % 列间距更小
\renewcommand{\arraystretch}{0.92}  % 行高更紧凑（可选）
\resizebox{\linewidth}{!}{
\scriptsize % 整体字号更小（可改为 \footnotesize 或 \tiny）
\begin{tabular}{@{}p{1.15cm} p{1.8cm} c c c c c c c c@{}}
\toprule
\multirow{2}{*}{\textbf{Method}} & \multirow{2}{*}{\textbf{Opt.}} &
\multicolumn{2}{c}{\textbf{CL Performance ($\uparrow$)}} &
\multicolumn{2}{c}{\textbf{Sharpness ($\downarrow$)}} &
\multicolumn{4}{c}{\textbf{Curvature ($\downarrow$)}} 
% \multicolumn{2}{c}{\textbf{Emp Fisher ($\downarrow$)}}
\\
\cmidrule(lr){3-4}\cmidrule(lr){5-6}\cmidrule(lr){7-10}
 & & $\operatorname{Acc}_{20}^{\mathrm{cls}}$  %&$\operatorname{ACC}^{\mathrm{cls}}$ 
&$\operatorname{Acc}_{20}^{\mathrm{nme}}$ 
 & AS$^{(0)}(\rho)$  & AS$^{(1)}(\rho)$
 & $\lambda_{\max}(H)$ & $\mathrm{tr}(H)$  
 & $\lambda_{\max}(F)$ & $\mathrm{tr}(F)$   \\
\midrule
% \multirow{1}{*}{\textbf{PerTaskLP}}
% & SGD   & $55.22^{\pm 0.31}$ & $59.36^{\pm 0.19}$ & $-8.35^{\pm 0.75}$ & 0.072   & 0.072 &     &   &    &   \\
% \midrule
% \multirow{2}{*}{\textbf{SeqFT}}
% & SGD   &$50.72^{\pm 2.20}$ & $59.89^{\pm 1.32}$ & $-38.94^{\pm 1.87}$ &   &  &   &  &   & \\
% & SAM   &$56.07^{\pm 2.60}$ & $66.47^{\pm 1.53}$ & $-33.04^{\pm 4.71}$&   &  &   &  &   &  \\
%  &GAM  &$59.78^{\pm 0.74}$ & $67.31^{\pm 1.18}$ & $-32.70^{\pm 0.26}$ &   &  &   &  &   &   \\
% \midrule
\multirow{4}{*}{{SeqLoRA}}
& SGD    &$63.14^{\pm 2.17}$ & $72.82^{\pm 1.27}$  & 0.142  & 0.114 & 47.37 & 1502.55 & 35.51 & 1401.15 \\\
& + $\mathrm{RS}^{0}_S(W,\rho)$  &$64.25^{\pm 0.87}$ & $73.01^{\pm 1.11}$ &0.121 &0.108 &47.29 &1444.96 &30.35 &1338.16\\
& + $\mathrm{AS}^{0}_S(W,\rho)$    &$67.04^{\pm 0.71}$ & $74.53^{\pm 0.60}$ & 0.120  & 0.093 & 24.41   & 996.59   & 21.39  & 821.45   \\
 & + $\mathrm{AS}^{(1)}_S(W,\rho)$   &$73.31^{\pm 0.52}$ & $75.74^{\pm 0.15}$  & 0.060  &0.058  & 4.81    & 339.49   & 3.72   & 240.91   \\
\midrule
\multirow{4}{*}{{IncLoRA}}
& SGD    &$65.77^{\pm 2.84}$  &$75.41^{\pm 0.26}$ & 0.076    & 0.062 & 15.46 & 969.34  & 8.89  & 836.91 \\
& + $\mathrm{RS}^{0}_S(W,\rho)$ & $67.42^{\pm 1.23}$ & $75.79^{\pm 0.42}$  &0.075 &0.073 &7.54 &917.35 &8.82 &774.56\\
& + $\mathrm{AS}^{(0)}_S(W,\rho)$    &$68.86^{\pm 0.92}$ & $75.95^{\pm 0.07}$  & 0.053   & 0.052 & 3.72    & 712.88   & 4.35   & 526.49   \\
 &+ $\mathrm{AS}^{(1)}_S(W,\rho)$  &$72.86^{\pm 0.23}$ & $76.72^{\pm 0.54}$  & 0.035   & 0.035 & 2.69    & 281.22   & 2.09   & 192.25   \\
\midrule
\multirow{4}{*}{{OLoRA}}
& SGD    &$66.64^{\pm 0.44}$ & $75.10^{\pm 0.41}$  & 0.078   & 0.058 & 6.94    & 938.93   & 9.00   & 831.23   \\
& + $\mathrm{RS}^{0}_S(W,\rho)$  & $67.53^{\pm 1.04}$ & $75.71^{\pm 0.27}$  &0.076 &0.073 &7.18 &918.33 &8.97 &813.62\\
& + $\mathrm{AS}^{(0)}_S(W,\rho)$     &$68.53^{\pm 0.78}$ & $75.73^{\pm 0.32}$  & 0.066   & 0.048 & 4.40    & 775.40   & 7.41   & 628.12   \\
 & + $\mathrm{AS}^{(1)}_S(W,\rho)$  &$72.44^{\pm 0.36}$ & $76.56^{\pm 0.53}$  & 0.036   & 0.036 & 2.22    & 261.47   & 2.08   & 198.30\\
% \midrule
% \multirow{1}{*}{\textbf{l2p}}
% & Adam   & $71.35^{\pm 0.05}$ & $74.84^{\pm 0.73}$ & $-6.34^{\pm 0.50}$  &-  &- &- &- &- &- \\
% \midrule
% \multirow{1}{*}{\textbf{dualprompt}}
% & Adam   & $71.59^{\pm 0.68}$ & $74.80^{\pm 1.18}$ & $-3.38^{\pm 0.45}$ &-  &- &- &- &- &- \\
\bottomrule
\end{tabular}
}
\end{table*}
} 

This clarifies how perturbation \emph{type} acts in PECL: since adapter update direction is constrained to $\mathcal W_{\Delta,t}$, perturbation types induce distinct curvature-aligned covariances $\Sigma_{\Delta,t}$, leading to different training trajectories and sharpness of the solutions. 
% different { perturbation types} allocate variance $\Sigma_{\Delta,t }$ in different ways relative to the restricted curvature $C_{\Delta,t}$, leading to different trajectories and  different flatness solutions.
% , by allocating variance differently within $\mathcal W_{\Delta,t}$.
% This clarifies how different optimizers act in PECL. Standard sharpness criteria can be re-expressed in the task-permissible geometry by restricting perturbations to $\mathcal W_{\Delta,t}$.
% Since parameter-efficient continual learning restricts adaptation to a task-admissible set of update directions, it is natural to focus on perturbations supported on the adapter-induced subspace $\mathcal W_{\Delta,t}$ when reasoning about training dynamics and flatness.
% Under this fixed scope, different optimizers mainly differ by the \emph{geometry of the perturbation distribution} within $\mathcal W_{\Delta,t}$, namely by how $\Sigma_t$ concentrates variance inside the same subspace.
% This clarifies how perturbation \emph{type} acts in PECL: since parameter-efficient continual learning restricts adaptation to a task-admissible set of update directions, only perturbations supported on the same adapter-induced subspace $\mathcal W_{\Delta,t}$ can materially alter the optimization trajectory and the resulting flatness; under this fixed scope, different optimizers mainly differ by the \emph{geometry of the perturbation distribution} within $\mathcal W_{\Delta,t}$, namely by how $\Sigma_t$ concentrates variance (or, for adversarial criteria, how the neighborhood is searched) inside the same subspace.
SGD induces an implicit stochastic perturbation on the trainable parameters, which can be viewed as an effective noise covariance on $\mathcal W_{\Delta,t}$\cite{pmlr-v97-zhu19e}. Random zeroth-order sharpness $\mathrm{RS}^{(0),\Delta}(W,\Sigma)$ corresponds to an explicit Gaussian perturbation $\varepsilon\sim\mathcal N(0,\Sigma)$ supported on $\mathcal W_{\Delta,t}$, and measures the average-case loss elevation. In contrast, adversarial criteria concentrate the perturbation on the same permissible subspace: $\mathrm{AS}^{(0),\Delta}$ optimize the worst-case loss increase over the $\rho$-ball in $\mathcal W_{\Delta,t}$ ~\cite{foretSharpnessAwareMinimizationEfficiently2021}, while $\mathrm{AS}^{(1),\Delta}$ further considers the neighborhood gradient norm and is more sensitive to dominant modes of the restricted curvature $C_{\Delta,t}$~\cite{Zhang_GAM}. 
SGD induces an implicit stochastic perturbation on the trainable parameters, which can be viewed as an effective noise covariance on $\mathcal W_{\Delta,t}$\cite{pmlr-v97-zhu19e}. Random zeroth-order sharpness $\mathrm{RS}^{(0),\Delta}(W,\Sigma)$ corresponds to an explicit Gaussian perturbation $\varepsilon\sim\mathcal N(0,\Sigma)$ supported on $\mathcal W_{\Delta,t}$, and measures the average-case loss elevation. In contrast, adversarial criteria concentrate the perturbation on the same permissible subspace: $\mathrm{AS}^{(0),\Delta}$ optimize the worst-case loss increase over the $\rho$-ball in $\mathcal W_{\Delta,t}$ ~\cite{foretSharpnessAwareMinimizationEfficiently2021}, while $\mathrm{AS}^{(1),\Delta}$ further considers the neighborhood gradient norm and is more sensitive to dominant modes of the restricted curvature $C_{\Delta,t}$~\cite{Zhang_GAM}. 


% Under a fixed scope, these choices correspond to increasingly concentrated perturbation geometries within $\mathcal W_{\Delta,t}$~\cite{mollenhoffSAMOptimalRelaxation2022,Zhang_GAM},
% $A$
% and therefore yield progressively stronger control of the sharpness-dependent empirical term in Theorem~\ref{thm:pecl-bound}.
% This mechanism is reflected in Fig.~\ref{fig:scope_type2}. under a full-parameter 1D probe along the leading curvature direction, the trained solutions exhibit clearly separated loss profiles in different order and the order is aligned with perturabtion  definiton. 


Under a fixed scope, these choices correspond to increasingly concentrated perturbation geometries within $\mathcal W_{\Delta,t}$, which induces an ordering of strength under a comparable radius (or covariance scale)~\cite{mollenhoffSAMOptimalRelaxation2022,Zhang_GAM}:
$ \mathrm{AS}^{(1)}_S(W,\rho)\ \ge \mathrm{AS}^{(0)}_S(W,\rho)\ \ge \mathrm{RS}^{(0)}_S(W,\Sigma_t).$
% {%
% \setlength{\abovedisplayskip}{1pt}
% \setlength{\belowdisplayskip}{1pt}
% \setlength{\abovedisplayshortskip}{1pt}
% \setlength{\belowdisplayshortskip}{1pt}
% \begin{equation*}
%     \mathrm{AS}^{(1)}_S(W,\rho)\ \ge \mathrm{AS}^{(0)}_S(W,\rho)\ \ge \mathrm{RS}^{(0)}_S(W,\Sigma_t).
% \end{equation*}
% }
Accordingly, more adversarial objectives impose a stronger penalty on sharp neighborhoods and thus exert progressively stronger control of the sharpness-dependent empirical term in Theorem~\ref{thm:pecl-bound}.
This mechanism is reflected in Fig.~\ref{fig:scope_type2}: under the full-parameter probe along the leading eigen-direction of the full curvature, the trained solutions exhibit clearly separated loss profiles, indicating that perturbation type, although imposed within $\mathcal W_{\Delta,t}$ during training, can translate into different global sensitivity of $W^{\mathrm{froz}}_t+\Delta W_t$.

% and therefore yield progressively stronger control of the sharpness-dependent empirical term in Theorem~\ref{thm:pecl-bound}. This mechanism is reflected in Fig.~\ref{fig:scope_type2}. Under the full-parameter probe along the leading eigen-direction of the full curvature, the trained solutions exhibit clearly separated loss profiles, indicating that different types of perturbation, although imposed inside $\mathcal W_{\Delta,t}$ during training, can translate into different global sensitivity of $W_0+\Delta W_t$.


% This mechanism is  reflected in Fig.~\ref{fig:scope_type2}: under the full-parameter probe, the  solutions exhibit clearly separated profiles.  These observations indicate that different subspace perturbation geometries, while acting inside $\mathcal W_{\Delta,t}$ during training, translate into different global sensitivity of $W_0+\Delta W_t$.


% These definitions imply an ordering in strength within the same subspace and radius constraint~\cite{mollenhoffSAMOptimalRelaxation2022,Zhang_GAM}:
% $\mathrm{AS}^{(1)}_S(W,\rho)\ \ge \mathrm{AS}^{(0)}_S(W,\rho)\ \ge \mathrm{RS}^{(0)}_S(W,\Sigma_t).$
% % {%
% % % \setlength{\abovedisplayskip}{3pt}
% % % \setlength{\belowdisplayskip}{3pt}
% % % \setlength{\abovedisplayshortskip}{1pt}
% % % \setlength{\belowdisplayshortskip}{1pt}
% % \begin{equation*}
% %     \mathrm{AS}^{(1)}_S(W,\rho)\ \ge \mathrm{AS}^{(0)}_S(W,\rho)\ \ge \mathrm{RS}^{(0)}_S(W,\Sigma_t),
% % \end{equation*}
% % }
% Accordingly, optimizing increasingly adversarial perturbation objectives typically yields solutions with smaller sharpness within task subspace. And as the perturbation becomes increasingly concentrated and aligned with the dominant high-curvature direction, it more effectively promotes escape from sharp minima and biases the trajectory toward flatter regions~\cite{pmlr-v97-zhu19e,chen2026sharpnessflatnessdecompositionframework}.



% These definitions imply an ordering in strength within the same subspace and radius constraint~\cite{mollenhoffSAMOptimalRelaxation2022,Zhang_GAM}:
% {%
% \setlength{\abovedisplayskip}{3pt}
% \setlength{\belowdisplayskip}{3pt}
% \setlength{\abovedisplayshortskip}{1pt}
% \setlength{\belowdisplayshortskip}{1pt}
% \begin{equation*}
%     \mathrm{AS}^{(1)}_S(W,\rho)\ \ge \mathrm{AS}^{(0)}_S(W,\rho)\ \ge \mathrm{RS}^{(0)}_S(W,\Sigma_t),
% \end{equation*}
% }
% % $$
% Accordingly, optimizing increasingly adversarial perturbation objectives typically yields solutions with smaller sharpness within task subspace. And as the perturbation becomes increasingly concentrated and aligned with the dominant high-curvature direction, it more effectively promotes escape from sharp minima and biases the trajectory toward flatter regions~\cite{pmlr-v97-zhu19e,chen2026sharpnessflatnessdecompositionframework}.
% This mechanism is directly reflected in Fig.~\ref{fig:scope_type2}: under the full-parameter probe, the same trained solutions exhibit clearly separated profiles. A one-dimensional probe within the LoRA subspace can produce seemingly similar loss curves (top of Fig.~\ref{fig:scope_type2}), however, when the resulting models are evaluated under full-parameter perturbations (bottom of Fig.~\ref{fig:scope_type2}), the profiles separate markedly, revealing that different subspace perturbation geometries lead to different full-model sensitivity. These observations indicate that different subspace perturbation geometries, while acting inside $\mathcal W_{\Delta,t}$ during training, translate into different global sensitivity of $W_0+\Delta W_t$.

% A one-dimensional probe within the LoRA subspace can produce seemingly similar loss curves (top of Fig.~\ref{fig:scope_type2}), however, when the resulting models are evaluated under full-parameter perturbations (bottom of Fig.~\ref{fig:scope_type2}), the profiles separate markedly, revealing that different subspace perturbation geometries lead to different full-model sensitivity.


% the same trained solutions separate markedly when evaluated under \emph{full-parameter} perturbations, indicating that different subspace perturbation geometries translate into different global sensitivity of $W_0+\Delta W_t$




% This clarifies how different optimizers act in PECL. More generally, standard sharpness criteria can be re-expressed in the task-admissible geometry by restricting perturbations to $\mathcal W_{\Delta,t}$. Standard SGD injects an implicit stochastic perturbation with an effective covariance $\Sigma_t^{\mathrm{SGD}}$ determined by minibatch noise\cite{pmlr-v97-zhu19e}. Random zeroth-order sharpness $\mathrm{RS}^{(0)}$ corresponds to explicit isotropic Gaussian noise $\Sigma=\sigma^2 I$, and measures the expected loss elevation under average-case perturbations supported on the admissible subspace. In contrast, adversarial criteria replace the average-case elevation with a worst-case concentration over the same adapter subspace. The adversarial zeroth-order sharpness $\mathrm{AS}^{(0),\Delta}$ chooses the worst-case perturbation that maximizes the local loss increase over $\Sigma=\rho\frac{\nabla L}{|\nabla L|_2}$
% % , yielding a gradient-aligned direction\cite{foretSharpnessAwareMinimizationEfficiently2021}.
% The adversarial first-order sharpness $\mathrm{AS}^{(1),\Delta}$ maximizes the neighborhood gradient norm
% $\Sigma=\rho\frac{\nabla|\nabla L|_2}{|\nabla|\nabla L|_2|_2}$,
% which emphasizes directions directions aligned with dominant eigenmodes of $C_{\Delta,t}$~\cite{zhangGradientNormAware2023}. These definitions imply an ordering in strength within the same subspace and radius constraint~\cite{mollenhoffSAMOptimalRelaxation2022,Zhang_GAM}:
% $\mathrm{AS}^{(1)}_S(W,\rho)\ \ge \mathrm{AS}^{(0)}_S(W,\rho)\ \ge \mathrm{RS}^{(0)}_S(W,\Sigma_t).$
% % {%
% % % \setlength{\abovedisplayskip}{3pt}
% % % \setlength{\belowdisplayskip}{3pt}
% % % \setlength{\abovedisplayshortskip}{1pt}
% % % \setlength{\belowdisplayshortskip}{1pt}
% % \begin{equation*}
% %     \mathrm{AS}^{(1)}_S(W,\rho)\ \ge \mathrm{AS}^{(0)}_S(W,\rho)\ \ge \mathrm{RS}^{(0)}_S(W,\Sigma_t),
% % \end{equation*}
% % }
% Accordingly, optimizing increasingly adversarial perturbation objectives typically yields solutions with smaller sharpness within task subspace. And as the perturbation becomes increasingly concentrated and aligned with the dominant high-curvature direction, it more effectively promotes escape from sharp minima and biases the trajectory toward flatter regions~\cite{pmlr-v97-zhu19e,chen2026sharpnessflatnessdecompositionframework}.



\noindent\textbf{LoRA Organization as a Structural Prior.}\quad 
To make the hyper process concrete in PECL, we represent each task-level prior by a minimal \emph{hyper-state}
$
\phi_t = \bigl(\mathcal W_{\Delta,t},\mu_{P,t},\Sigma_{P,t}\bigr),
$
where $\mathcal W_{\Delta,t}$ specifies the permissible adapter update scope and $(\mu_{P,t},\Sigma_{P,t})$ parameterize a task prior on these coordinates.
Sampling $P_t\sim\mathcal P_{t-1}$ can then be interpreted as drawing $\phi_t\sim\mathcal P_{t-1}$ and instantiating the corresponding task prior. Under this view, SeqLoRA, IncLoRA, and constraint-augmented designs correspond to different evolution rules for $\phi_t$, and thus shape both the restricted sharpness channel and the KL complexity in Theorem~\ref{thm:pathwise_pac}.




SeqLoRA keeps an approximately fixed scope $\mathcal W_{\Delta,t}\equiv\mathcal W_\Delta$ and concentrates successive updates in the same geometry, which tends to induce stronger direction competition across tasks.
% under the deterministic choice $P_t^\Delta=Q_{t-1}^\Delta$, the complexity reduces to $\sum_t\mathrm{KL}(Q_t^\Delta\Vert Q_{t-1}^\Delta)$.
% This shared scope limit admissible directions for new task and amplify cross-task interference for previous task.
% A fixed scope concentrates successive updates in the same geometry, which tends to induce stronger direction competition across tasks.
IncLoRA relaxes the fixed-scope assumption by expanding $\mathcal W_{\Delta,t}$ over time, which redistribute adaptation across additional degrees of freedom and reduced subspace overwrite, while the increased flexibility may also admit more directions with projected curvature.
% . At the same time, a larger freedom may admit more directions with projected curvature.
% This provides extra admissible directions and reduce the influence of previous task direction. 
Constraint-augmented designs further control how new degrees of freedom are allocated and initialized by imposing structure on $\mathcal W_{\Delta,t}$, for example via orthogonality regularization~\cite{wang2023orthogonal}, gradient projection~\citep{liangInfLoRAInterferenceFreeLowRank2024a,yang2024parametercollisionhinderingcontinual}.
It reduce repeated overwrite and stabilize the evolution of the permissible geometry.











% \section{Flatness in the Adapter Subspace Matters for Continual Learning: A PAC-Bayesian Perspective}
% \label{section PAC}


% % \subsection{Motivating Observation: LoRA-Only Flatness Improves CL Performance}



% \subsection{PAC–Bayesian setup for CL}

% % \begin{figure}[htbp]
% % \centering
% % \begin{minipage}{0.9\linewidth}
% % \centering
% % \includegraphics[width=0.9\linewidth]{figures_TRML/HyperPAC_2.pdf}
% % \caption{Hierarchical PAC-Bayesian CL process }
% % \label{fig:Hierarchical PAC-Bayesian CL process}
% % \end{minipage}
% % \end{figure}


% % To analyze the role of sharpness in CL, we adopt a hierarchical PAC-Bayesian framework.
% % For each task \(t \in \{1, \dots, T\}\), let \(S_t = \{z_{tj}\}_{j=1}^{m_t}\) be i.i.d. draws from \(\mathcal{D}_t\), and let \(\ell(w, z) \in [0,1]\) be the loss function. Define the empirical risk and expected risk for a model \(W\) as:
% % \[
% % \hat{L}_{S_t}(W) = \frac{1}{m_t} \sum_{j=1}^{m_t} \ell(W, z_{tj}), \quad
% % L_{\mathcal{D}_t}(W) = \mathbb{E}_{z \sim \mathcal{D}_t} \ell(W, z).
% % \]
% To analyze the role of sharpness in CL, we treat the prior $P$ itself as a random variable of $M(\mathcal H)$. A \emph{hyper-posterior} is a probability measure $\mathcal P \in \mathcal M(\mathcal M(\mathcal H))$ over such priors. We write $\mathcal F_t = \sigma(S_1,\dots,S_t)$ for the $\sigma$-algebra generated by the first $t$ datasets and consider the filtration $\mathcal F_0 \subseteq \mathcal F_1 \subseteq \cdots \subseteq \mathcal F_T$. At the beginning of task $t$, before observing $S_t$, the learner holds a task-level hyper-posterior $\mathcal P_{t-1}$ that is $\mathcal F_{t-1}$-measurable, which encodes all information extracted from previous tasks. The learning process at task $t$ is then described in two stages. First, a task-level prior $P_t$ is sampled from the current hyper-posterior,
% $
% P_t \sim \mathcal P_{t-1}.
% $
% Second, after observing $S_t$, the learner updates from $P_t$ to a task-level posterior $Q_t(\cdot \mid P_t,S_t)$ on $\mathcal H$, and the hyper-posterior is updated to $\mathcal P_t$. We assume that these updates are absolutely continuous, so that the associated KL divergences are well defined.
% % $
% % \mathcal P_t \ll \mathcal P_{t-1},
% % \quad
% % Q_t(\cdot \mid P,S_t) \ll P(\cdot)
% % \quad \text{for } \mathcal P_{t-1}\text{-almost every } P.
% % $
% The corresponding task-level risks under the hierarchical posterior are obtained by averaging over the sampling of priors and hypotheses,
% $$
% \mathcal L_t
% =
% \mathbb E_{P_t \sim \mathcal P_t}
% \,
% \mathbb E_{W_t \sim Q_t(\cdot \mid P_t,S_t)}
% \big[
% L_{\mathcal D_t}(W_t)
% \big],
% \qquad
% \hat{\mathcal L}_t
% =
% \mathbb E_{P_t \sim \mathcal P_t}
% \,
% \mathbb E_{W_t \sim Q_t(\cdot \mid P_t,S_t)}
% \big[
% \hat L_{S_t}(W_t)
% \big].
% $$








% % Before observing \(S_t\), a task-level prior \(P_t\) is drawn from the hyper-posterior \(\mathcal{P}_{t-1}\), i.e., \(P_t \sim \mathcal{P}_{t-1}\), where \(\mathcal{P}_{t-1}\) is \(\mathcal{F}_{t-1}\)-measurable.   \textit{hyper-posterior} $\mathcal{P}\sim \mathcal{M}(\mathcal{M}(\mathcal{H}))$ i.e., a distribution over priors $P$, into \textit{hyper-posterior} $\mathcal{Q}\sim \mathcal{M}(\mathcal{M}(\mathcal{H}))$. Here \(\mathcal{F}_{t-1} = \sigma(S_1, \dots, S_{t-1})\) is the \(\sigma\)-algebra encapsulating all historical information before task \(t\) and $\mathcal{F}_0\subseteq\mathcal{F}_1\subseteq\cdots\subseteq\mathcal{F}_{t-1}$. After training on \(S_t\), the learner obtains the posterior \(Q_t(\cdot \mid P_t, S_t)\), and the hyper-posterior is updated to \(\mathcal{P}_t\), with absolute continuity:
% % $
% % \mathcal{P}_t \ll \mathcal{P}_{t-1}, \quad
% % Q_t(\cdot \mid P, S_t) \ll P(\cdot) \quad \text{for } \mathcal{P}_{t-1}\text{-a.e. } P.
% % $
% % The expected risk and empirical risk on task \(t\) are then defined as:
% % \[
% % \mathcal{L}_t = \mathbb{E}_{P_t \sim \mathcal{P}_{t}} \mathbb{E}_{W_t \sim Q_t(\cdot \mid P, S_t)} L_{\mathcal{D}_t}(W_t), \quad
% % \hat{\mathcal{L}}_t = \mathbb{E}_{P_t \sim \mathcal{P}_{t}} \mathbb{E}_{W_t \sim Q_t(\cdot \mid P, S_t)} \hat{L}_{S_t}(W_t).
% % \]


% In the PAC-Bayesian framework, rather than analyzing a single deterministic parameter vector, we consider a posterior distribution over parameters.  A common and tractable choice is to model the posterior at task \(t\) as a Gaussian
%  $ Q_t = \mathcal{N}(\hat{W}_t, \Sigma_t),$
%  where \(\hat{W}_t\) is the trained parameter vector after task \(t\), and \(\Sigma_t\) is  a covariance matrix capturing local uncertainty around \(\hat{W}_t\). Sampling  from \(W \sim Q_t\) is equivalent to 
%  $ W = \hat{W}_t + \varepsilon,
%  \quad
%  \varepsilon \sim \mathcal{N}(0, \Sigma_t).$
% Under this posterior view, the expected increase in empirical loss under such perturbations serves as a Bayesian measure of local sensitivity.
% From a robustness perspective, the expected sharpness $\mathrm{E\text{-}Sh}$ quantifies the model's sensitivity to small parameter perturbations within the support of the posterior, linking flatness to Bayesian theory as discussed in prior work~\citep{mllenhoff2023sam}. Specifically, SAM has been proven to be an optimal convex relaxation that upper bounds the true negative-loss objective in Bayesian inference~\citep{mllenhoff2023sam}. As a prominent technique for achieving flat minima, we adopt SAM to empirically validate our theoretical findings.


% % The following PAC-Bayesian theorem is established specifically for the continual learning scenario, with explicit consideration given to the dynamic evolution of hyper-posteriors across tasks and its connection to a flatness term. The proof is provided in the Appendix.

% \begin{theorem}[PAC-Bayes with time-varying hyper-posterior for CL]\label{thm:pathwise-pacbayes}
% For any $\delta\in(0,1]$, with probability at least $1-\delta$ over $S_{1:T}$,
% \begin{equation}\label{eq:main-bound}
% \begin{aligned} 
% & \frac{1}{T} \sum_{t=1}^T \underbrace{\mathbb{E}_{P_t \sim \mathcal{P}_t} \mathbb{E}_{W \sim Q_t(\cdot \mid P_t, S_t)} L_{\mathcal{D}_t}(W)}_{\text{test risk at step } t} 
% \le \ \frac{1}{T} \sum_{t=1}^T \underbrace{\mathbb{E}_{P_t \sim \mathcal{P}_t} \mathbb E_{\varepsilon\sim\mathcal N(0,\Sigma_t)}\hat L_{S_t}(\hat W_t+\varepsilon)}_{\text{empirical risk with flatness at step } t} \\ 
% & + \sqrt{\frac{ \sum_{t=1}^T \mathbb{E}_{P_t \sim \mathcal{P}_t} \mathrm{KL}\big(Q_t(\cdot \mid P_t, S_t) \,\|\, P_t\big) + \sum_{t=1}^T \mathrm{KL}(\mathcal{P}_t \,\|\, \mathcal{P}_{t-1}) + \log\frac{1}{\delta} }{2 \sum_{t=1}^T (m_t - 1)}}. 
% \end{aligned} 
% \end{equation}
% \end{theorem}

% % This theorem provides a PAC-Bayesian bound that is tailored to task-level continual learning. 
% % In contrast to classical single-task results, it introduces two quantities that are specific to the CL setting. 
% This bound strictly generalizes classical single-task PAC-Bayesian inequalities and differs from existing hierarchical bounds in two structural ways.
% First, the term $\sum_{t=1}^T \mathrm{KL}(\mathcal{P}_t \,\|\, \mathcal{P}_{t-1})$ captures the evolution of the hyper-posterior over priors as tasks accumulate, and therefore quantifies how much the task-level prior distribution drifts over time. Second, the empirical term on the right-hand side is evaluated under Gaussian perturbations, which embeds an explicit flatness component into the bound through the expected sharpness of the empirical loss.
% Our formulation contributes a PAC-Bayesian bound that is explicitly designed for task-level continual learning, where the complexity decomposes into task-wise posterior divergences and a hyper-posterior drift term and this structure prepares the ground for the subspace-aware PECL specialization.


% % This theorem provides a PAC-Bayesian bound tailored to continual learning, which explicitly incorporates a flatness term through Gaussian perturbations in the empirical risk. Compared with classical results, it introduces two CL-specific ingredients: a sum of KL divergences between successive hyper-posteriors, which captures the evolution of task-level priors, and an expected sharpness term that quantifies the sensitivity of empirical risk to Gaussian perturbations.

% % This theorem provides a PAC-Bayesian generalization bound specifically designed for continual learning, which explicitly incorporates a flatness term through Gaussian perturbations in the empirical risk. Unlike standard PAC-Bayes bounds, our approach accounts for the sequential nature of CL through two key components: (i)  the dynamic evolution of hyper-posteriors across tasks, captured by the sum of KL divergences between consecutive hyper-posteriors, and  (ii) the task-adaptive posterior distributions that enable knowledge transfer. Intuitively, the expected sharpness term captures the increase of empirical risk under small Gaussian perturbations in parameter space. Controlling this term in PECL therefore amounts to stabilizing the learned solution against parameter drift across tasks.

% % The bound reveals that in continual learning, flat minima (as captured by the perturbation term) can help stabilize the learning process and mitigate catastrophic forgetting by controlling the complexity of both task-specific solutions and the evolving meta-knowledge.

% % \subsection{Main Result: A Sharpness-Aware PAC-Bayes Bound for PECL}


% % \subsection{KL Decomposition and Sharpness in the PECL Parameter Space}
% \subsection{Subspace Decomposition of Complexity and Sharpness}
% % \subsubsection{KL Decomposition and Sharpness in the PECL Parameter Space}


% % the overall model complexity, measured by $\mathrm{KL}(Q^{\mathrm{full}} \| P^{\mathrm{full}})$, is solely determined by the distribution of the updated parameters in the adapter subspace $\mathcal{W}_\Delta$, formalized as $\mathrm{KL}(Q^{\Delta} \| P^{\Delta})$, while the frozen backbone contributes no divergence. 

% \begin{figure}[htbp]
% \centering
% \begin{minipage}{0.9\linewidth}
% \centering
% \includegraphics[width=0.8\linewidth]{figures_TRML/shiyitu668.pdf}
% \caption{Illustration of the KL divergence decomposition and coordinate equivalence. 
%     (1) \textbf{KL Decomposition}: In PECL, the overall model complexity, measured by $\mathrm{KL}(Q^{\mathrm{full}} \| P^{\mathrm{full}})$, is solely determined by the distribution of the updated parameters in the adapter subspace $\mathcal{W}_\Delta$, formalized as $\mathrm{KL}(Q^{\Delta} \| P^{\Delta})$, while the frozen backbone contributes no divergence. 
%     (2) \textbf{Two equivalent views}: The optimization in the low-rank factor coordinates $(A, B)$ is theoretically isomorphic to operating directly in the update subspace $\mathcal{W}_\Delta$. This equivalence preserves geometric properties such as flatness and curvature, which allows sharpness-aware optimization to be confined to the factorized LoRA parameters.}
% \label{fig: EFM_seqFT}
% \end{minipage}
% \end{figure}


% % \subsubsection{Subspace-Restricted Sharpness in PECL}

% In PECL, the model parameters at task $t$ decompose as
% $
% W_t=W_{\mathrm {frozen }}+\Delta W_t,
% $
% where $W_{\mathrm {frozen }}$ includes the pretrained backbone and all previously learned adapters (kept fixed during task $t$ ), and $\Delta W_t $ is the new low-rank LoRA update trained only on task $t$. This reflects the core PECL constraint: past knowledge is preserved via freezing, while adaptation is confined to the newly introduced LoRA block. 
% To formalize the LoRA subspace rigorously, we first define the mathematical context of the LoRA subspace. Let $\mathcal W_\Delta\subseteq\mathbb R^d$ be a linear subspace and let $W+\mathcal W_\Delta$ be its affine translate. Equip both sets with the trace Borel $\sigma$-algebra induced from $\mathbb R^d$. $
% \mathcal B(\mathcal W_\Delta)\ :=\{B\cap \mathcal W_\Delta:\ B\in \mathcal B(\mathbb R^d)\},\quad
% \mathcal B(W+\mathcal W_\Delta)\ :=\{B\cap (W+\mathcal W_\Delta):\ B\in \mathcal B(\mathbb R^d)\}.
% $
% These coincide with the Borel $\sigma$-algebras generated by the subspace topologies on $\mathcal W_\Delta$ and $W+\mathcal W_\Delta$, respectively. In particular, both are standard Borel spaces.


% To formalize the probabilistic relationship between the LoRA subspace and the full model space, we define the translation map:
% \[
% T_W:\ (\mathcal W_\Delta,\mathcal B(\mathcal W_\Delta))\ \to\ (W+\mathcal W_\Delta,\mathcal B(W+\mathcal W_\Delta)),\quad
% T_W(\Delta)=W+\Delta,
% \]
% which is a Borel isomorphism between the subspace and affine space.  The full-space prior and posterior are defined as pushforwards onto the affine subspace by
% $
% P^{\mathrm{full}}:=(T_W)_{\text{\#}} P^{\Delta},\quad
% Q^{\mathrm{full}}:=(T_W)_{\text{\#}} Q^{\Delta}.
% $
% where $P_t^{\Delta}$ and $Q_t^{\Delta}$ are probability measures on $(\mathcal{W}\Delta, \mathcal{B}(\mathcal{W}\Delta))$.
% For any measurable $F$, the pushforward $F_{\text{\#}}\mu$ is defined by $(F_{\#}\mu)(A)=\mu(F^{-1}(A))$. This pushforward construction leads to a fundamental simplification of the KL divergence between full-model distributions, as established in the following lemma.


% \begin{lemma}[KL Decomposition with Frozen Backbone]
% \label{lem:kl-decomposition}
% Let $P_t^{\Delta}, Q_t^{\Delta}$ be probability measures on $\mathcal{W}_\Delta$ with $Q_t^{\Delta} \ll P_t^{\Delta}$. Then $Q^{\mathrm{full}}\ll P^{\mathrm{full}}$ and
% \[
% \mathrm{KL}\left(Q^{\mathrm{full}}\middle|P^{\mathrm{full}}\right)
% =\mathrm{KL}\left(Q^{\Delta}\middle|P^{\Delta}\right).
% \]
% \end{lemma}
% % Equivalently, identifying ${W}\times\mathcal W_\Delta$ with $W+\mathcal W_\Delta$ via $\Phi$, then
% % $P^{\mathrm{full}}=\Phi_{\text{\#}}(\delta_W\otimes P^{\Delta})$ and
% % $Q^{\mathrm{full}}=\Phi_{\text{\#}}(\delta_W\otimes Q^{\Delta})$, and the same equality holds.

% Lemma~\ref{lem:kl-decomposition} shows that, in PECL, the KL complexity depends only on the distribution of the adapter update $\Delta W \in W_\Delta$. The frozen backbone induces no additional divergence, so any posterior covariance can be chosen within $W_\Delta$ without loss of generality.

% % Lemma \ref{lem:kl-decomposition} demonstrates that the PAC-Bayesian complexity term in PECL arises solely from adaptation in the LoRA subspace, as the frozen component contributes no uncertainty. This simplification directly informs the structure of posterior uncertainty: since only $\Delta W_t$ contributes to the divergence, the perturbation covariance $\Sigma_t$ must be confined to $\mathcal{W}_\Delta$. This foundational insight enables deriving a sharpness-aware generalization bound specific to the adapter subspace, where flatness optimization can be focused precisely where uncertainty exists.

% We consider a Gaussian posterior \( Q_t^{\Delta} = \mathcal{N}(\hat{\Delta W}_t, \Sigma_t) \) on the measurable space \( (\mathcal{W}_\Delta, \mathcal{B}(\mathcal{W}_\Delta)) \), where \( \hat{\Delta W}_t \in \mathcal{W}_\Delta \) is the mean and \( \Sigma_t \succeq 0 \) is a positive semidefinite covariance operator acting on the adapter subspace \( \mathcal{W}_\Delta \). When viewed in the ambient space \( \mathbb{R}^d \), this covariance is supported entirely within \( \mathcal{W}_\Delta \). Sampling from \( \Delta W \sim Q_t^{\Delta} \) is equivalent to 
% $
% \Delta W = \hat{\Delta W}_t + \varepsilon, \quad \varepsilon \sim \mathcal{N}(0, \Sigma_t) \ \text{on } \mathcal{W}_\Delta.
% $ This Gaussian posterior formulation enables a natural decomposition of the empirical risk into nominal performance and subspace-restricted sharpness, as formalized below.

% \begin{lemma}
% \label{lemma: subspace-full}
%     The posterior expected empirical loss decomposes as
% \begin{equation}
%     \mathbb{E}_{W\sim Q_t^{\mathrm{full}}}\big[\hat{L}_{S_t}(W)\big]
% = \mathbb{E}_{\varepsilon\sim\mathcal N(0,\Sigma_t)}
%     \big[\hat{L}_{S_t}(W_\mathrm {frozen }+\hat{\Delta W}_t+\varepsilon)\big]
% = \hat{L}_{S_t}(W_\mathrm {frozen }+\hat{\Delta W}_t)
%   + \mathrm{E\text{-}Sh}_t^{\Delta}(\hat{\Delta W}_t,\Sigma_t),
% \end{equation}
% where the (subspace) expected sharpness is
% $$
% \mathrm{E\text{-}Sh}_t^{\Delta}(\hat{\Delta W}_t,\Sigma_t)
% := \mathbb{E}_{\varepsilon\sim\mathcal N(0,\Sigma_t)}
%    \big[\hat{L}_{S_t}(W_\mathrm {frozen }+\hat{\Delta W}_t+\varepsilon)
%        - \hat{L}_{S_t}(W_\mathrm {frozen }+\hat{\Delta W}_t)\big],\quad
% \varepsilon\sim\mathcal N(0,\Sigma_t)\ \text{on } \mathcal W_\Delta$$
% \end{lemma}

% In contrast to full-space analyses, this decomposition shows that, once the backbone is frozen, only perturbations within the update subspace $\mathcal{W}_{\Delta}$ contribute to the posterior expected empirical loss. The expected sharpness term $\mathrm{E\text{-}Sh}_t^{\Delta}(\hat{\Delta W}_t,\Sigma_t)$ measures the increase in empirical loss under Gaussian perturbations supported on $\mathcal{W}_{\Delta}$, and the frozen backbone plays no role in this quantity. This is a structural property of the PECL regime: because the operative model at task $t$ is \(W_{{frozen}}+\Delta W_t\) with $W_{{frozen}}$ fixed, any admissible perturbation of the model parameters can be represented as a perturbation of $\Delta W_t$ within $\mathcal{W}_{\Delta}$.

% This perspective differs from the full-parameter viewpoint advocated by Flat-LoRA. Our result shows that, within the PECL hypothesis where only adapters are updated and priors/hyperpriors are supported on $W_\Delta$, the PAC-Bayesian bound can be expressed entirely in terms of adapter-subspace quantities.
% % under the PECL assumption of a frozen backbone, such global perturbations are not required by the PAC-Bayesian analysis.
% The contribution of flatness to the bound can be fully expressed through perturbations confined to $\mathcal{W}_{\Delta}$, since the backbone does not carry posterior uncertainty or task-specific updates in this regime.

% % Here \(\Sigma_t \succeq 0\) denotes a positive semidefinite covariance operator on the adapter subspace \(\mathcal W_\Delta\).

% % This analysis reveals a fundamental distinction from the Flat-LoRA perspective~\citep{li2025flatloralowrankadaptationflat}. 
% % While Flat-LoRA correctly identifies the importance of flatness in the merged full-model weights for generalization, it assumes this must be achieved by perturbing all parameters—including the frozen backbone.
% % However, our theoretical framework demonstrates that such full-parameter perturbation is unnecessary and potentially redundant in PECL. Because the operative model at task \(t\) is \(W_{{frozen}}+\Delta W_t\) with \(W_{\text{frozen}}\) fixed, the flatness of the merged model along all admissible directions is determined by the flatness of \(\Delta W_t\). Consequently, optimizing for flatness by perturbing only the adapter subspace inherently controls the geometry of the full model—enabled by the additive nature of LoRA updates—while preserving the stability of frozen, knowledge-preserving parameters.

% % Lemma~\ref{lemma: subspace-full}identifies the flatness term in our framework: the increase of empirical risk under Gaussian perturbations is exactly captured by the subspace expected sharpness $\mathbb{E}\text{-Sh}t^\Delta$. This term is entirely confined to $W\Delta$, since the perturbations act only on the adapter parameters.


% \subsection{A Sharpness-Aware PAC-Bayesian Bound for PECL}
% % \subsubsection{Sharpness-aware PAC-Bayes bound for PECL}

% % This insight, that subspace-restricted sharpness control is both necessary and sufficient in PECL, is formally captured in our main generalization bound, which explicitly confines both the sharpness measure and complexity terms to the adapter subspace.

% We now specialize Theorem \ref{thm:pathwise-pacbayes} to a PECL model with a frozen backbone and Gaussian posteriors supported on the LoRA update subspace $\mathcal{W}_{\Delta}$. In this setting, both the flatness and complexity terms in the bound can be expressed entirely in terms of quantities restricted to $\mathcal{W}_{\Delta}$.
% \begin{theorem}[Sharpness-aware PAC--Bayes bound for PECL]
% \label{thm:pecl-bound}
% For any $\delta\in(0,1]$, with probability at least $1-\delta$ over $S_1,\ldots,S_T$,
% \begin{equation}\label{eq:pecl-main}
% \begin{aligned}
% &\frac{1}{T}\sum_{t=1}^T
% \mathbb{E}_{P_t\sim\mathcal P_{t-1}}\mathbb{E}_{W\sim Q_t}L_{\mathcal{D}_t}(W)
% \ \le\
% \ \frac{1}{T}\sum_{t=1}^T
% \mathbb{E}_{P_t^\Delta\sim\mathcal P_{t-1}}
% \mathbb{E}_{\varepsilon \sim \mathcal{N}(0, \Sigma_t)}[\hat{L}_{S_t}(W_0 + \hat{\Delta W}_t + \varepsilon)
% \\
% &\quad +\sqrt{\frac{
% \displaystyle \sum_{t=1}^T \mathbb{E}_{P_t^\Delta\sim\mathcal P_{t-1}}\!\Big[\mathrm{KL}\!\big(Q_t^\Delta\|P_t^\Delta\big)\Big]
% \;+\;\displaystyle \sum_{t=1}^T \mathrm{KL}\!\big(\mathcal P_t\|\mathcal P_{t-1}\big)
% \;+\;\log\!\frac{1}{\delta}
% }{2\,\displaystyle\sum_{t=1}^T (m_t-1)}}, \qquad \varepsilon\sim\mathcal N(0,\Sigma_t)\ \text{on } \mathcal W_\Delta.
% \end{aligned}
% \end{equation}
% \end{theorem}

% % Theorem~\ref{thm:pecl-bound} specializes the hierarchical PAC-Bayesian framework to PECL. It shows that the generalization error decomposes into an empirical risk term and a complexity term entirely confined to the adapter subspace, regularized by the subspace expected sharpness $\mathbb{E}\text{-Sh}t^\Delta$. In other words, both curvature and posterior uncertainty are evaluated within $W\Delta$. Section 5 empirically examines how controlling this subspace sharpness affects continual learning performance.

% Theorem~\ref{thm:pecl-bound} specialises the hierarchical PAC-Bayesian framework to a PECL regime with a frozen backbone and posteriors supported on the adapter subspace $\mathcal{W}_\Delta$. In this setting, the empirical term is evaluated under Gaussian perturbations \( \varepsilon \in \mathcal W_\Delta \), which gives rise to the adapter-subspace expected sharpness \( \mathrm{E\text{-}Sh}^{\Delta}_t(\hat{\Delta W_t},\Sigma_t) \), while the complexity term decomposes into adapter-level divergences \( \mathrm{KL}(Q^\Delta_t \Vert P^\Delta_t) \) and a hyper-posterior drift \( \mathrm{KL}(P_t \Vert P_{t-1}) \). Since the backbone is frozen, neither sharpness nor complexity depends on directions outside \( \mathcal W_\Delta \), so the bound singles out adapter-subspace flatness as the relevant notion of flatness in PECL and theoretically justifies focusing flatness-aware optimisation on the adapter subspace.



% % The bound decomposes the average generalization risk into an empirical term, evaluated under Gaussian perturbations in $\mathcal{W}_\Delta$, and a complexity term that depends on the adapter-space divergences $\mathrm{KL}(Q_t^\Delta \Vert P_t^\Delta)$ and the hyper-posterior drift $\mathrm{KL}(\mathcal P_t \Vert \mathcal P_{t-1})$. Curvature enters the bound through the adapter-subspace expected sharpness $E\mbox{-}\mathrm{Sh}_\Delta$, which quantifies how sensitive the perturbed empirical risk is to small Gaussian perturbations confined to $\mathcal{W}_\Delta$. Under these modelling assumptions, controlling  sharpness in $\mathcal{W}_\Delta$ will tighten the PAC-Bayesian guarantee. In Section~5, we empirically study how variations in this subspace sharpness correlate with continual-learning metrics.


% % This bound explicitly decomposes the generalization error into an empirical risk term regularized by the expected sharpness within the LoRA subspace, and a complexity term confined to the adapter's parameter distribution. It thus provides a theoretical justification for our empirical finding: optimizing for flatness only in the low-rank adapter effectively controls the overall generalization performance in continual learning. 
% % \paragraph{What changes from full parameters to LoRA?}
% % (i) Task KL: $\mathrm{KL}(Q_t\|P_t)\;\to\;\mathrm{KL}(Q_t^{\Delta}\|P_t^{\Delta})$ (support restricted to the LoRA subspace);
% % (ii) Sharpness: $\mathrm{E\text{-}Sh}_t\;\to\;\mathrm{E\text{-}Sh}^{\Delta}_t$ (perturbations only along trainable directions);
% % (iii) Hyper-drift: $\sum_t \mathrm{KL}(\mathcal{P}_t\|\mathcal{P}_{t-1})$ persists, reflecting sequential knowledge accumulation at the hyper level.

% % This bound provides an explicit decomposition of the generalization error into an empirical risk term regularized by the expected sharpness within the LoRA subspace, and a complexity term confined to the adapter's parameter distribution. This formulation offers a theoretical justification for our empirical observation that optimizing for flatness solely in the low-rank adapter effectively controls overall generalization performance in continual learning. Compared to the full-parameter setting, our approach introduces three key adaptations: (i) the task-specific KL divergence now becomes \(\mathrm{KL}(Q_t^{\Delta}\|P_t^{\Delta})\), whose support is restricted to the LoRA subspace; (ii) sharpness is measured as \(\mathrm{E\text{-}Sh}^{\Delta}_t\), capturing perturbations only along trainable (low-rank) directions; and (iii) the hyper-drift term \(\sum_t \mathrm{KL}(\mathcal{P}_t\|\mathcal{P}_{t-1})\) remains unchanged, reflecting the sequential accumulation of knowledge at the hyperparameter level.


% % Notably, when a single adapter is updated sequentially with \(P_t^{\Delta} := Q_{t-1}^{\Delta}\) (conditioning on realized priors), the hyper-drift term vanishes completely. This occurs because the sequential adapter update eliminates the need for separate hyper-posterior distributions \(\mathcal{P}_t\)—each task's prior is directly defined by the previous task's posterior within the adapter subspace. Consequently, the complexity term simplifies to \(\sum_{t=1}^T \mathrm{KL}(Q_t^{\Delta}\|Q_{t-1}^{\Delta})\), capturing only the sequential accumulation of knowledge through adapter updates without additional hyper-level distribution shifts. This yields the following more tractable bound:

% In the SeqLoRA regime we update a single adapter sequentially and choose, for each task $t$, the adapter prior as the previous task posterior,
% $
% P_t^{\Delta} := Q_{t-1}^{\Delta}
% $,
% conditionally on the realized history $S_{1:t-1}$. From the hyper-distribution viewpoint, this update rule is deterministic: the random prior $P_t^{\Delta}$ is a measurable function of $Q_{t-1}^{\Delta}$ with the same law. Consequently the hyper-posteriors satisfy
% $
% \mathcal{P}_t = \mathcal{P}_{t-1}
% $
% for all $t$, and the hyper drift term in Theorem~\ref{thm:pecl-bound} vanishes,
% $
% \sum_{t=1}^{T} \mathrm{KL}\bigl(\mathcal{P}_t \,\Vert\, \mathcal{P}_{t-1}\bigr) = 0.
% $
% The bound therefore simplifies as follows.

% \begin{corollary}[SeqLoRA]\label{cor:seq-lora}
% If a single adapter is updated sequentially with priors $P_t^{\Delta} := Q_{t-1}^{\Delta}$, then with probability at least $1-\delta$,
% \[
% \frac1T\sum_{t=1}^T \mathbb{E}_{W\sim Q_t} L_{\mathcal{D}_t}(W)
% \ \le\
% \frac1T\sum_{t=1}^T \mathbb{E}_{\epsilon \sim \mathcal{N}(0,\Sigma_t)} \widehat L_{S_t}(W_t + \epsilon)
% \;+\;
% \sqrt{\frac{
% \sum_{t=1}^T \mathrm{KL}\bigl(Q_t^{\Delta}\,\Vert\,Q_{t-1}^{\Delta}\bigr)
% +\log(1/\delta)}
% {2\,\sum_{t=1}^T (m_t-1)}}.
% \]
% \end{corollary}


% % In the SeqLoRA setting, we consider a single adapter updated sequentially, where each task prior is fixed to the previous task posterior, that is, $P_t^{\Delta}:=Q_{t-1}^{\Delta}$ conditional on the realized history $S_{1:t-1}$. From the hyper distribution perspective, the hyper update is deterministic, so
% %  $
% %  \mathcal{P}_{t-1}=\delta_{Q_{t-1}^{\Delta}},
% %  \quad
% %  \mathcal{P}_{t}=\delta_{Q_{t}^{\Delta}}.
% %  $
% %  The pathwise change of measure at the hyper level is therefore
% %  $
% %  \sum_{t=1}^{T}\mathrm{KL}\big(\mathcal{P}_t\|\mathcal{P}_{t-1}\big)
% % =\sum_{t=1}^{T}\mathrm{KL}\big(\delta_{Q_t^{\Delta}}\|\delta_{Q_{t-1}^{\Delta}}\big)=0,$.
% % This yields the following more tractable bound:
% % \begin{corollary}[SeqLoRA]\label{cor:seq-lora}
% % If a single adapter is updated sequentially and $P_t^{\Delta}:=Q_{t-1}^{\Delta}$, then the hyper-term $\mathcal{P}_{t-1}=\delta_{Q_{t-1}^{\Delta}},
% %  \mathcal{P}_{t}=\delta_{Q_{t}^{\Delta}}$ and $\mathrm{KL}\big(\mathcal{P}_t\|\mathcal{P}_{t-1}\big)=0$ . With probability at least $1-\delta$,
% % \[
% % \frac1T\sum_{t=1}^T \mathbb{E}_{w\sim Q_t}L_{\mathcal{D}_t}(W)
% % \ \le\
% % \frac1T\sum_{t=1}^T \mathbb{E}_{\epsilon \sim \mathcal{N}(0,\sum)}\widehat L_{S_t}(W_0+\Delta W+\epsilon)
% % \;+\;
% % \sqrt{\frac{
% % \sum_{t=1}^T \mathrm{KL}(Q_t^{\Delta}\|Q_{t-1}^{\Delta})
% % +\log(1/\delta)}
% % {2\,\sum_{t=1}^T (m_t-1)}}.
% % \]
% % \end{corollary}



% % \section{From Weight Space Flatness to Feature Space Robustness}

% % % \subsection{From Abstract Subspaces to Concrete LoRA Implementations}

% % Section~3 shows that, in PECL, the PAC-Bayesian complexity depends only on flatness within the trainable update subspace $W_\Delta$, since curvature and posterior uncertainty are confined to this subspace rather than the full parameter space. A natural question is how such subspace flatness in weight space translates into robustness in feature space, and how this abstract view relates to concrete adapters such as LoRA that implement $W_\Delta$ in practice.
% % In this section we address these questions by expressing the flatness term in a form restricted to $W_\Delta$ and by making explicit how this update subspace is instantiated through LoRA parameterisations. 


% %In this section we show that the LoRA factor space is geometrically equivalent to the abstract subspace $W_\Delta$, so that flatness optimization in LoRA coordinates directly shapes the curvature term in the PAC-Bayesian bound and, as a result, both tightens the bound and stabilizes the learned representations.


% % The PAC-Bayesian analysis in Section~3 establishes a key conclusion for parameter efficient continual
% % learning (PECL). It is sufficient to control flatness within the trainable update subspace $W_\Delta$ in
% % order to improve generalization and the stability plasticity tradeoff. The full parameter space is not
% % required in the complexity term of the bound. Instead, the relevant curvature and posterior uncertainty are
% % strictly confined to $W_\Delta$.

% % A fundamental question remains. How does flatness in weight space, restricted to $W_\Delta$, translate into
% % robustness against catastrophic forgetting in feature space? Addressing this question requires a bridge
% % between the abstract theoretical framework developed in Section~3 and the concrete objects studied
% % throughout this section, namely low rank adapters such as LoRA that realize $W_\Delta$ in practice.

% % Our framework begins at the level of an abstract subspace $W_\Delta$, where both optimization and
% % generalization guarantees are formulated. In practical PECL systems such as LoRA, this subspace is
% % implemented by a specific low rank parameterization. In this section we show that this concrete
% % implementation does not compromise the theoretical guarantees. On the contrary, there is a precise
% % geometric equivalence between the abstract subspace view and the factor space on which LoRA is optimized.
% % Flatness optimization in LoRA factor coordinates shapes the geometry of $W_\Delta$ itself and therefore
% % operates directly on the curvature term that appears in the PAC-Bayesian bound.

% % This link allows us to explain why promoting flatness inside a low rank LoRA subspace simultaneously
% % tightens the generalization bound and stabilizes the learned representations, which in turn improves
% % continual learning performance.

% % \subsection{Realizing the Update Subspace via LoRA}\
% \subsection{Equivalence: LoRA and Update Subspaces}

% In PECL, the PAC-Bayesian complexity depends only on flatness within the trainable update subspace $W_\Delta$, since curvature and posterior uncertainty are confined to this subspace rather than the full parameter space.
% % A natural question is how such subspace flatness in weight space translates into final performance, and how this abstract view relates to concrete adapters such as LoRA that implement $W_\Delta$ in practice.
% % In this section we address these questions by expressing the flatness term in a form restricted to $W_\Delta$ and by making explicit how this update subspace is instantiated through LoRA parameterisations. 
% To understand how flatness in $W_\Delta$ influences model behaviour, we decompose the network into a
% shared representation mapping and a linear classifier. Let
% $
% \phi(x; W) \in \mathbb{R}^D
% $
% denote the representation produced by the backbone together with LoRA adapters, and let
% $W_{\mathrm{cls}} \in \mathbb{R}^{D \times K}$ denote the classification head. The predictive distribution is
% $
% p(y \mid x; W, W_{\mathrm{cls}})
% =
% \mathrm{softmax}\big(W_{\mathrm{cls}}^\top \phi(x; W)\big).
% $
% In PECL the model parameters decompose as
% $
% W = W_{\mathrm{froz}} + \Delta W,
% \quad
% \Delta W \in \mathcal{W}_\Delta,
% $
% where $W_{\mathrm{froz}}$ is the frozen backbone and $\Delta W$ is the only component updated on new tasks.

% Consider a small parameter perturbation $\Delta W \in \mathcal{W}_\Delta$. The induced feature shift can be
% approximated as
% $
% \Delta \phi(x)
% \approx
% \frac{\partial \phi(x; W)}{\partial W} \,\Delta W.
% $
% The sensitivity of the predictive distribution to this feature shift can be quantified through the local
% quadratic approximation of the KL divergence
% $$
% \mathrm{KL}\big(
% p(y \mid \phi + \Delta \phi; W_{\mathrm{cls}})
% \;\big\|\;
% p(y \mid \phi; W_{\mathrm{cls}})
% \big)
% \approx
% \Delta \phi(x)^\top E_\phi(x) \,\Delta \phi(x),
% $$
% where $
% E_\phi(x)
% :=
% \mathbb{E}_{y \sim p(\cdot \mid x;W,W_{\mathrm{cls}})}
% \Big[
% \nabla_\phi \log p(y \mid x; W, W_{\mathrm{cls}})
% \nabla_\phi \log p(y \mid x; W, W_{\mathrm{cls}})^\top
% \Big]
% $
% is the local feature matrix. For a linear classifier followed by a softmax, there is an analytic formulation
% $  E_\phi(x)
% =
% W_{\mathrm{cls}}
%  \Big(
%  \mathrm{Diag}(p) - p p^\top
%  \Big)
%  W_{\mathrm{cls}}^\top.
%  $
% % It measures how the current representation
% % $\phi(x)$ is aligned with the discriminative directions defined by $W_{\mathrm{cls}}$.
% % $E_\phi(x)$ is the local feature matrix. This matrix measures how the current representation
% % $\phi(x)$ is aligned with the discriminative subspace defined by $W_{\mathrm{cls}}$. It can be written as
% % $
% % E_\phi(x)
% % :=
% % \mathbb{E}_{y \sim p(\cdot \mid x;W,W_{\mathrm{cls}})}
% % \Big[
% % \nabla_\phi \log p(y \mid x; W, W_{\mathrm{cls}})
% % \nabla_\phi \log p(y \mid x; W, W_{\mathrm{cls}})^\top
% % \Big].
% % $
% % By the chain rule,
% % $
% % \nabla_W \log p(y \mid x; W, W_{\mathrm{cls}})
% % =
% % \Big(
% % \frac{\partial \phi(x; W)}{\partial W}
% % \Big)^\top
% % \nabla_\phi \log p(y \mid x; W, W_{\mathrm{cls}}).
% % $
% % For a small parameter perturbation $\Delta W$, the change in predictive distribution admits the quadratic
% % approximation
% % $$
% % \mathrm{KL}\big(
% % p(\cdot \mid x; W + \Delta W)
% % \;\big\|\;
% % p(\cdot \mid x; W)
% % \big)
% % \approx
% % \frac{1}{2}\,
% % \Delta W^\top F_W(x) \,\Delta W,
% % $$
% % where
% % $
% % F_W(x)
% % =
% % \Big(
% % \frac{\partial \phi(x; W)}{\partial W}
% % \Big)^\top
% % E_\phi(x)
% % \frac{\partial \phi(x; W)}{\partial W}
% % $
% % is the \textbf{Fisher information matrix at $x$}. Averaging over a dataset $S$ yields the empirical Fisher
% % matrix
% % $
% % F
% % =
% % \frac{1}{|S|}
% % \sum_{x \in S} F_W(x),
% % $
% % which serves as the curvature measure that controls the flatness term in the PAC-Bayesian bound.
% % In the PECL, the update only occurs in the lora part, so the fisher  information matrix is exactly the fisher information in the lora space. 
% For a small parameter perturbation $\Delta W$, the change in predictive distribution admits the quadratic
% approximation
% $$
% \mathrm{KL}\big(
% p(\cdot \mid x; W + \Delta W)
% \;\big\|\;
% p(\cdot \mid x; W)
% \big)
% \approx
% \frac{1}{2}\,
% \Delta W^\top F_W(x) \,\Delta W,
% $$
% where
% $
% F_W(x)
% =
% \Big(
% \frac{\partial \phi(x; W)}{\partial W}
% \Big)^\top
% E_\phi(x)
% \frac{\partial \phi(x; W)}{\partial W}
% $
% is the Fisher information matrix at $x$. Averaging over a dataset $S$ yields the empirical Fisher
% matrix
% $
% F
% =
% \frac{1}{|S|}
% \sum_{x \in S} F_W(x),
% $
% which serves as the curvature measure that controls the flatness term in the PAC--Bayesian bound.
% In the PECL regime, only the LoRA adapters are updated, so the relevant curvature is given by the restriction of $F_W(x)$ to the corresponding update subspace $\mathcal W_\Delta$, i.e.\ by its quadratic forms $\Delta W^\top F_W(x)\Delta W$ for perturbations $\Delta W \in \mathcal W_\Delta$.



% % Restricting attention to perturbations in the update subspace, write $W_\Delta = \mathrm{span}(U)$ with
% % $U \in \mathbb{R}^{d \times r}$ and parameterize $\Delta W = U \alpha$ for $\alpha \in \mathbb{R}^r$. The
% % quadratic approximation becomes
% % $$
% % \mathrm{KL}\big(
% % p(\cdot \mid x; W + \Delta W)
% % \;\big\|\;
% % p(\cdot \mid x; W)
% % \big)
% % \approx
% % \frac{1}{2}\,
% % \alpha^\top F_\Delta(x) \,\alpha,
% % $$
% % with  \textbf{the subspace Fisher matrix}
% % $
% % F_\Delta(x)
% % :=
% % U^\top F_W(x) U
% % =
% % U^\top
% % \Big(
% % \frac{\partial \phi(x; W)}{\partial W}
% % \Big)^\top
% % E_\phi(x)
% % \frac{\partial \phi(x; W)}{\partial W}
% % U
% % \in \mathbb{R}^{r \times r}.
% % $
% % This restricted matrix captures curvature along admissible directions in $W_\Delta$. 
% % Since the flatness term in Theorem~\ref{thm:pecl-bound} is given by
% % $
% % \mathbb{E}\text{-Sh}_t^\Delta
% % \;\approx\;
% % \tfrac{1}{2}\,\mathbb{E}_{x \in S_t}\bigl[\mathrm{tr}\bigl(F_\Delta(x)\,\Sigma_t\bigr)\bigr],
% % $
% % this chain of relations shows that spectral control of $F_\Delta(x)$ simultaneously tightens the PAC-Bayesian bound and limits the sensitivity of $\phi(x; W)$ to perturbations in $W_\Delta$.


% % In Section~3 the
% % subspace expected sharpness term in the PAC-Bayesian bound is expressed in terms of quadratic forms
% % associated with $F_\Delta(x)$ under Gaussian perturbations supported on $W_\Delta$. Thus, the spectrum
% % of $F_\Delta(x)$ governs both the flatness term in the bound and the sensitivity of the learned features to
% % updates in the adapter subspace.

% % \subsection{Geometric Equivalence Between LoRA and the Abstract Subspace $W_\Delta$}

% % In practice, LoRA does not parameterize $W_\Delta$ directly as a free linear subspace. Instead, it realizes
% % $W_\Delta$ through a structured low rank decomposition. 

% % \subsection{LoRA as a realization of the update subspace}
% In practice, PECL does not operate on an abstract update subspace but on concrete low-rank adapters.
% For a layer $\ell$ with weight matrix $W_\ell \in \mathbb R^{d_o^{(\ell)} \times d_i^{(\ell)}}$, a LoRA module introduces an update
% $
% \Delta W_\ell = B_\ell A_\ell,
% $
% where $B_\ell \in \mathbb R^{d_o^{(\ell)} \times r_\ell}$ and $A_\ell \in \mathbb R^{r_\ell \times d_i^{(\ell)}}$ are trainable factors. Let $(\widehat{A}_\ell, \widehat{B}_\ell)$ denote the current factors and consider first-order perturbations $(\delta A_\ell, \delta B_\ell)$.
% % By the product rule,
% % $
% % \delta(\Delta W_\ell)
% % =
% % \widehat{B}_\ell \,\delta A_\ell
% % +
% % \delta B_\ell \,\widehat{A}_\ell.
% % $
% % Vectorising and stacking the perturbations across all layers yields a global
% Writing
% $
% \theta
% =
% \big[
% \mathrm{vec}(\delta A_1);\,
% \mathrm{vec}(\delta B_1);\,
% \dots;\,
% \mathrm{vec}(\delta A_L);\,
% \mathrm{vec}(\delta B_L)
% \big]
% \in \mathbb R^q,
% $
% small changes in the factor coordinates induce weight updates
% $
% \mathrm{vec}(\Delta W)
% \approx
% J_{\mathrm{LoRA}} \,\theta,
% $
% where $J_{\mathrm{LoRA}} \in \mathbb R^{d \times q}$, $d$ is the dimension of the full parameter vector and $q$ is the total number of factor parameters.


% Variations of the LoRA factor parameters $\theta$ span exactly the weight update subspace $\mathcal W_\Delta \subset \mathbb R^d$. Optimisation in factor space is therefore equivalent to optimisation over perturbations confined to $\mathcal W_\Delta$, and any posterior on $\theta$ induces a posterior on weight updates supported on this subspace. In particular, a Gaussian posterior over $\theta$ corresponds to a Gaussian posterior over $\mathcal W_\Delta$, so the quantities in Theorem~2, such as $\mathrm{E\text{-}Sh}_t^{\Delta}(\hat{\Delta W}_t,\Sigma_t)$ and $\mathrm{KL}(Q^\Delta_t \Vert P^\Delta_t)$, can be evaluated either in factor space or on $\mathcal W_\Delta$. This correspondence makes explicit that promoting flatness in LoRA coordinates amounts to controlling the subspace sharpness that governs the PAC-Bayesian bound.






% % The key geometric insight is that the set of weight-space perturbations generated in $\theta$,  coincides exactly with the update subspace $W_\Delta$. Although optimisation is carried out in the factor coordinates $\theta$, the set of reachable perturbations in parameter space is exactly the linear subspace $W_\Delta$ assumed in the PAC--Bayesian analysis. Thus, a posterior over LoRA factors simply induces a posterior over weight updates supported on $W_\Delta$, and flatness-oriented optimisation in LoRA coordinates can be understood as controlling the curvature of the loss restricted to this update subspace.
% % From the PAC-Bayesian viewpoint, choosing a Gaussian posterior over the LoRA factor coordinates $\theta$ is therefore equivalent to choosing a Gaussian posterior over weight updates supported on $\mathcal{W}_{\Delta}$ . All quantities that appear in Theorem 2, such as $\mathrm{E\text{-}Sh}_t^{\Delta}(\hat{\Delta W}_t,\Sigma_t)$ and  \( \mathrm{KL}(Q^\Delta_t \Vert P^\Delta_t) \), can thus be computed either in factor space or in the abstract update subspace. This formalises the intuition that flatness promotion in LoRA coordinates directly controls the subspace sharpness that governs the bound.




% % For a layer $\ell$ with weight matrix
% % $W_\ell \in \mathbb{R}^{d_o^{(\ell)} \times d_i^{(\ell)}}$, the LoRA update is
% % $
% % \Delta W_\ell = B_\ell A_\ell,
% % $
% % where $B_\ell \in \mathbb{R}^{d_o^{(\ell)} \times r_\ell}$ and $A_\ell \in \mathbb{R}^{r_\ell \times d_i^{(\ell)}}$
% % are the trainable factors. Optimization algorithms such as SAM act directly on these factors rather than on
% % $W_\ell$ itself.

% % To connect theory and practice, we analyze how flatness should be understood under this specific
% % parameterization. Let $(\widehat{A}_\ell, \widehat{B}_\ell)$ denote the current factors and consider
% % first order perturbations $(\delta A_\ell, \delta B_\ell)$. By the product rule,
% % $
% % \delta(\Delta W_\ell)
% % =
% % \widehat{B}_\ell \,\delta A_\ell
% % +
% % \delta B_\ell \,\widehat{A}_\ell.
% % $
% % Vectorizing and stacking the perturbations across all layers yields a global LoRA Jacobian
% % $J_{\mathrm{LoRA}} \in \mathbb{R}^{d \times q}$, where $d$ is the dimension of the full parameter vector and
% % $q$ is the total number of factor parameters. Let
% % $
% % \theta
% % =
% % \big[
% % \mathrm{vec}(\delta A_1);\,
% % \mathrm{vec}(\delta B_1);\,
% % \dots;\,
% % \mathrm{vec}(\delta A_L);\,
% % \mathrm{vec}(\delta B_L)
% % \big]
% % \in \mathbb{R}^q
% % $
% % collect all factor perturbations. Then small changes in $\theta$ induce parameter updates
% % $
% % \mathrm{vec}(\Delta W)
% % \approx
% % J_{\mathrm{LoRA}} \,\theta.
% % $

% % The key geometric insight is that the image of $J_{\mathrm{LoRA}}$ in parameter space coincides with the
% % LoRA update subspace,
% % $
% % \mathrm{Im}(J_{\mathrm{LoRA}}) = W_\Delta.
% % $
% % Although optimization is carried out in the factor coordinates $\theta$, the set of reachable perturbations
% % in weight space is exactly the abstract subspace $W_\Delta$ that appears in the PAC-Bayesian analysis.
% % The factor space Fisher matrix
% % $
% % F_\theta(x)
% % =
% % J_{\mathrm{LoRA}}^\top F_W(x) J_{\mathrm{LoRA}}
% % $
% % is therefore related to the restricted Fisher $F_\Delta(x)$ by a change of coordinates. More precisely, there
% % exists an invertible matrix $R$ such that
% % $
% % F_\theta(x)
% % =
% % R^\top F_\Delta(x) R
% % \quad
% % \text{for all } x,
% % $
% % whenever the columns of $J_{\mathrm{LoRA}}$ span $W_\Delta$. In other words, $F_\theta(x)$ and
% % $F_\Delta(x)$ are congruent matrices with the same eigenvalues up to multiplicity.
% % This isomorphism validates the application of Theorem~\ref{thm:pecl-bound} to LoRA. Optimizing flatness in factor coordinates, for example
% % by using SAM to reduce the spectral norm of $F_\theta(x)$, is exactly shaping the geometry of the abstract
% % subspace $W_\Delta$. The concrete LoRA implementation is not an ad hoc approximation. It is a
% % parameter efficient realization of the subspace geometry required by the PAC-Bayesian framework in
% % Section~3.



% \subsection{Mechanism: From Weight Space Flatness to Feature Stability}

% % The equivalence between $F_\Delta(x)$ and $F_\theta(x)$ allows us to describe precisely how flatness
% % optimization improves feature stability. 



% In the PECL setting, Theorem~\ref{thm:pecl-bound} bounds the task-averaged generalisation error by an empirical term plus a subspace flatness term and two subspace complexity terms. All three terms are defined on, or induced by, the update subspace $\mathcal W_\Delta$.
% On the loss side, the definition of the measures Sh$(0)$, Sh$(1)$ and $\mathrm{E\text{-}Sh}$ together with the quadratic approximation \eqref{eq: approx} shows that subspace sharpness is governed by the curvature of the loss along $\mathcal W_\Delta$, as captured by the Fisher or Hessian restricted to this subspace.
% Large eigenvalues along $\mathcal W_\Delta$ mean that small adapter perturbations can cause substantial loss increases, which directly inflates $\mathrm{E\text{-}Sh}^{\Delta}_t$ and loosens the bound. When flatness-aware optimisation such as SAM, GAM, drives the solution into regions where the Fisher spectrum along $\mathcal W_\Delta$ is compressed, the same posterior covariance yields a smaller sharpness contribution, and the empirical term of the bound becomes tighter.
% On the complexity side, curvature along $\mathcal W_\Delta$ also constrains the two KL terms. In a highly curved subspace, keeping the loss small requires the adapter posterior to be either very concentrated or significantly shifted away from its prior, which inflates $\mathrm{KL}(Q^\Delta_t \Vert P^\Delta_t)$. When the subspace is flatter, the same loss control can be achieved with a less concentrated and more prior-aligned posterior, so the adapter KL remains moderate. A similar effect appears at the hyper level: if successive tasks can be fitted within a relatively flat region of $\mathcal W_\Delta$, the preferred adapter priors evolve more smoothly and the hyper-posterior drift $\mathrm{KL}(P_t \Vert P_{t-1})$ stays small. Subspace flatness therefore relaxes the trade-off between expected sharpness and both KL components of the bound and improves all curvature-dependent contributions simultaneously.
% % On the complexity side, curvature along $\mathcal W_\Delta$ also constrains the two KL terms. For the adapter-level divergence $\mathrm{KL}(Q^\Delta_t \Vert P^\Delta_t)$, a highly curved subspace forces the posterior to concentrate strongly or to move far from its prior in order to keep the loss small, both of which increase the KL contribution. Subspace flatness, by contrast, allows one to maintain a less concentrated and more prior-aligned posterior while still controlling $\mathrm{E\text{-}Sh}^{\Delta}*t$, thereby keeping the adapter KL moderate. An analogous effect holds at the hyper level: when successive tasks can be accommodated within a relatively flat region of $\mathcal W*\Delta$, the preferred adapter priors do not need to move aggressively, so the hyper-posterior drift $\mathrm{KL}(P_t \Vert P_{t-1})$ remains small. In this way, subspace flatness relaxes the trade-off between expected sharpness and both KL components of the bound and simultaneously improves all curvature-dependent terms.


% Finally, the feature-space perspective links these subspace curvature quantities to representation dynamics. The Fisher matrix can be written in terms of the Jacobian of the feature map $\phi(x;W)$, so its eigenvalues restricted to $\mathcal W_\Delta$ quantify how sensitive the representations are to LoRA updates. Reducing this subspace curvature stabilises features and class prototypes across tasks, preserves inter-task similarity structure, and leads to smoother decision boundaries. The improvements we observe in NME accuracy, prototype drift, CKA, and EFM spectra therefore provide a feature-level manifestation of the same mechanism: adapter-subspace flatness tightens the PAC–Bayesian bound and, at the same time, promotes feature stability in continual learning.



% % When flatness-aware optimisation drives the solution into regions where the Fisher spectrum along $\mathcal W*\Delta$ is compressed, the same posterior covariance yields a smaller sharpness contribution, and the empirical term of the bound becomes tighter. In our experiments, the measures Sh$(0)$, Sh$(1)$ and $\mathrm{E\text{-}Sh}$ are precisely different numerical approximations to this subspace curvature, so their systematic reduction under SAM, GAM, and related methods can be interpreted as a direct tightening of the loss-dependent part of the bound.






% % The PAC-Bayesian bound decomposes the generalization error in PECL into an empirical risk term and a flatness term restricted to the adapter subspace $W_\Delta$. 




% % For a Gaussian posterior \( Q_\Delta^t \), the adapter-subspace expected sharpness \( \mathrm{E\mbox{-}Sh}_t^\Delta \) measures the loss variation under parameter perturbations in \( W_\Delta \). Using a second-order Taylor approximation, this sharpness term approximates to \( \frac{1}{2} \mathbb{E}_{x \in S_t} [\operatorname{tr}(F_W(x) \Sigma_t)] \), where \( F_W(x) \) is the Fisher information matrix. Since \( \Sigma_t \) has support in \( W_\Delta \), the trace depends on the eigenvalues of \( F_W(x) \) along update directions. Large eigenvalues of $F_W(x)$ therefore correspond to directions in $W_\Delta$ along which small parameter changes can induce large increases in the loss.

% % Conversely, when flatness promotion reduces the dominant eigenvalues of $F_W(x)$ along $W_\Delta$, the same update $\Delta W$ produces a smaller loss increase on each task. In sequential learning, this reduces the interference with the loss of older tasks when adapting to new tasks, thereby alleviating catastrophic forgetting.

% % For a task level Gaussian posterior $Q_t^\Delta = \mathcal{N}(\hat{\Delta}W_t, \Sigma_t)$ supported on $W_\Delta$, the adapter-subspace expected sharpness 
% % can be written as
% % \[
% % \mathrm{E\mbox{-}Sh}_t^\Delta
% % =
% % \mathbb{E}_{\varepsilon \sim \mathcal{N}(0,\Sigma_t)}
% % \Big[
% % \widehat{L}_{S_t}(W_{\mathrm{froz}} + \widehat{\Delta W}_t + \varepsilon)
% % -
% % \widehat{L}_{S_t}(W_{\mathrm{froz}} + \widehat{\Delta W}_t)
% % \Big],
% % \qquad
% % \varepsilon \in W_\Delta.
% % \]


% % Under a second order approximation of the loss around $W_{\mathrm{froz}} + \widehat{\Delta W}_t$, this term admits a Fisher form
% % $
% % \mathrm{E\mbox{-}Sh}_t^\Delta
% % \;\approx\;
% % \frac{1}{2}\,
% % \mathbb{E}_{x \in S_t}
% % \big[
% % \mathrm{tr}\big(F_W(x)\,\Sigma_t\big)
% % \big],
% % $
% % where $F_W(x)$ is the Fisher information matrix defined and $\Sigma_t$ has support in $W_\Delta$. Since $\Sigma_t$ only explores directions within $W_\Delta$, this trace depends on the eigenvalues of $F_W(x)$ along update directions in $W_\Delta$. For positive semidefinite matrices the trace inequality
% % $
% % \mathrm{tr}\big(F_W(x)\,\Sigma_t\big)
% % \;\le\;
% % \lambda_{\max}\big(F_W(x)\big)\,\mathrm{tr}\big(\Sigma_t\big)
% % $
% % shows that the contribution of the flatness term is dominated by the largest eigenvalues of $F_W(x)$ in the update directions and by the overall spread of the posterior covariance $\Sigma_t$ inside $W_\Delta$.




% % This mechanism has a natural interpretation in feature space. we have
% % $
% % F_W(x)
% % =
% % \Big(
% % \frac{\partial \phi(x; W)}{\partial W}
% % \Big)^\top
% % E_\phi(x)
% % \frac{\partial \phi(x; W)}{\partial W},
% % $
% % so the largest eigenvalues of $F_W(x)$ quantify how sensitive the representation $\phi(x;W)$ is to infinitesimal changes in the trainable directions of $W$. Since PECL restricts updates to $W_\Delta$, these eigenvalues effectively measure the sensitivity of the features to LoRA adapter updates. Suppressing the dominant eigenvalues of $F_W(x)$ therefore limits how quickly decision boundaries and feature prototypes can move when $\Delta W$ is updated for new tasks. In other words, subspace flatness in $W_\Delta$ interpreted as a small spectrum of $F_W(x)$ along the update directions implies smaller loss changes under the parameter updates induced by new tasks and hence less destructive interference with previously learned solutions together with more stable feature representations in PECL.


% % The effect of these quantities on continual learning can be understood through a Taylor expansion. For a small update $\Delta W \in W_\Delta$, the loss on an earlier task $s$ satisfies
% % \[
% % L_{\mathcal{D}_s}(W_{\mathrm{froz}} + \widehat{\Delta W}_t + \Delta W)
% % \;\approx\;
% % L_{\mathcal{D}_s}(W_{\mathrm{froz}} + \widehat{\Delta W}_t)
% % +
% % \frac{1}{2}\,\Delta W^\top H_s\,\Delta W,
% % \]
% % where $H_s$ denotes the Hessian of $L_{\mathcal{D}_s}$ with respect to $W$ evaluated at $W_{\mathrm{froz}} + \widehat{\Delta W}_t$, and we only consider its action on directions in $W_\Delta$. 

% % In expectation over the data, this projected Hessian is well approximated by $F_W(x)$ along the same directions. 












% % the subspace expected sharpness can be expressed as
% % $
% % \mathrm{E\text{-}Sh}_t^\Delta
% % \approx
% % \frac{1}{2}\,
% % \mathbb{E}_{x \in S_t}
% % \big[
% % \mathrm{tr}\big(F_W\Sigma_t\big)
% % \big],\quad
% % \varepsilon \in W_\Delta.
% % $
% % Under a second order approximation of the loss around $W_
% % $ where  $F_W$ is obtained by projecting the parameter Fisher $F_W(x)$ onto $W_\Delta$. 
% % If $U$ denotes an orthonormal basis of $W_\Delta$, then
% % $
% % F_\Delta(x)
% % =
% % U^\top F_W(x)\,U.
% % $ 
% % For positive semidefinite matrices, the trace inequality
% % $
% % \mathrm{tr}\bigl(F_W(x)\,\Sigma_t\bigr)
% % \;\le\;
% % \lambda_{\max}\bigl(F_W(x)\bigr)\,\mathrm{tr}(\Sigma_t)
% % $
% % shows that the contribution of the flatness term is dominated by the largest eigenvalues of $F_W(x)$ and by the overall variance encoded in $\Sigma_t$.

% % In practice, flatness aware optimizers such as SAM construct, at each iteration, a worst case perturbation in parameter space that maximizes the empirical loss within a small neighbourhood. When optimization is confined to the LoRA factors, these perturbations align with leading eigenvectors of the local curvature in full space. Through repeated updates, this procedure suppresses the dominant eigenvalues of the factor space Fisher $F_W(x)$.


% % This mechanism can be interpreted in feature space: the eigenvalues of \( F_W(x) \) reflect the sensitivity of representations \( \phi(x; W) \) to adapter updates. By promoting flatness, PECL stabilizes feature representations and decision boundaries, further ensuring minimal disruption to previous tasks.

% % This spectral contraction has two complementary roles in PECL.

% % \begin{itemize}
% %     \item \textbf{Tightening the generalization bound.} For a fixed posterior covariance $\Sigma_t$, reducing $\lambda_{\max}\bigl(F_\Delta(x)\bigr)$ directly decreases $E\text{-Sh}_t^\Delta$ and thus the flatness term on the right hand side of the PAC-Bayesian bound. This yields a tighter generalization guarantee for the adapter subspace.
% %     \item \textbf{Stabilizing feature representations.} Suppressing the spectrum of $F_W(x)$ therefore limits how sensitive $\phi(x; W)$ is to updates in $W_\Delta$. This constrains the deformation of previously learned decision boundaries induced by new low rank updates and mitigates catastrophic forgetting.
% % \end{itemize}



% % For positive semidefinite matrices, the trace inequality
% % $
% % \mathrm{tr}\big(F_\theta(x)\Sigma_t\big)
% % \le
% % \lambda_{\max}\big(F_\theta(x)\big)\,
% % \mathrm{tr}(\Sigma_t)
% % $
% % shows that the contribution of flatness to the bound is dominated by the largest eigenvalue of the restricted
% % Fisher matrix.

% % Sharpness aware optimizers such as SAM iteratively search for perturbations that maximize the loss in a
% % small neighbourhood. These perturbations align with the leading eigenvectors of the local curvature in
% % factor space, and repeated updates have the effect of reducing the largest eigenvalues of $F_\theta(x)$. By
% % geometric equivalence, this also reduces $\lambda_{\max}(F_\theta(x))$.

% % This spectral contraction has two complementary roles in PECL.

% % \begin{enumerate}
% %     \item \textbf{Tightening the generalization bound.}
% %     For a fixed posterior covariance $\Sigma_t$, reducing $\lambda_{\max}(F_\theta(x))$ directly decreases
% %     $\mathrm{E\text{-}Sh}_t^\Delta$, which is the flatness term on the right hand side of the PAC-Bayesian
% %     bound. This leads to a tighter generalization guarantee for the adapter subspace.
% %     \item \textbf{Stabilizing feature representations.}
% %     % Since
% %     % $
% %     % F_\Delta(x)
% %     % =
% %     % U^\top
% %     % \Big(
% %     % \frac{\partial \phi(x; W)}{\partial W}
% %     % \Big)^\top
% %     % E_\phi(x)
% %     % \frac{\partial \phi(x; W)}{\partial W}
% %     % U,
% %     % $
% %     Suppressing the spectrum of $F_\theta(x)$ limits the sensitivity of $\phi(x; W)$ to updates in $W_\Delta$.
% %     This constrains how much new low rank updates can deform previously learned decision boundaries and
% %     thereby mitigates catastrophic forgetting.
% % \end{enumerate}







% \section{Experiment}



% \subsection{Experiment Settings}
% \label{sec:exp-settings}


% \textbf{Datasets and Evaluation Metric}
% We evaluate PECL methods on distribution drift benchmarks derived from ImageNet. First, we adopt ImageNet-R~\cite{ImagenetR_Hendrycks_2021_ICCV} as our main testbed, which is introduced to PECL by existing work~\cite{liangInfLoRAInterferenceFreeLowRank2024a} and has become a standard benchmark. We further adopt Tiny ImageNet-C~\cite{hendrycks2018benchmarking} and Tiny ImageNet-P~\cite{hendrycks2018benchmarking} to evaluate robustness on common corruptions and perturbations. Tiny ImageNet-C is a corrupted variant of Tiny ImageNet with 200 classes containing 15 types of corruption. Tiny ImageNet-P provides short perturbation sequences for each image.


% % ImageNet-R is our main testbed for class incremental learning. It contains natural images rendered into diverse artistic styles while preserving the original semantics, and is widely adopted in continual learning studies. To assess robustness under corruption and perturbation, we further use Tiny ImageNet-C and Tiny ImageNet-P. Tiny ImageNet-C applies common corruptions of fifteen types and five severity levels to Tiny ImageNet validation images, while Tiny ImageNet-P provides short perturbation sequences for each validation image.


% % ImageNet-R contains about 30k artistic renditions of 200 ImageNet classes, and represents a real-world distribution shift where the semantic labels remain fixed but the image style changes.

% % In addition, we evaluate on a 15-task text benchmark~\cite{razdaibiedina2023progressivepromptscontinuallearning} in NLP domains, which includes datasets from CL, GLUE~\cite{wang2019gluemultitaskbenchmarkanalysis}, SuperGLUE~\cite{NEURIPS2019_4496bf24}, and IMDB~\cite{maas2011learning}. 

% % This multi-task setup provides a comprehensive framework for assessing generalization across diverse NLP domains.
% % For computational efficiency, we adopt only "gaussian noise", "impulse noise", "shot noise" and one strength level for Tiny ImageNet-C. Also, we uniformly sample three frames from each sequence as static inputs to approximate the full perturbation trajectory for Tiny ImageNet-P.

% % After finishing task $t$, we evaluate the current model on all previously seen
% % tasks $j \in \{1,\ldots,t\}$ to get $ACC_t=\frac{1}{t}\sum_{j=1}^{t}a_{tj}$, where $a_{tj}$ denotes the accuracy of the $j$-th task once the model has learned the $t$-th task. At the end of the sequence we report standard
% % continual-learning metrics, namely the final accuracy ($\mathrm{ACC}_T$), the averaged accuracy $\mathrm{AAA}=\frac{1}{T} \sum_{i=1}^T \mathrm{ACC}_i$, and backward transfer $\mathrm{BWT} = \frac{1}{T - 1} \sum_{j=1}^{T - 1} \left(a_{Tj} - a_{jj}\right)$

% % As standard CL metrics, we report final average accuracy $\mathrm{ACC}_T$, overall averaged accuracy $\mathrm{AAA}$, and backward transfer $\mathrm{BWT}$. Besides, we also report the corresponding \textbf{Nearest Mean of Exemplars (NME)} accuracy, which assesses the quality of the learned representation independently of the classifier head. 
% % As standard CL metrics, we report accuracy $\operatorname{ACC}_T^{\mathrm{cls}}$, average accuracy over time $\operatorname{AAA}^{\mathrm{cls}}$, and backward transfer $\operatorname{BWT}^{\mathrm{cls}}$, computed from the task accuracy matrix obtained with the learned classifier head. In addition, we introduce the corresponding prototype based metrics $(\operatorname{ACC}_T^{\mathrm{nme}}, \operatorname{AAA}^{\mathrm{nme}}, \operatorname{BWT}^{\mathrm{nme}})$, where each class is represented by the mean feature of its exemplars and predictions are made by nearest prototype search. 
% % The former family measures end to end classification performance, while the latter isolates the quality and stability of the learned representation.
% % For each class $c$ with training samples $\{x_i : y_i = c\}$, we compute a prototype $\mu_c = \mathrm{normalize}\Big( \frac{1}{n_c} \sum_{i : y_i = c} \frac{\phi(x_i)}{\lVert \phi(x_i) \rVert_2} \Big),$ and classify test samples by nearest prototype $\hat{y} = \arg\min_c \lVert \frac{\phi(x)}{\lVert \phi(x) \rVert_2} - \mu_c \rVert_2^2$. 
% % To comprehensively evaluate the flatness of the learned model, we employ several geometric metrics.

% % In the class incremental setting, after training task $t$ we evaluate the current model on all previously seen tasks $j \in \{1,\ldots,t\}$ and record the accuracies $a_{tj}$. 

% We define the accuracy at step $t$ as
% $
% \operatorname{ACC}_t = \frac{1}{t} \sum_{j=1}^{t} a_{tj},
% $
% and report three standard CL metrics:
% $
% \operatorname{ACC}_T,$
% $\operatorname{AAA} =\frac{1}{T} \sum_{t=1}^{T} \operatorname{ACC}_t, $
% $\operatorname{BWT} =\frac{1}{T - 1} \sum_{j=1}^{T - 1} \bigl(a_{Tj} - a\bigr).$
% As standard CL metrics, we report final accuracy $\operatorname{ACC}_T^{\mathrm{cls}}$, average accuracy over time $\operatorname{AAA}^{\mathrm{cls}}$, and backward transfer $\operatorname{BWT}^{\mathrm{cls}}$, which uses the trained classification head. In addition, we also introduce the corresponding prototype based metrics $(\operatorname{ACC}_T^{\mathrm{nme}}, \operatorname{AAA}^{\mathrm{nme}}, \operatorname{BWT}^{\mathrm{nme}})$, which are computed in the same way from the Nearest Mean of Exemplars (NME) accuracy obtained by nearest prototype classification, where each class is represented by the mean feature of its exemplars.



% For loss landscape flatness, we compute the maximum eigenvalue $\lambda_{\max}\big(H(\hat{W})\big)$ and trace $\mathrm{tr}\big(H(\hat{W})\big)$ of the Hessian matrix, alongside $\lambda_{\max}\big(F(\hat{W})\big)$ and $\mathrm{tr}\big(F(\hat{W})\big)$ derived from the empirical Fisher information matrix, using the Lanczos~\citep{ghorbaniInvestigationNeuralNet2019} algorithms follow~\citep{foretSharpnessAwareMinimizationEfficiently2021}. For feature robustness, we measure {the L2 prototype drift}to quantify representation shift, and employ linear $\text{CKA}(X_0, X_T) = \frac{\| \bar{X}_0^T \bar{X}_T \|_F^2}{\| \bar{X}_0^T \bar{X}_0 \|_F \cdot \| \bar{X}_T^T \bar{X}_T \|_F}$ to assess feature structural consistency. Additionally, we analyze the empirical feature matrix through its trace and effective rank $r_{\text{eff}} = \frac{(\sum_i \lambda_i)^2}{\sum_i \lambda_i^2}$.

% % ,  providing a multi-faceted perspective on representation stability and flatness throughout training.

%  % $\Delta_c^{L2} = \|\mu_c^{(T)} - \mu_c^{(0)}\|_2$ 

% % and anisotropy index $\text{ani} = \frac{\rho(E)}{\mathrm{tr}(E)/d}$,

% \textbf{Architecture and Methods}  
% We employ a ViT-B/16 backbone ~\cite{dosovitskiy2020image}, pre-trained in a supervised manner on ImageNet-21K, as the base model for image-class incremental learning. 
% % For NLP experiments, we adopt the T5 model~\cite{raffelExploringLimitsTransfer2023}, following the setups of ~\cite{wang2023orthogonal}.
% To probe the intrinsic representation ability of the backbone, we employ a PerTask Linear Probe that trains an independent linear classifier for each task while keeping the backbone frozen. 

% % For PECL we consider several adapter organizations. SeqFT sequentially fine tunes all backbone parameters across tasks and serves as a strong full parameter baseline. SeqLoRA attaches a single shared LoRA adapter to each layer and updates it across all tasks, which corresponds to a compact and fully shared adapter subspace. IncLoRA extends this design by progressively expanding the adapter subspace as tasks arrive, thereby fusing task specific and shared directions. OLoRA augments the shared subspace with orthogonality promoting constraints that encourage task specific adapters to explore complementary directions within a controlled geometry.



% To study the effect of flatness promotion, we integrate several perturbation based optimizers. SAM minimises the worst case empirical loss within an $\ell_2$ neighbourhood of radius $\rho$ and corresponds to the zero-order sharpness $Sh^{(0)}_S(W,\rho)$. RWP construct noise weighted perturbations that target expected sharpness, GAM optimize the first-order flatness and C-Flat combines elements of SAM and GAM. Together, these methods span different relaxation levels of the sharpness notions $Sh^{(0)}_S(W,\rho)$, $E\text{-Sh}_t(W,\Sigma_t)$, and $Sh^{(1)}_S(W,\rho)$ introduced in Section~3. Besides, we also evaluate the Flat-LoRA method which apply RWP on the whole model rather than only LoRA part.


% % We also employ SeqFT, which sequentially fine-tunes all backbone parameters across tasks. To systematically analyze parameter-sharing behaviors among task adapters, we implement SeqLoRA, IncLoRA, and OLoRA, corresponding to fully shared subspaces, partially fused subspaces, and constraint-augmented sharing, respectively. Furthermore, to investigate various flatness-promoting objectives, we integrate SAM~\cite{foretSharpnessAwareMinimizationEfficiently2021}, RWP~\cite{liRevisitingRandomWeight2024}, and GAM~\cite{Zhang_GAM}—each aligned with one of the three sharpness notions $\mathrm{Sh}^{(0)}_S(W,\rho)$, $\mathrm{E\text{-}Sh}_t(\hat{W}_t, \Sigma_t)$, and $\mathrm{Sh}^{(1)}_S(W,\rho)$. Additionally, we include C-Flat~\citep{liCFlatMoreEfficient2025}, which combines elements of both SAM and GAM.

% %  We use PerTaskFT to evaluate the backbone's ablility, which train a separate linear classfier for each task.  We adapot SeqFT  train the full backbone parameters on a sequence of tasks .
% %  To fully understand task adapters share a parameter space, we adopt the SeqLoRA, IncLoRA and OLoRA to represent the Fully shared subspace, Partially fused subspaces, and Constraint-augmented sharing respectly. 
% % Besides, to check different flatness prometing method, we adopt SAM, RWP, GAM which related to the the three definition $\mathrm{Sh}^{(0)}_S(W,\rho) ,\mathrm{E\text{-}Sh}_t(\hat{W}_t, \Sigma_t),\mathrm{Sh}^{(1)}_S(W,\rho)$. Addtionaly, we also adopt C-Flat which is mixture of SAM and GAM.  




 

% \textbf{Implementation Details}

% Unless stated otherwise, we use SGD as the base optimizer, train each session for $20$ epochs and SAM radius $\rho = 0.05$. LoRA based variants are trained with learning rate $0.01$ without momentum or weight decay. The same training configuration is used across SeqLoRA, IncLoRA, and OLoRA, so that differences in performance can be attributed to the geometry of the adapter subspace and the choice of flatness optimizer rather than to changes in training budget.

% % For NLP experiments with T5-Large, we fine-tune using LoRA-based adapters optimized by AdamW in bfloat16 precision. Models are trained for 5 epochs with a constant learning rate of \(1\times 10^{-4}\), effective batch size of 32 (16 per device across 2 GPUs). All LoRA variants use rank \(r=8\), while OLoRA incorporates constraint weights \(\lambda_1=0.5, \lambda_2=0\). 


% \subsection{Main Result}

% % In this section we empirically stress test the central claim of our sharpness-aware PAC-Bayesian analysis. 
% % Theorem~2 states that, in PECL with a frozen backbone, generalization is controlled by flatness and complexity restricted to the adapter subspace $W_\Delta$. 
% % Experiments E1 to E4 are organized to investigate this claim. 
% % E1 examines whether promoting flatness in $W_\Delta$ improves continual learning metrics and feature stability on ImageNet-R. 
% % E2 verifies that the same relationship holds under on Tiny ImageNet-C and Tiny ImageNet-P. 
% % E3 orders  optimisers by their strength in suppressing subspace curvature. 
% % E4 contrasts adapter-only with global perturbations, directly challenging the global flatness hypothesis advocated by Flat-LoRA and testing the structural assumptions behind our PAC-Bayesian analysis.

% % \subsection{Implications and empirical hypotheses}
% % \label{sec:implications-q1q3}

% % The sharpness aware PAC-Bayesian bound in Section~3 yields three central implications for parameter efficient continual learning.

% % First, the complexity term depends only on the adapter subspace $W_\Delta$. Lemma~1 shows that the Kullback–Leibler divergence between full model distributions reduces to the divergence between their restrictions on $W_\Delta$, because the frozen backbone contributes no posterior uncertainty. As a result, the posterior covariance $\Sigma_t$ can and should be confined to $W_\Delta$, and the flatness term in the bound is fully captured by the subspace expected sharpness $E\text{-Sh}_t^\Delta$.

% % Second, for a Gaussian posterior $Q_t^\Delta = \mathcal{N}(\hat{\Delta}W_t, \Sigma_t)$ supported on $W_\Delta$, the expected sharpness term admits the approximation
% % $$
% % E\text{-Sh}_t^\Delta
% % \;\approx\;
% % \frac{1}{2}\,\mathbb{E}_{x \in S_t}
% % \bigl[\mathrm{tr}\bigl(F_\Delta(x)\,\Sigma_t\bigr)\bigr],
% % $$
% % where $F_\Delta(x)$ is the Fisher information restricted to $W_\Delta$. This term is controlled by the spectrum of $F_\Delta(x)$ and the posterior covariance $\Sigma_t$. In particular, the trace inequality
% % $$
% % \mathrm{tr}\bigl(F_\Delta(x)\,\Sigma_t\bigr)
% % \;\le\;
% % \lambda_{\max}\bigl(F_\Delta(x)\bigr)\,\mathrm{tr}(\Sigma_t)
% % $$
% % shows that the dominant eigenvalues of $F_\Delta(x)$ govern the contribution of flatness to the generalization bound.

% % Third, in the SeqLoRA setting where a single adapter is updated sequentially with $P_t^\Delta := Q_{t-1}^\Delta$, the hyper drift term $\sum_t \mathrm{KL}(P_t \,\Vert\, P_{t-1})$ vanishes by construction. In this regime the bound simplifies: the generalization behaviour across tasks is controlled by the empirical risks on each task together with the adapter subspace expected sharpness $E\text{-Sh}_t^\Delta$ and the task level Kullback–Leibler divergences $\mathrm{KL}(Q_t^\Delta \,\Vert\, P_t^\Delta)$.

% % These observations lead to concrete, testable hypotheses.

% % \begin{itemize}
% %     \item[(H1)] For a fixed adapter organization and training budget, promoting flatness within $W_\Delta$ by reducing $E\text{-Sh}_t^\Delta$ should improve standard continual learning metrics such as final accuracy $\mathrm{ACC}_T$, averaged accuracy $\mathrm{AAA}$, and backward transfer $\mathrm{BWT}$.
% %     \item[(H2)] Improvements in subspace flatness should manifest primarily through feature level stability. In particular, reductions in the spectrum of $F_\Delta(x)$ are expected to coincide with more stable class prototypes, higher nearest mean of exemplars (NME) accuracy in the frozen feature space, and more consistent feature geometry across tasks.
% %     \item[(H3)] The quantitative gains from flatness promotion should depend on the geometry of $W_\Delta$ and on the choice of flatness optimizer. Larger or more entangled adapter subspaces can facilitate fitting in distribution data but may also admit more unstable directions in $F_\Delta(x)$ and $\Sigma_t$, which influences the attainable trade off between flexibility and stability.
% % \end{itemize}

% % In Section~5 we design experiments that directly address these hypotheses. Experiment~E1 focuses on (H1) and (H2) on ImageNet-R, while E2–E4 analyze the roles of adapter geometry, the choice of flatness optimizer, and local versus global perturbations, thereby addressing (H3).



% \begin{figure}[thbp]
%     \centering
%     \begin{minipage}{0.495\textwidth}
%         \centering
%         \includegraphics[width=0.9\linewidth]{figures_TRML/AAA_curve_Imagenet-R_summary_inc10_rank16—point2.pdf}
%         % {figures_TRML/AAA_curve_Imagenet-R_summary_inc10_rank16_replot2.pdf}
%         \subcaption{Accuracy $\operatorname{ACC}_t^{\mathrm{cls}}$ on ImageNet-R}
%         \label{fig:curves_imageR_1_AAA}
%     \end{minipage}
%     \begin{minipage}{0.495\textwidth}
%         \centering
%         \includegraphics[width=0.9\linewidth]{figures_TRML/NME_curve_Imagenet-R_summary_inc10_rank16_point2.pdf}
%         % {figures_TRML/NME_curve_Imagenet-R_summary_inc10_rank16_replot2.pdf}
%         \subcaption{NME-based accuracy $\operatorname{ACC}_t^{\mathrm{nme}}$ on ImageNet-R}
%         \label{fig:curves_imageR_2_NME}
%     \end{minipage}
%     \caption{
%     % Summary curves of accuracy with the classifier head $\operatorname{ACC}_t^{\mathrm{cls}}$ and prototype-based accuracy $\operatorname{ACC}_t^{\mathrm{nme}}$ on ImageNet-R ($T=20, r=16$). 
%     %     The curves report performance after each task for different class incremental methods and optimization configurations. 
%         Summary curves of classifier-head accuracy $\operatorname{ACC}_t^{\mathrm{cls}}$ and prototype-based accuracy $\operatorname{ACC}_t^{\mathrm{nme}}$ on ImageNet-R ($T=20, r=16$) over the task sequence, comparing different class-incremental methods and optimization configurations.
%         Flatness-promoting optimizers such as SAM and GAM consistently outperform SGD, and their advantage becomes more pronounced on later tasks, which suggests more stable representations and improved performance.
%         % Flatness promotion optimizers (SAM and GAM) consistently outperform SGD, and their advantage tends to increase over later tasks, indicating more stable representations and better performance.
%     }
%     \label{fig:combined_curves_imageR}
% \end{figure}



% \begin{table}[htpb]
% \centering
% \caption{
% Class incremental performance and weight space flatness metrics on ImageNet-R ($T=20$, $R=16$). 
% The table reports standard class incremental metrics  together with three sharpness indices and empirical Hessian and Fisher statistics in the weight space.
% }
% \label{table:q1_ci_with_weight_flatness_reorg}
% \setlength{\tabcolsep}{4pt}
% \resizebox{\linewidth}{!}{
% \begin{tabular}{l l ccc ccccc c}
% \toprule
% \multirow{2}{*}{\textbf{Method}} & \multirow{2}{*}{\textbf{Opt.}} &
% \multicolumn{3}{c}{\textbf{CL ($\uparrow$)}} &
% \multicolumn{2}{c}{\textbf{Sharpness ($\downarrow$)}} &
% \multicolumn{2}{c}{\textbf{Emp Hessian ($\downarrow$)}} &
% \multicolumn{2}{c}{\textbf{Emp Fisher ($\downarrow$)}} \\
% \cmidrule(lr){3-5}\cmidrule(lr){6-7}\cmidrule(lr){8-9}\cmidrule(lr){10-11}
%  & & $\operatorname{ACC}_{20}^{\mathrm{cls}}$  &$\operatorname{AAA}^{\mathrm{cls}}$   & $\operatorname{BWT}^{\mathrm{cls}}$
%  & Sh$^{(0)}(\rho)$  & Sh$^{(1)}(\rho)$
%  & $\lambda_{\max}(H)$ & $\mathrm{tr}(H)$  
%  & $\lambda_{\max}(F)$ & $\mathrm{tr}(F)$   \\
% \midrule
% \multirow{1}{*}{\textbf{PerTaskLP}}
% & SGD   & $55.22^{\pm 0.31}$ & $59.36^{\pm 0.19}$ & $-8.35^{\pm 0.75}$ & 0.072   & 0.072 & 5.31    & 1182.49  & 8.47   & 1015.07  \\
% \midrule
% \multirow{2}{*}{\textbf{SeqFT}}
% & SGD   &$50.72^{\pm 2.20}$ & $59.89^{\pm 1.32}$ & $-38.94^{\pm 1.87}$ & 0.582  & 0.517 & 4692.68 & 25621.40 & 538.79 & 29703.66 \\
% & SAM   &$56.07^{\pm 2.60}$ & $66.47^{\pm 1.53}$ & $-33.04^{\pm 4.71}$ & 0.259  & 0.252 & 474.25  & 19747.59 & 75.43  & 5881.32  \\
% \rowcolor{LightPurple} &GAM  &$59.78^{\pm 0.74}$ & $67.31^{\pm 1.18}$ & $-32.70^{\pm 0.26}$ & 0.177  & 0.176 & 877.97  & 18879.15 & 32.00  & 3273.08  \\
% \midrule
% \multirow{2}{*}{\textbf{SeqLoRA}}
% & SGD    &$63.14^{\pm 2.17}$ & $69.07^{\pm 1.33}$ & $-11.91^{\pm 1.82}$ & 0.142  & 0.114 & 47.37 & 1502.55 & 35.51 & 1401.15 \\
% & SAM    &$67.04^{\pm 0.71}$ & $71.42^{\pm 0.71}$ & $-9.26^{\pm 0.66}$ & 0.120  & 0.093 & 24.41   & 996.59   & 21.39  & 821.45   \\
% \rowcolor{LightPurple} &GAM  &$73.31^{\pm 0.52}$ & $76.80^{\pm 0.56}$ & $-9.37^{\pm 0.61}$ & 0.060  &0.058  & 4.81    & 339.49   & 3.72   & 240.91   \\
% \midrule
% \multirow{2}{*}{\textbf{IncLoRA}}
% & SGD    &$65.77^{\pm 2.84}$ & $71.29^{\pm 0.78}$ & $-4.47^{\pm 2.85}$ & 0.076    & 0.062 & 15.46 & 969.34  & 8.89  & 836.91 \\
% & SAM    &$68.86^{\pm 0.92}$ & $72.59^{\pm 0.39}$ & $-3.51^{\pm 0.50}$ & 0.053   & 0.052 & 3.72    & 712.88   & 4.35   & 526.49   \\
% \rowcolor{LightPurple} &GAM  &$72.86^{\pm 0.23}$ & $76.52^{\pm 0.48}$ & $-6.41^{\pm 0.48}$ & 0.035   & 0.035 & 2.69    & 281.22   & 2.09   & 192.25   \\
% \midrule
% \multirow{2}{*}{\textbf{OLoRA}}
% & SGD    &$66.64^{\pm 0.44}$ & $71.41^{\pm 0.47}$ & $-3.52^{\pm 0.43}$ & 0.078   & 0.058 & 6.94    & 938.93   & 9.00   & 831.23   \\
% & SAM    &$68.53^{\pm 0.78}$ & $72.75^{\pm 0.76}$ & $-3.16^{\pm 0.68}$ & 0.066   & 0.048 & 4.40    & 775.40   & 7.41   & 628.12   \\
% \rowcolor{LightPurple} &GAM &$72.44^{\pm 0.36}$ & $76.05^{\pm 0.52}$ & $-7.12^{\pm 0.88}$ & 0.036   & 0.036 & 2.22    & 261.47   & 2.08   & 198.30\\
% \midrule
% \multirow{1}{*}{\textbf{l2p}}
% & Adam   & $71.35^{\pm 0.05}$ & $74.84^{\pm 0.73}$ & $-6.34^{\pm 0.50}$  &-  &- &- &- &- &- \\
% \midrule
% \multirow{1}{*}{\textbf{dualprompt}}
% & Adam   & $71.59^{\pm 0.68}$ & $74.80^{\pm 1.18}$ & $-3.38^{\pm 0.45}$ &-  &- &- &- &- &- \\
% \bottomrule
% \end{tabular}
% }
% \end{table}



% % \begin{table}[htpb]
% % \centering
% % \caption{
% % Class incremental performance and weight space flatness metrics on ImageNet-R ($T=20$, $R=16$). 
% % The table reports standard class incremental metrics  together with three sharpness indices and empirical Hessian and Fisher statistics in the weight space.
% % }
% % \label{table:q1_ci_with_weight_flatness_reorg}
% % \setlength{\tabcolsep}{4pt}
% % \resizebox{\linewidth}{!}{
% % \begin{tabular}{l l ccc cccccc c}
% % \toprule
% % \multirow{2}{*}{\textbf{Method}} & \multirow{2}{*}{\textbf{Opt.}} &
% % \multicolumn{3}{c}{\textbf{CL ($\uparrow$)}} &
% % \multicolumn{3}{c}{\textbf{Sharpness ($\downarrow$)}} &
% % \multicolumn{2}{c}{\textbf{Emp Hessian ($\downarrow$)}} &
% % \multicolumn{2}{c}{\textbf{Emp Fisher ($\downarrow$)}} \\
% % \cmidrule(lr){3-5}\cmidrule(lr){6-8}\cmidrule(lr){9-10}\cmidrule(lr){11-12}
% %  & & $\operatorname{ACC}_{20}^{\mathrm{cls}}$  &$\operatorname{AAA}^{\mathrm{cls}}$   & $\operatorname{BWT}^{\mathrm{cls}}$
% %  & Sh$^{(0)}(\rho)$ & E-Sh$(Q)$ & Sh$^{(1)}(\rho)$
% %  & $\lambda_{\max}(H)$ & $\mathrm{tr}(H)$  
% %  & $\lambda_{\max}(F)$ & $\mathrm{tr}(F)$   \\
% % \midrule
% % \multirow{1}{*}{\textbf{PerTaskLP}}
% % & SGD   & 55.22^{\pm 0.31 & 59.36^{\pm 0.19 & -8.35^{\pm 0.75 & 0.072 & 0.125  & 0.072 & 5.31    & 1182.49  & 8.47   & 1015.07  \\
% % \midrule
% % \multirow{2}{*}{\textbf{SeqFT}}
% % & SGD   &50.72 ^{\pm 2.20 & 59.89 ^{\pm 1.32 & -38.94 ^{\pm 1.87 & 0.582 & -0.105 & 0.517 & 4692.68 & 25621.40 & 538.79 & 29703.66 \\
% % & SAM   &56.07^{\pm 2.60 & 66.47 $\pm$1.53 & -33.04^{\pm 4.71 & 0.259 & -0.003 & 0.252 & 474.25  & 19747.59 & 75.43  & 5881.32  \\
% % \rowcolor{LightPurple} &GAM  &59.78^{\pm 0.74 & 67.31^{\pm 1.18 & -32.70^{\pm 0.26 & 0.177 & -0.036 & 0.176 & 877.97  & 18879.15 & 32.00  & 3273.08  \\
% % \midrule
% % \multirow{2}{*}{\textbf{SeqLoRA}}
% % & SGD    & 63.14^{\pm 2.17 & 69.07^{\pm 1.33 & -11.91^{\pm 1.82 & 0.142 & -0.006 & 0.114 & 47.37 & 1502.55 & 35.51 & 1401.15 \\
% % & SAM    & 67.04^{\pm 0.71 & 71.42^{\pm 0.71 & -9.26 ^{\pm 0.66 & 0.120 & -0.123 & 0.093 & 24.41   & 996.59   & 21.39  & 821.45   \\
% % \rowcolor{LightPurple} &GAM  &73.31^{\pm 0.52 & 76.80^{\pm 0.56 & -9.37  $\pm$0.61 & 0.060 & 0.080  & 0.058 & 4.81    & 339.49   & 3.72   & 240.91   \\
% % \midrule
% % \multirow{2}{*}{\textbf{IncLoRA}}
% % & SGD    & 65.77$\pm$2.84 & 71.29^{\pm 0.78 & -4.47^{\pm 2.85 & 0.076 & 0.071  & 0.062 & 15.46 & 969.34  & 8.89  & 836.91 \\
% % & SAM    & 68.86^{\pm 0.92 & 72.59^{\pm 0.39 & -3.51 ^{\pm 0.50 & 0.053 & 0.058  & 0.052 & 3.72    & 712.88   & 4.35   & 526.49   \\
% % \rowcolor{LightPurple} &GAM  &72.86^{\pm 0.23 & 76.52^{\pm 0.48 & -6.41 ^{\pm 0.48 & 0.035 & 0.052  & 0.035 & 2.69    & 281.22   & 2.09   & 192.25   \\
% % \midrule
% % \multirow{2}{*}{\textbf{OLoRA}}
% % & SGD    & 66.64^{\pm 0.44 & 71.41^{\pm 0.47 & -3.52 ^{\pm 0.43 & 0.078 & 0.020  & 0.058 & 6.94    & 938.93   & 9.00   & 831.23   \\
% % & SAM    & 68.53^{\pm 0.78 & 72.75^{\pm 0.76 & -3.16 ^{\pm 0.68 & 0.066 & 0.026  & 0.048 & 4.40    & 775.40   & 7.41   & 628.12   \\
% % \rowcolor{LightPurple} &GAM &72.44^{\pm 0.36 & 76.05^{\pm 0.52 & -7.12 ^{\pm 0.88 & 0.036 & 0.051  & 0.036 & 2.22    & 261.47   & 2.08   & 198.30\\
% % \bottomrule
% % \end{tabular}
% % }
% % \end{table}


% \subsubsection{(E1)Empirical validation of adapter-subspace flatness}



% Experiment~E1 provides the main empirical on ImageNet-R. We train PerTask Linear Probe, SeqFT, SeqLoRA, IncLoRA, and OLoRA under identical schedules, using SGD as the base optimizer and optionally applying SAM or GAM to the adapter parameters. We then evaluate continual learning metrics, parameter space sharpness, and feature space diagnostics as described in Section~\ref{sec: flat minima} and~\ref{sec:exp-settings}. The results are summarised in Table \ref{table:q1_ci_with_weight_flatness_reorg} and Table~\ref{tab:q1_lp_feat_reorg}, and the evolution of $\operatorname{ACC}_t^{\mathrm{cls}}$ and $\operatorname{ACC}_t^{\mathrm{nme}}$ over tasks is shown in Figure~\ref{fig:combined_curves_imageR}.


% From the PECL perspective, flatness promotion confined to $W_\Delta$ consistently improves performance. For all three LoRA organizations, replacing SGD with SAM increases final accuracy $\mathrm{ACC}_{20}$ and averaged accuracy $\mathrm{AAA}$, while making $\mathrm{BWT}$ less negative. The $\operatorname{ACC}_t^{\mathrm{cls}}$ curves in Figure~\ref{fig:curves_imageR_1_AAA} show that these improvements are not restricted to the last task. The SAM and GAM variants maintain higher accuracy throughout the sequence, and the gap relative to SGD typically widens for later tasks. The $\operatorname{ACC}_t^{\mathrm{nme}}$ curves in Figure~\ref{fig:curves_imageR_2_NME} follow the same pattern, which indicates that the benefits are not solely due to changes in the classifier head but extend to the quality of the shared representation.


% \begin{table}[htpb]
% \centering
% \caption{
% NME based performance and feature space diagnostics on ImageNet-R ($T=20$, $R=16$). 
% The table reports NME accuracy metrics together with prototype drift, CKA similarity to the initial representation, and spectral properties of the empirical feature matrix.
% }
% \label{tab:q1_lp_feat_reorg}
% \setlength{\tabcolsep}{4pt}
% \renewcommand{\arraystretch}{0.95} % 可选：稍微压缩行距
% \resizebox{0.9\linewidth}{!}{
% \begin{tabular}{l l ccc  cc cc}
% \toprule
% \multirow{2}{*}{\textbf{Method}} & \multirow{2}{*}{\textbf{Opt.}} &
% \multicolumn{3}{c}{\textbf{NME} ($\uparrow$)} &
% \multicolumn{2}{c}{\textbf{Feature Stability}} &
% \multicolumn{2}{c}{\textbf{Feature Matrix}} \\
% \cmidrule(lr){3-5} \cmidrule(lr){6-7} \cmidrule(lr){8-9}
% &  & $ \operatorname{ACC}_{20}^{\mathrm{nme}}$  & $\operatorname{AAA}^{\mathrm{nme}}$  & $\operatorname{BWT}^{\mathrm{nme}}$ 
%  & L2 drift & CKA$_{0}$
%  & $\mathrm{tr}(E)$ & eff \\ 
% \midrule
% \multirow{1}{*}{\textbf{PerTaskLP}}
%   & SGD &$62.96^{\pm0.12}$ & $66.58^{\pm0.37}$ & $-7.66^{\pm1.27}$
%      & 0        & 1        & 0.82 & 128.11  \\
% \midrule
% \multirow{2}{*}{\textbf{SeqFT}}
%  & SGD     &$69.49^{\pm0.16}$ & $75.10^{\pm0.57}$ & $-16.11^{\pm0.23}$ & 29.57 & 0.78 & 0.58 & 93.53  \\
%  & SAM     &$72.36^{\pm2.46}$ & $78.31^{\pm0.50}$ & $-14.04^{\pm3.30}$ & 17.04 & 0.86 & 0.74 & 127.28 \\
%  \rowcolor{LightPurple}&GAM  &$73.74^{\pm0.88}$ & $78.32^{\pm0.94}$ & $-13.36^{\pm0.20}$ & 13.80 & 0.85 & 0.78 & 127.87 \\
% \midrule
% \multirow{2}{*}{\textbf{SeqLoRA}}
%  & SGD     &$72.82^{\pm1.27}$ & $76.62^{\pm0.68}$ & $-6.75^{\pm1.29}$  & 16.58 & 0.73 & 0.77 & 125.50 \\
%  & SAM     &$74.53^{\pm0.60}$ & $77.62^{\pm0.43}$ & $-5.54^{\pm0.72}$  & 14.82 & 0.73 & 0.81 & 135.38 \\
% \rowcolor{LightPurple} &GAM   &$75.74^{\pm0.15}$ & $78.86^{\pm0.53}$ & $-6.19^{\pm0.33}$  & 11.28 & 0.78 & 0.78 & 133.96 \\
% \midrule
% \multirow{2}{*}{\textbf{IncLoRA}}
%  & SGD      &$75.41^{\pm0.26}$ & $78.37^{\pm0.26}$ & $-4.32^{\pm0.59}$   & 16.87 & 0.69 & 0.67 & 121.18 \\
%  & SAM     &$75.95^{\pm0.07}$ & $78.68^{\pm0.25}$ & $-3.87^{\pm0.03}$ & 15.14 & 0.68 & 0.70 & 127.16 \\
%  \rowcolor{LightPurple}&GAM  &$76.72^{\pm0.54}$ & $79.42^{\pm0.40}$ & $-3.63^{\pm0.67}$  & 12.15 & 0.75 & 0.73 & 119.68 \\
% \midrule
% \multirow{2}{*}{\textbf{OLoRA}}
%  & SGD     &$75.10^{\pm0.41}$ & $78.65^{\pm0.67}$ & $-3.90^{\pm0.54}$ & 17.23 & 0.69 & 0.68 & 120.89 \\
%  & SAM     &$75.73^{\pm0.32}$ & $78.70^{\pm0.67}$ & $-3.47^{\pm0.54}$ & 15.17 & 0.72 & 0.70 & 123.76 \\
% \rowcolor{LightPurple} &GAM  &$76.56^{\pm0.53}$ & $79.16^{\pm0.43}$ & $-3.37^{\pm0.68}$  & 12.16 & 0.75 & 0.74 & 119.36\\
% \bottomrule
% \end{tabular}
% }
% \end{table}




% The feature diagnostics in Table~\ref{tab:q1_lp_feat_reorg} and Figure~\ref{fig:combined_curves_imageR} further support the mechanism above. NME accuracy in the frozen feature space increases under SAM and GAM, which reflects a better alignment between class labels and feature prototypes. Prototype drift across tasks is reduced, while CKA between the initial backbone features and the final features is higher, indicating that flatness aware training stabilises the global structure of the representation. At the same time, the spectrum of the empirical feature matrix becomes more isotropic: the trace is moderated and the effective rank increases, so that sensitivity is distributed more evenly across feature directions rather than concentrating on a few highly curved modes. Taken together, these observations confirm that controlling sharpness inside $W_\Delta$ is sufficient to improve continual learning metrics and that the effect is mediated by enhanced feature stability.




% % The preceding theoretical analysis yields a clear prediction: in PECL, the stability-plasticity trade off is governed primarily by the flatness of the adapter subspace. Our empirical results, summarized in Table \ref{table:q1_ci_with_weight_flatness_reorg} and visualized in Figures \ref{fig:combined_curves_imageR} and corroborate this prediction. 
% % Table~\ref{table:q1_ci_with_weight_flatness_reorg} shows that, across all LoRA organizations, SAM improves the continual learning metrics when restricted to the adapter parameters. Final accuracy $\mathrm{ACC}_{20}$ and averaged accuracy AAA consistently increase, whereas backward transfer $\mathrm{BWT}$ becomes less negative. The AAA curves in Figure~\ref{fig:combined_curves_imageR} confirm that the gain is not limited to the final task. SAM maintains a higher accuracy throughout the sequence and the gap often widens in later tasks. The NME-based AAA curves in Figure ~\ref{fig:curves_imageR_2_NME} exhibit the same pattern, which indicates that the improvement is not only due to the classifier head but also in the feature space.

% % This subsection addresses Q1 (sufficiency): does enforcing flatness only on the LoRA parameters already suffice to improve CL performance?


% % \begin{table}[t]
% % \centering
% % \caption{Class-incremental learning on ImageNet-R (T=20) using ViT-B/16-IN21K; means^{\pm sample std over 5 seeds with shuffled class orders. We set \( \rho=0.05 \) for both \( \mathrm{Sh}^{(0)}(\rho) \)  and \( \mathrm{Sh}^{(1)}(\rho) \) ; \( \mathrm{E\text{-}Sh}(Q) \) is estimated using 100 Gaussian perturbations \( \delta\sim\mathcal N(0,\sigma^2 I) \) with \( \sigma=\rho/\sqrt{d} \), where $d$ is the dimensionality of the parameter vector and table entries report \(\mathrm{E\text{-}Sh}(Q)\times10^{4}\) for readability.
% % }
% % % .All methods are implemented under a unified codebase.~\citep{sun2025pilot,zhou2024class,zhou2024continual}
% % \label{table:q1_ci_with_weight_flatness_reorg}
% % \setlength{\tabcolsep}{4pt}
% % \resizebox{\linewidth}{!}{
% % \begin{tabular}{l l ccc cccccc c}
% % \toprule
% % \multirow{2}{*}{\textbf{Opt.}} & \multirow{2}{*}{\textbf{Method}} &
% % \multicolumn{3}{c}{\textbf{CL ($\uparrow$)}} &
% % \multicolumn{3}{c}{\textbf{Sharpness ($\downarrow$)}} &
% % \multicolumn{2}{c}{\textbf{Emp Hessian ($\downarrow$)}} &
% % \multicolumn{2}{c}{\textbf{Emp Fisher ($\downarrow$)}} \\
% % \cmidrule(lr){3-5}\cmidrule(lr){6-8}\cmidrule(lr){9-10}\cmidrule(lr){11-12}
% %  & & ${\text{ACC}}_{20}$  & $\overline{\text{ACC}}_{20}$    &BWT
% %  % & BWT
% %  & Sh$^{(0)}(\rho)$ & E-Sh$(Q)$ & Sh$^{(1)}(\rho)$
% % & $\lambda_{\max}(H)$ & $\mathrm{tr}(H)$  
% % & $\lambda_{\max}(F)$ & $\mathrm{tr}(F)$   \\
% % \midrule
% % % \multirow{6}{*}{\textbf{without SAM}}
% % % \multirow{2}{*}{\centering{\textbf{Backbone}}}
% % % & MulTaskLP  \\
% % % & MulTaskFT   \\
% % % & MulTaskLoRA   \\
% % % \midrule
% % % without SAM \rotatebox[origin=c]{90}
% % \multirow{1}{*}{\centering{\textbf{Backbone}}}
% % & PerTaskLP   &55.22$\pm$0.31  &59.36$\pm$0.19  &-8.35$\pm$0.74  &0.0717  &0.102  &0.0747  &6.40  &1166.77  &7.75  &995.16  \\
% % \midrule
% % \multirow{4}{*}{\centering{\textbf{SGD}}}
% % & SeqFT   &52.30$\pm$2.34 &60.43$\pm$1.00  &-37.68$\pm$2.26  &0.556  &-0.0286  &0.528  &3252.05  & 25285.24 &525.58  &30321.04  \\
% % \cmidrule(lr){2-12}
% %  & SeqLoRA    &63.67$\pm$1.97  & 69.41$\pm$1.11 & -11.55$\pm$1.63 &0.134  &-0.0814  &0.104  &51.98  &1415.11  &33.20  &1350.37  \\
% % \cmidrule(lr){2-12}
% %  & IncLoRA    &65.69$\pm$2.79  &71.51$\pm$0.91  &-4.49$\pm$2.83    &0.0821  &0.0125  &0.0556   &16.12  &963.40  &10.107 &865.42  \\
% %   % & InfLoRA \\
% % % & MOS        &  &  &  &  &  &  &  &\\
% %  \cmidrule(lr){2-12}
% %  & OLoRA      &66.60$\pm$0.44  &71.79$\pm$0.13  &-3.44$\pm$0.68  &0.0907  &-0.00657  &0.051  &7.84  &1012.26  &11.74  &931.91  \\
% %  % & TUNA        &  &  &  &  &  &  &  &\\
% %  % with 
% % \midrule
% % \multirow{4}{*}{\centering{\textbf{SAM}}}
% % % & PerTaskLP   &55.80$\pm$0.30  &60.00$\pm$0.15  &-8.32$\pm$0.56  &0.0708  &0.106  &0.0738  &6.37  &1140.97  &7.73  &978.06  \\
% % & SeqFT   &58.17$\pm$2.86 &66.9$\pm$1.04  &-31.73$\pm$3.13  &0.251  &0.00646  &0.259  &369.13  &19823.97  &91.59  &6280.24  \\
% % \cmidrule(lr){2-12}
% %  & SeqLoRA    &67,43$\pm$1.06  &71.37$\pm$0.63  &-9.10$\pm$0.66  &0.120  &-0.123  &0.0928  &24.41  &996.59  &21.39  &821.45  \\
% % \cmidrule(lr){2-12}
% %  & IncLoRA    &68.86$\pm$0.92  &72.59$\pm$0.39  &-3.51$\pm$0.50   &0.0529  & 0.0581 &0.0522  &3.723 &712.88 &4.35 &526.49  \\
% %   % & InfLoRA \\
% % % & MOS        &  &  &  &  &  &  &  &\\
% %  \cmidrule(lr){2-12}
% %  & OLoRA      &68.18$\pm$0.67  &73.07$\pm$0.72  &-3.19$\pm$0.95  &0.0766  &-0.000273  &0.0431  &5.136  &794.64  &10.02  &689.33  \\
% %  % & TUNA        &  &  &  &  &  &  &  &\\
% % \bottomrule
% % \end{tabular}
% % }
% % \end{table}



% % \begin{table}[t]
% % \centering
% % \caption{Class-incremental learning on ImageNet-R (T=20) using ViT-B/16-IN21K; means^{\pm sample std over 5 seeds with shuffled class orders. We set \( \rho=0.05 \) for both \( \mathrm{Sh}^{(0)}(\rho) \)  and \( \mathrm{Sh}^{(1)}(\rho) \) ; \( \mathrm{E\text{-}Sh}(Q) \) is estimated using 100 Gaussian perturbations \( \delta\sim\mathcal N(0,\sigma^2 I) \) with \( \sigma=\rho/\sqrt{d} \), where $d$ is the dimensionality of the parameter vector and table entries report \(\mathrm{E\text{-}Sh}(Q)\times10^{4}\) for readability.
% % }
% % % .All methods are implemented under a unified codebase.~\citep{sun2025pilot,zhou2024class,zhou2024continual}
% % \label{table:q1_ci_with_weight_flatness_reorg}
% % \setlength{\tabcolsep}{4pt}
% % \resizebox{\linewidth}{!}{
% % \begin{tabular}{l l ccc cccccc c}
% % \toprule
% % \multirow{2}{*}{\textbf{Opt.}} & \multirow{2}{*}{\textbf{Method}} &
% % \multicolumn{3}{c}{\textbf{CL ($\uparrow$)}} &
% % \multicolumn{3}{c}{\textbf{Sharpness ($\downarrow$)}} &
% % \multicolumn{2}{c}{\textbf{Emp Hessian ($\downarrow$)}} &
% % \multicolumn{2}{c}{\textbf{Emp Fisher ($\downarrow$)}} \\
% % \cmidrule(lr){3-5}\cmidrule(lr){6-8}\cmidrule(lr){9-10}\cmidrule(lr){11-12}
% %  & & ${\text{ACC}}_{20}$  &  AAA   &BWT
% %  % & BWT
% %  & Sh$^{(0)}(\rho)$ & E-Sh$(Q)$ & Sh$^{(1)}(\rho)$
% % & $\lambda_{\max}(H)$ & $\mathrm{tr}(H)$  
% % & $\lambda_{\max}(F)$ & $\mathrm{tr}(F)$   \\
% % \midrule
% % % $\overline{\text{ACC}}_{20}$
% % % \multirow{6}{*}{\textbf{without SAM}}
% % % \multirow{2}{*}{\centering{\textbf{Backbone}}}
% % % & MulTaskLP  \\
% % % & MulTaskFT   \\
% % % & MulTaskLoRA   \\
% % % \midrule
% % % without SAM \rotatebox[origin=c]{90}
% % \multirow{1}{*}{\centering{\textbf{Backbone}}}
% % & PerTaskLP    & 55.22^{\pm 0.31 & 59.36^{\pm 0.19 & -8.35^{\pm 0.75 & 0.072 & 0.125  & 0.072 & 5.31    & 1182.49  & 8.47   & 1015.07  \\
% % \midrule
% % \multirow{4}{*}{\centering{\textbf{SGD}}}
% % & SeqFT   & 51.48^{\pm 2.03 & 60.18 $\pm$1.06 & -38.66^{\pm 1.41 & 0.562 & -0.027 & 0.502 & 3261.86 & 26125.32 & 590.76 & 29588.55 \\
% % \cmidrule(lr){2-12}
% %  & SeqLoRA    & 63.56^{\pm 2.25 & 69.39^{\pm 1.28 & -11.55^{\pm 1.89 & 0.134 & -0.081 & 0.104 & 51.98   & 1415.11  & 33.20  & 1350.38  \\
% % \cmidrule(lr){2-12}
% %  & IncLoRA    & 65.69$\pm$ 2.79 & 71.51^{\pm 0.91 & -4.49 ^{\pm 2.83 & 0.082 & 0.012  & 0.056 & 16.12   & 963.40   & 10.11  & 865.42   \\
% %   % & InfLoRA \\
% % % & MOS        &  &  &  &  &  &  &  &\\
% %  \cmidrule(lr){2-12}
% %  & OLoRA       & 66.48^{\pm 0.37 & 71.48^{\pm 0.55 & -3.55 ^{\pm 0.52 & 0.078 & 0.020  & 0.058 & 6.94    & 938.93   & 9.00   & 831.23   \\

% %  % & TUNA        &  &  &  &  &  &  &  &\\
% %  % with 
% % \midrule
% % \multirow{4}{*}{\centering{\textbf{SAM}}}
% % % & PerTaskLP   &55.80$\pm$0.30  &60.00$\pm$0.15  &-8.32$\pm$0.56  &0.0708  &0.106  &0.0738  &6.37  &1140.97  &7.73  &978.06  \\
% % & SeqFT   & 57.53 ^{\pm 3.12 & 66.82 ^{\pm 1.24 & -32.02 ^{\pm 3.76 & 0.253 & -0.011 & 0.248 & 474.53  & 20157.25 & 82.17  & 6039.48  \\
% % \cmidrule(lr){2-12}
% %  & SeqLoRA     & 67.04 ^{\pm 0.71 & 71.42^{\pm 0.71 & -9.26  ^{\pm 0.66 & 0.120 & -0.123 & 0.093 & 24.41   & 996.59   & 21.39  & 821.45   \\

% % \cmidrule(lr){2-12}
% %  & IncLoRA    & 68.86 ^{\pm 0.92 & 72.59 ^{\pm 0.39 & -3.51  ^{\pm 0.50 & 0.053 & 0.058  & 0.052 & 3.72    & 712.88   & 4.35   & 526.49   \\
 
% %   % & InfLoRA \\
% % % & MOS        &  &  &  &  &  &  &  &\\
% %  \cmidrule(lr){2-12}
% %  & OLoRA      & 68.53 ^{\pm 0.78 & 72.75 ^{\pm 0.76 & -3.16  ^{\pm 0.68 & 0.066 & 0.026  & 0.048 & 4.40    & 775.40   & 7.41   & 628.12  
% %  \\
% %  % & TUNA        &  &  &  &  &  &  &  &\\
% % \bottomrule
% % \end{tabular}
% % }
% % \end{table}





% % \begin{table}[t]
% % \centering
% % \caption{Class-incremental learning on ImageNet-R (T=20) using ViT-B/16-IN21K; means^{\pm sample std over 5 seeds with shuffled class orders. We set \( \rho=0.05 \) for both \( \mathrm{Sh}^{(0)}(\rho) \)  and \( \mathrm{Sh}^{(1)}(\rho) \) ; \( \mathrm{E\text{-}Sh}(Q) \) is estimated using 100 Gaussian perturbations \( \delta\sim\mathcal N(0,\sigma^2 I) \) with \( \sigma=\rho/\sqrt{d} \), where $d$ is the dimensionality of the parameter vector and table entries report \(\mathrm{E\text{-}Sh}(Q)\times10^{4}\) for readability.
% % }
% % % .All methods are implemented under a unified codebase.~\citep{sun2025pilot,zhou2024class,zhou2024continual}
% % \label{table:q1_ci_with_weight_flatness_reorg}
% % \setlength{\tabcolsep}{4pt}
% % \resizebox{\linewidth}{!}{
% % \begin{tabular}{l l ccc cccccc c}
% % \toprule
% % \multirow{2}{*}{\textbf{Opt.}} & \multirow{2}{*}{\textbf{Method}} &
% % \multicolumn{3}{c}{\textbf{CL ($\uparrow$)}} &
% % \multicolumn{3}{c}{\textbf{Sharpness ($\downarrow$)}} &
% % \multicolumn{2}{c}{\textbf{Emp Hessian ($\downarrow$)}} &
% % \multicolumn{2}{c}{\textbf{Emp Fisher ($\downarrow$)}} \\
% % \cmidrule(lr){3-5}\cmidrule(lr){6-8}\cmidrule(lr){9-10}\cmidrule(lr){11-12}
% %  & & ${\text{ACC}}_{20}$  &  AAA   &BWT
% %  % & BWT
% %  & Sh$^{(0)}(\rho)$ & E-Sh$(Q)$ & Sh$^{(1)}(\rho)$
% % & $\lambda_{\max}(H)$ & $\mathrm{tr}(H)$  
% % & $\lambda_{\max}(F)$ & $\mathrm{tr}(F)$   \\
% % \midrule
% % % $\overline{\text{ACC}}_{20}$
% % % \multirow{6}{*}{\textbf{without SAM}}
% % % \multirow{2}{*}{\centering{\textbf{Backbone}}}
% % % & MulTaskLP  \\
% % % & MulTaskFT   \\
% % % & MulTaskLoRA   \\
% % % \midrule
% % % without SAM \rotatebox[origin=c]{90}
% % \multirow{1}{*}{\centering{\textbf{Backbone}}}
% % & PerTaskLP    & 55.22^{\pm 0.31 & 59.36^{\pm 0.19 & -8.35^{\pm 0.75 & 0.072 & 0.125  & 0.072 & 5.31    & 1182.49  & 8.47   & 1015.07  \\
% % \midrule
% % \multirow{4}{*}{\centering{\textbf{SGD}}}
% % & SeqFT   & 51.48^{\pm 2.03 & 60.18 $\pm$1.06 & -38.66^{\pm 1.41 & 0.562 & -0.027 & 0.502 & 3261.86 & 26125.32 & 590.76 & 29588.55 \\
% % \cmidrule(lr){2-12}
% %  & SeqLoRA    & 63.56^{\pm 2.25 & 69.39^{\pm 1.28 & -11.55^{\pm 1.89 & 0.134 & -0.081 & 0.104 & 51.98   & 1415.11  & 33.20  & 1350.38  \\
% % \cmidrule(lr){2-12}
% %  & IncLoRA    & 65.69$\pm$ 2.79 & 71.51^{\pm 0.91 & -4.49 ^{\pm 2.83 & 0.082 & 0.012  & 0.056 & 16.12   & 963.40   & 10.11  & 865.42   \\
% %   % & InfLoRA \\
% % % & MOS        &  &  &  &  &  &  &  &\\
% %  \cmidrule(lr){2-12}
% %  & OLoRA       & 66.48^{\pm 0.37 & 71.48^{\pm 0.55 & -3.55 ^{\pm 0.52 & 0.078 & 0.020  & 0.058 & 6.94    & 938.93   & 9.00   & 831.23   \\

% %  % & TUNA        &  &  &  &  &  &  &  &\\
% %  % with 
% % \midrule
% % \multirow{4}{*}{\centering{\textbf{SAM}}}
% % % & PerTaskLP   &55.80$\pm$0.30  &60.00$\pm$0.15  &-8.32$\pm$0.56  &0.0708  &0.106  &0.0738  &6.37  &1140.97  &7.73  &978.06  \\
% % & SeqFT   & 57.53 ^{\pm 3.12 & 66.82 ^{\pm 1.24 & -32.02 ^{\pm 3.76 & 0.253 & -0.011 & 0.248 & 474.53  & 20157.25 & 82.17  & 6039.48  \\
% % \cmidrule(lr){2-12}
% %  & SeqLoRA     & 67.04 ^{\pm 0.71 & 71.42^{\pm 0.71 & -9.26  ^{\pm 0.66 & 0.120 & -0.123 & 0.093 & 24.41   & 996.59   & 21.39  & 821.45   \\

% % \cmidrule(lr){2-12}
% %  & IncLoRA    & 68.86 ^{\pm 0.92 & 72.59 ^{\pm 0.39 & -3.51  ^{\pm 0.50 & 0.053 & 0.058  & 0.052 & 3.72    & 712.88   & 4.35   & 526.49   \\
 
% %   % & InfLoRA \\
% % % & MOS        &  &  &  &  &  &  &  &\\
% %  \cmidrule(lr){2-12}
% %  & OLoRA      & 68.53 ^{\pm 0.78 & 72.75 ^{\pm 0.76 & -3.16  ^{\pm 0.68 & 0.066 & 0.026  & 0.048 & 4.40    & 775.40   & 7.41   & 628.12  
% %  \\
% %  % & TUNA        &  &  &  &  &  &  &  &\\
% % \bottomrule
% % \end{tabular}
% % }
% % \end{table}


% % \begin{table}[t]
% % \centering
% % \caption{Class-incremental learning on ImageNet-R (T=20) using ViT-B/16-IN21K; means^{\pm sample std over 5 seeds with shuffled class orders. We set \( \rho=0.05 \) for both \( \mathrm{Sh}^{(0)}(\rho) \)  and \( \mathrm{Sh}^{(1)}(\rho) \) ; \( \mathrm{E\text{-}Sh}(Q) \) is estimated using 100 Gaussian perturbations \( \delta\sim\mathcal N(0,\sigma^2 I) \) with \( \sigma=\rho/\sqrt{d} \), where $d$ is the dimensionality of the parameter vector and table entries report \(\mathrm{E\text{-}Sh}(Q)\times10^{4}\) for readability.}
% % \label{table:q1_ci_with_weight_flatness_reorg}
% % \setlength{\tabcolsep}{4pt}
% % \resizebox{\linewidth}{!}{
% % \begin{tabular}{l l ccc cccccc c}
% % \toprule
% % \multirow{2}{*}{\textbf{Method}} & \multirow{2}{*}{\textbf{Opt.}} &
% % \multicolumn{3}{c}{\textbf{CL ($\uparrow$)}} &
% % \multicolumn{3}{c}{\textbf{Sharpness ($\downarrow$)}} &
% % \multicolumn{2}{c}{\textbf{Emp Hessian ($\downarrow$)}} &
% % \multicolumn{2}{c}{\textbf{Emp Fisher ($\downarrow$)}} \\
% % \cmidrule(lr){3-5}\cmidrule(lr){6-8}\cmidrule(lr){9-10}\cmidrule(lr){11-12}
% %  & & ${\text{ACC}}_{20}$  &  AAA   & BWT
% %  & Sh$^{(0)}(\rho)$ & E-Sh$(Q)$ & Sh$^{(1)}(\rho)$
% %  & $\lambda_{\max}(H)$ & $\mathrm{tr}(H)$  
% %  & $\lambda_{\max}(F)$ & $\mathrm{tr}(F)$   \\
% % \midrule
% % \multirow{1}{*}{\textbf{PerTaskLP}}
% % & SGD   & 55.22^{\pm 0.31 & 59.36^{\pm 0.19 & -8.35^{\pm 0.75 & 0.072 & 0.125  & 0.072 & 5.31    & 1182.49  & 8.47   & 1015.07  \\
% % \midrule
% % \multirow{2}{*}{\textbf{SeqFT}}
% % & SGD   & 51.48^{\pm 2.03 & 60.18^{\pm 1.06 & -38.66^{\pm 1.41 & 0.562 & -0.027 & 0.502 & 3261.86 & 26125.32 & 590.76 & 29588.55 \\
% % & SAM   & 57.53^{\pm 3.12 & 66.82^{\pm 1.24 & -32.02^{\pm 3.76 & 0.253 & -0.011 & 0.248 & 474.53  & 20157.25 & 82.17  & 6039.48  \\
% % &GAM  \\
% % \midrule
% % \multirow{2}{*}{\textbf{SeqLoRA}}
% % & SGD    & 63.56^{\pm 2.25 & 69.39^{\pm 1.28 & -11.55^{\pm 1.89 & 0.134 & -0.081 & 0.104 & 51.98   & 1415.11  & 33.20  & 1350.38  \\
% % & SAM    & 67.04^{\pm 0.71 & 71.42^{\pm 0.71 & -9.26 ^{\pm 0.66 & 0.120 & -0.123 & 0.093 & 24.41   & 996.59   & 21.39  & 821.45   \\
% % &GAM  \\
% % \midrule
% % \multirow{2}{*}{\textbf{IncLoRA}}
% % & SGD    & 65.69^{\pm 2.79 & 71.51^{\pm 0.91 & -4.49 ^{\pm 2.83 & 0.082 & 0.012  & 0.056 & 16.12   & 963.40   & 10.11  & 865.42   \\
% % & SAM    & 68.86^{\pm 0.92 & 72.59^{\pm 0.39 & -3.51 ^{\pm 0.50 & 0.053 & 0.058  & 0.052 & 3.72    & 712.88   & 4.35   & 526.49   \\
% % &GAM  \\
% % \midrule
% % \multirow{2}{*}{\textbf{OLoRA}}
% % & SGD    & 66.48^{\pm 0.37 & 71.48^{\pm 0.55 & -3.55 ^{\pm 0.52 & 0.078 & 0.020  & 0.058 & 6.94    & 938.93   & 9.00   & 831.23   \\
% % & SAM    & 68.53^{\pm 0.78 & 72.75^{\pm 0.76 & -3.16 ^{\pm 0.68 & 0.066 & 0.026  & 0.048 & 4.40    & 775.40   & 7.41   & 628.12   \\
% % &GAM  \\
% % \bottomrule
% % \end{tabular}
% % }
% % \end{table}


%  % We set \( \rho=0.05 \) for both \( \mathrm{Sh}^{(0)}(\rho) \)  and \( \mathrm{Sh}^{(1)}(\rho) \) ; \( \mathrm{E\text{-}Sh}(Q) \) is estimated using 100 Gaussian perturbations \( \delta\sim\mathcal N(0,\sigma^2 I) \) with \( \sigma=\rho/\sqrt{d} \), where $d$ is the dimensionality of the parameter vector and table entries report \(\mathrm{E\text{-}Sh}(Q)\times10^{4}\) for readability.




% % \begin{figure}[htbp]
% % \centering
% % \begin{minipage}{0.9\linewidth}
% % \centering
% % \includegraphics[width=0.9\linewidth]{figures_TRML/2D_lossLandscape_task0_lora_opt_sam_vs_sgd_3d_compare_withxyz.pdf}
% % \caption{Visualization of the 3D loss landscape for a ViT backbone trained for incremental learning with LoRA adapters, comparing SGD and SAM optimizers. The SAM optimizer not only discovers a wider and flatter basin but also achieves a lower loss value, indicating a more robust and superior solution.}
% % \label{fig:3d_loss_landscape}
% % \end{minipage}

% % \end{figure}


% % {figures_TRML/AAA_curve_Imagenet-R_summary_inc10_rank16.pdf}

% %The performance gap typically widens on later tasks, indicating that SAM learns more stable representations with reduced forgetting.




% % The preceding theoretical analysis yields a clear prediction: in PECL, the stability-plasticity trade-off is governed primarily by the flatness of the adapter subspace. 



% % These CL gains coincide with systematic reductions in sharpness. Table~\ref{table:q1_ci_with_weight_flatness_reorg} reports three definitions of sharpness, together with spectral statistics of empirical Hessian and Fisher. Across all LoRA organizations, SAM lowers the sharpness metrics and shrinks the dominant eigenvalues and traces of both curvature matrices. The 3D loss landscape visualization in figure~\ref{fig:3d_loss_landscape} aligns with it.  


% % Taken together, these phenomena show that, in PECL with a frozen backbone, flatness-aware optimization confined to the LoRA adapter subspace is already sufficient to improve the stability–plasticity trade-off. The LoRA-restricted sharpness term predicted by our PAC–Bayesian analysis decreases, and the standard CL metrics improve in a consistent way. 


%  % (zeroth-order sharpness $\mathrm{Sh}^{(0)}$, expected sharpness $\mathrm{E\text{-}Sh}(Q)$, and first-order sharpness $\mathrm{Sh}^{(1)}$)




% % Across all LoRA configurations, SAM-optimized adapters consistently achieve higher final accuracy $\mathrm{ACC}_{20}$ and significantly mitigate catastrophic forgetting, as reflected by less negative backward transfer $\mathrm{BWT}$. For instance, in the challenging SeqLoRA setting, SAM elevates final accuracy from $63.56$ to $67.04$ while reducing negative backward transfer from $-11.55$ to $-9.26$. This advantage is not transient. The  and NME curves in Figure \ref{fig:combined_curves_imageR} show that SAM's performance gain is sustained throughout the task sequence and often widens on later tasks, which indicates that flatter adapter solutions accumulate less interference and exhibit greater long term stability.


% % Importantly, the consistent improvement in NME accuracy, which evaluates feature quality independently of the classifier head, suggests that the gains cannot be attributed solely to a better tuned linear probe. Rather, they point to a systematic enhancement of the learned representations. A detailed analysis of this representation-level mechanism is deferred to Section~4. Here we focus on the phenomenological conclusion: flatness-aware optimization restricted to the adapter subspace is sufficient to improve standard continual learning metrics, even when the backbone and other training conditions are held fixed.


% % Our empirical results first demonstrate that optimizing for adapter subspace flatness directly translates to superior continual learning performance. As shown in Table \ref{table:q1_ci_with_weight_flatness_reorg}, SAM-optimized adapters consistently achieve higher final accuracy ($\text{ACC}_{20}$) and significantly mitigate catastrophic forgetting, evidenced by improved backward transfer (BWT). For instance, in the challenging SeqLoRA setting, SAM elevates final accuracy from $63.56$ to $67.04$ while reducing negative backward transfer from $-11.55$ to $-9.26$. This advantage is not static but evolves over time. The AAA and NME curves in Fig \ref{fig:combined_curves_imageR} reveal that SAM's performance gain is not only sustained but often widens during later tasks, indicating that flatter adapter solutions accumulate less interference and exhibit greater long-term stability. Crucially, the consistent improvement in NME accuracy—which assesses feature quality independently of the classifier head—confirms that the benefit stems from more robust and stable representations in the feature space, rather than merely from a better-tuned classification head. This points to an intrinsic enhancement in the feature learning process.


% % The performance gains are systematically corroborated by a suite of sharpness metrics and geometric visualizations. SAM solutions exhibit consistently lower values across all evaluated sharpness measures. This pattern of reduced sensitivity to parameter perturbations is further validated by spectral analysis of the Hessian ($\lambda_{\max}(H)$, $\text{tr}(H)$) and the empirical Fisher information matrix ($\lambda_{\max}(F)$, $\text{tr}(F)$). The consistent findings across these distinct but theoretically aligned definitions and proxies provide multi-faceted evidence for the formation of flatter minima. The 3D loss landscape visualization in Figure \ref{fig:3d_loss_landscape} offers direct geometric confirmation: compared to SGD, SAM discovers solutions residing in visibly wider and flatter basins, while concurrently achieving a lower loss value. This visual evidence solidifies the quantitative measurements, confirming that flatness-aware optimization within the adapter subspace effectively steers the solution toward more robust and generalizable regions of the loss landscape.


% % These performance gains are systematically corroborated by a suite of weight-space flatness diagnostics. SAM solutions exhibit consistently lower values across all evaluated sharpness measures. This reduced sensitivity to parameter perturbations is further supported by spectral analysis of the empirical Hessian and Fisher information matrices, where both the dominant eigenvalues $\lambda_{\max}$ and the traces are smaller under SAM. The 3D loss landscape visualization in Figure~\ref{fig:3d_loss_landscape} provides geometric confirmation: compared to SGD, SAM discovers solutions that reside in visibly wider and flatter basins while achieving a lower training loss. Taken together, these observations answer Q1 in the affirmative. Enforcing flat minima in the LoRA subspace alone already yields substantial improvements in the stability-plasticity balance. The remaining question, addressed in Section~4, is why such subspace flatness translates into more stable feature representations and robustness under distribution shift.



% % This finding reframes the principle of flatness in continual learning: inducing a flat minimum within the low-dimensional adapter suffices. However, a critical question remains regarding the underlying mechanism: why do flatter adapter solutions exhibit reduced interference? In Section 5, we transition from the weight space to the feature space to analyze how subspace flatness promotes more stable representations and mitigates task interference.



% % \section{From Weight-Space Flatness to Feature-Space Robustness: Theory and Validation}

% % The PAC-Bayesian bounds established in Section \ref{section PAC} demonstrate that subspace-restricted sharpness control suffices for generalization in PECL. However, a fundamental question still remains: why do flatter adapter solutions exhibit reduced catastrophic forgetting? To answer this, we now transition from the weight-space analysis to the feature-space perspective, examining how flatness in the LoRA subspace propagates to feature stability.
% %  Specifically, we ask: does the benefit stem merely from better-aligned classifier heads, or does flatness-enhancing training fundamentally stabilize the shared representations? Our analysis reveals the mechanism through feature-space robustness
 
 
% %  % Our analysis reveals the mechanism  through feature-space robustness analyse.
 
% %  % formally captured by the Fisher information matrix.


% %  \subsection{Bridging Weight-Space Flatness and Feature Sensitivity}

 
% % %  To understand how this parameter-space flatness translates to feature stability, we decompose the model into a shared representation extractor and task-specific classifier. Let $\phi(x;W) $ denote the feature representation, where $W \in \mathbb{R}^d$ includes both the frozen backbone and the trainable LoRA adapter in $\mathcal{W}$. The linear classifier $W_{\text{cls}}$ maps features to logits:
% % %  \[
% % %  f(x;W,W_{\mathrm{cls}}) = W_{\mathrm{cls}}^\top \phi(x;W),
% % %  \quad
% % %  p(y\mid x;W,W_{\mathrm{cls}})=\mathrm{softmax}(f(x;W,W_{\mathrm{cls}})).
% % %  \]
% % % This decomposition allows us to precisely trace how weight perturbations affect predictive distributions through their impact on features, which is the missing link between our PAC-Bayesian theory and empirical performance gains.



% % % LoRA updates act only on the low-rank adapter component inside $W$ with the task specific head $W_{\mathrm{cls}}$. We first ignore the aligned problem of classifer and foucus on the difference on the feature layer caused by the weight. To ensure the target prediction
% % % results have a comparable basis we fix the input $x$  to investigate how different training strategies designed to
% % % enhance flatness (e.g., SAM) alter feature representations and consequently predictive distributions.


% % % Then we fix the input $x$ to ensure the target prediction
% % % results have a comparable basis and further investigate how different training strategies designed to
% % % enhance flatness (e.g., SAM) alter feature representations and consequently predictive distributions.

% % % Then we fix the input $x$ and examine how parameter perturbations—specifically those confined to the LoRA subspace $\mathcal{W}\Delta$—alter feature representations and consequently predictive distributions.



% % % Averaging over inputs with the $S$,  we obtain the empirical Fisher information matrix 
% % % $
% % % F= \frac{1}{|S|}\sum_{x\in S} J_\phi(x)^\top E_\phi(x)  J_\phi(x)
% % % $
% % % which bridges feature-space robustness to parameter-space flatness, exactly corresponding to the curvature surrogate that governs the expected sharpness in our PAC-Bayesian bounds.



% % % \paragraph{Interpreting the new symbol and its curvature role}
% % % The matrix $E_\phi(x)$ quantifies the sensitivity of the predictive law to infinitesimal
% % % feature perturbations at $x$. Via the chain rule with the representation Jacobian
% % % $J_\phi(x)=\partial\phi(x;W)/\partial W$, the parameter-level pullback
% % % $F_W(x)=J_\phi(x)^\top E_\phi(x) J_\phi(x)$ measures how small parameter changes
% % % alter the predictive distribution. For $\Delta W$ sufficiently small,
% % % \[
% % % \mathrm{KL}\!\Big(p(\cdot\mid x;W+\Delta W)\,\big\|\,p(\cdot\mid x;W)\Big)
% % % \;\approx\;\tfrac12\,\Delta W^\top F_W(x)\,\Delta W,
% % % \]
% % % so $F_W(x)$ acts as a local, PSD curvature metric in probability space. Under
% % % cross-entropy with a softmax output, $F_W(x)$ coincides with the generalized Gauss-Newton
% % % (GGN) matrix and closely tracks the Hessian near well-specified minima.
% % % When updates are confined to the LoRA subspace $\mathcal W_\Delta=\mathrm{span}(U)$,
% % % writing $\Delta W=U\alpha$ yields the restricted metric
% % % $F_\Delta(x)=U^\top F_W(x)U$ and the quadratic form
% % % $\tfrac12\,\alpha^\top F_\Delta(x)\alpha$, i.e., a local curvature restricted to
% % % admissible directions. This is precisely the object controlled by our subspace
% % % expected sharpness and the PAC-Bayesian analysis that follows.

% % % To analyse flatness in the LoRA parameter subspace to stability of the learned features, we factor the model into a shared representation map and a linear classification head. 
% % To relate flatness in the LoRA parameter subspace to the stability of the learned features, we factor the model into a shared representation map and a linear classification head.
% % Let $\phi(x;W)\in\mathbb{R}^D$ denote the representation produced by the backbone with LoRA where $W=W_{\text{froz}}+\Delta W$ and $\Delta W\in\mathcal{W}_\Delta$, and the head $W_{\mathrm{cls}}\in\mathbb{R}^{D\times K}$ maps features to logits. 
% % \[
% % f(x;W,W_{\mathrm{cls}})=W_{\mathrm{cls}}^\top \phi(x;W),\qquad
% % p(y\mid x;W,W_{\mathrm{cls}})=\mathrm{softmax}\!\big(f(x;W,W_{\mathrm{cls}})\big).
% % \]
% % In PECL, LoRA updates act solely on the low-rank adapter component within $W$ and the task-specific classifier head $W_{\mathrm{cls}}$. We focus on the differences in the feature layer induced by weights. To ensure that the target prediction results have a comparable basis, we fix the input $x$ to investigate how different training strategies designed to enhance flatness alter feature representations and consequently affect predictive distributions.

% % % and ignore the alignment issue of the classifier

% % The KL divergence between the perturbed and unperturbed predictive laws gives the local quadratic approximation
% % \begin{equation}
% % \label{eq: pro-feature}
% %  \mathrm{KL}\Big(  \ p(y \mid  \phi + \Delta\phi; W_{\mathrm{cls}})  \Vert  p(y \mid \phi; W_{\mathrm{cls}}) \  \big)
% % =
% % \Delta\phi(x)^\top
% %  E_\phi(x)
% %  \Delta\phi(x) + \mathcal{O}(\|\Delta\phi(x)\|^3)
% % \end{equation}
% % where the {local feature matrix}
% %  $
% %  E_\phi(x)
% %  \triangleq
% %  \mathbb{E}_{y \sim p(\cdot \mid x; W, W_{\mathrm{cls}})}
% %  \Big[
% %  \nabla_\phi \log p(y \mid x; W, W_{\mathrm{cls}})
% %  \nabla_\phi \log p(y \mid x; W, W_{\mathrm{cls}})
% %  ^\top
% %  \Big]
% %  $
% %  % Since the representation shift is caused by parameter perturbation rather than linear classifier $W_{\mathrm{cls}}$, 
% % measures predictive sensitivity to feature perturbations ~\citep{EFC2024}.
% % For a linear classifier followed by a softmax, there is an analytic formulation
% % $  E_\phi(x)
% % =
% % W_{\mathrm{cls}}
% %  \Big(
% %  \mathrm{Diag}(p) - p p^\top
% %  \Big)
% %  W_{\mathrm{cls}}^\top.
% %  $
% %  It quantifies the alignment strength between the current feature representation and the discriminative subspace spanned by the linear classifier, and does so in a boundary-aware manner.




% % % Let $J_\phi(x)=\partial \phi(x;W)/\partial W$ and 

% % % $\frac{\partial \phi(x;W)}{\partial W}$

% % The chain rule gives
% %  $
% %  \nabla_W \log p(y \mid x; W, W_{\mathrm{cls}})={\frac{\partial \phi(x;W)}{\partial W}}^\top  \nabla_\phi \log p(y \mid x; W, W_{\mathrm{cls}}).
% %  $
% % For small parameter perturbations $\Delta W \in \mathcal{W}_\Delta$, the induced representation shift is $\Delta\phi(x) \approx \frac{\partial \phi(x;W)}{\partial W}\Delta W$, yielding :
% % \[
% % \mathrm{KL}\!\Big(p(\cdot\mid x;W+\Delta W)\,\big\|\,p(\cdot\mid x;W)\Big)
% % \;\approx \tfrac12\,\Delta W^\top{\frac{\partial \phi(x;W)}{\partial W}}^\top E_\phi(x){\frac{\partial \phi(x;W)}{\partial W}}\Delta W =\;\tfrac12\,\Delta W^\top\,F_W(x)\,\Delta W
% % \]
% % where $F_W(x)={\frac{\partial \phi(x;W)}{\partial W}}^\top E_\phi(x) {\frac{\partial \phi(x;W)}{\partial W}}$ is the \textbf{local parameter Fisher information matrix at $x$}.  $F_W(x)$ measures how
% % small parameter changes alter the predictive distribution and acts as a local, curvature metric in probability space. Under
% % cross-entropy with a softmax output, $F_W(x)$ coincides with the generalized Gauss-Newton
% % (GGN) matrix and closely tracks the Hessian near well-specified minima. Averaging over inputs with the $S$,  we obtain the empirical Fisher information matrix 
% % $
% % F= \frac{1}{|S|}\sum_{x\in S} F_W(x).
% % $

% % % PSD

% % % When updates are confined to the LoRA subspace $\mathcal W_\Delta=\mathrm{span}(U)$,
% % % writing $\Delta W=U\alpha$ yields the restricted metric
% % % $F_\Delta(x)=U^\top F_W(x)U$ and the quadratic form
% % % $\tfrac12\,\alpha^\top F_\Delta(x)\alpha$, i.e., a local curvature restricted to
% % % admissible directions. This is precisely the object controlled by our subspace
% % % expected sharpness and the PAC-Bayesian analysis that follows.




% % \subsection{Expected Sharpness and PAC-Bayesian Coupling in LoRA SubSpace}

% % % \paragraph{From feature sensitivity to LoRA-restricted curvature}
% % % We factor the model into a shared representation and a linear head:
% % % $f(x;W,W_{\mathrm{cls}})=W_{\mathrm{cls}}^\top\phi(x;W)$, 
% % % $p(y\mid x;W,W_{\mathrm{cls}})=\mathrm{softmax}(f)$, with $W=W_{t-1}+\Delta W$ and $\Delta W\in\mathcal W_\Delta$.
% % % Fix $x$ and let $E_\phi(x)$ be the feature-level Fisher. With $J_\phi(x)=\partial\phi(x;W)/\partial W$,
% % % the parameter-level (pullback) Fisher is
% % % \[
% % % F_W(x)\;=\;J_\phi(x)^\top\,E_\phi(x)\,J_\phi(x).
% % % \]
% % % For sufficiently small $\Delta W$,
% % % \[
% % % \mathrm{KL}\!\Big(p(\cdot\mid x;W+\Delta W)\,\big\|\,p(\cdot\mid x;W)\Big)
% % % \;\approx\;\tfrac12\,\Delta W^\top\,F_W(x)\,\Delta W,
% % % \]
% % % so $F_W(x)$ acts as a local, PSD curvature metric in probability space. Under cross-entropy with softmax,
% % % $F_W(x)$ coincides with the generalized Gauss-Newton and closely tracks the Hessian near well-specified minima.

% % % We analyse LoRA geometry from two complementary angles that answer different needs.
% % % \begin{itemize}
% % %     \item {Update-subspace view}: PAC-Bayesian quantities (e.g., expected sharpness) live on the set of {admissible updates}; this view makes explicit that we only need curvature along directions that can actually change in PECL.
% % %     \item {Factor-coordinate view}: optimisation acts on the trainable factors $(A_\ell,B_\ell)$; this view exposes how curvature is seen by the optimiser and clarifies invariances due to the non-uniqueness of the factorisation.
% % % \end{itemize}
% % % We therefore introduce both views and then prove their equivalence, which ensures that all subsequent conclusions are \emph{coordinate-free} and depend only on the LoRA subspace.

% % We analyse LoRA flatness through two complementary routes.
% % % \emph{Update-subspace view:} PAC-Bayes quantities depend only on directions that can actually change during PECL, so we should measure curvature within the admissible update subspace.
% % % \emph{Factor-coordinate view:} optimisation operates on the trainable factors $(A_\ell,B_\ell)$, and algorithmic choices (preconditioning, SAM-style steps, gauge invariance) are most transparent in the factor coordinates.
% % \begin{itemize}
% %     \item {Update-subspace view:} PAC-Bayes quantities depend only on directions that can actually change during PECL, so we should measure curvature within the admissible update subspace.
% %     \item {factor coordinate view:} Optimization operates on trainable factors $(A_\ell,B_\ell)$, and algorithmic choices are most transparent in factor coordinates.
% % \end{itemize}
% % We present both views and then show they are equivalent, which yields coordinate-free conclusions.


% % % (i) \emph{Update-subspace view}: PAC--Bayesian quantities (e.g., expected sharpness) live on the set of \emph{admissible updates}; this view makes explicit that we only need curvature along directions that can actually change in PECL.
% % % (ii) \emph{Factor-coordinate view}: optimisation acts on the trainable factors $(A_\ell,B_\ell)$; this view exposes how curvature is seen by the optimiser and clarifies invariances due to the non-uniqueness of the factorisation.


% % % \subsubsection{View I: update subspace}
% % \textbf{View I: update subspace}

% % % $\mathcal{W}_\Delta$ 

% % % Lemma \ref{lem:kl-decomposition} shows that the KL complexity is entirely borne by the adapter component. 


% % Let \(\mathcal W_\Delta=\mathrm{span}(U)\) with \(U\in\mathbb R^{d\times r}\) and write \(\Delta W=U\alpha\). For a fixed input \(x\), the quadratic KL approximation along any admissible perturbation $\Delta W \in \mathcal{W}_\Delta$ becomes
% %  \[
% %  \mathrm{KL}\big(p(\cdot\mid W+\Delta W)|p(\cdot\mid W)\big)
% %  \approx\tfrac12\Delta W^\top F_W(x)\Delta W
% %  =\tfrac12\alpha^\top\underbrace{U^\top F_W(x)U}_{\displaystyle F_\Delta(x)}\alpha.
% %  \]
% % The \textbf{local LoRA parameter Fisher}
% % $
% % F_\Delta(x):=U^\top F_W(x)U=U^\top {\frac{\partial \phi(x;W)}{\partial W}}^\top E_\phi(x){\frac{\partial \phi(x;W)}{\partial W}}U \in\ \mathbb R^{r\times r}  
% % $ act as the pullback of \(F_W(x)\) onto \(\mathcal W_\Delta\) and is exactly the curvature along admissible LoRA directions. It follows the fact that the KL complexity is entirely borne by the adapter component and aligns with Lemma \ref{lem:kl-decomposition}.

% % % If we block-partition \(F_W\) by \((\text{frozen},\Delta)\), then for feasible \(\delta_\theta=(0,\alpha)\),
% % %  \[
% % %  \delta_\theta^\top F_W\delta_\theta=\alpha^\top F_{\Delta}\alpha
% % %  \]

% %  % $F_W(x)={\frac{\partial \phi(x;W)}{\partial W}}^\top E_\phi(x) {\frac{\partial \phi(x;W)}{\partial W}}$


% % % Assume the task-$t$ posterior on $\mathcal W_\Delta$ is Gaussian, $Q_t^\Delta=\mathcal N(\hat{\Delta W}_t,\Sigma_t)$ with $\Sigma_t\in\mathbb R^{r\times r}$, then the subspace expected sharpness in our PAC-Bayes bound is
% % % \[
% % % \operatorname{E-Sh}_t^\Delta
% % % \;\approx\;\tfrac12\,\mathbb E_{x\in S_t}\!\Big[\mathrm{tr}\big(F_\Delta(x)\,\Sigma_t\big)\Big],
% % % \]
% % % this quantity inherits the classifier-feature alignment encoded by \(E_\phi(x)\) and retains only the component viewed through the LoRA-admissible directions. Consequently, controlling flatness {within} \(\mathcal W_\Delta\)is both necessary (only these directions can change in PECL) and sufficient (it directly reduces \(\operatorname{ESh}_t^\Delta\) and tightens the PAC-Bayes bound).




% % % namely the data-averaged inner product between the LoRA-restricted curvature $F_\Delta$ and the posterior covariance $\Sigma_t$.









% % % Since the posterior places a Dirac mass on the frozen block and all admissible perturbations satisfy \(\Delta W\in\mathcal W_\Delta\), the relevant curvature is the pullback of the parameter-level Fisher onto the LoRA update subspace. Concretely, letting \(\mathcal W_\Delta=\mathrm{span}(U)\) and \(\Delta W=U\alpha\), we define the local LoRA parameter Fisher by*\
% % %  \[
% % %  F_\Delta(x);:=;U^\top F_W(x),U,\qquad
% % %  \mathrm{KL}\ \approx\ \tfrac12,\alpha^\top F_\Delta(x),\alpha.
% % %  \]







% % % Let $\mathcal W_\Delta=\mathrm{span}(U)$ be the LoRA update subspace with $U\in\mathbb R^{d\times r}$ and write $\Delta W=U\alpha$.
% % % Pulling back the parameter-level Fisher $F_W(x)$ onto $\mathcal W_\Delta$ gives the \textbf{local LoRA parameter Fisher}
% % % \[
% % % F_\Delta(x)\;:=\;U^\top F_W(x)\,U\ \in\ \mathbb R^{r\times r},
% % % \qquad
% % % \mathrm{KL}\ \approx\ \tfrac12\,\alpha^\top F_\Delta(x)\,\alpha,
% % % \]
% % % which is exactly the curvature along admissible LoRA directions. This means that perturbations of features are first seen through the classifier geometry and then restricted to the directions allowed by LoRA.














% % % % \paragraph{LoRA subspace Fisher}
% % % Let $\mathcal W_\Delta=\mathrm{span}(U)$ denote the LoRA update subspace (columns of $U\in\mathbb R^{d\times r}$).
% % % Writing $\Delta W=U\alpha$ yields the \textbf{local LoRA parameter Fisher}
% % % \[
% % % F_\Delta(x)\;:=\;U^\top\,F_W(x)\,U\ \in\ \mathbb R^{r\times r},
% % % \qquad
% % % \mathrm{KL}\ \approx\ \tfrac12\,\alpha^\top F_\Delta(x)\,\alpha,
% % % \]
% % % which is precisely the curvature along admissible (LoRA) directions. This two-step pullback shows that the curvature we control in PECL is not arbitrary curvature of the full model, but exactly the curvature that arises when feature perturbations are projected first onto the classifier space and then onto the admissible LoRA directions.
% % % A Gaussian posterior $Q_t^\Delta=\mathcal N(\hat{\Delta W}_t,\Sigma_t)$ supported on $\mathcal W_\Delta$ satisfies
% % % \[
% % % \operatorname{E-Sh}_t^\Delta
% % % \;\approx\;\tfrac12\,\mathbb E_{x\in S_t}\!\Big[\mathrm{tr}\!\big(F_\Delta(x)\,\Sigma_t\big)\Big],
% % % \]
% % % so the subspace expected sharpness used in our PAC-Bayesian bound is the data-averaged inner product between $F_\Delta$ and the posterior covariance $\Sigma_t$.

% % % \subsubsection{View II: factor coordinates}
% % \textbf{View II: factor coordinates}

% % % \(\theta\)


% % For an adapted layer $\ell$ with $W_\ell\in\mathbb R^{d_o\times d_i}$, LoRA constrains the update to $\Delta W_\ell=B_\ell A_\ell$ with $B_\ell\in\mathbb R^{d_o\times r_\ell}$ and $A_\ell\in\mathbb R^{r_\ell\times d_i}$.
% % Fix current estimates $(\hat B_\ell,\hat A_\ell)$ and consider first-order perturbations $(\delta B_\ell,\delta A_\ell)$.
% % By the product rule,
% % $
% % \delta(\Delta W_\ell)\;=\;\hat B_\ell\,\delta A_\ell\;+\;\delta B_\ell\,\hat A_\ell.
% % $
% % Vectorization with $\mathrm{vec}(U D V)=(V^\top\!\otimes U)\mathrm{vec}(D)$ yields the linear relation
% % \[
% % \mathrm{vec}\big(\delta(\Delta W_\ell)\big)
% % \;=\;
% % \underbrace{(I_{d_i^{(\ell)}}\otimes \hat B_\ell)}_{J_{A,\ell}}\mathrm{vec}(\delta A_\ell)
% % \;+\;
% % \underbrace{(\hat A_\ell^\top\otimes I_{d_o^{(\ell)}})}_{J_{B,\ell}}\mathrm{vec}(\delta B_\ell).
% % \]






% % % Linearising around $(\hat B_\ell,\hat A_\ell)$ yields
% % % \[
% % % \delta_{(\Delta W_\ell)}=\hat B_\ell\,\delta_{A_\ell}+\delta_{B_\ell}\,\hat A_\ell,
% % % \quad
% % % \mathrm{vec}\big(\delta(\Delta W_\ell)\big)
% % % =(I_{d_i}\otimes \hat B_\ell)\mathrm{vec}(\delta A_\ell)
% % % +(\hat A_\ell^\top\otimes I_{d_o})\mathrm{vec}(\delta B_\ell).
% % % \]


% % Let $\Pi_\ell$ embed $\mathrm{vec}(\Delta W_\ell)$ into the global parameter vector $\mathrm{vec}(\Delta W)$.
% % Define the factor coordinate
% % $
% % \theta := \big[\mathrm{vec}(\delta A_1);\mathrm{vec}(\delta B_1);\ldots;\mathrm{vec}(\delta A_L);\mathrm{vec}(\delta B_L)\big] \in \mathbb R^{q}
% % $, where
% % $
% % q:=\sum_{\ell=1}^L r_\ell\big(d_i^{(\ell)}+d_o^{(\ell)}\big).
% % $
% % The global LoRA Jacobian $J_{\mathrm{LoRA}}\in\mathbb R^{d\times q}$ with $d:=\sum_\ell d_o^{(\ell)}d_i^{(\ell)}$ is the block concatenation $
% % J_{\mathrm{LoRA}} :=
% % \big[\ \Pi_1J_{A,1}\ \ \Pi_1J_{B,1}\ \ \cdots\ \ \Pi_LJ_{A,L}\ \ \Pi_LJ_{B,L}\ \big]
% % $, and $
% % \mathrm{vec}(\Delta W)\ \approx\ J_{\mathrm{LoRA}}\theta.
% % $
% % By construction $\mathrm{Im}(J_{\mathrm{LoRA}})$ equals the LoRA update subspace $\mathcal W_\Delta\subset\mathbb R^d$ reachable by factor perturbations at $(\hat B,\hat A)$. When the intended LoRA directions are indeed spanned by the columns of $\hat B_\ell$ and the rows of $\hat A_\ell$, $J_{\mathrm{LoRA}}$ is full column rank and thus provides a linear isomorphism between {factor coordinates} and {admissible parameter updates}.


% % % Stacking all layers, define the factor coordinate
% % % \[
% % % \theta=[\mathrm{vec}(\delta A_1);\mathrm{vec}(\delta B_1);\ldots;\mathrm{vec}(\delta A_L);\mathrm{vec}(\delta B_L)]\in\mathbb R^{q},
% % % \quad
% % % q=\sum_{\ell=1}^L r_\ell\big(d_i^{(\ell)}+d_o^{(\ell)}\big),
% % % \]
% % % and the global LoRA Jacobian $J_{\mathrm{LoRA}}\in\mathbb R^{d\times q}$, so that
% % % \[
% % % \mathrm{vec}(\Delta W)\ \approx\ J_{\mathrm{LoRA}}\theta,
% % % \qquad
% % % \mathrm{Im}(J_{\mathrm{LoRA}})=\mathcal W_\Delta.
% % % \]
% % % The curvature experienced by the optimiser in factor space is
% % % $
% % % F_\theta(x):=J_{\mathrm{LoRA}}^\top\,F_W(x)\,J_{\mathrm{LoRA}}\ \in\ \mathbb R^{q\times q}.
% % % $
% % %  For small factor perturbations,
% % % \[
% % % \mathrm{KL}\!\Big(p(\cdot\mid x;W+\mathrm{unvec}(J_{\mathrm{LoRA}}\theta))\ \big\|\ p(\cdot\mid x;W)\Big)
% % % \;\approx\;\tfrac12\,\theta^\top F_\theta(x)\,\theta.
% % % \]
% % % Thus the factor-space curvature $J_{\mathrm{LoRA}}^\top F_W J_{\mathrm{LoRA}}$ is simply a coordinate representation of the LoRA-restricted Fisher $F_\Delta$.
% % % Spectral quantities such as $\lambda_{\max}$ and $\mathrm{tr}$ coincide up to a change of basis. Hence spectral control in factor space is exactly the same as control of $F_\Delta$ in the update subspace.

% % % Let the task-$t$ posterior in factor space be Gaussian, $\theta\sim\mathcal N(\hat\theta_t,\Sigma_{\theta,t})$, which induces the parameter covariance
% % % \[
% % % \Sigma_t\;=\;J_{\mathrm{LoRA}}\Sigma_{\theta,t}J_{\mathrm{LoRA}}^\top
% % % \;\in\;\mathbb R^{d\times d},\qquad
% % % \text{with support in }\ \mathcal W_\Delta.
% % % \]
% % % The subspace expected sharpness appearing in our PAC--Bayes bound admits the factor-space representation
% % % \[
% % % \operatorname{E-Sh}_t^\Delta
% % % \;\approx \tfrac12\,\mathrm{tr}\!\big(\frac{1}{|S_t|}\sum_{x\in S_t}F_\theta(x),\Sigma_{\theta,t}\big).
% % % \]

% % % \subsubsection{Equivalence and consequences}
% % \textbf{Equivalence and consequences}

% % \begin{lemma}[Factor--subspace curvature equivalence]
% % \label{lem:lora-factor-subspace}
% % If $\mathrm{Im}(J_{\mathrm{LoRA}})=\mathcal W_\Delta$, then there exists an invertible matrix $R$ such that
% % \[
% % F_\theta(x)\;=\;J_{\mathrm{LoRA}}^\top F_W(x) J_{\mathrm{LoRA}}
% % \;=\;R^\top \big(U^\top F_W(x)U\big) R
% % \;=\;R^\top F_\Delta(x)\,R\quad\text{for all }x.
% % \]
% % In particular, if the columns of $J_{\mathrm{LoRA}}$ form the basis $U$, then $F_\theta(x)=F_\Delta(x)$.
% % \end{lemma}
% % % \begin{proof}[Proof sketch]
% % % Because $\mathrm{Im}(J_{\mathrm{LoRA}})=\mathcal W_\Delta=\mathrm{span}(U)$, there exists invertible $R$ with $J_{\mathrm{LoRA}}=UR$. Substitute $J_{\mathrm{LoRA}}=UR$ into $F_\theta(x)$ to obtain $F_\theta(x)=R^\top U^\top F_W(x)U R$.
% % % \end{proof}
% % % While the individual eigenvalues of $F_\theta$ may change, 

% % The LoRA factorization admits gauge transformations $(B_\ell,A_\ell)\mapsto(B_{\ell} R_\ell,R_\ell^{-1}A_\ell)$ with invertible $R_\ell$,which induce an invertible change of basis $\theta\mapsto R\theta$ on the factor coordinates. Under this transformation, 
% % $
% % F_\theta\ \mapsto\ R^\top F_\theta R
% % $
% % , and $
% % \Sigma_{\theta,t}\ \mapsto\ R^{-1}\Sigma_{\theta,t}(R^{-1})^\top.
% % $
% % These congruences leave the quadratic form
% % $
% % \operatorname{tr}(F_\theta \Sigma_{\theta,t})
% % =
% % \mathbb{E}_{\varepsilon \sim \mathcal{N}(0,\Sigma_{\theta,t})}
% % \bigl[\varepsilon^\top F_\theta \varepsilon\bigr]
% % $
% % remains invariant. Hence, the sharpness term that enters our PAC-Bayesian bound is a coordinate-free property of the LoRA subspace geometry, independent of the particular factorization of the adapters. Consequently, algorithmic interventions in the factor coordinates $\theta$ can be interpreted as geometric regularization in the adapter update subspace $\mathcal{W}_\Delta$. 


% % In summary, subspace flatness reduces the dominant eigenvalues of the parameter-space Fisher, which in turn suppresses the feature-space Fisher and the empirical feature matrix. This chain explains why lower sharpness is consistently accompanied by more isotropic and stable features in our diagnostics.




% % % In the remainder of this section, we therefore focus on how such subspace-level spectral control translates into reduced expected sharpness and more stable feature representations. 





% % % Crucially, the energy term $\mathrm{tr}( F_\theta\Sigma_{\theta})$ remains invariant under these transformations. This implies that both the curvature spectrum and the sharpness term are \textbf{coordinate-free} intrinsic properties of the LoRA subspace geometry. Consequently, algorithmic interventions in the factor coordinates $\theta$ (e.g., via SAM) directly translate to geometric regularizations in the adapter update subspace $\mathcal{W}_\Delta$, mathematically validating the sufficiency of subspace-restricted optimization.





% % \subsection{From Spectral Control to Feature Stability: Mechanism and Verification}

% % \subsubsection{Theoretical Mechanism: Spectral Suppression Tightens the Bound.}

% % The mathematical equivalence established above bridges the gap between optimization algorithms and generalization theory. We now analyze how this mechanism tightens the PAC-Bayes bound and empirically validate its impact on feature stability.

% % Our PAC-Bayesian analysis decomposes the generalization error into empirical risk and a subspace-restricted sharpness term. Considering task-wise Gaussian posteriors $Q_t^\Delta = \mathcal{N}(\hat{\Delta W}_t, \Sigma_t)$, the expected sharpness is given by
% % $
% % \operatorname{E-Sh}_t^\Delta \approx \tfrac12 \mathbb E_{x\in S_t}\left[\mathrm{tr}\left(F_\Delta(x)\Sigma_t\right)\right].
% % $
% % By the congruence of \(F_\theta\) and \(F_\Delta\) (Lemma~\ref{lem:lora-factor-subspace}), and applying PSD trace inequalities, we derive:
% % \begin{equation}
% % \mathrm{tr}\left(F_\Delta(x)\Sigma_t\right) \le \lambda_{\max}\left(F_\Delta(x)\right) \mathrm{tr}(\Sigma_t).
% % \end{equation}
% % This inequality reveals that generalization is governed by the spectral norm of the curvature matrix. SAM updates, which iteratively seek the worst-case perturbation, align with the leading eigenvector of the local curvature in $\theta$, thereby effectively minimizing $\lambda_{\max}(F_\theta)$ and, by equivalence, $\lambda_{\max}(F_\Delta)$.

% % This spectral contraction serves two critical purposes in PECL:
% % \begin{enumerate}
% %     \item \textbf{Tightening the Bound:} It directly reduces the expected sharpness term $\operatorname{E-Sh}_t^\Delta$ for a fixed posterior variance. $\mathrm{E\text{-}Sh}_t^\Delta$ is exactly the subspace sharpness term that appears in the right-hand side of the bound in Theorem~2, reducing $\lambda_{\max}(F_\Delta)$ tightens the PAC-Bayesian generalization guarantee for a fixed posterior variance. 
% %     \item \textbf{Stabilizing Features:} Since $F_\Delta = U^\top J_\phi^\top E_\phi J_\phi U$,suppressing the spectrum of $F_\Delta$ limits the sensitivity of the representation $\phi(x)$ to parameter updates along admissible LoRA directions. This restricts how far new low-rank updates can move previously learned decision boundaries and thereby mitigates catastrophic forgetting.
% % \end{enumerate}



% % % The curvature matrices \(F_\theta\) and \(F_\Delta\) are congruent (Lemma~\ref{lem:lora-factor-subspace}). Consequently, spectral regularization in the factor space—such as diminishing \(\lambda_{\max}\) or \(\mathrm{tr}(\cdot)\) via SAM,  is exactly equivalent to controlling the LoRA-restricted curvature \(F_\Delta\). Together with a Gaussian posterior on the adapter subspace $\mathcal W_\Delta$, this equivalence implies that flatness regularization confined to $\mathcal W_\Delta$ tightens the expected sharpness term in the PAC Bayes bound and improves stability and generalization in PECL, independently of the specific LoRA parameterization. Controlling flatness within $\mathcal W_\Delta$ suffices to govern the effective curvature encountered during training, which stabilizes shared features along the chain $  F_{\Delta} \to F_W \to E_\phi$.


% % % Consider task-wise Gaussian posteriors over the LoRA subspace, defined as \(Q_t^\Delta = \mathcal{N}(\hat{\Delta W}_t, \Sigma_t)\). The sharpness contribution in our bound is given by:
% % % $
% % % \operatorname{E-Sh}_t^\Delta \approx \tfrac12 \mathbb E_{x\in S_t}\left[\mathrm{tr}\left(F_\Delta(x)\Sigma_t\right)\right],
% % % $ where \(F_\Delta(x) = U^\top J_\phi(x)^\top E_\phi(x) J_\phi(x)U\). Increased flatness acts upon this term and the KL complexities in two distinct ways. By the congruence of \(F_\theta\) and \(F_\Delta\), and applying PSD trace inequalities, we derive:
% % % $$
% % % \mathrm{tr}\left(F_\Delta(x)\Sigma_t\right) \le \lambda_{\max}\left(F_\Delta(x)\right) \mathrm{tr}(\Sigma_t).
% % % $$
% % % % \quad\text{and}\quad
% % % % \mathrm{tr}\left(F_\Delta(x)\Sigma_t\right) \le \mathrm{tr}\left(F_\Delta(x)\right) \lambda_{\max}(\Sigma_t)
% % % Therefore, any method that reduces \(\lambda_{\max}(F_\theta)\) or \(\mathrm{tr}(F_\theta)\) tightens the bound on \(\operatorname{E-Sh}_t^\Delta\) for a fixed \(\Sigma_t\). Furthermore, flatter subspaces enable the same empirical risk to be achieved with reduced posterior variance and smaller mean shifts, thereby diminishing the KL divergence \(\mathrm{KL}(Q_t^\Delta \mid P_t^\Delta)\).

% % % In summary, At each step, SAM seeks a worst-case perturbation of the loss within \(\mathcal{W}_\Delta\) and subsequently updates the parameters against it. Under a local quadratic approximation, the inner maximizer aligns with the top eigenvector of \(F_\theta\), and hence of \(F_\Delta\). Iterating this mechanism suppresses \(\lambda_{\max}(F_\Delta)\). The suppression of \(\lambda_{\max}\) directly controls the most sensitive admissible direction, while improved conditioning redistributes curvature from a few sharp axes to many flat ones. Since \(F_\Delta = U^\top J_\phi^\top E_\phi J_\phi U\) accentuates boundary-dominated, classifier-aligned directions, this spectral contraction limits the extent to which new low-rank updates can distort previously learned margins. The result is a reduction in \(\operatorname{E-Sh}_t^\Delta\), lower KL penalties for a comparable empirical fit, and diminished interference across tasks, all while preserving plasticity through optimization within \(\mathcal{W}_\Delta\).




% % % The LoRA factorization admits gauge transformations $(B_\ell,A_\ell)\mapsto(B_\ellR_\ell,R_\ell^{-1}A_\ell)$ with invertible $R_\ell$, which act on $\theta$ by an invertible change of basis $\theta\mapsto R\theta$. Under this transformation
% % % \[
% % % F_\theta\ \mapsto\ R^\top F_\theta R,
% % % \qquad
% % % \Sigma_{\theta,t}\ \mapsto\ R^{-1}\Sigma_{\theta,t}(R^{-1})^\top,
% % % \]
% % % so the energy $\mathrm{tr}(\bar F_\theta\Sigma_{\theta,t})$ is invariant. Hence both the curvature and the sharpness term are \emph{coordinate-free} properties of the LoRA subspace geometry.





% % % \begin{lemma}[Equivalence of factor and subspace curvature]
% % % \label{lem:coord-equiv}
% % % If $\mathrm{Im}(J_{\mathrm{LoRA}})=\mathcal W_\Delta$, then for every $x$ there exists an invertible $R$ such that
% % % \[
% % % J_{\mathrm{LoRA}}^\top F_W(x) J_{\mathrm{LoRA}}
% % % \;=\;R^\top\!\big(U^\top F_W(x)U\big)R
% % % \;=\;R^\top F_\Delta(x)R.
% % % \]
% % % In particular, if the columns of $J_{\mathrm{LoRA}}$ form the basis $U$, then $F_\theta(x)=F_\Delta(x)$.
% % % \end{lemma}

% % % \noindent


% % % With dataset averages $\bar F_\Delta:=|S_t|^{-1}\sum_{x\in S_t}F_\Delta(x)$ and a factor-space posterior $\theta\sim\mathcal N(\hat\theta_t,\Sigma_{\theta,t})$ satisfying $J_{\mathrm{LoRA}}\Sigma_{\theta,t}J_{\mathrm{LoRA}}^\top=U\Sigma_tU^\top$, the PAC-Bayes coupling used in Theorem~\ref{thm:pecl-bound} can be written as
% % % \[
% % % \mathbb E_{\varepsilon\sim\mathcal N(0,\Sigma_t)}\!\left[\hat L_{S_t}(W_{t-1}+\hat{\Delta W}_t+\varepsilon)-\hat L_{S_t}(W_{t-1}+\hat{\Delta W}_t)\right]
% % % \;\approx\;\tfrac12\,\big\langle \bar F_\Delta,\ \Sigma_t\big\rangle
% % % \;=\;\tfrac12\,\big\langle \bar F_\theta,\ \Sigma_{\theta,t}\big\rangle.
% % % \]


% % % For an adapted layer $\ell$ with weight $W_\ell\in\mathbb R^{d_o\times d_i}$, LoRA constrains the update to a rank-$r_\ell$ factorization
% % % $\Delta W_\ell = B_\ell A_\ell$ with $B_\ell\in\mathbb R^{d_o\times r_\ell}$ and $A_\ell\in\mathbb R^{r_\ell\times d_i}$.
% % % Linearizing around $(\hat B_\ell,\hat A_\ell)$ yields
% % % \[
% % % \delta(\Delta W_\ell)=\hat B_\ell\,\delta A_\ell+\delta B_\ell\,\hat A_\ell.
% % % \]
% % % Vectorization gives
% % % \[
% % % \mathrm{vec}\big(\delta(\Delta W_\ell)\big)
% % % =\underbrace{(I_{d_i}\otimes \hat B_\ell)}_{J_{A,\ell}}\mathrm{vec}(\delta A_\ell)
% % % +\underbrace{(\hat A_\ell^\top\otimes I_{d_o})}_{J_{B,\ell}}\mathrm{vec}(\delta B_\ell).
% % % \]
% % % Let $\Pi_\ell$ place $\mathrm{vec}(\Delta W_\ell)$ into the global parameter vector $\mathrm{vec}(\Delta W)$ and define the factor coordinate
% % % \[
% % % \theta\ :=\ [\mathrm{vec}(\delta A_1);\mathrm{vec}(\delta B_1);\ldots;\mathrm{vec}(\delta A_L);\mathrm{vec}(\delta B_L)]
% % % \ \in\ \mathbb R^{q},\qquad
% % % q=\sum_{\ell=1}^L r_\ell\big(d_i^{(\ell)}+d_o^{(\ell)}\big).
% % % \]
% % % The global LoRA Jacobian $J_{\mathrm{LoRA}}\in\mathbb R^{d\times q}$ is
% % % \[
% % % J_{\mathrm{LoRA}}=\big[\ \Pi_1J_{A,1}\ \ \Pi_1J_{B,1}\ \ \cdots\ \ \Pi_LJ_{A,L}\ \ \Pi_LJ_{B,L}\ \big],
% % % \qquad
% % % \mathrm{vec}(\Delta W)\ \approx\ J_{\mathrm{LoRA}}\theta,
% % % \]
% % % with image $\mathrm{Im}(J_{\mathrm{LoRA}})=\mathcal W_\Delta$. When the intended LoRA directions are spanned by the columns of $\hat B_\ell$ and rows of $\hat A_\ell$, $J_{\mathrm{LoRA}}$ is full column rank and provides a linear isomorphism between factor coordinates and admissible updates.












% % % If $U\in\mathbb R^{d\times r}$ is a basis of $\mathcal W_\Delta$, the LoRA-restricted Fisher is
% % % \[
% % % F_\Delta(x):=U^\top F_W(x)U\in\mathbb R^{r\times r},\qquad
% % % \mathrm{KL}\ \approx\ \tfrac12\,\alpha^\top F_\Delta(x)\alpha\quad(\Delta W=U\alpha).
% % % \]
% % % The curvature that the optimizer perceives in factor coordinates is
% % % \[
% % % F_\theta(x):=J_{\mathrm{LoRA}}^\top F_W(x)J_{\mathrm{LoRA}}\in\mathbb R^{q\times q}.
% % % \]



% % % LoRA imposes a factorized, layer-wise update structure: for each weight matrix $W_\ell \in \mathbb{R}^{d_o \times d_i}$, the update takes the form $\Delta W_\ell = B_\ell A_\ell$ with $B_\ell \in \mathbb{R}^{d_o \times r}$ and $A_\ell \in \mathbb{R}^{r \times d_i}$. Around current estimates $(\hat{B}_\ell, \hat{A}_\ell)$, the first-order variation satisfies:
% % % $
% % % \delta(\Delta W_\ell) = \hat{B}_\ell \delta A_\ell + \delta B_\ell \hat{A}_\ell.
% % % $
% % % In vectorized form:
% % % $
% % % \mathrm{vec}(\delta(\Delta W_\ell)) = (I_{d_i} \otimes \hat{B}_\ell)\mathrm{vec}(\delta A_\ell) + (\hat{A}_\ell^\top \otimes I_{d_o})\mathrm{vec}(\delta B_\ell).
% % % $
% % % Stacking across layers yields a global linear embedding from LoRA factor coordinates $\theta = [\mathrm{vec}(\delta A_1); \mathrm{vec}(\delta B_1); \dots]$ to the adapter subspace $\mathcal{W}_\Delta$:
% % % \[
% % % \mathrm{vec}(\Delta W) \approx J_{\mathrm{LoRA}} \theta,
% % % \]
% % % where $J_{\mathrm{LoRA}}$ is block-structured and full column rank, establishing an isomorphism between the factor space and admissible directions in $\mathcal{W}_\Delta$.

% % % % \paragraph{LoRA factor coordinates with $F_\Delta$}

% % % LoRA uses layerwise factors $\Delta W_\ell=B_\ell A_\ell$. Linearizing around $(\hat B_\ell,\hat A_\ell)$ gives
% % % $\mathrm{vec}(\delta\Delta W)\approx J_{\mathrm{LoRA}}\theta$, where $\theta$ stacks $\{\mathrm{vec}(\delta A_\ell),\mathrm{vec}(\delta B_\ell)\}$ and $J_{\mathrm{LoRA}}$ is block-structured and full column rank.

% % % \begin{lemma}[Coordinate equivalence of factor and subspace curvature]
% % % \label{lem:coord-equiv}
% % % If $\mathrm{Im}(J_{\mathrm{LoRA}})=\mathcal W_\Delta$, then for every $x$ there exists an invertible change of basis $R$
% % % such that
% % % \[
% % % J_{\mathrm{LoRA}}^\top\,F_W(x)\,J_{\mathrm{LoRA}}
% % % \;=\;R^\top\,F_\Delta(x)\,R.
% % % \]
% % % In particular, if the columns of $J_{\mathrm{LoRA}}$ form the basis $U$, then
% % % $J_{\mathrm{LoRA}}^\top F_W(x)J_{\mathrm{LoRA}}=F_\Delta(x)$.
% % % \end{lemma}

% % % \noindent
% % % Thus the “factor-space curvature” $G_{\mathrm{fact}}:=J_{\mathrm{LoRA}}^\top F_W J_{\mathrm{LoRA}}$ is nothing but a coordinate representation of the LoRA-restricted Fisher $F_\Delta$.
% % % Its spectrum (e.g., $\lambda_{\max}$, $\mathrm{tr}$) coincides with that of $F_\Delta$ up to a basis change, so optimizing flatness \emph{in the LoRA subspace} controls the effective curvature experienced by training.


% % % Combining this with the Gaussian posterior on $\mathcal W_\Delta$ yields the PAC-Bayesian coupling used in Theorem~\ref{thm:pecl-bound}:
% % % \[
% % % \mathbb E_{\varepsilon\sim\mathcal N(0,\Sigma_t)}
% % % \!\left[\hat L_{S_t}\big(W_{t-1}+\hat{\Delta W}_t+\varepsilon\big)-\hat L_{S_t}\big(W_{t-1}+\hat{\Delta W}_t\big)\right]
% % % \;\approx\;\tfrac12\,\Big\langle \ \bar F_\Delta\ ,\ \Sigma_t\ \Big\rangle,
% % % \quad
% % % \bar F_\Delta:=\frac1{|S_t|}\sum_{x\in S_t}F_\Delta(x).
% % % \]
% % % This explicitly ties the empirical sharpness penalty to the LoRA-restricted Fisher introduced above, ensuring that all subsequent analyses only invoke $F_W$ and $F_\Delta$.






% % % Building on the Gaussian posterior formulation $Q_t^{\Delta} = \mathcal{N}(\hat{\Delta W}_t, \Sigma_t)$ on $\mathcal{W}_\Delta$ from Section 3.3.2, we now parameterize the adapter subspace through its factorized LoRA structure. This refined coordinate system allows us to explicitly trace how parameter perturbations propagate through the model's geometric hierarchy.


% % % LoRA imposes a factorized, layer-wise update structure: for each weight matrix $W_\ell \in \mathbb{R}^{d_o \times d_i}$, the update takes the form $\Delta W_\ell = B_\ell A_\ell$ with $B_\ell \in \mathbb{R}^{d_o \times r}$ and $A_\ell \in \mathbb{R}^{r \times d_i}$. Around current estimates $(\hat{B}_\ell, \hat{A}_\ell)$, the first-order variation satisfies:
% % % $
% % % \delta(\Delta W_\ell) = \hat{B}_\ell \delta A_\ell + \delta B_\ell \hat{A}_\ell.
% % % $
% % % In vectorized form:
% % % $
% % % \mathrm{vec}(\delta(\Delta W_\ell)) = (I_{d_i} \otimes \hat{B}_\ell)\mathrm{vec}(\delta A_\ell) + (\hat{A}_\ell^\top \otimes I_{d_o})\mathrm{vec}(\delta B_\ell).
% % % $
% % % Stacking across layers yields a global linear embedding from LoRA factor coordinates $\theta = [\mathrm{vec}(\delta A_1); \mathrm{vec}(\delta B_1); \dots]$ to the adapter subspace $\mathcal{W}_\Delta$:
% % % \[
% % % \mathrm{vec}(\Delta W) \approx J_{\mathrm{LoRA}} \theta,
% % % \]
% % % where $J_{\mathrm{LoRA}}$ is block-structured and full column rank, establishing an isomorphism between the factor space and admissible directions in $\mathcal{W}_\Delta$.


% % % Substituting this chart into the parameter Fisher gives the \textbf{projected Fisher in LoRA factor space}
% % % \begin{equation}
% % % \label{eq:lora-fisher}
% % % \mathrm{KL}\!\left(p(\cdot\mid W+\mathrm{unvec}(J_{\mathrm{LoRA}}\theta))\,\big\|\,p(\cdot\mid W)\right)
% % % \;=\; \theta^\top G_{\mathrm{fact}}\,\theta,
% % % \qquad
% % % G_{\mathrm{fact}}:=J_{\mathrm{LoRA}}^\top\,F\,J_{\mathrm{LoRA}}.
% % % \end{equation}

% % % Geometrically, the feature-space sensitivity ellipsoids governed by $E_\phi(x)$ are first pulled back to parameter space through $J_\phi(x)$ to form $F$, and then restricted to the admissible LoRA directions through $J_{\mathrm{LoRA}}$. The spectrum of $G_{\mathrm{fact}}$ therefore captures the curvature that the optimizer effectively experiences in the LoRA subspace: $\lambda_{\max}(G_{\mathrm{fact}})$ identifies the most sensitive direction for prediction stability and $\mathrm{tr}(G_{\mathrm{fact}})$ measures the aggregate sensitivity.


% % % Because $F_\theta$ and $F_\Delta$ are congruent (Lemma~\ref{lem:lora-factor-subspace}), spectral control in factor space (e.g., shrinking $\lambda_{\max}$ or $\mathrm{tr}$ via SAM/natural-gradient preconditioning) is \emph{exactly} equivalent to controlling the LoRA-restricted curvature $F_\Delta$.
% % % Combined with the Gaussian posterior above, this shows that flatness regularisation confined to $\mathcal W_\Delta$ reduces the expected sharpness term in the PAC--Bayes bound and thereby improves stability and generalisation within PECL, independently of the particular LoRA coordinates used.
% % % This shows that optimizing flatness {only} inside the LoRA subspace is sufficient to control the effective curvature seen by training, which in turn stabilizes the shared features through the chain 
% % % $E_\phi \;\to\; F \;\to\; G_{\mathrm{fact}}$.







% % \subsubsection{Empirical Validation: Feature Stability and Performance Gains}


% % % \begin{table}[t]
% % % \centering
% % % \caption{Class-Incremental Learning Performance and Feature-Space Diagnostics on ImageNet-R (ViT-B/16). Reports Nearest-Mean-Exemplars (NME) accuracy, Linear Probe (LP) accuracy, L2 prototype drift $\Delta^{L2}_{0}$, Centered Kernel Alignment (CKA), and spectral properties of the feature matrix (trace, effective rank, anisotropy).
% % % }.
% % % \label{tab:q1_lp_feat_reorg}
% % % \setlength{\tabcolsep}{4.5pt}
% % % \resizebox{\linewidth}{!}{
% % % \begin{tabular}{l l ccc c cc ccc}
% % % \toprule
% % % \multirow{2}{*}{\textbf{Opt.}} & \multirow{2}{*}{\textbf{Method}} &
% % % \multicolumn{3}{c}{\textbf{NEM} ($\uparrow$)} &
% % % \multicolumn{1}{c}{\textbf{LP} ($\uparrow$)} &
% % % \multicolumn{2}{c}{\textbf{Feature Stability}} &
% % % \multicolumn{2}{c}{\textbf{Feature Matrix}} & \\
% % % \cmidrule(lr){3-5} \cmidrule(lr){6-6} \cmidrule(lr){7-8}\cmidrule(lr){9-11}
% % % &  & ${\text{ACC}}_{20}$  & AAA  & BWT 
% % %  & ${\text{ACC}}_{20}$ 
% % %  & L2 drift &CKA$_{0}$
% % %  & $\mathrm{tr}(E)$ & eff  & ani\\ 
% % % % \midrule
% % % % \multirow{3}{*}{\centering{\textbf{Backbone}}}
% % % % & MulTaskFT   \\
% % % % & MulTaskLP  \\
% % % % & MulTaskLoRA  \\
% % % % $\overline{\text{ACC}}_{20}$ 
% % % \midrule
% % % \multirow{1}{*}{\centering{\textbf{Backbone}}}
% % %   & PerTaskLP &62.96$\pm$0.12	&66.58$\pm$0.37	&-7.66$\pm$1.27
% % %  & 69.58     & 0        & 1        & 0.82 & 128.11  &18.85\\
% % % \midrule
% % % \multirow{7}{*}{\centering{\textbf{SGD}}} %\rotatebox[origin=c]{90}
% % %  & SeqFT     &70.37$\pm$1.26   & 75.27$\pm$0.50 & -15.75$\pm$1.13& 77.65    & 29.94 & 0.79 & 0.56 & 97.29 &23.38 \\ 
% % %  %  & SeqLP &62.96$\pm$0.12	&66.58$\pm$0.37	&-7.66$\pm$1.27
% % %  % & 69.58     & 0        & 1        & 0.82 & 128.11  &18.85\\
% % % \cmidrule(lr){2-11}
% % %  & SeqLoRA    &72.82$\pm$1.27   & 76.62$\pm$0.68 & -6.75$\pm$1.29 & 77.34 & 16.58 & 0.73 & 0.77 & 125.50 &21.20 \\
% % %  \cmidrule(lr){2-11}
% % %  & IncLoRA    &75.41$\pm$0.26   & 78.37$\pm$0.26 & -4.32$\pm$0.58 & 78.97 & 16.58 & 0.70 & 0.68 & 119.41 &17.68\\
% % %   % & MOS             &  &  &  &  &  &  &  \\
% % % \cmidrule(lr){2-11}
% % %  & OLoRA   &75.10$\pm$0.41  & 78.65$\pm$0.67 & -3.90$\pm$0.54 & 79.47    & 16.67 & 0.73 & 0.68 & 117.33  &18.54\\

% % %   % & InfLoRA  &  &  &  &  &  &  &  \\
% % %  % & TUNA            &  &  &  &  &  &  &  \\
% % % \midrule
% % % \multirow{7}{*}{\centering{\textbf{SAM}}} %\rotatebox[origin=c]{90}
% % %  & SeqFT     &73.50$\pm$2.01 & 78.58$\pm$0.43 & -13.45$\pm$2.19 & 79.43 & 17.42& 0.86 & 0.72 & 128.31  &16.09\\
% % % % & SeqLP &62.96$\pm$0.12	&66.58$\pm$0.37	&-7.66$\pm$1.27     & 69.58 & 0        & 1        & 0.81 & 125.40 &19.05\\
% % % \cmidrule(lr){2-11}
% % %  & SeqLoRA    &74.53$\pm$0.60  & 77.62$\pm$0.43 & -5.54$\pm$0.72  & 77.87 & 14.82  & 0.73 & 0.81 & 135.38  &13.50\\
% % %  \cmidrule(lr){2-11}
% % %  & IncLoRA     &75.95$\pm$0.073	&78.68$\pm$0.25	&-3.87$\pm$0.029
% % %  &79.60 &15.14 &0.68 &0.70  &127.16  &15.52  \\
% % % % & MOS             &  &  &  &  &  &  &  \\
% % % \cmidrule(lr){2-11}
% % %  & OLoRA    &75.73$\pm$0.32  & 78.70$\pm$0.67 & -3.47$\pm$0.54 & 79.75    & 15.13 & 0.74 & 0.72 & 121.23  &15.94\\
% % %   % & InfLoRA  &  &  &  &  &  &  &  \\
% % %  % & TUNA            &  &  &  &  &  &  &  \\
% % % % \midrule
% % % % \multicolumn{2}{r}{\textbf{Avg $\Delta$}} &  &  &  &  &  &  &  \\
% % % \bottomrule
% % % \end{tabular}
% % % }
% % % \end{table}

% % \begin{table}[htpb]
% % \centering
% % \caption{Class-Incremental Learning Performance and Feature-Space Diagnostics on ImageNet-R (ViT-B/16). Reports Nearest-Mean-Exemplars (NME) accuracy, Linear Probe (LP) accuracy, L2 prototype drift $\Delta^{L2}_{0}$, Centered Kernel Alignment (CKA), and spectral properties of the feature matrix (trace, effective rank, anisotropy).}
% % \label{tab:q1_lp_feat_reorg}
% % \setlength{\tabcolsep}{4.5pt}
% % \resizebox{\linewidth}{!}{
% % \begin{tabular}{l l ccc  cc cc}
% % \toprule
% % \multirow{2}{*}{\textbf{Method}} & \multirow{2}{*}{\textbf{Opt.}} &
% % \multicolumn{3}{c}{\textbf{NME} ($\uparrow$)} &
% % \multicolumn{2}{c}{\textbf{Feature Stability}} &
% % \multicolumn{2}{c}{\textbf{Feature Matrix}} \\
% % \cmidrule(lr){3-5} \cmidrule(lr){6-7} \cmidrule(lr){8-9}
% % &  & ${\text{ACC}}_{20}$  & AAA  & BWT 
% %  & L2 drift & CKA$_{0}$
% %  & $\mathrm{tr}(E)$ & eff \\ 
% % \midrule
% % \multirow{1}{*}{\textbf{PerTaskLP}}
% %   & SGD & 62.96$\pm$0.12 & 66.58$\pm$0.37 & -7.66$\pm$1.27
% %      & 0        & 1        & 0.82 & 128.11  \\
% % \midrule
% % \multirow{2}{*}{\textbf{SeqFT}}
% %  & SGD     & 69.49$\pm$0.16 & 75.10$\pm$0.57 & -16.11$\pm$0.23 & 29.57 & 0.78 & 0.58 & 93.53  \\
% %  & SAM     & 72.36$\pm$2.46 & 78.31$\pm$0.50 & -14.04$\pm$3.30 & 17.04 & 0.86 & 0.74 & 127.28 \\
% %  &GAM  & 73.74$\pm$0.88 & 78.32$\pm$0.94 & -13.36$\pm$0.20 & 13.80 & 0.85 & 0.78 & 127.87 \\
% % \midrule
% % \multirow{2}{*}{\textbf{SeqLoRA}}
% %  & SGD     & 72.82$\pm$1.27 & 76.62$\pm$0.68 & -6.75$\pm$1.29  & 16.58 & 0.73 & 0.77 & 125.50 \\
% %  & SAM     & 74.53$\pm$0.60 & 77.62$\pm$0.43 & -5.54$\pm$0.72  & 14.82 & 0.73 & 0.81 & 135.38 \\
% %  &GAM   & 75.74$\pm$0.15 & 78.86$\pm$0.53 & -6.19$\pm$0.33  & 11.28 & 0.78 & 0.78 & 133.96 \\
% % \midrule
% % \multirow{2}{*}{\textbf{IncLoRA}}
% %  & SGD      & 75.41$\pm$0.26 & 78.37$\pm$0.26 & -4.32$\pm$0.59  & 16.87 & 0.69 & 0.67 & 121.18 \\
% %  & SAM     & 75.95$\pm$0.07 & 78.68$\pm$0.25 & -3.87$\pm$0.03 & 15.14 & 0.68 & 0.70 & 127.16 \\
% %  &GAM  & 76.72$\pm$0.54 & 79.42$\pm$0.40 & -3.63 $\pm$0.67  & 12.15 & 0.75 & 0.73 & 119.68 \\
% % \midrule
% % \multirow{2}{*}{\textbf{OLoRA}}
% %  & SGD     & 75.10$\pm$0.41 & 78.65$\pm$0.67 & -3.90$\pm$0.54 & 17.23 & 0.69 & 0.68 & 120.89 \\
% %  & SAM     & 75.73$\pm$0.32 & 78.70$\pm$0.67 & -3.47$\pm$0.54 & 15.17 & 0.72 & 0.70 & 123.76 \\
% %  &GAM  & 76.56$\pm$0.53 & 79.16$\pm$0.43 & -3.37 $\pm$0.68  & 12.16 & 0.75 & 0.74 & 119.36\\
% % \bottomrule
% % \end{tabular}
% % }
% % \end{table}







% % % \begin{figure}[htbp]
% % % \centering
% % % \begin{minipage}{0.495\linewidth}
% % % \centering
% % % \includegraphics[width=0.9\linewidth]{figures_TRML/finetune_42_sam-sgd_efm_spectrum_compare2.pdf}
% % % \caption{Finetune EFM Matrix}
% % % \label{fig: EFM_seqFT}
% % % \end{minipage}
% % % \begin{minipage}{0.495\linewidth}
% % % \centering
% % % \includegraphics[width=0.9\linewidth]{figures_TRML/seqlora_42_sam-sgd_efm_spectrum_compare2.pdf}
% % % \caption{SeqLoRA EFM Matrix}
% % % \label{fig: EFM_SeqLoRa}
% % % \end{minipage}
% % % \end{figure}


% % % \begin{figure}[htbp]
% % % \centering
% % % \begin{minipage}{0.495\linewidth}
% % % \centering
% % % \includegraphics[width=0.9\linewidth]{figures_TRML/finetune_42_sam-sgd_efm_spectrum_compare2.pdf}
% % % \subcaption{Eigenvalue Spectrum in Sequential Fine-tuning (SeqFT)}
% % % \label{fig:EFM_seqFT}
% % % \end{minipage}
% % % \begin{minipage}{0.495\linewidth}
% % % \centering
% % % \includegraphics[width=0.9\linewidth]{figures_TRML/seqlora_42_sam-sgd_efm_spectrum_compare2.pdf}
% % % \subcaption{Eigenvalue Spectrum in Sequential LoRA (SeqLoRA)}
% % % \label{fig:EFM_SeqLoRA}
% % % \end{minipage}
% % % \caption{
% % % Comparison of EFM eigenvalue spectra in the feature space across incremental learning methods. The horizontal axis represents eigenvalue rank (sorted by magnitude), while the vertical axis shows eigenvalue magnitude. Only the top 200 eigenvalues are non-zero, with remaining dimensions at zero. SAM consistently produces flatter, heavier-tailed spectra compared to SGD across both configurations, indicating reduced energy concentration in the leading principal directions. This spectral profile suggests that SAM learns more isotropic representations with higher effective dimensionality.
% % % }
% % % \label{fig:EFM_spectra_comparison}
% % % \end{figure}



% % To analyse the mechanism behind the gains observed in Section~3.4, 
% % we complement CL metrics and weight-space sharpness with feature-space diagnostics that directly probe representation stability.


% % We first evaluate Nearest-Mean-of-Exemplars (NME) accuracy, which performs classification in the frozen feature space using class prototypes and therefore isolates the quality of the representation from the classifier head. To quantify temporal stability, we measure prototype drift as the $\ell_2$ distance between class prototypes across tasks and use Centered Kernel Alignment (CKA) to assess how much the learned features deviate from the initial backbone representation. Finally, we study the empirical feature matrix (EFM) constructed from normalized features. Its spectrum provides three complementary indicators: the trace measures overall feature energy, and the effective rank captures how many directions carry significant variance.

% % The results, summarized in Table~\ref{tab:q1_lp_feat_reorg}, reveal a coherent mechanism. First, SAM consistently improves NME accuracy across SeqLoRA, IncLoRA, and OLoRA, which shows that adapter-subspace flatness enhances the geometry of the learned features rather than merely tuning the classifier. Second, SAM reduces prototype drift and increases CKA stability, indicating that the shared representation changes more smoothly across tasks. Third, the EFM spectra become less anisotropic: the leading eigenvalues shrink relative to the average eigenvalue, the anisotropy index decreases, and the effective rank increases. These patterns match the theoretical picture in Section~4.3.1, where spectral suppression of the LoRA-restricted Fisher $F_\Delta$ tightens the expected sharpness term and yields more isotropic, stable features. 

% % % Section~\ref{sec:main-results} (E2) will distill these observations into a concise empirical statement for Q2.



% % % First, we focus on the NME-based evaluation, which isolates the quality of the learned representation from the classifier head. The consistent performance gains observed in the NME curves in Figure~\ref{fig:combined_curves_imageR} and in the NME accuracies reported in Table~\ref{tab:q1_lp_feat_reorg} show that SAM improves feature quality across SeqLoRA, IncLoRA, and OLoRA. The enhanced NME performance indicates that SAM fosters more robust and stable feature embeddings that are less susceptible to interference from subsequent tasks. This is precisely in line with our theoretical picture: optimizing for flatness within the LoRA subspace, promotes smoother and more stable representations and thereby mitigates catastrophic forgetting at the representation level.


% % % To further dissect the mechanism behind feature-space stability, we analyze feature drift and the eigenvalue spectrum.  As reported in Table~\ref{tab:q1_lp_feat_reorg}, SAM consistently reduces the L2 prototype drift $\Delta L^2_c$ and improves structural consistency of features, as reflected in higher CKA stability values. These metrics confirm that features learned with SAM deviate less from their initial configuration and retain their structure more faithfully throughout the task sequence. The eigenvalue spectra in Figure~\ref{fig:EFM_spectra_comparison} provide a complementary geometric view: compared to SGD, SAM-optimized models exhibit less concentrated leading eigenvalues and heavier spectral tails. In other words, energy is less concentrated in a few sharp directions and more evenly spread across many moderately curved directions, which corresponds to a more isotropic representation geometry. This behaviour matches our theoretical prediction that SAM reduces the spectral norm of the curvature and redistributes it within the LoRA subspace.




% % % Taken together, the representation-level evidence in this subsection and the continual-learning metrics in Section~3.4 paint a coherent picture. Subspace spectral control through SAM reduces the dominant curvature of the LoRA-restricted Fisher, which tightens the expected sharpness and KL terms in our PAC-Bayesian bound. Empirically, this manifests as more isotropic and stable feature representations, lower prototype drift, higher CKA stability, and flatter EFM spectra. These changes in turn account for the improvements in $\mathrm{ACC}_{20}$, $\mathrm{BWT}$, AAA, and NME observed in Section~3.4. This closes Q2: the benefits of LoRA subspace flatness are mediated by feature-space robustness rather than by a purely classifier-level effect.



% % % Complementing this, the EFM eigenvalue spectra in Figure 5 reveal that SAM produces flatter, more heavy-tailed distributions compared to SGD, indicating a more isotropic allocation of feature energy across dimensions rather than concentration in a few dominant directions. This spectral profile—characterized by higher effective rank ($r_{\text{eff}}$) and reduced anisotropy—corroborates that SAM-learned representations possess greater dimensionality and robustness. The consequent stabilization of the feature base directly supports the linear classifier head, enabling it to maintain performance over time with reduced fine-tuning effort. These empirical findings resonate with our PAC-Bayesian analysis: flatter minima in the LoRA subspace restrict disruptive updates along sensitive feature directions ($F_\Delta$), thereby jointly enhancing both representation robustness and classification stability.





% % % Combined with the Gaussian posterior above, this shows that flatness regularisation confined to $\mathcal W_\Delta$ reduces the expected sharpness term in the PAC--Bayes bound and thereby improves stability and generalisation within PECL, independently of the particular LoRA coordinates used. 
% % % This shows that optimizing flatness {only} inside the LoRA subspace is sufficient to control the effective curvature seen by training, which in turn stabilizes the shared features through the chain 
% % % $E_\phi \;\to\; F \;\to\; G_{\mathrm{fact}}$.


% % % We consider task-wise Gaussian posteriors on the LoRA subspace, \(Q_t^\Delta=\mathcal N(\hat{\Delta W}_t,\Sigma_t)\). The sharpness contribution in our bound is
% % %  \[
% % %  \operatorname{E-Sh}_t^\Delta \approx \tfrac12\mathbb E_{x\in S_t}\big[\mathrm{tr}\big(F_\Delta(x)\Sigma_t\big)\big],
% % %  \]
% % %  with \(F_\Delta(x)=U^\top J_\phi(x)^\top E_\phi(x) J_\phi(x)U\). Increasing flatness acts on this term and on the KL complexities in two precise ways. By congruence \(F_\theta \sim F_\Delta\) and PSD trace inequalities,
% % %  \[
% % %  \mathrm{tr}\big(F_\Delta(x)\Sigma_t\big)\le \lambda_{\max}!\big(F_\Delta(x)\big),\mathrm{tr}(\Sigma_t)
% % %  \quad\text{and}\quad
% % %  \mathrm{tr}\big(F_\Delta(x)\Sigma_t\big)\le \mathrm{tr}\big(F_\Delta(x)\big),\lambda_{\max}(\Sigma_t).
% % %  \]
% % %  Any method that reduces \(\lambda_{\max}(F_\theta)\) or\(\mathrm{tr}(F_\theta)\) therefore tightens \(\operatorname{E-Sh}_t^\Delta\) for fixed \(\Sigma_t\).  Flatter subspaces require smaller posterior variance and smaller mean shifts to achieve the same empirical risk, which reduces \(\mathrm{KL}\big(Q_t^\Delta|P_t^\Delta\big)\). 



% %   % In each step SAM searches, within \(\mathcal W_\Delta\), for a worst-case perturbation of the loss and then updates against it. In the local quadratic regime the inner maximizer aligns with the top eigenvector of \(F_\theta\) and hence of \(F_\Delta\). Iterating this mechanism suppresses \(\lambda_{\max}(F_\Delta)\) and improves conditioning \(\kappa(F_\Delta)\). The reduction of \(\lambda_{\max}\) controls the single most sensitive admissible direction, while better conditioning redistributes curvature from a few sharp axes to many gentler ones. Because \(F_\Delta=U^\top J_\phi^\top E_\phi J_\phi U\) emphasizes boundary-dominated, classifier-aligned directions, this spectral contraction limits how much new low-rank updates can distort previously learned margins. The outcome is smaller \(\operatorname{E-Sh}_t^\Delta\), smaller KL penalties for comparable empirical fit, and reduced interference across tasks, with plasticity preserved by operating inside \(\mathcal W_\Delta\).


 

% % % Finally, we can reinterpret SAM through this lens. SAM, when applied only to the LoRA parameters, works by explicitly searching (at each update) for the perturbation inside the LoRA subspace that maximizes the loss, and then updating in a way that counters that perturbation. Geometrically, that adversarial perturbation aligns with the top eigenvector of $G_{\text {LoRA }}$, i.e., the most dangerous direction for prediction stability.
% % % Thus, SAM-in-LoRA has two direct spectral effects:

% % % It suppresses $\lambda_{\max}\left(G_{\mathrm{fact}}\right)$, reducing the strength of the single worst ("catastrophic") direction.

% % % It reduces Condition number $\kappa\left(G_{\mathrm{fact}}\right)$, distributing risk more evenly across multiple gentler directions, rather than letting one direction dominate.

% % % For continual learning, this is exactly what we want. Future tasks will continue to update their own LoRA blocks in low-rank subspaces. If the shared representation has already been shaped so that no single LoRA direction can easily shatter previous features, then new tasks can adapt without erasing old ones. In Section 5.5, we empirically confirm this story: applying SAM only in the LoRA subspace systematically shrinks $\lambda_{\max}\left(G_{\mathrm{fact}}\right)$, improves the effective rank of $E_{\phi}$, stabilizes cross-task representations, and ultimately reduces forgetting.



% % To validate this mechanism, we complement standard continual learning metrics with feature space
% % diagnostics. Nearest Mean of Exemplars (NME) accuracy evaluates classification performance in the frozen
% % feature space using class prototypes, which isolates representation quality from the classifier head. Prototype
% % drift measures the $\ell_2$ distance between class prototypes across tasks. Centered Kernel Alignment (CKA)
% % quantifies the structural similarity between the learned features and the initial backbone representation.
% % Finally, the spectrum of the empirical feature matrix (EFM) reveals how the sensitivity of the predictive
% % distribution is distributed across feature directions through its trace, effective rank, and anisotropy.

% % On ImageNet-R, these diagnostics exhibit a consistent pattern. Flatness optimization within the LoRA
% % subspace, through methods such as SAM and GAM, increases NME accuracy, reduces prototype drift, and
% % produces higher CKA values relative to SGD. At the same time, the trace of the EFM is moderated and its
% % effective rank increases, which indicates that sensitivity is more evenly distributed instead of concentrating
% % on a few highly curved directions. These trends are robust across different LoRA organizations.

% % These empirical observations align with the theoretical mechanism above. Spectral suppression of the
% % restricted Fisher matrix $F_\Delta$ controls the sensitivity of the representation, tightens the PAC-Bayesian
% % bound, and improves feature stability at the same time.














% \subsubsection{(E2) ImageNet-C/P: robustness and the role of LoRA geometry}

% Experiment~E2 examines robustness under corruptions and perturbations. Using the same training configuration, we evaluate SeqLoRA, IncLoRA, and OLoRA with SGD and SAM on corruption and perturbation benchmarks, namely Tiny ImageNet-C and Tiny ImageNet-P. Figure~\ref{fig:robustness_curves_comparison} reports $\operatorname{ACC}_t^{\mathrm{cls}}$ and $\operatorname{ACC}_t^{\mathrm{nme}}$ over the task sequence for each configuration.



% % Figure~\ref{fig:robustness_curves_comparison} reports AAA and NME on corruption and perturbation benchmarks. Across ImageNet-C and ImageNet-P, SAM consistently outperforms SGD over the task sequence in both metrics and for all LoRA organizations. This shows that adapter subspace flatness improves robustness of both the classifier and the feature space under distribution shift.

% % \begin{figure}[htbp]
% %     \centering
% %     % 第一行：ImageNet-C
% %     \begin{minipage}{0.495\textwidth}
% %         \centering
% %         \includegraphics[width=0.9\linewidth]{figures_TRML/AAA_curve_Imagenet-C_summary_inc10_rank1--3.pdf}
% %         \subcaption{Average Accuracy with task on ImageNet-C}
% %         \label{fig:aaa_imagenet_c}
% %     \end{minipage}%
% %     \begin{minipage}{0.495\textwidth}
% %         \centering
% %         \includegraphics[width=0.9\linewidth]{figures_TRML/NME_curve_Imagenet-C_summary_inc10_rank16--3.pdf}
% %         \subcaption{Nearest Mean of Exemplars with task  on ImageNet-C}
% %         \label{fig:nme_imagenet_c}
% %     \end{minipage}
    
% %     % 第二行：ImageNet-P  
% %     \begin{minipage}{0.495\textwidth}
% %         \centering
% %         \includegraphics[width=0.9\linewidth]{figures_TRML/AAA_curve_Imagenet-P_summary_inc10_rank16--3.pdf}
% %         \subcaption{Average Accuracy  with task on ImageNet-P}
% %         \label{fig:aaa_imagenet_p}
% %     \end{minipage}%
% %     \begin{minipage}{0.495\textwidth}
% %         \centering
% %         \includegraphics[width=0.9\linewidth]{figures_TRML/NME_curve_Imagenet-P_summary_inc10_rank16--3.pdf}
% %         \subcaption{Nearest Mean of Exemplars with task on ImageNet-P}
% %         \label{fig:nme_imagenet_p}
% %     \end{minipage}
% %     \caption{
% %         Robustness of class incremental representations evaluated on ImageNet-C and ImageNet-P. 
% %         % The panels show AAA and NME curves under corruption robustness (ImageNet-C) and perturbation robustness (ImageNet-P) for the same class incremental configurations as in Figure~\ref{fig:combined_curves_imageR}. 
% %         Flatness optimizers (SAM and GAM) consistently outperform SGD also on these robustness benchmarks, which indicates that the benefits of promoting flat minima in the adapter transfer to corruptions and perturbations. %temporal
% %     }
% %     \label{fig:robustness_curves_comparison}
% % \end{figure}

% \begin{figure}[htbp]
%     \centering
%     % ImageNet-C: cls / nme
%     \begin{minipage}{0.495\textwidth}
%         \centering
%         \includegraphics[width=0.95\linewidth]{figures_TRML/AAA_curve_Imagenet-C_summary_inc10_rank16_ponit2.pdf}
%         \subcaption{ImageNet-C: $\operatorname{ACC}_t^{\mathrm{cls}}$}
%         \label{fig:aaa_imagenet_c}
%     \end{minipage}\hfill
%     \begin{minipage}{0.495\textwidth}
%         \centering
%         \includegraphics[width=0.95\linewidth]{figures_TRML/NME_curve_Imagenet-C_summary_inc10_rank16_point2.pdf}
%         \subcaption{ImageNet-C: $\operatorname{ACC}_t^{\mathrm{nme}}$}
%         \label{fig:nme_imagenet_c}
%     \end{minipage}\hfill
%     % ImageNet-P: cls / nme
%     \begin{minipage}{0.495\textwidth}
%         \centering
%         \includegraphics[width=0.95\linewidth]{figures_TRML/AAA_curve_Imagenet-P_summary_inc10_rank16_point2.pdf}
%         \subcaption{ImageNet-P: $\operatorname{ACC}_t^{\mathrm{cls}}$}
%         \label{fig:aaa_imagenet_p}
%     \end{minipage}\hfill
%     \begin{minipage}{0.495\textwidth}
%         \centering
%         \includegraphics[width=0.95\linewidth]{figures_TRML/NME_curve_Imagenet-P_summary_inc10_rank16_point2.pdf}
%         \subcaption{ImageNet-P: $\operatorname{ACC}_t^{\mathrm{nme}}$}
%         \label{fig:nme_imagenet_p}
%     \end{minipage}
%     \caption{
%         Robustness of accuracy with the classifier head ($\operatorname{ACC}_t^{\mathrm{cls}}$) and the NME-based ($\operatorname{ACC}_t^{\mathrm{nme}}$) on ImageNet-C and ImageNet-P. 
%         % The four panels report performance after each task for the same class incremental methods and optimization configurations as in Figure~\ref{fig:combined_curves_imageR}. 
%          SAM and GAM consistently outperform SGD on both corruption robustness (ImageNet-C) and perturbation robustness (ImageNet-P), which indicates that promoting flat minima in the adapter improves robustness of the learned representations.
%     }
%     \label{fig:robustness_curves_comparison}
% \end{figure}



% Across both benchmarks and all LoRA organizations, SAM consistently outperforms SGD in AAA and NME, which shows that adapter subspace flatness improves robustness of both the classifier and the feature space under a wide range of corruptions and perturbations. The joint evolution of AAA and NME remains closer and more stable under SAM, indicating that the representations and the classifier stay better aligned when the adapter subspace is flattened.
% % The relative behaviour of different LoRA organizations changes in the presence of strong corruptions. 
% On ImageNet-C and ImageNet-P, SeqLoRA becomes the most robust configuration, while IncLoRA and OLoRA suffer larger accuracy drops. This is consistent with the dependence of the effective sharpness term on the spectrum of $F_W$ restricted to the adapter subspace. Larger update subspaces admit more directions along which curvature and shift induced drift can accumulate, whereas the compact SeqLoRA subspace constrains both $\Sigma_t$ and $F_\Delta$ to a smaller region and therefore exposes fewer unstable directions. 


% % This unified advantage is evident in both linear probe classification (AAA) and feature-based evaluation (NME), demonstrating that the flatter minima discovered by SAM yield representations with enhanced stability and reduced sensitivity to distribution shifts. 

% % On ImageNet-R, IncLoRA and OLoRA slightly outperform SeqLoRA, which reflects the benefit of a richer update subspace in clean conditions.



% % The relative behaviour of different LoRA organizations changes once strong corruptions are introduced. On ImageNet-C and ImageNet-P, SeqLoRA becomes the most robust configuration, whereas IncLoRA and OLoRA suffer larger accuracy drops. From the PAC-Bayesian perspective, this is consistent with the dependence of the effective sharpness term on the spectrum of $F_\Delta$ restricted to the adapter subspace. Larger update subspaces provide more directions along which curvature and shift-induced drift can accumulate, while the compact SeqLoRA subspace constrains both the posterior covariance and $F_\Delta$ to a smaller region and therefore admits fewer unstable directions.

% % The joint evolution of AAA and NME further links robustness to feature-space geometry. Under SAM, AAA and NME remain closer and more stable, which is consistent with the view that adapter-subspace flatness primarily stabilises the feature geometry rather than only improving the classifier head.


% % Under SGD, NME typically degrades faster than AAA on ImageNet-C and ImageNet-P, indicating that class prototypes become less representative when features are corrupted

% % Table~\ref{table:q1_ci_with_weight_flatness_reorg} and Table~\ref{tab:q1_lp_feat_reorg}, together with Figures~\ref{fig:combined_curves_imageR} and \ref{fig:EFM_spectra_comparison}, compare SGD and SAM on ImageNet-R. Across SeqFT, SeqLoRA, IncLoRA and OLoRA, SAM consistently improves $\mathrm{ACC}_{20}$, $\mathrm{AAA}$ and $\mathrm{BWT}$, while simultaneously reducing sharpness, as well as the dominant eigenvalues and traces of the empirical Hessian and Fisher. NME accuracies and feature diagnostics further show more stable class prototypes, higher CKA, and flatter EFM spectra. This block provides the primary evidence that enforcing flat minima in the LoRA subspace is already sufficient to improve PECL performance and answers Q1 in the affirmative.
% % % These results answer Q1: enforcing flat minima only in the LoRA adapter subspace already suffices to improve the stability-plasticity balance in PECL.

% % \subsubsection{(E2) ImageNet-C/P: robustness and the role of LoRA geometry.}
% % Figure~\ref{fig:robustness_curves_comparison} extends the analysis to corruption and perturbation benchmarks.Across both ImageNet-C and ImageNet-P, SAM maintains a consistent advantage over SGD in AAA and NME, which shows that adapter subspace flatness improves robustness of both classifier and features under distribution shift. 
% % The relative ranking of LoRA organizations changes between ImageNet-R and ImageNet-C/P: configurations with richer update subspaces, such as IncLoRA and OLoRA, are strongest in domain, whereas the more compact SeqLoRA subspace becomes comparatively more robust under strong corruptions and perturbations. This pattern indicates that the size and organization of $\mathcal W_\Delta$ modulate the effect of flatness regularization and therefore addresses Q3 from a robustness perspective.






% % Figure~\ref{fig:robustness_curves_comparison} extends the ImageNet-R analysis to corrupted (ImageNet-C) and perturbed (ImageNet-P) inputs. Across all LoRA organizations, SAM maintains a clear advantage over SGD in both AAA and NME, and this gap persists along the task sequence. Adapter-subspace flatness therefore improves not only in-domain performance but also robustness under distribution shift.


% % The behaviour of different LoRA organizations changes once strong corruptions are introduced. On ImageNet-R, IncLoRA and OLoRA slightly outperform SeqLoRA, which reflects the benefit of a richer update subspace in clean conditions. On ImageNet-C and ImageNet-P, SeqLoRA becomes the most robust adapter configuration, while IncLoRA and OLoRA suffer larger accuracy drops. This is consistent with the PAC-Bayesian view in which the effective sharpness term depends on the spectrum of $F_\Delta$ in the adapter subspace. Larger subspaces provide more directions along which curvature and shift-induced drift can accumulate, whereas the compact SeqLoRA subspace constrains both the posterior covariance and $F_\Delta$ to a smaller region and therefore admits fewer unstable directions.


% % The relative evolution of AAA and NME further links these robustness results to the feature-space analysis. Under SGD, NME tends to degrade faster than AAA on ImageNet-C and ImageNet-P, indicating that class prototypes become less representative when features are corrupted. Under SAM, the gap between AAA and NME is smaller and more stable, which is consistent with the interpretation that adapter-subspace flatness primarily acts by stabilising the feature geometry rather than only improving the classifier head.

% % SAM~\cite{foretSharpnessAwareMinimizationEfficiently2021}, RWP~\cite{liRevisitingRandomWeight2024}, and GAM~\cite{liRevisitingRandomWeight2024}


% \subsubsection{(E3) Different flatness optimizers on ImageNet-R}

% Experiment~E3 compares several perturbation based flatness optimizers on ImageNet-R while keeping the adapter organization fixed. We consider RWP, SAM, C-Flat, and GAM in addition to the SGD baseline, and evaluate $\operatorname{AAA}^{\mathrm{nme}}$ for SeqLoRA, IncLoRA, and OLoRA. The results are reported in Table~\ref{tab:aaa_nme_lora_org} as absolute scores and as gains relative to SGD.



% \begin{table}[thpb]
% \centering
% \caption{The $\operatorname{AAA}^{\mathrm{cls}}$ and $\operatorname{AAA}^{\mathrm{nme}}$ performance across LoRA organizations under different flatness optimizers. 
% % Rows below SGD report the change relative to the SGD baseline.
% }
% \label{tab:aaa_nme_lora_org}
% % \setlength{\tabcolsep}{5pt}
% % \resizebox{\linewidth}{!}
% \setlength{\tabcolsep}{3pt}
% \renewcommand{\arraystretch}{0.8} % 可选：稍微压缩行距
% \resizebox{0.75\linewidth}{!}{
% \begin{tabular}{l cc cc cc}
% \toprule
% \multirow{3}{*}{\textbf{Opt.}} &
% \multicolumn{2}{c}{\textbf{SeqLoRA}} &
% \multicolumn{2}{c}{\textbf{IncLoRA}} &
% \multicolumn{2}{c}{\textbf{OLoRA}} \\
% \cmidrule(lr){2-3} \cmidrule(lr){4-5} \cmidrule(lr){6-7}
% % & $\overline{\text{ACC}}_{20}$ (CL) & $\overline{\text{ACC}}_{20}$ (NME) & $\overline{\text{ACC}}_{20}$ (CL) & $\overline{\text{ACC}}_{20}$ (NME) & $\overline{\text{ACC}}_{20}$ (CL) & $\overline{\text{ACC}}_{20}$ (NME) \\
% % \cmidrule(lr)
% % &  CL &  NME & CL &  NME &  CL &  NME \\
% % & $\overline{\text{ACC}}_{20}$ & $\overline{\text{ACC}}_{20}$  & $\overline{\text{ACC}}_{20}$  & $\overline{\text{ACC}}_{20}$  & $\overline{\text{ACC}}_{20}$  & $\overline{\text{ACC}}_{20}$  \\
% & $\operatorname{AAA}^{\mathrm{cls}}$ &  $\operatorname{AAA}^{\mathrm{nme}}$ & $\operatorname{AAA}^{\mathrm{cls}}$ & $\operatorname{AAA}^{\mathrm{nme}}$ & $\operatorname{AAA}^{\mathrm{cls}}$ & $\operatorname{AAA}^{\mathrm{nme}}$ \\
% \midrule
% % $\overline{\text{ACC}}_{20}$
% SGD  & 69.07 & 76.62 & 71.29 & 78.29 & 71.41 & 78.47 \\
% \midrule
% $\Delta$ RWP   &0.70  & 0.29  & 0.86  & 0.53  & 0.28  & 0.12   \\
% $\Delta$ SAM   & 2.35  & 1.00  & 1.29  & 0.40  & 1.08  & 0.22  \\
% $\Delta$ CFLAT & 2.35  & 1.04  & 2.67  & 1.26  & 1.07  & 0.12  \\
% $\Delta$ GAM  & 7.73  & 2.24  & 5.23  & 1.13  & 4.64  & 0.68   \\
% \bottomrule
% \end{tabular}}
% \end{table}


% All four methods improve AAA and NME over SGD, but the magnitude of the gain follows a clear and consistent hierarchy: RWP yields only modest improvements, SAM and C-FLAT provide moderate and reliable boosts, and GAM achieves the largest gains across all three LoRA organisations. When combined with the curvature statistics in Table~\ref{table:q1_ci_with_weight_flatness_reorg}, this ordering reveals a parallel hierarchy on the geometric side: optimisers that more aggressively suppress subspace sharpness and the dominant eigenvalues of the LoRA-restricted Fisher $F_\Delta$ also achieve better continual learning performance. In other words, the trend appears simultaneously in the flatness diagnostics and in the CI metrics. This monotone relationship is precisely what Theorem~\ref{thm:pecl-bound} predicts: for a fixed LoRA geometry and training budget, methods that reduce the adapter-subspace expected sharpness $\mathrm{E\text{-}Sh}^\Delta_t$ tighten the empirical term of the PAC–Bayesian bound and allow the posterior and hyper-posterior KL terms to remain moderate, which manifests empirically as higher AAA and NME.

% The dependence on LoRA organisation is also systematic. Under SGD, IncLoRA and OLoRA slightly outperform SeqLoRA on ImageNet-R, reflecting the benefit of richer adapter subspaces in fitting the distribution drift. Under stronger flatness promotion, particularly with GAM, SeqLoRA exhibits the largest improvement over its own SGD baseline, and the performance gap between others organisations narrows. This pattern is consistent with the view that aggressive spectral suppression is most effective when perturbations are confined to a more compact update subspace $\mathcal W_\Delta$ in which curvature can be controlled uniformly, whereas larger adapter manifolds benefit more from capacity under weak regularisation but are harder to flatten.  


% % It reinforces our central conclusion that, in the PECL regime, controlling flatness within the LoRA update subspace $\mathcal W_\Delta$ is sufficient to steer the behaviour of the full model: perturbation-based optimizers that more effectively flatten $\mathcal W_\Delta$ consistently deliver stronger continual learning performance.


% % All four methods improve AAA and NME over SGD, but the magnitude of the gain follows a clear hierarchy. RWP yields modest improvements, SAM and C-Flat provide moderate and reliable boosts, and GAM achieves the largest gains, especially in classifier AAA. This ordering matches the intended strength of spectral suppression on the LoRA restricted curvature $F_\Delta$: GAM most aggressively reduces the largest eigenvalues of $F_\Delta$, SAM and C-Flat achieve intermediate suppression, and RWP provides the weakest effect. The empirical hierarchy in Table~3 therefore offers a controlled test of the sharpness term in the bound.

% % The dependence on LoRA organization is also systematic. Under SGD, IncLoRA and OLoRA slightly outperform SeqLoRA on ImageNet-R, reflecting the benefit of larger adapter subspaces in fitting the in distribution data. Under stronger flatness regularisation, particularly with GAM, SeqLoRA exhibits the largest improvement over its own SGD baseline. This pattern is consistent with the view that aggressive spectral suppression is more effective when perturbations are confined to a compact subspace in which curvature can be controlled. E3 thus refines hypothesis (H3) by showing that both the choice of optimizer and the geometry of $W_\Delta$ jointly determine the quantitative benefits of flatness promotion.



% % Table~\ref{tab:aaa_nme_lora_org} compares four perturbation based flatness optimizers on ImageNet-R: RWP~\cite{liRevisitingRandomWeight2024}, SAM~\cite{foretSharpnessAwareMinimizationEfficiently2021}, C-FLAT~\citep{liCFlatMoreEfficient2025}, and GAM~\cite{liRevisitingRandomWeight2024}. All methods improve AAA and NME over the SGD baseline, and the gains follow a clear hierarchy. RWP yields only small improvements, SAM and C-FLAT provide moderate and stable gains, and GAM achieves the largest gains, especially in classifier $\mathrm{AAA}$. This ordering is consistent with the intended strength of spectral suppression on the LoRA restricted curvature $F_\Delta$, and thus provides a controlled test of the sharpness term in our bound.

% % The dependence on the LoRA organization is also systematic. Under SGD, IncLoRA and OLoRA slightly outperform SeqLoRA, which reflects the advantage of a richer adapter subspace for fitting data. Under stronger flatness regularisation, particularly with GAM, SeqLoRA shows the largest improvement over its own SGD baseline. In this case, aggressive spectral suppression reduces the effective sharpness term most efficiently because perturbations are restricted to a compact subspace where curvature can be controlled more precisely.


% % Taken together, these results indicate that RWP provides relatively mild flatness regularisation, SAM and C-FLAT serve as robust general purpose choices, and GAM is most beneficial when the adapter subspace is compact and more aggressive spectral shaping is acceptable.



% % Table~\ref{tab:aaa_nme_lora_org} ompares four perturbation based optimizers on ImageNet-R: RWP, SAM, C-FLAT, and GAM. All four methods improve $\mathrm{AAA}$ and NME over the SGD baseline. The gains follow a clear hierarchy. RWP yields only modest improvements, SAM and C-FLAT provide moderate yet reliable gains, and GAM produces the largest improvements. This empirical ordering mirrors the strength with which each method suppresses the spectrum of the LoRA restricted curvature and provides a controlled test of the sharpness term in our bound.

% % Table~\ref{tab:aaa_nme_lora_org} compares four perturbation based optimisers on ImageNet-R. All methods outperform SGD in terms of AAA and NME, but the gains follow a clear hierarchy. RWP provides only modest improvements, SAM and C-FLAT yield moderate and reliable gains, and GAM achieves the largest improvements, especially for $\mathrm{ACC}_{20}$. This ordering matches the intended strength of spectral control on the LoRA restricted curvature $F_\Delta$.

% % The interaction with the LoRA organization is systematic. Under SGD, IncLoRA and OLoRA slightly outperform SeqLoRA, reflecting the benefit of a richer adapter subspace in fitting the training distribution. Under stronger flatness regularisation, particularly with GAM, SeqLoRA obtains the largest improvement over its own SGD baseline and nearly closes the gap to IncLoRA and OLoRA. This is consistent with the PAC-Bayesian analysis: when $W_\Delta$ is low dimensional and shared, aggressive spectral suppression can reduce the sharpness term more effectively, since injected perturbations are confined to a compact subspace where curvature is easier to control.

% % In practical terms, these results suggest that RWP offers a light-weight but relatively weak form of flatness regularisation, SAM and C-FLAT provide a robust default choice, and GAM is most beneficial when the adapter subspace is compact and one can afford more aggressive spectral shaping.



% % The interaction with LoRA geometry is also structured: under pure SGD, IncLoRA and OLoRA outperform SeqLoRA; under stronger flatness regularization, SeqLoRA exhibits the largest relative improvement and nearly closes the gap. This pattern matches the prediction that methods which more aggressively suppress the spectrum of $F_\Delta$ yield larger continual learning gains, and that the magnitude of the gain is modulated by the size and structure of $\mathcal W_\Delta$.

% \subsubsection{(E4) Local versus global perturbations}


% \begin{table}[thpb]
% \centering
% \caption{$\operatorname{AAA}^{\mathrm{cls}}$ and $\operatorname{AAA}^{\mathrm{nme}}$ on ImageNet-R under local global (full model) versus (LoRA only).}
% \label{tab:lora——or_full}
% \setlength{\tabcolsep}{3pt}
% \renewcommand{\arraystretch}{0.8}
% \resizebox{0.75\linewidth}{!}{
% \begin{tabular}{c c cc cc cc}
% \toprule
% \multicolumn{2}{c}{\textbf{Opt.}} &
% \multicolumn{2}{c}{\textbf{SeqLoRA}} &
% \multicolumn{2}{c}{\textbf{IncLoRA}} &
% \multicolumn{2}{c}{\textbf{OLoRA}} \\
% \cmidrule(lr){1-2} \cmidrule(lr){3-4} \cmidrule(lr){5-6} \cmidrule(lr){7-8}
% \textbf{strength} & \textbf{range} &
% $\operatorname{AAA}^{\mathrm{cls}}$ & $\operatorname{AAA}^{\mathrm{nme}}$ &
% $\operatorname{AAA}^{\mathrm{cls}}$ & $\operatorname{AAA}^{\mathrm{nme}}$ &
% $\operatorname{AAA}^{\mathrm{cls}}$ & $\operatorname{AAA}^{\mathrm{nme}}$ \\
% \midrule
% \multirow{1}{*}{$0$} &Base &
% 69.07 & 76.62 & 71.29 & 78.29 & 71.41 & 78.47 \\
% \midrule
% % 组 1：\rho_1
% \multirow{2}{*}{$0.05$}
% &full 
% &-3.5  & -1.54 & -2    & -0.38 & -2.62 & -0.8  \\ 
% & lora &0.27  & 0.02  & 0.79  & 0.51  & 0.12  & 0.16  \\ 
% \midrule
% % 组 2：\rho_2
% \multirow{2}{*}{$0.10$} & full &-4.82 & -1.29 & -4.74 & -0.74 & -5.53 & -1.18 \\
% & lora &0.7   & 0.29  & 0.86  & 0.53  & 0.28  & 0.12  \\
% \midrule
% % 组 3：\rho_3
% \multirow{2}{*}{$0.15$} & full &-6.72 & -1.69 & -6.84 & -0.98 & -6.49 & -1.29 \\
% & lora &0.41  & 0.08  & 0.77  & 0.53  & 0.27  & 0.11 \\
% \bottomrule
% \end{tabular}}
% \end{table}


% This experiment is not intended to challenge the conclusions of Flat-LoRA in its original single task setting. Instead, we ask whether promoting global flatness remains beneficial in PECL regime. We compare the original Flat-LoRA configuration, which applies RWP to all model parameters, with a PECL compatible variant that restricts RWP to the LoRA branch only.

% Table~\ref{tab:lora——or_full} reports $\operatorname{AAA}^{\mathrm{cls}}$ and $\operatorname{AAA}^{\mathrm{nme}}$ on ImageNet-R for three LoRA organizations under different perturbation strengths. For all strengths and all three architectures, global perturbations applied to the full model consistently degrade performance. In contrast, perturbations confined to the adapter subspace yield small but systematic gains over the base configuration. These results indicate that, in the PECL setting, promoting flatness through local perturbations within the LoRA subspace is beneficial, whereas enforcing global perturbations on all parameters is uniformly harmful in this experiment.



% % This experiment does not intend to refute Flat-LoRA in its original single-task setting. Its goal is to examine whether global flatness promote remain beneficial under a strict PECL regime with a frozen backbone. We compare  the Flat-LoRA method which apply RWP on the whole model with the basic setting RWP only on the lora part in PECL.


% % Table \ref{tab:lora——or_full} reports $\operatorname{AAA}^{\mathrm{nme}}$ on ImageNet R for three LoRA organizations under different range. For all strengths and all three architectures, global perturbations (full) consistently hurt performance. In contrast, adapter only perturbations (lora) yield small but systematic gains over the base configuration. These results indicate that, in the PECL setting, promoting flatness through local perturbations confined to the LoRA subspace is beneficial, while enforcing global perturbations on all parameters is uniformly detrimental in this class incremental experiment.


% % Experiment~E4 investigates the role of local versus global perturbations in a long sequence NLP benchmark with T5 adapters. We compare three configurations: a baseline that uses Adam without explicit flatness promotion, RWP-LoRA, which restricts perturbations to the adapter subspace $W_\Delta$, and RWP-Full, which applies RWP to all parameters of the model. All methods are trained under matched budgets and evaluated using $\mathrm{ACC}_{15}$ and $\mathrm{BWT}$.

% % Across all adapter organizations and task orders, RWP-LoRA consistently improves over the Adam baseline in both final accuracy and backward transfer. In contrast, RWP-Full is unstable. In relatively benign regimes it can sometimes match or slightly exceed Adam, but it is consistently weaker than the corresponding RWP-LoRA configuration and can even degrade performance. In particular, global perturbations tend to amplify catastrophic forgetting, which suggests that perturbing the frozen backbone violates the assumptions under which the subspace based PAC-Bayesian bound is valid.

% % These findings confirm that the PECL regime in which the backbone is frozen and only adapters are updated is precisely the regime where subspace perturbations provide reliable benefits. Restricting perturbations and flatness optimisation to $W_\Delta$ aligns with the theoretical framework in Section~3 and Section~4, whereas global perturbations fall outside this regime and need not improve, and may even harm, continual learning performance.


% % Table~\ref{tab:longseq_nlp} presents results on a long sequence NLP benchmark with T5 adapters. Across all LoRA organizations and task orders, RWP-LoRA, which restricts perturbations to $\mathcal W_\Delta$, consistently improves over the Adam baseline in both $\mathrm{ACC}_{15}$ and $\mathrm{BWT}$. In contrast, RWP-Full, which perturbs all parameters, is unstable: it sometimes matches or slightly exceeds Adam in relatively benign regimes, but it is always weaker than the corresponding RWP-LoRA configuration and can even degrade performance. These findings confirm that the PECL regime in which the backbone is frozen and only adapters are updated is precisely the regime where subspace perturbations provide reliable benefits, while global perturbations fall outside the assumptions of our theory.

% % Taken together, these four experimental blocks deliver a coherent picture across modalities, datasets, and algorithms. Controlling sharpness inside the LoRA adapter is sufficient to improve continual learning performance (Q1), operates primarily through feature space stability quantified in Section~4 (Q2), and its quantitative effect is governed by both the geometry of $\mathcal W_\Delta$ and the choice of flatness promoting optimizer (Q3). The following subsections return to each experimental block in more detail and analyze robustness, optimizer behaviour, and local versus global perturbations beyond these summary conclusions.

% % \subsubsection{(E5) Exps in NLP}



% % \subsection{Detailed Discussion}
% % \textbf{Robustness under corruption and perturbation.}

% % \begin{figure}[htbp]
% % \centering
% % \begin{minipage}{0.495\linewidth}
% % \centering
% % \includegraphics[width=0.9\linewidth]{figures_TRML/AAA_curve_Imagenet-C_summary_inc10_rank16.pdf}
% % \caption{$\mathrm{AAA}$ Curve in Imagenet-C}
% % \label{fig: 1_AAA}
% % \end{minipage}
% % \begin{minipage}{0.495\linewidth}
% % \centering
% % \includegraphics[width=0.9\linewidth]{figures_TRML/NME_curve_Imagenet-C_summary_inc10_rank16.pdf}
% % \caption{$\mathrm{NME}$ Curve in Imagenet- C }
% % \label{fig: 2_NME}
% % \end{minipage}
% % \end{figure}

% % \begin{figure}[htbp]
% %     \centering
% %     \begin{minipage}{0.495\textwidth}
% %         \centering
% %         \includegraphics[width=0.9\linewidth]{figures_TRML/AAA_curve_Imagenet-C_summary_inc10_rank16.pdf}
% %         \subcaption{Average Accuracy (AAA) on ImageNet-C}
% %         \label{fig:aaa_imagenet_c}
% %     \end{minipage}%
% %     \begin{minipage}{0.495\textwidth}
% %         \centering
% %         \includegraphics[width=0.9\linewidth]{figures_TRML/NME_curve_Imagenet-C_summary_inc10_rank16.pdf}
% %         \subcaption{Nearest Mean of Exemplars (NME) on ImageNet-C}
% %         \label{fig:nme_imagenet_c}
% %     \end{minipage}
% %     \caption{
% %         AAA and NME summary curves on ImageNet-C ($T=20, r=16$). Under common corruptions, SAM maintains its advantage over SGD across various PECL configurations, achieving higher accuracy in both evaluation protocols. The consistent performance gap demonstrates that the flatter solutions found by SAM yield more robust representations that are less sensitive to input corruptions.
% %     }
% %     \label{fig:curves_imagenet_c}
% % \end{figure}


% % \begin{figure}[htbp]
% %     \centering
% %     \begin{minipage}{0.495\textwidth}
% %         \centering
% %         \includegraphics[width=0.9\linewidth]{figures_TRML/AAA_curve_Imagenet-P_summary_inc10_rank16.pdf}
% %         \subcaption{Average Accuracy (AAA) on ImageNet-P}
% %         \label{fig:aaa_imagenet_p}
% %     \end{minipage}%
% %     \begin{minipage}{0.495\textwidth}
% %         \centering
% %         \includegraphics[width=0.9\linewidth]{figures_TRML/NME_curve_Imagenet-P_summary_inc10_rank16.pdf}
% %         \subcaption{Nearest Mean of Exemplars (NME) on ImageNet-P}
% %         \label{fig:nme_imagenet_p}
% %     \end{minipage}
% %     \caption{
% %         AAA and NME summary curves on ImageNet-P ($T=20, r=16$). When facing input perturbations, SAM consistently outperforms SGD across various PECL configurations. The sustained performance advantage underscores that the flatter minima discovered by SAM confer greater stability to the learned representations, making them more resilient to input variations.
% %     }
% %     \label{fig:curves_imagenet_p}
% % \end{figure}



% % \begin{figure}[htbp]
% % \centering
% % \begin{minipage}{0.495\linewidth}
% % \centering
% % \includegraphics[width=0.9\linewidth]{figures_TRML/AAA_curve_Imagenet-P_summary_inc10_rank16.pdf}
% % \caption{$\mathrm{AAA}$ Curve in Imagenet-P }
% % \label{fig: 1_AAA}
% % \end{minipage}
% % \begin{minipage}{0.495\linewidth}
% % \centering
% % \includegraphics[width=0.9\linewidth]{figures_TRML/NME_curve_Imagenet-P_summary_inc10_rank16.pdf}
% % \caption{$\mathrm{NME}$ Curve in Imagenet-P }
% % \label{fig: 2_NME}
% % \end{minipage}
% % \end{figure}





% % Figure~\ref{fig:robustness_curves_comparison} extends our analysis to robustness benchmarks with input corruptions (ImageNet-C) and perturbations (ImageNet-P). 
% % First, SAM consistently outperforms SGD across all configurations in both AAA and NME accuracy, demonstrating that flatter adapters enhance both classifier robustness and feature stability.
% % Second, the systematic gap between linear probe accuracy and NME accuracy across robustness benchmarks reveals how much of the continual learning degradation originates from the classification head. In all three tables that report both metrics, NME is consistently higher than the corresponding $\mathrm{AAA}$. Besides, the NME of the frozen backbone is often above the performance of the PerTaskLinearProbe clarify it. 
% % Third, the robustness benchmarks highlight how the organization and size of the LoRA update subspace interact with flatness and drift. On ImageNet-R, IncLoRA and OLoRA slightly outperform SeqLoRA, which indicates that enlarging or enriching the adapter subspace $\mathcal{W}_\Delta$ helps to exploit domain specific structure. On ImageNet-C and ImageNet-P, however, all PECL variants experience a noticeable drop relative to the backbone and the ranking between organizations changes: SeqLoRA becomes the strongest LoRA variant under severe corruptions and perturbations. This behaviour aligns with the fact that SeqLoRA trains a single shared adapter subspace across all tasks,less over specialized update subspace.  So both the posterior covariance $\Sigma_t$ and the curvature $F_\Delta$ are constrained in a smaller region of parameter space, stochastic gradient noise and input driven perturbations have less freedom to amplify along unstable directions. By contrast, IncLoRA and OLoRA distribute task specific information over a richer family of adapters, which increases the coverage of $\mathcal{W}_\Delta$ and improves in domain performance but also creates more directions along which curvature can concentrate and drift can accumulate when the distribution shifts. The reversal of the ranking across ImageNet-R and ImageNet-C/P therefore provides a further, and structurally interpretable, validation of our subspace based PAC-Bayesian analysis.





% % First, SAM consistently outperforms SGD across all configurations in both AAA and NME accuracy, demonstrating that flatter adapters enhance both classifier robustness and feature stability.





% % Second, the systematic gap between linear probe (AAA) and NME accuracy reveals that performance degradation often stems from classifier drift, not representation collapse. Third, LoRA organization interacts distinctly with robustness: while IncLoRA/OLoRA excel on in-domain data (ImageNet-R), SeqLoRA often becomes strongest under severe corruption. This reversal reflects their underlying geometry












% % They probe the effect of adapter subspace flatness in two complementary regimes of feature drift and performance drift, and thereby provide an additional validation of our theoretical predictions.

% % First, the qualitative conclusion regarding SAM remains stable across datasets. For all LoRA organizations and on both ImageNet-C and ImageNet-P, SAM consistently dominates SGD in terms of both $\mathrm{AAA}$ and $\mathrm{NME}$. This holds for models with a trained linear head and for models evaluated by NME, which depends only on backbone representations. In other words, increasing flatness within the LoRA subspace improves robustness of the learned features as well as robustness of the task specific classifier. From the perspective of our bound in Theorem~2, this behaviour is exactly what one would expect when the subspace sharpness term $\mathrm{E\text{-}Sh}_t^\Delta$ is reduced by suppressing the dominant eigenvalues of $F_\Delta$. Lower expected sharpness tightens the bound for clean data and at the same time reduces the sensitivity of the predictive distribution to input corruptions and perturbations.

% % Second, the systematic gap between linear probe accuracy and NME accuracy across robustness benchmarks reveals how much of the continual learning degradation originates from the classification head. In all three tables that report both metrics, NME is consistently higher than the corresponding $\mathrm{AAA}$, and the NME of the frozen backbone is often above the performance of the PerTaskLinearProbe model trained in a sequential fashion. This phenomenon is particularly clear in Figure~\ref{fig:robustness_curves_comparison}, where several PECL configurations with a trained head fall below the backbone baseline on ImageNet-C and ImageNet-P, while their NME curves remain closer to or above the backbone. This pattern is fully consistent with our subspace based analysis. The LoRA parameters are confined to the update subspace $\mathcal{W}_\Delta$, where SAM enforces flatness and reduces $\mathrm{E\text{-}Sh}_t^\Delta$ as well as the KL divergence between successive posteriors. The linear head, in contrast, is not constrained to remain within this regularized subspace. Its parameters drift under non stationary supervision and can accumulate curvature along unstable directions that are not controlled by the LoRA specific Fisher $F_\Delta$. As a consequence, head level drift can degrade performance even when the backbone representation remains relatively robust. The fact that SAM improves both $\mathrm{AAA}$ and $\mathrm{NME}$, while the NME gains are more stable and closer to the backbone, provides a second and independent confirmation that our theoretical conclusions primarily act at the representation level.

% % Third, the robustness benchmarks highlight how the organization of the LoRA subspace interacts with flatness and drift, again in a manner that is consistent with our theoretical framework. On ImageNet-R, IncLoRA and OLoRA slightly outperform SeqLoRA, which suggests that accumulating or orthogonalizing task specific adapters in a shared subspace can be beneficial when the test distribution matches the training domains. On ImageNet-C and ImageNet-P, however, all PECL variants experience a noticeable drop relative to the backbone, and the ranking between organizations changes. SeqLoRA often becomes the strongest LoRA variant under severe corruptions and perturbations. From the subspace viewpoint, this reversal can be understood as follows. IncLoRA and OLoRA explicitly combine task specific updates into a shared structure, which sharpens the associated curvature $F_\Delta = U^\top J_\phi^\top E_\phi J_\phi U$ along directions that are highly adapted to the original domains. This concentrated structure is advantageous for reusing knowledge on clean data, but it also makes the model more sensitive when the input distribution drifts far from the domains that shaped $F_\Delta$. SeqLoRA, by maintaining more task specific and less entangled adapters, induces a more fragmented and less over specialized update subspace. The resulting $F_\Delta$ is less tightly aligned with any single domain and therefore less sensitive to domain specific corruptions, which can translate into better robustness at the cost of slightly weaker in domain performance.


% % Third, the robustness benchmarks highlight how the organization and size of the LoRA update subspace interact with flatness and drift, again in a manner that is consistent with our theoretical framework. On ImageNet-R, IncLoRA and OLoRA slightly outperform SeqLoRA, which indicates that enlarging or enriching the adapter subspace $\mathcal{W}_\Delta$ helps to exploit domain specific structure when the test distribution matches the training data. In terms of our bound, these organizations effectively increase the dimensionality and expressiveness of $\mathcal{W}_\Delta$, which allows the learner to reduce the empirical risk more aggressively at the price of a larger KL term and a potentially sharper curvature $F_\Delta$ that must be controlled by SAM.

% % On ImageNet-C and ImageNet-P, however, all PECL variants experience a noticeable drop relative to the backbone and the ranking between organizations changes: SeqLoRA often becomes the strongest LoRA variant under severe corruptions and perturbations. This behaviour aligns with the fact that SeqLoRA trains a single shared adapter subspace across all tasks. Its update subspace $\mathcal{W}_\Delta$ is lower dimensional and less fragmented, so both the posterior covariance $\Sigma_t$ and the curvature $F_\Delta$ are constrained in a smaller region of parameter space. As a result, stochastic gradient noise and input driven perturbations have less freedom to amplify along unstable directions, and the effective subspace sharpness $\mathrm{E\text{-}Sh}_t^\Delta$ remains smaller when the input distribution drifts far from the source domain. The same structural constraint that limits SeqLoRA's capacity on clean ImageNet-R therefore becomes an advantage on ImageNet-C and ImageNet-P: the model is less expressive but also less exposed to noise, which translates into comparatively stronger robustness.

% % By contrast, IncLoRA and OLoRA distribute task specific information over a richer family of adapters, which increases the coverage of $\mathcal{W}_\Delta$ and improves in domain performance but also creates more directions along which curvature can concentrate and drift can accumulate when the distribution shifts. In terms of Theorem~2, SeqLoRA corresponds to a smaller and more tightly regularized subspace, with reduced $\mathrm{E\text{-}Sh}_t^\Delta$ and KL penalties under large shifts, while IncLoRA and OLoRA correspond to larger subspaces that achieve lower empirical loss on clean data but are more fragile under corruption and perturbation. The reversal of the ranking across ImageNet-R and ImageNet-C/P therefore provides a further, and structurally interpretable, validation of our subspace based PAC-Bayesian analysis.








% % Viewed together with the ImageNet-R results, the robustness experiments therefore validate our theory from two complementary sides. In regimes where continual adapters can surpass the backbone (clean ImageNet-R), SAM reduces subspace sharpness and keeps drift small, which lifts performance above the baseline. In regimes where continual training inevitably introduces additional drift and performance falls below the backbone (corrupted and perturbed ImageNet-C and ImageNet-P), SAM still moderates this drift and consistently closes the gap between adapter based models and the backbone. Both behaviours match the predictions of Theorem~2. Subspace spectral control through SAM reduces $\mathrm{E\text{-}Sh}_t^\Delta$ and the associated KL penalties, and this manifests empirically as more robust representations, more stable NME performance, and a systematic dependence of robustness on the geometry of the LoRA update subspace across SeqLoRA, IncLoRA, and OLoRA.



% % \textbf{Different flatness method }






% % compares different flatness optimizers across three LoRA organizations. For every structure, all four methods improve both $\overline{\mathrm{ACC}}_{20}$ and NME over the SGD baseline. The magnitude of the gains follows a clear hierarchy: RWP brings only modest improvements, SAM and C-FLAT provide stable but moderate boosts, and GAM achieves the largest gains, especially on $\overline{\mathrm{ACC}}_{20}$. This ordering matches the strength of the underlying perturbation schemes. In our configuration, RWP injects only isotropic Gaussian noise into the LoRA parameters and mixes noisy and clean gradients, which yields relatively weak subspace regularization. SAM and C-FLAT explicitly minimize flatness; C-FLAT further incorporates a first order correction, whereas GAM directly targets first order flatness by manipulating gradient norms and suppressing directions associated with large curvature in $F_\Delta$.

% % The interaction with the LoRA organization shows a complementary effect. Under SGD, IncLoRA and OLoRA slightly outperform SeqLoRA, reflecting the advantage of a richer update subspace $\mathcal{W}_\Delta$ for fitting the clean training distribution. With stronger flatness control, particularly under GAM, SeqLoRA exhibits the largest relative improvement and almost closes the gap, while IncLoRA and OLoRA still attain the highest NME scores. SeqLoRA operates in a single shared, lower dimensional subspace, so both the posterior covariance $\Sigma_t$ and the curvature $F_\Delta$ remain confined to a smaller region of parameter space, which makes perturbation based regularization more effective and limits drift. IncLoRA and OLoRA span a larger and more structured $\mathcal{W}_\Delta$, which increases expressiveness but also introduces more directions along which curvature and drift can accumulate when the distribution or task changes. Together with the robustness results on ImageNet R/C/P, these patterns support the central prediction of our subspace based PAC-Bayesian analysis: more aggressive reduction of subspace sharpness yields larger continual learning gains, and the size and geometry of $\mathcal{W}_\Delta$ modulate the extent of this improvement.




% % \textbf{Local vs. Global in NLP}







% % We evaluate the CL performance of the model using the final accuracy ($\mathrm{ACC}_T$), the averaged accuracy $\overline{\mathrm{ACC}}_T=$, and backward transfer $\mathrm{BWT}_T$. We also show the result related  NEM perforamnce. And train a liner probe using all previous dataset. 

% % For evaluate the flatness, we using  $\lambda_{\max}\big(H(\hat{W})\big)$ and $\mathrm{tr}\big(H(\hat{W})\big)$ related Hessian matrix, as well as  experimental Fisher  $\lambda_{\max}\big(F(\hat{W})\big)$ and $\mathrm{tr}\big(F(\hat{W})\big)$. We use the Power Iteration algorithm~\citep{kaddourWhenFlatMinima2022}
% % the Lanczos algorithm ~\citep{ghorbaniInvestigationNeuralNet2019} and iteration 100 times follow.

% % For evaluate the feature flatness, we measure the L2 drfit and CKA for the first task'class . Besides, we calculate the Feature Matrix and the 


% % We evaluate PECL methods on multiple robustness benchmarks based on ImageNet. Our primary testbed is ImageNet-R, which has been established as a standard benchmark in PECL by prior work. ImageNet-R serves as a benchmark for domain shift, comprising artistic renditions of ImageNet classes where semantic labels are conserved across significant stylistic variations. 
% % We further employ Tiny ImageNet-C and Tiny ImageNet-P to assess model performance under common corruptions and perturbations. Tiny ImageNet-C is a corrupted version of Tiny ImageNet across 200 classes, incorporating 15 types of corruption at 5 severity levels. ImageNet-P provides short MP4 perturbation sequences for each validation image. For computational efficiency, we only include three corruption types in Tiny ImageNet-C "gaussian noise", "impulse noise", and "shot noise" at a single severity level. For Tiny ImageNet-P, we uniformly sample three frames from each video (beginning, middle, and end) as static inputs to approximate the full perturbation trajectory.




% %  We use ImageNet-R as a standard benchmark for continual learning methods based on PEFT. Imagenet-R is generated through artistic processing of 200 classes from ImageNet, which represents real-world distribution shift
% % datasets consisting of changes in image style.
% % Besides, we also use IMAGENET-C and IMAGENET-P to evaluates performance on common corruptions and perturbations.  Tiny ImageNet-C has 200 classes with images of size 64x64 on results on all 15 corruptions and .  ImageNet-P sequences are MP4s to  spatter perturbation sequence is a validation sequence. And for each we sample 3 photo to represect the begining to the end transfore.




% % \begin{table*}[thbp]

% % \resizebox{\textwidth}{!}{
% % \begin{tabular}{cc|cc|cc|cc|cc}
% % \toprule
% % & & \multicolumn{2}{c|}{\textbf{Order4}}
% %   & \multicolumn{2}{c|}{\textbf{Order5}}
% %   & \multicolumn{2}{c|}{\textbf{Order6}}
% %   & \multicolumn{2}{c}{\textbf{Average (Avg)}} \\
% % \cmidrule(lr){3-4} \cmidrule(lr){5-6} \cmidrule(lr){7-8} \cmidrule(lr){9-10}
% % \textbf{Baseline} & \textbf{Method}
% %   & ${\text{ACC}}_{15}$ &  $\overline{\text{ACC}}_{20}$
% %   & ${\text{ACC}}_{15}$ & $\overline{\text{ACC}}_{20}$
% %   & ${\text{ACC}}_{15}$ & $\overline{\text{ACC}}_{20}$
% %  & ${\text{ACC}}_{15}$  & $\overline{\text{ACC}}_{20}$\\
% % \midrule
% % \multirow{4}{*}{{SeqLoRA}}  %\rotatebox[origin=c]{90}
% % & Adam      & 70.44 & 66.89 & 68.11 & 61.37 & 70.84 & 68.63 & 69.80 & 65.63   \\
% % & RWP-Full  & 67.74 & 66.25 & 65.11 & 64.67 & 73.11 & 69.90 & 68.65 & 66.94   \\

% % & RWP-LoRA  & 68.08 & 67.42 & 67.64 & 65.67 & 72.80 & 70.09 & 69.51 & 67.73   \\
% % \midrule
% % \multirow{4}{*}{{IncLoRA}} 
% % & Adam       & 70.13 & 67.07 & 70.13 & 64.90 & 43.16 & 64.80 & 61.14 & 65.59   \\

% % & RWP-Full   & 69.48 & 68.58 & 69.17 & 61.89 & 67.52 & 67.07 & 68.72 & 65.85  \\
% % % & LoRA-SAM        & -10.72  & 68.69 
% % %                   & -16.45  & 68.52 
% % %                   & \textbf{-9.45}  & 70.69 
% % %                   & -12.21  & 69.30 \\
% % & RWP-LoRA   & 70.44 & 71.77 & 71.07 & 62.41 & 65.16 & 67.50 & 68.89 & 67.23   \\
% % \midrule
% % \multirow{4}{*}{{OLoRA}} 
% % & Adam      & 66.95 & 73.21 & 66.91 & 68.92 & 38.85 & 66.25 & 57.57 & 69.46   \\
% % & RWP-Full  & 67.64 & 73.16 & 65.03 & 67.12 & 31.83 & 64.91 & 54.83 & 68.40   \\
% % % & LoRA-SAM        & -6.55  & \textbf{68.04} 
% % %                   & -9.22  & \textbf{66.99} 
% % %                   & -4.00  & 71.68 
% % %                   & -6.59  & \textbf{68.90} \\
% % & RWP-LoRA  & 66.41 & 73.52 & 65.21 & 67.36 & 27.79 & 64.15 & 53.14 & 68.34   \\
% % \bottomrule
% % \end{tabular}
% % }
% % \caption{CL results on the Long Sequence Benchmark (Orders 4-6), using the same evaluation metrics as \ref{tab:all_orders3}. LoRA-specific perturbations remain effective, with LoRA-SAM showing strong performance under stricter structural constraints.}
% % \label{tab:all_orders_456}
% % \vspace{-10pt}
% % \end{table*}


% % \begin{table*}[thbp]

% % \resizebox{\textwidth}{!}{
% % \begin{tabular}{cc|cc|cc|cc|cc}
% % \toprule
% % & & \multicolumn{2}{c|}{\textbf{Order4}}
% %   & \multicolumn{2}{c|}{\textbf{Order5}}
% %   & \multicolumn{2}{c|}{\textbf{Order6}}
% %   & \multicolumn{2}{c}{\textbf{Average (Avg)}} \\
% % \cmidrule(lr){3-4} \cmidrule(lr){5-6} \cmidrule(lr){7-8} \cmidrule(lr){9-10}
% % \textbf{Baseline} & \textbf{Method}
% %   & ${\text{ACC}}_{15}$ & BWT
% %   & ${\text{ACC}}_{15}$ & BWT
% %   & ${\text{ACC}}_{15}$ & BWT
% %  & ${\text{ACC}}_{15}$  & BWT\\
% % \midrule
% % \multirow{4}{*}{{SeqLoRA}}  %\rotatebox[origin=c]{90}
% % & Adam          & \textbf{70.44} & -14.42
% %                     & 68.11 & -22.57
% %                     & 38.85 & -16.61
% %                    & 59.13  & -17.87 \\
% % & RWP-Full     & 67.74  & -16.96
% %                     & 65.11  & -18.40
% %                   & 27.79   & -19.50 
% %                     & 53.55 & -18.29\\
% % % & LoRA-SAM        & \textbf{-11.02}  & 70.17 
% % %                   & -19.27  & 66.87 
% % %                   & -19.54  & 18.11 
% % %                   & \textbf{-16.61}  & 51.72 \\
% % & RWP-LoRA       & 68.99 & -14.89
% %                    & 66.03  & -20.69
% %                     & \textbf{45.07} & \textbf{-15.56}
% %                     & \textbf{60.03}  & -17.05 \\
% % \midrule
% % \multirow{4}{*}{{IncLoRA}} 
% % & Adam       & -12.23  & 70.13 
% %                   & 70.13  & -14.60
% %                    & 43.16 & -13.77
% %                     & 61.14 & -13.53\\

% % & RWP-Full     & 69.48 & -12.87
% %                    & 69.17 & -18.67
% %                     & 65.16 & -9.80
% %                    & 67.94 & -13.78\\
% % % & LoRA-SAM        & -10.72  & 68.69 
% % %                   & -16.45  & 68.52 
% % %                   & \textbf{-9.45}  & 70.69 
% % %                   & -12.21  & 69.30 \\
% % & RWP-LoRA       & 69.29 & \textbf{-10.50}
% %                     & \textbf{72.52} & \textbf{-12.44}
% %                    & \textbf{72.20}  & -9.46
% %                    & \textbf{71.34}  & \textbf{-10.80}\\
% % \midrule
% % \multirow{4}{*}{{OLoRA}} 
% % & Adam         & 66.95 & -6.27
% %                     & 66.91 & \textbf{-7.70}
% %                     & 70.84 & -5.14
% %                     & 68.23 & -6.37\\
% % & RWP-Full     & 67.64 & \textbf{-5.19}
% %                     & 65.03 & -10.02
% %                    & 72.80  & -4.61
% %                    & 68.49  & -6.61  \\
% % % & LoRA-SAM        & -6.55  & \textbf{68.04} 
% % %                   & -9.22  & \textbf{66.99} 
% % %                   & -4.00  & 71.68 
% % %                   & -6.59  & \textbf{68.90} \\
% % & RWP-LoRA       & 66.43 & -6.15
% %                    & 66.00 & -10.22 
% %                    & 72.17 & -3.85
% %                     & 68.20 & -6.74\\
% % \bottomrule
% % \end{tabular}
% % }
% % \caption{CL results on the Long Sequence Benchmark (Orders 4-6), using the same evaluation metrics as \ref{tab:all_orders3}. LoRA-specific perturbations remain effective, with LoRA-SAM showing strong performance under stricter structural constraints.}
% % \label{tab:all_orders_456}
% % \vspace{-10pt}
% % \end{table*}


% % \begin{table*}[htbp]
% % \caption{CL results on the Long Sequence Benchmark. LoRA-specific perturbations remain effective, with LoRA-SAM showing strong performance under stricter structural constraints.}
% % \resizebox{\textwidth}{!}{
% % \begin{tabular}{cc|cc|cc|cc|cc}
% % \toprule
% % & & \multicolumn{2}{c|}{\textbf{Order1}}
% %   & \multicolumn{2}{c|}{\textbf{Order2}}
% %   & \multicolumn{2}{c|}{\textbf{Order3}}
% %   & \multicolumn{2}{c}{\textbf{Average (Avg)}} \\
% % \cmidrule(lr){3-4} \cmidrule(lr){5-6} \cmidrule(lr){7-8} \cmidrule(lr){9-10}
% % \textbf{Baseline} & \textbf{Method}
% %   & ${\text{ACC}}_{15}$ & BWT
% %   & ${\text{ACC}}_{15}$ & BWT
% %   & ${\text{ACC}}_{15}$ & BWT
% %  & ${\text{ACC}}_{15}$  & BWT\\
% % \midrule
% % \multirow{4}{*}{{SeqLoRA}}  %\rotatebox[origin=c]{90}
% % & Adam      & 70.44 & -11.45 & 68.11 & -14.88 & 38.85 & -45.16 & 59.13 & -23.83 \\
% % & Flat-LoRA(Full)  & 67.74 & -15.04 & 65.11 & -17.52 & 27.79 & -57.47 & 53.55 & -30.01 \\
% % % & LoRA-SAM        & \textbf{-11.02}  & 70.17 
% % %                   & -19.27  & 66.87 
% % %                   & -19.54  & 18.11 
% % %                   & \textbf{-16.61}  & 51.72 \\
% % & RWP(LoRA)  & 68.99 & -13.13 & 67.64 & -15.25 & 45.07 & -38.29 & 60.56 & -22.22 \\
% % \midrule
% % \multirow{4}{*}{{IncLoRA}} 
% % & Adam       & 70.13 & -10.67 & 70.13 & -11.67 & 43.16 & -38.91 & 61.14 & -20.42 \\

% % & RWP-Full(Full)     & 69.24 & -11.73 & 69.17 & -12.31 & 65.16 & -15.90 & 67.86 & -13.31 \\
% % % & LoRA-SAM        & -10.72  & 68.69 
% % %                   & -16.45  & 68.52 
% % %                   & \textbf{-9.45}  & 70.69 
% % %                   & -12.21  & 69.30 \\
% % & RWP(LoRA)    & 70.44 & -10.94 & 72.52 & -8.94  & 72.20 & -8.04  & 71.72 & -9.31  \\
% % \midrule
% % \multirow{4}{*}{{OLoRA}} 
% % & Adam       & 66.95 & -8.99  & 66.91 & -8.57  & 70.84 & -5.49  & 68.23 & -7.68  \\
% % & RWP-Full(Full)  & 67.64 & -7.66  & 65.03 & -11.24 & 72.80 & -4.84  & 68.49 & -7.91  \\
% % % & LoRA-SAM        & -6.55  & \textbf{68.04} 
% % %                   & -9.22  & \textbf{66.99} 
% % %                   & -4.00  & 71.68 
% % %                   & -6.59  & \textbf{68.90} \\
% % & RWP(LoRA)  & 66.43 & -9.47  & 66.00 & -10.50 & 73.11 & -3.37  & 68.51 & -7.78 \\
% % \bottomrule
% % \end{tabular}
% % }
% % \label{tab:longseq_nlp}
% % \end{table*}




% % \begin{table*}[htbp]
% % \caption{CL results on the Long Sequence Benchmark. LoRA-specific perturbations remain effective, with LoRA-SAM showing strong performance under stricter structural constraints.}
% % \resizebox{\textwidth}{!}{
% % \begin{tabular}{cc|cc|cc|cc|cc}
% % \toprule
% % & & \multicolumn{2}{c|}{\textbf{Order1}}
% %   & \multicolumn{2}{c|}{\textbf{Order2}}
% %   & \multicolumn{2}{c|}{\textbf{Order3}}
% %   & \multicolumn{2}{c}{\textbf{Average (Avg)}} \\
% % \cmidrule(lr){3-4} \cmidrule(lr){5-6} \cmidrule(lr){7-8} \cmidrule(lr){9-10}
% % \textbf{Baseline} & \textbf{Method}
% %   & ${\text{ACC}}_{15}$ & BWT
% %   & ${\text{ACC}}_{15}$ & BWT
% %   & ${\text{ACC}}_{15}$ & BWT
% %  & ${\text{ACC}}_{15}$  & BWT\\
% % \midrule
% % \multirow{4}{*}{{SeqLoRA}}  %\rotatebox[origin=c]{90}
% % & Adam      & 70.44 & -11.45 & 68.11 & -14.88 & 38.85 & -45.16 & 59.13 & -23.83 \\
% % & Flat-LoRA(Full)  & 67.74 & -15.04 & 65.11 & -17.52 & 27.79 & -57.47 & 53.55 & -30.01 \\
% % % & LoRA-SAM        & \textbf{-11.02}  & 70.17 
% % %                   & -19.27  & 66.87 
% % %                   & -19.54  & 18.11 
% % %                   & \textbf{-16.61}  & 51.72 \\
% % & RWP(LoRA)  & 68.99 & -13.13 & 67.64 & -15.25 & 45.07 & -38.29 & 60.56 & -22.22 \\
% % \midrule
% % \multirow{4}{*}{{IncLoRA}} 
% % & Adam       & 70.13 & -10.67 & 70.13 & -11.67 & 43.16 & -38.91 & 61.14 & -20.42 \\

% % & RWP-Full(Full)     & 69.24 & -11.73 & 69.17 & -12.31 & 65.16 & -15.90 & 67.86 & -13.31 \\
% % % & LoRA-SAM        & -10.72  & 68.69 
% % %                   & -16.45  & 68.52 
% % %                   & \textbf{-9.45}  & 70.69 
% % %                   & -12.21  & 69.30 \\
% % & RWP(LoRA)    & 70.44 & -10.94 & 72.52 & -8.94  & 72.20 & -8.04  & 71.72 & -9.31  \\
% % \midrule
% % \multirow{4}{*}{{OLoRA}} 
% % & Adam       & 66.95 & -8.99  & 66.91 & -8.57  & 70.84 & -5.49  & 68.23 & -7.68  \\
% % & RWP-Full(Full)  & 67.64 & -7.66  & 65.03 & -11.24 & 72.80 & -4.84  & 68.49 & -7.91  \\
% % % & LoRA-SAM        & -6.55  & \textbf{68.04} 
% % %                   & -9.22  & \textbf{66.99} 
% % %                   & -4.00  & 71.68 
% % %                   & -6.59  & \textbf{68.90} \\
% % & RWP(LoRA)  & 66.43 & -9.47  & 66.00 & -10.50 & 73.11 & -3.37  & 68.51 & -7.78 \\
% % \bottomrule
% % \end{tabular}
% % }
% % \label{tab:longseq_nlp}
% % \end{table*}

% % Table~\ref{tab:all_orders_456} 


% % reports results on the long sequence benchmark in NLP and serves two purposes. First, it tests our theory in a different modality and under more complex task orders. Across all three organizations and orders 4-6, RWP-LoRA consistently improves over the Adam baseline in both $\mathrm{ACC}_{15}$ and BWT (most prominently for IncLoRA, where the average $\mathrm{ACC}_{15}$ rises from $61.14$ to $71.72$ and BWT improves from $-20.42$ to $-9.31$). These gains confirm that restricting perturbations to the adapter subspace $\mathcal{W}_\Delta$ remains effective beyond vision and under long-range sequence modelling, in line with our PAC-Bayesian analysis that assumes a fixed backbone and task-dependent posteriors only on $\mathcal{W}_\Delta$.

% % Second, this setup allows us to revisit the tension between local and global perturbations discussed in LoRA-SAM and Flat-LoRA. Our bound is derived under the PECL assumption that the pretrained backbone is frozen; global perturbations violate this assumption and therefore fall outside the scope of the theorem. The RWP-Full rows make this explicit. In SeqLoRA, RWP-Full systematically degrades both $\mathrm{ACC}_{15}$ and BWT relative to Adam, and is clearly dominated by RWP-LoRA. In IncLoRA and OLoRA, RWP-Full can sometimes help compared with Adam (for example, average $\mathrm{ACC}_{15}$ improves from $61.14$ to $67.86$ in IncLoRA), but it never surpasses the corresponding RWP-LoRA configuration. This suggests that full-parameter RWP only works in benign regimes where the injected Gaussian noise is small and the backbone is already robust, yet even in these cases it remains inferior to subspace-restricted perturbations. Overall, the NLP results reinforce our main conclusion: in PECL, perturbing only the LoRA subspace yields more reliable gains than perturbing the full model, and occasional improvements from RWP-Full should be interpreted as side effects of small, well tolerated noise rather than as evidence against subspace-based flatness. A deeper understanding of these rare cases will likely require analysing the alignment between perturbation directions and the LoRA update subspace, which we leave for future work.


% % The long-sequence NLP benchmark with T5 adapters allows us to revisit perturbation strategies in a different modality. Table~\ref{tab:longseq_nlp} compares Adam, RWP-LoRA and RWP-Full across several task orders. Restricting perturbations to the adapter subspace (RWP-LoRA) consistently improves both $\mathrm{ACC}_{15}$ and $\mathrm{BWT}$ over Adam, which confirms that subspace perturbations remain effective when the backbone is a large language model. In contrast, RWP-Full is unstable: it seldom improves over Adam and is always weaker than the corresponding RWP-LoRA configuration when results are averaged across task orders.

% % This discrepancy aligns with the theoretical assumptions of PECL. The PAC-Bayesian bound is derived under a frozen backbone, with task-dependent posteriors supported on $W_\Delta$. Subspace perturbations modify the posterior covariance and the LoRA restricted curvature $F_\Delta$ without altering the backbone prior and therefore reduce the expected sharpness term in a controlled manner. Global perturbations implicitly modify both backbone and 。，蜜f【, fall outside this regime, and can introduce additional hyper-drift across tasks. Occasional benign cases where RWP-Full matches Adam can be interpreted as induced small-scale regularisation on an already robust backbone, rather than as evidence in favour of full parameter perturbations.  Local perturbations in $W_\Delta$ provide consistent and theoretically grounded benefits, whereas global RWP is less predictable. 

% % \subsection{Discussion: Closing the Loop Between Theory and Practice}

% % Sections~3 and~4 together form a closed loop between theory and practice. Section~3 shows that in PECL
% % the generalization behaviour is governed by flatness in the update subspace $W_\Delta$. Section~4 explains
% % how this subspace level flatness translates into concrete performance gains through feature stability and
% % demonstrates that LoRA implements exactly the geometry required by the theory.

% % Different LoRA organizations, such as SeqLoRA, IncLoRA, and OLoRA, correspond to different choices of
% % $W_\Delta$. SeqLoRA uses a single compact adapter subspace, whereas IncLoRA and OLoRA provide
% % richer and higher dimensional update subspaces. On clean in distribution data, the larger subspaces of
% % IncLoRA and OLoRA can offer an advantage in fitting the training distribution. Under strong noise and
% % distribution shift, however, as in ImageNet-C and ImageNet-P, the compact SeqLoRA subspace tends to be
% % more robust, since it restricts the number of directions along which curvature and posterior variance can
% % accumulate.

% % This phenomenon is consistent with the PAC-Bayesian perspective. Enlarging $W_\Delta$ increases the
% % capacity for modelling but also increases the potential for unstable directions in both $F_\Delta$ and
% % $\Sigma_t$. A smaller, well regularized subspace provides a more favourable balance between flexibility and
% % stability, especially under distribution shift.

% % In summary, flatness oriented optimization in the update subspace $W_\Delta$ is not a heuristic trick. It is a
% % direct intervention on the key term in the PAC-Bayesian bound. By operating in LoRA factor coordinates, we
% % implement this intervention in a parameter efficient manner while preserving pre trained knowledge. This
% % theoretically guided practice improves final accuracy and systematically enhances feature stability, providing
% % a principled route to robust parameter efficient continual learning.


% % \subsection{Discussion: the Link Between Theory and Practice}

% % The experiments in Section~5 provide a stress test of the PAC-Bayesian view developed in Sections~3 and~4. Theorem~2 shows that the continual learning risk is controlled by four quantities: an empirical term that averages task-wise training losses under Gaussian perturbations, the subspace expected sharpness $E\mbox{-}\mathrm{Sh}_\Delta$, the adapter-level divergence $\mathrm{KL}(Q_\Delta \Vert P_\Delta)$, and the hyper-posterior drift $\sum_t \mathrm{KL}(P_t \Vert P_{t-1})$. Section~4 further links $E\mbox{-}\mathrm{Sh}_\Delta$ to the leading eigenvalues of the LoRA-restricted Fisher $F_\Delta$ and to the stability of feature representations.

% % Experiments E1 and E2 directly probe the relationship between adapter-subspace flatness and both continual learning performance and feature stability. Across ImageNet-R and its corrupted and perturbed variants, reducing $E\mbox{-}\mathrm{Sh}_\Delta$ and the spectrum of $F_\Delta$ consistently improves $\mathrm{ACC}_T$, $\overline{\mathrm{ACC}}_T$, and NME, while limiting prototype drift and preserving CKA similarity. E3 then varies the strength of flatness promotion along a hierarchy of perturbation-based optimizers and confirms that stronger spectral suppression in $F_\Delta$ yields predictable gains, especially when the update subspace is compact. Finally, E4 contrasts local and global perturbations and shows that perturbations confined to $W_\Delta$ are beneficial and stable, whereas global perturbations that move the frozen backbone can degrade performance and amplify forgetting, which lies outside the regime where the subspace-aware bound is valid.

% % Taken together, these results support a coherent picture. In PECL, operating in LoRA factor coordinates is not merely a convenient parameterization. It realizes in practice the subspace geometry assumed by the PAC-Bayesian analysis, and flatness-oriented optimization in $W_\Delta$ acts as a direct intervention on the sharpness term in the bound. The observed improvements in accuracy and feature robustness therefore do not only indicate empirical gains, but provide concrete evidence that the sharpness-aware PAC-Bayesian interpretation of adapter-subspace flatness is predictive of continual learning behaviour.



% % \subsection{Discussion: Closing the Loop Between Theory and Practice}


% % The experiments in Section~5 instantiate and test the decomposition given by Theorem~2. 
% % The bound splits the continual learning test risk into an empirical term, a subspace expected sharpness term $E\text{-}Sh_\Delta$ that depends on the spectrum of the adapter-space Fisher $F_\Delta$, a sum of adapter-space complexity terms $\sum_t \mathrm{KL}(Q^\Delta_t \Vert P^\Delta_t)$, and a hyper-drift term $\sum_t \mathrm{KL}(P_t \Vert P_{t-1})$. 
% % E1 and E2 jointly probe the link between the empirical component and the flatness term. 
% % They show that reductions in subspace sharpness and in the dominant eigenvalues of $F_\Delta$ systematically coincide with improvements in $\mathrm{ACC}_T$, $\overline{\mathrm{ACC}}_T$, $\mathrm{BWT}_T$, and feature stability diagnostics such as prototype drift, CKA, and the effective rank of the empirical feature matrix. 
% % E3 acts on the same sharpness term by varying the strength of spectral suppression in $W_\Delta$ through different optimisers and demonstrates a consistent hierarchy of gains that matches their curvature control strength. 
% % E4 targets the structural part of the bound. 
% % It confirms that the favourable decomposition of complexity into adapter-space KL and hyper-drift terms relies on restricting perturbations to $W_\Delta$; once perturbations are extended to the frozen backbone, the empirical benefits deteriorate and catastrophic forgetting increases.

% % Across LoRA organisations, enlarging $W_\Delta$ increases modelling capacity but also creates more directions along which curvature and posterior variance may accumulate. 
% % The experiments indicate that compact subspaces such as SeqLoRA can be more robust under strong flatness promotion and distribution shift, which is consistent with the PAC-Bayesian view that both $F_\Delta$ and the posterior covariance $\Sigma_t$ scale with the geometry of $W_\Delta$.

% % Overall, flatness-oriented optimisation in the update subspace $W_\Delta$ is not a heuristic. 
% % It is a direct intervention on the key terms that appear in the sharpness-aware PAC-Bayesian bound. 
% % By operating in LoRA factor coordinates, we realise this intervention in a parameter-efficient manner while preserving pre-trained knowledge. 
% % Therefore, the experiments not only show performance gains but also provide direct empirical support for the sharpness-aware PAC-Bayesian interpretation proposed in Section~3.


% % The experiments in Section~5 confirm the picture developed in Sections~3 and~4. The PAC-Bayesian analysis shows that, in PECL, generalization is governed by flatness within the update subspace $W_\Delta$ and by the associated subspace expected sharpness $E\text{-Sh}_t^\Delta$.The feature-level analysis then links spectral suppression of the LoRA-restricted Fisher $F_\Delta$ to more stable representations. Experiments E1–E4 confirm this picture across modalities and optimizers.

% % Across LoRA organisations, enlarging $W_\Delta$ increases modelling capacity but also introduces more directions along which curvature and posterior variance may accumulate. The experiments suggest that compact subspaces such as SeqLoRA can be more robust under strong flatness promotion and distribution shift, which aligns with the PAC-Bayesian view that both $F_\Delta$ and the posterior covariance $\Sigma_t$ scale with the geometry of $W_\Delta$.

% % Overall, flatness oriented optimization in the update subspace $W_\Delta$ is not a heuristic. It is a direct intervention on the key term in the PAC-Bayesian bound. By operating in LoRA factor coordinates, we implement this intervention in a parameter efficient manner while preserving pre trained knowledge. This theoretically guided practice improves final accuracy and systematically enhances feature stability, providing a principled route to robust parameter efficient continual learning.


% \subsection{Ablation Studies}



% \begin{figure}[thbp]
%     \centering
%     \begin{minipage}{0.495\textwidth}
%         \centering
%         \includegraphics[width=0.9\linewidth]{figures_TRML/length_NME_T_avg_by_method_opt_imagenetr.pdf}
%         \subcaption{Different Length}
%         \label{fig:curves_imageR_length}
%     \end{minipage}%
%     \begin{minipage}{0.495\textwidth}
%         \centering
%         \includegraphics[width=0.9\linewidth]{figures_TRML/rank_NME_T_avg_by_method_opt_imagenetr.pdf}
%         \subcaption{Different Rank}
%         \label{fig:curves_imageR_rank}
%     \end{minipage}
%     \caption{
%         Effect of task sequence length and adapter rank on NME-based average accuracy on ImageNet-R..
%     }
%     \label{fig:ablation}
% \end{figure}

% We further vary the task sequence length and the adapter rank on ImageNet-R and report NME-based average accuracy in Figure~\ref{fig:ablation}. Across all values of $T$ and $r$, the relative behavior observed remains stable: flatness aware optimizers consistently outperform or SGD counterparts for SeqLoRA, IncLoRA, and OLoRA. This consistency across task lengths and adapter capacities indicates that the benefit of controlling flatness in the adapter subspace is systematic rather than tied to a specific configuration.










% % \textbf{LoRA rank}

% % \begin{figure}[htbp]
% %     \centering
% %     \begin{minipage}{0.495\textwidth}
% %         \centering
% %         \includegraphics[width=0.9\linewidth]{}
% %         \subcaption{Average Accuracy (ACC)  on ImageNet-R}
% %         \label{fig:curves_imageR_1_AAA}
% %     \end{minipage}%
% %     \begin{minipage}{0.495\textwidth}
% %         \centering
% %         \includegraphics[width=0.9\linewidth]{}
% %         \subcaption{Nearest Mean of Exemplars (NME) on ImageNet-R}
% %         \label{fig:curves_imageR_2_NME}
% %     \end{minipage}
% %     \caption{
% %         AAA and NME summary curves on ImageNet-R ($T=20, r=16$) with different rank. ons with reduced forgetting.
% %     }
% %     \label{fig:combined_curves_imageR}
% % \end{figure}



% % \textbf{Local vs Global perturbation Strength}



% \section{Conclusion}

% % \section{Conclusion}

% This work revisits parameter efficient continual learning through the geometry of LoRA adapter subspaces. Within a hierarchical PAC Bayesian framework, we derive a generalization bound in which the continual learning risk is controlled by an empirical term under Gaussian perturbations, an adapter subspace expected sharpness term $E\mbox{-}\mathrm{Sh}_\Delta$, the divergence $\mathrm{KL}(Q_\Delta \Vert P_\Delta)$, and the hyper posterior drift across tasks. By restricting updates to the LoRA manifold $W_\Delta$ and working in factor coordinates, the analysis identifies subspace sharpness as the relevant notion of flatness and links it to the spectrum of a LoRA restricted Fisher matrix and to the stability of feature representations.

% The framework is deliberately aligned with a PECL regime where the pretrained backbone and previously learned adapters remain frozen. It does not cover partial or periodic finetuning of the backbone or richer perturbation models beyond Gaussian noise. Extending the analysis to such settings is a natural direction for future work and may reveal additional interactions between subspace geometry and global updates. Even within its current scope, the results suggest several practical guidelines for LoRA based PECL: flatness control should primarily act in the trainable adapter subspace $W_\Delta$, the organisation of adapters can be viewed as a structural prior that shapes the Fisher spectrum and stability plasticity trade offs, and subspace sharpness together with Fisher based diagnostics provides a useful lens for interpreting and comparing parameter efficient continual learning methods.





% % The framework has limitations. Theoretical guarantees are established under the assumption that the pretrained backbone and previously learned adapters remain frozen, which aligns with common PECL practice but excludes finetuning regimes where part of the backbone is updated. The bound also focuses on a specific class of Gaussian perturbations and on a fixed adapter parameterization. Extending the analysis to settings with partial or periodic finetuning, dynamically expanding adapter subspaces, or alternative noise models would be a natural direction for future work and may reveal additional interactions between subspace geometry and global parameter updates.

% % Despite these restrictions, the results offer several conceptual lessons for the LoRA and PECL communities. First, they identify adapter subspace flatness as a principled target for regularization, suggesting that flatness control should primarily act in $W_\Delta$ rather than on the full parameter space. Second, they highlight the geometry of the LoRA organization itself as a meaningful design dimension: different adapter layouts induce different Fisher spectra and stability properties, and can therefore be seen as structural priors on how knowledge is organised over tasks. Third, they point to subspace sharpness and related Fisher based diagnostics as tools for interpreting and comparing PECL methods. Taken together, these insights deepen the theoretical understanding of LoRA based continual learning and provide a basis for more geometry aware designs in future parameter efficient architectures.




% % This work investigates parameter efficient continual learning through the geometry of LoRA adapter subspaces. Within a hierarchical PAC Bayesian framework, we derive a generalization bound in which the continual learning risk is controlled by an empirical term under Gaussian perturbations, an adapter subspace expected sharpness term $E\mbox{-}\mathrm{Sh}_\Delta$, the divergence $\mathrm{KL}(Q_\Delta \Vert P_\Delta)$, and the hyper posterior drift across tasks. By restricting updates to $W_\Delta$ and working in LoRA factor coordinates, we link this sharpness term to the spectrum of a LoRA restricted Fisher matrix and to representation stability measured in feature space.

% % % Experiments on ImageNet R and its corrupted and perturbed variants, provide a systematic stress test of this analysis. Across LoRA organisations, promoting flat minima within $W_\Delta$ improves final and averaged accuracy, reduces catastrophic forgetting, and stabilises features, while stronger suppression of the spectrum of $F_\Delta$ yields predictable gains, especially in compact adapter subspaces. In contrast, global perturbations that move the frozen backbone can degrade performance, which supports the view that flatness should be controlled primarily in the trainable adapter manifold. These findings suggest practical guidelines for LoRA based PECL: treat subspace sharpness metrics as informative diagnostics, design adapter geometry with both capacity and curvature in mind, and focus flatness oriented optimisation on the update subspace rather than the full parameter space.

% %  The analysis and experiments focus on LoRA adapters for vision transformers. Extending the framework to other parameter efficient mechanisms, more diverse architectures and task streams, and algorithms that explicitly regulate both $\mathrm{KL}(Q_\Delta \Vert P_\Delta)$ and hyper drift is an important direction for future work. Even so, the present results indicate that geometry aware flatness in the adapter subspace offers a useful organising principle for understanding and improving parameter efficient continual learning.



% % \subsubsection*{Author Contributions}
% % If you'd like to, you may include a section for author contributions as is done
% % in many journals. This is optional and at the discretion of the authors. Only add
% % this information once your submission is accepted and deanonymized. 

% % \subsubsection*{Acknowledgments}
% % Use unnumbered third level headings for the acknowledgments. All
% % acknowledgments, including those to funding agencies, go at the end of the paper.
% % Only add this information once your submission is accepted and deanonymized. 

\section{Experiments and Analysis}

% \subsection{Experimental Setup}



\noindent\textbf{Datasets}\quad
We evaluate PECL methods on vision benchmarks derived from ImageNet. ImageNet-R~\cite{ImagenetR_Hendrycks_2021_ICCV} serves as our main class-incremental benchmark following prior work~\cite{liangInfLoRAInterferenceFreeLowRank2024a}.
%and is split into $T=20$ tasks over $200$ classes
% To assess robustness under corruptions and perturbations, we use the Tiny ImageNet variants of ImageNet-C and ImageNet-P (200 classes)~\cite{hendrycks2018benchmarking}. 
% For brevity, we refer to these datasets as “ImageNet-C” and “ImageNet-P”.
To assess robustness under corruptions and perturbations, we use the Tiny ImageNet variants of ImageNet-C and ImageNet-P (200 classes)~\cite{hendrycks2018benchmarking}.
%, where the former applies 15 types of corruption types and the latter provides short perturbation sequences per image

% We further adopt Tiny ImageNet-C~\cite{hendrycks2018benchmarking} and Tiny ImageNet-P~\cite{hendrycks2018benchmarking} to evaluate robustness on common corruptions and perturbations. Tiny ImageNet-C is a corrupted variant of Tiny ImageNet with 200 classes containing 15 types of corruption. Tiny ImageNet-P provides short perturbation sequences for each image
%throughout the rest of the paper
% To assess robustness under corruptions and perturbations, we further use ImageNet-C and ImageNet-P~\cite{hendrycks2018benchmarking}. 
% ImageNet-C applies common corruptions to ImageNet validation images, while ImageNet-P provides short perturbation sequences per image.

{%
\setlength{\intextsep}{4pt}
\setlength{\textfloatsep}{4pt}
\setlength{\abovecaptionskip}{6pt}
\setlength{\belowcaptionskip}{4pt}
\captionsetup{skip=10pt}
\begin{figure*}[t]
\centering
\includegraphics[width=1\linewidth]{figures_TRML/weight_overlap_5fig.pdf}
\vspace{-15pt}
\caption{Evolution of the curvature projection ratio $c_t$ (left), amplification factor $\alpha_t$ (middle) and loss with perturbation $\hat{L}(W+\epsilon)$ and the average accuracy $\operatorname{Acc}_t$ (right) in seqLoRA.
% across the task sequence, together with the corresponding average accuracy $\operatorname{Acc}_t$ (dashed) in seqLoRA.
% under SGD, $\mathrm{AS}^{(0)}$, and $\mathrm{AS}^{(1)}$
% Evolution of the amplification factor $\alpha_t$ (top) and curvature projection ratio $c_t$ (bottom) under SGD and apply perturbation $\mathrm{AS}^{(1)}$ and $\mathrm{AS}^{(1)}$ in seqLoRA.
% Compared with SGD, $\mathrm{AS}^{(1)}$ produces consistently larger and increasing $\alpha_t$ and $c_t$ over tasks, indicating progressively stronger adapter-induced amplification and a growing concentration of task-conditioned curvature in the adapter-induced geometry.
$\alpha_t$ and $c_t$ are computed from the last layer of ViT-B/16.
% Evolution of the amplification factor $A_t$ and curvature projection ratio $C_t$ across tasks in SeqLoRA. The top panel shows $A_t$, measuring how strongly the adapter update $\Delta W_t$ amplifies directions relative to the backbone restricted to the same rank-$r$ subspace. The bottom panel shows $C_t$, comparing the concentration of dominant task-conditioned curvature modes in the adapter-induced subspace versus the backbone subspace. The weight matrices are taken from the  11th layer of ViT-B/16.
}
\label{fig:evolute}
\vspace{-15pt}
\end{figure*}
}


\noindent\textbf{Models}\quad 
All experiments use a ViT-B/16 backbone~\cite{dosovitskiy2020image} pre-trained on ImageNet-21K.
We attach LoRA adapters to each transformer block and consider three organizations: SeqLoRA, IncLoRA, and OLoRA~\cite{wang2023orthogonal}.

% Unless stated otherwise, PECL updates and perturbations are applied only to
% to the trainable LoRA coordinates of the current task, while the backbone and past adapters remain frozen. 





% To study perturbation {scope}, we consider a Flat-LoRA-style~\cite{li2024flatlora} as a \emph{full-space} baseline, which applies $\mathrm{RS}^{(0)}_S(W,\Sigma)$ perturbations over all model parameters.
To study perturbation \emph{scope}, we adopt {full-parameter} setting~\cite{li2024flatlora} by injecting isotropic Gaussian perturbation $\mathrm{RS}_S^{(0)}({W}, \sigma^2 
I)$ into all model parameters.
% , instead of restricting perturbations to the adapter subspace.
% To study perturbation {type}, we compare SAM~\cite{foretSharpnessAwareMinimizationEfficiently2021} ($\mathrm{AS}^{(0)}_S(W,\rho)$), RWP~\cite{pmlr-v151-bisla22a} ($\mathrm{RS}^{(0)}_S(W,\Sigma)$), and GAM~\cite{zhangGradientNormAware2023} ($\mathrm{AS}^{(1)}_S(W,\rho)$), all acting on the trainable adapters in PECL.
To study perturbation \emph{type}, we compare SAM~\cite{foretSharpnessAwareMinimizationEfficiently2021}, RWP~\cite{pmlr-v151-bisla22a}, and GAM~\cite{Zhang_GAM}, which optimize different sharpness objectives: SAM minimizes $\mathrm{AS}^{(0)}_S(W,\rho)$, RWP targets $\mathrm{RS}^{(0)}_S(W,\Sigma)$, and GAM targets $\mathrm{AS}^{(1)}_S(W,\rho)$. For a controlled comparison in PECL, all methods apply their perturbations only to the trainable LoRA coordinates.
% $\theta_\Delta$; the induced weight perturbation is $\Delta W(\theta_\Delta+\xi)-\Delta W(\theta_\Delta)$.
%all methods apply their perturbations only to the trainable adapters.
Unless otherwise stated, we fix the perturbation radius to $\rho=0.05$, following the standard empirical setting used in SAM, and use a dedicated radius sweep only when diagnosing the scope effect.




% based on the task accuracy matrix $a_{tj}$, where $a_{tj}$ denotes the accuracy on task $j$ after learning task $t$. 
\noindent\textbf{Metrics}\quad
The average class-incremental accuracy at step $t$ is
$
\operatorname{Acc}_t = \frac{1}{t} \sum_{j=1}^{t} a_{tj}
$,  where $a_{tj}$ denotes the accuracy on task $j$ after learning task $t$. 
We report accuracy $\operatorname{Acc}_t^{\mathrm{cls}}$ using the linear classifier and
% , and also evaluate representation quality using Nearest Class Mean (NCM) classifier, denoted by $\operatorname{Acc}_t^{\mathrm{nme}}$.
compute the corresponding Nearest-Class-Mean Classifier (NCM) results
% Nearest Mean of Exemplars (NME) metrics 
$(\operatorname{Acc}_t^{\mathrm{ncm}})$ using class prototypes in the feature space.
%, which isolates the feature representations from the classifier head
% We report class-incremental performance using the task-wise classification accuracy $\operatorname{Acc}_t^{\mathrm{cls}}$.
% In addition, we evaluate representation quality via Nearest Mean of Exemplars (NME), denoted by $\operatorname{Acc}_t^{\mathrm{nme}}$.
% This metric reduces the influence of the classifier head and provides a complementary view of how stable and transferable the learned representations are across tasks.
To quantify flatness, we estimate the largest eigenvalue $\lambda_{\max}$ and the trace $\mathrm{tr}(\cdot)$ of both the Hessian and the Fisher information matrix using Lanczos~\cite{ghorbaniInvestigationNeuralNet2019}.

\subsection{Result Analysis}





\noindent\textbf{Scope of perturbations.}\quad
To study perturbation scope in PECL, we inject isotropic Gaussian perturbation $\mathrm{RS}_S^{(0)}({W}, \sigma^2 
I)$ under two scopes: perturbations restricted to the task-permissible adapter coordinates versus perturbations defined over the full parameter space.
As shown in Table~\ref{tab:lora_or_full_cls—show}, restricting perturbations to the trainable adapters consistently improves average accuracy over the SGD baseline, whereas applying perturbations that include the frozen backbone substantially degrades performance, with the strongest drop observed for SeqLoRA.



% Let $U_{\Delta,t}$ and $ V_{\Delta,t}$ are the left- and right-singular matrices of the SVD decomposition of $\Delta W$, and $U_{\Delta,t}U_{\Delta,t}^\top W  V_{\Delta,t}^\top  V_{\Delta,t}$ gives the “projection” of $W $onto the subspace spanned by $\Delta W$. Following~\cite{hu2022lora}, we adopt the feature amplification factor
% $
% A_t =
% \frac{\left\lVert \Delta W_{t}\right\rVert_F}
% {\left\lVert U_{\Delta,t}^\top W  V_{\Delta,t}^\top\right\rVert_F}
% $
% to measure how much of task-specific directions are amplified by $\Delta W$. 
% Here $\|U_{\Delta,t}^\top W V_{\Delta,t}^\top\|_F$ measures the energy of $W$ restricted to the same bilinear subspace induced by $\Delta W_t$ and equals the Frobenius norm of the projected of $W$ onto the subspace spanned by $\Delta W$ $U_{\Delta,t}(U_{\Delta,t}^\top W V_{\Delta,t}^\top) V_{\Delta,t}$.
% where $|U_{\Delta,t}^\top W V_{\Delta,t}|_F$ measures the energy of $W$ restricted to the same bilinear subspace induced by $\Delta W_t$. Equivalently, since $U{\Delta,t}$ and $V_{\Delta,t}$ are orthonormal, this quantity equals the Frobenius norm of the projected weight $\Pi_{\Delta,t}(W)=U_{\Delta,t}(U_{\Delta,t}^\top W V_{\Delta,t})V_{\Delta,t}^\top$
% Here $(U_{\Delta,t},V_{\Delta,t})$ are the rank-$r$ left and right singular vector matrices of $\Delta W $
% left and right singular vectors of $\Delta W_t$ from its SVD~\cite{hu2022lora}.
% and $\Delta U_t \Delta U_t^\top W \Delta V_t\Delta V_t^\top$ corresponds to the projection of $W$ onto the subspace spanned by $\Delta W_t$.

{%
\setlength{\textfloatsep}{0pt}      % 页顶/页底 table 与正文的距离（最关键）
\setlength{\intextsep}{1pt}         % 若 table 被放在正文中时的上下距离（备用）
\setlength{\floatsep}{1pt}          % 两个浮动体之间的距离（备用）
\setlength{\abovecaptionskip}{2pt}  % caption 与表格之间
\setlength{\belowcaptionskip}{8pt}  % caption 下方到正文
\captionsetup{skip=4pt}           % 若你用 caption 包，可用它更细调
\begin{table}[t]
\vspace{-5pt}
\centering
\caption{Scope effect in PECL: reporting $\Delta\operatorname{Acc}^{\mathrm{cls}}$ (vs.\ SGD), adapter-only perturbations consistently improve accuracy, whereas full-parameter perturbations degrade performance.}
\label{tab:lora_or_full_cls—show}
\footnotesize
\setlength{\tabcolsep}{3pt}
\renewcommand{\arraystretch}{0.88}
\begin{tabular}{@{}c c c c c@{}}
\toprule
\textbf{Scope} & \textbf{Strength} & \textbf{SeqLoRA} & \textbf{IncLoRA} & \textbf{OLoRA} \\
\midrule
% Base & --   & 69.07 & 71.29 & 71.41 \\
% \midrule
\multirow{3}{*}{Full}
 & 0.05 & -3.50 & -2.00 & -2.62 \\
 & 0.10 & -4.82 & -4.74 & -5.53 \\
 & 0.15 & -6.72 & -6.84 & -6.49 \\
\midrule
\multirow{3}{*}{LoRA}
 & 0.05 & 0.27 & 0.79 & 0.12 \\
 & 0.10 & 0.70 & 0.86 & 0.28 \\
 & 0.15 & 0.41 & 0.77 & 0.27 \\
\bottomrule
\end{tabular}
\end{table}
\vspace{-5pt}
}% end local group


To understand why the perturbation scope is decisive, we examine how the learned update $\Delta W_t$ interacts with the frozen weight $W^\mathrm{froz}_t$ and how dominant task-conditioned curvature modes localize relative to the geometry induced by $\Delta W_t$.
% To understand why the perturbation scope is decisive, we further analyze how the learned adapter update $\Delta W_t$ interacts with the frozen backbone weight $W$ and the task-conditioned curvature.
% To elucidate why the perturbation scope is decisive, we examine where the dominant task-conditioned curvature modes localize relative to the bilinear subspaces induced by the learned update $\Delta W_t$ and by the frozen backbone weight $W$.
% To elucidate why the perturbation scope is decisive, we analyze how the learned adapter update $\Delta W_t$ interacts with the frozen backbone weight $W$ and the task-conditioned curvature.
% Let $\Delta W_t = U_{\Delta,t} S_{\Delta,t} V_{\Delta,t}^\top$ be the rank-$r$ truncated SVD, where
% Let $U_{\Delta,t}$ and $V_{\Delta,t}^\top$ are the left and right singular vectors of the SVD decomposition of $\Delta W$.
% and the "projection” of matrix $W$ onto the subspace spanned by $\Delta W$ is
% $U_{\Delta,t}(U_{\Delta,t}^\top W V_{\Delta,t})V_{\Delta,t}^\top.$
% Frobenius-orthogonal projection of $W$ onto the bilinear subspace induced by $(U_{\Delta,t},V_{\Delta,t})$ 
Following~\cite{hu2022lora}, we adopt the feature amplification factor
$
\alpha_t=\frac{\|\Delta W_t\|_F}{\|U_{\Delta,t}^\top W^{\mathrm{froz}}_t V_{\Delta,t}\|_F},
$ to measure how much of task-specific directions are amplified by $\Delta W$, where $U_{\Delta,t}$ and $V_{\Delta,t}^\top$ are the left and right singular vectors of the SVD decomposition of $\Delta W$.
Similarly, to quantify whether dominant task-conditioned curvature modes concentrate more on the subspace induced by the learned update $\Delta W_t$ than on a backbone reference, we define the curvature projection ratio
$
c_t=\frac{R_{\Delta,t}}{R_{W^{\mathrm{froz}_t},t}}.
$
Here $R_{\Delta,t}$ and $R_{W,t}$ are localization scores of the top-$k$ task-conditioned curvature modes on the update-induced and backbone-reference subspaces, respectively.
The precise definitions of $R_{\Delta,t}, R_{W,t}$ are provided in Appendix~\ref{app:weight_projection}.
Larger $c_t$ indicates stronger curvature localization on the update-induced subspace relative to the backbone reference.
% the projection operators, as well as 
% Similarly, to quantify whether dominant task-conditioned curvature modes concentrate more on the subspace induced by the learned update $\Delta W_t$ than on a backbone reference subspace, 
% % Similarly, to compare how strongly dominant task-conditioned curvature modes concentrate more onto the adapter-induced subspace than onto the backbone reference subspace,
% % the adapter-induced subspace,
% % direction project onto the adapter-induced subspace versus a backbone reference subspace,
% % to assess whether dominant curvature modes have larger projection onto the adapter-induced subspace than onto the backbone reference subspace, 
% we define the curvature projection ratio 
% % $
% % c_t=\frac{\sum_{i=1}^k \lambda_{t, i}\left\|U_{\Delta, t}^{\top} {v_{t, i}} V_{\Delta, t} \right\|_F^2}{\sum_{i=1}^k \lambda_{t, i}\left\|U_{W, t}^{\top} v_{t, i} V_{W, t} \right\|_F^2},
% % $
% $
% c_t=\frac{R_{\Delta,t}}{R_{W,t}}.
% $
% where $R_{\Delta,t}$ and $R_{W,t}$ summarize the eigenvalue-weighted localization of the top-$k$ task-conditioned curvature modes onto  the  subspace induced by the learned update $\Delta W_t$ and a backbone reference subspace.
% after mapping each curvature eigen-direction to its effective weight-space perturbation. 
% The precise definitions of $R_{\Delta,t}, R_{W,t}$ and the projection operators, as well as the implementation details, are provided in Appendix~\ref{app:weight_projection}. 
% Larger $c_t$ indicates that dominant curvature modes exhibit stronger projected energy within the subspace induced by $\Delta W_t$ than within the backbone reference subspace.
% where ${(\lambda_{t,i},v_{t,i})}_{i=1}^k$ are the top-$k$ curvature eigenpairs at task $t$ in the full parameter space (Appendix~\ref{app:weight_projection}). 
% Larger $c_t$ indicates that dominant curvature modes have higher projected energy within subspace induced by $\Delta W_t$ than within the backbone reference subspace. 
% Moreover, let $\{(\lambda_{t,i},v_{t,i})\}_{i=1}^k$ be the top-$k$ eigenpairs of a positive semidefinite curvature operator at task $t$ computed in the full parameter space. We quantify curvature localization by eigenvalue-weighted projected-energy ratios onto $\mathcal S_{\Delta,t}$ and onto a backbone reference subspace, and report their preference ratio $c_t$ (Appendix~\ref{weight projection}).

% Larger $C_t$ indicate that the dominant curvature modes align more with the adapter-induced subspace spanned by  $(U_{\Delta,t},V_{\Delta,t})$.

Fig.~\ref{fig:evolute} shows consistent trends across the task sequence and aligns closely with the evolution of accuracy.
As the task index increases, the curvature projection ratio $c_t$ increases over tasks and remains above the reference level $c_t\approx 1$, suggesting that dominant task-conditioned curvature modes are concentrated in the adapter-induced subspace rather than in the backbone subspace.
The amplification factor $\alpha_t$ also rises, indicating that the learned update $\Delta W_t$ progressively amplifies task-specific directions.
Crucially, this structural strengthening is accompanied by the evolution of the perturbation-aware loss $\hat L+\mathrm{RS}^{(0),\Delta}$ and aligns with the degradation of $\operatorname{Acc}_t$.



{%
% {%
%
% \setlength{\intextsep}{4pt}
\setlength{\textfloatsep}{4pt}
% \setlength{\abovecaptionskip}{6pt}
% \setlength{\belowcaptionskip}{pt}
% \captionsetup{skip=10pt}
\begin{figure}[ht]
\centering
\includegraphics[width=0.95\linewidth]{figures_TRML/ACC_dataset_pair_Robustness_to_Corruptions_on_ImageNet-C_Robustness_to_Perturbation_on_ImageNet-P2.pdf}
% \includegraphics[width=0.9\linewidth]{figures_TRML/hebing_robutness2.pdf}

\vspace{-6pt}
\caption{\textbf{Robustness of perturbation types under CL} Classification average accuracy $\operatorname{Acc}_t^{\mathrm{cls}}$ over tasks evaluated under distribution shift on ImageNet-C  and ImageNet-P.
Across tasks and LoRA organizations, both $\mathrm{AS}^{(0)}$ and $\mathrm{AS}^{(1)}$ consistently outperform SGD, and the gap widens in later tasks. Here ImageNet-C/P denote the Tiny ImageNet variants; see Experimental Setup.}
\label{fig:robustness_imagenet_cp}
\vspace{-10pt}
\end{figure}
}%

Taken together, these diagnostics offer a mechanism-level account of Table~\ref{tab:lora_or_full_cls—show}. As training progresses, both the effective adaptation directions and the task-conditioned curvature that governs sensitivity concentrate within the LoRA-permissible coordinates. Perturbations restricted to the adapter scope therefore act on the controllable geometry and can be counteracted by permissible LoRA updates, whereas full-scope perturbations introduce variations in the frozen model outside the permissible update subspace, creating an irreducible scope mismatch that accelerates performance degradation under the PECL constraint.
% Moreover, the perturbation type shapes the covariance geometry in LoRA coordinates, which steers the evolution of $\Delta W_t$ and biases the full-model solution $W_0+\Delta W_t$ toward flatter task-conditioned regions associated with better continual accuracy.

% Taken together, these diagnostics provide a mechanism-level explanation for Table~\ref{tab:lora_or_full_cls—show}:
% both the effective adaptation directions and the task-conditioned curvature that governs sensitivity become increasingly concentrated within the LoRA-admissible coordinates; therefore, perturbations restricted to the adapter scope act on the same controllable geometry whereas full-scope perturbations inject variations into frozen backbone that are outside the admissible update subspace, creating a scope mismatch that cannot be compensated by admissible LoRA updates under PECL constraint and accelerates performance degradation.

%  Under a fixed PECL scope, the perturbation design acts through the geometry within LoRA coordinates, which determines how $\Delta W_t$ evolves and, consequently, how the full-model solution $W_0+\Delta W_t$ bias toward flatter task-conditioned regions.


% Taken together, these diagnostics provide a mechanism-level explanation for Table~\ref{tab:lora_or_full_cls—show}: both adaptation and the curvature that governs sensitivity are increasingly localized within the LoRA-admissible coordinates in CL, so restricting perturbations to the adapter scope matches the controllable geometry, whereas full-scope perturbations inject variations into frozen backbone that are outside the admissible update subspace, creating a scope mismatch that cannot be compensated by admissible LoRA updates under PECL constraint and accelerates performance degradation.

% This strengthening of $\alpha_t$ and $c_t$ coincides with the growing performance gap: methods that maintain larger and increasing $\alpha_t$ and $c_t$ exhibit less loss and more stable $\operatorname{Acc}_t$,


% Fig.~\ref{fig:evolute} shows consistent trends across the task sequence: $\alpha_t$ increases with the task index, indicating that $\Delta W_t$ progressively strengthens task-specific directions. $C_t$ also increases over tasks and remains above the neutral reference level $c_t\approx 1$, suggesting that the dominant task-conditioned curvature becomes increasingly concentrated in the adapter-induced subspace rather than in the backbone subspace.






\noindent\textbf{Type of perturbations.}\quad
To isolate the effect of perturbation \emph{type} under PECL constraints, we apply three flatness optimizers to the \emph{trainable LoRA parameters only}:
$\mathrm{RS}^{(0)}_S(W,\Sigma)$, $\mathrm{AS}^{(0)}S(W,\rho)$, and $\mathrm{AS}^{(1)}S(W,\rho)$, with SGD as the baseline.
Across tasks on ImageNet-C and ImageNet-P (Fig.~\ref{fig:robustness_imagenet_cp}), both adversarial objectives consistently outperform SGD across the task sequence, and the gap widens in later tasks, indicating improved stability under continual adaptation. This trend is consistent with the final class-incremental results on ImageNet-R (Table~\ref{table:q1_ci_with_weight_flatness_reorg}), where moving from SGD to $\mathrm{RS}^{(0)}$, $\mathrm{AS}^{(0)}$, and then $\mathrm{AS}^{(1)}$ yields progressively higher $\operatorname{Acc}^{\mathrm{cls}}_{20}$ and $\operatorname{Acc}^{\mathrm{nme}}_{20}$.


% $\mathrm{AS}^{(0)}$ and $\mathrm{AS}^{(1)}$ outperform SGD throughout the sequence, and the advantage becomes more pronounced in later tasks, indicating improved stability under continual adaptation.
% This trend is consistent with the final class-incremental results in Table~\ref{table:q1_ci_with_weight_flatness_reorg}: across LoRA organizations, moving from SGD to $\mathrm{RS}^{(0)}$, $\mathrm{AS}^{(0)}$, and then $\mathrm{AS}^{(1)}$ yields progressively higher $\operatorname{ACC}{20}^{\mathrm{cls}}$ and $\operatorname{ACC}{20}^{\mathrm{nme}}$.
% Overall, stronger perturbation types lead to systematically better continual performance even when the update scope is fixed to LoRA, suggesting that the \emph{geometry} of perturbations inside the adapter coordinates matters beyond mere subspace restriction.


% To isolate the effect of perturbation \emph{type} under PECL constraints, we apply three representative flatness optimizers to the {trainable LoRA parameters only}: 
% $\mathrm{RS}^{(0)}_S(W,\Sigma)$, $\mathrm{AS}^{(0)}_S(W,\rho)$, and $\mathrm{AS}^{(1)}_S(W,\rho)$, with SGD as the baseline.
% Across tasks on ImageNet-C and ImageNet-P (Fig.~\ref{fig:robustness_imagenet_cp}), we observe a consistent performance ordering, where both $\mathrm{AS}^{(0)}$ and $\mathrm{AS}^{(1)}$ outperform SGD throughout the sequence, and the margin becomes more pronounced in later tasks, indicating improved stability under continual adaptation.
% % Fig.~\ref{fig:robustness_imagenet_cp} shows a consistent ordering across tasks on ImageNet-C and ImageNet-P: both $\mathrm{AS}^{(0)}$ and $\mathrm{AS}^{(1)}$ dominate SGD throughout the sequence, and the gap becomes more visible in later tasks, indicating improved stability under continual adaptation. 
% The same ordering is reflected by the final class-incremental performance and weight-space geometry in Table~\ref{table:q1_ci_with_weight_flatness_reorg}. Across LoRA configurations, moving from SGD to $\mathrm{RS}^{(0)}$, $\mathrm{AS}^{(0)}$, and then $\mathrm{AS}^{(1)}$ yields progressively higher $\operatorname{Acc}_{20}^{\mathrm{cls}}$ and $\operatorname{Acc}_{20}^{\mathrm{nme}}$, while simultaneously reducing sharpness and curvature diagnostics, including $\lambda{\max}(H)$ and $\operatorname{tr}(H)$, and similarly for Fisher. Overall, perturbation types that more strongly suppress adapter-subspace sharpness achieve better continual performance, consistent with the sharpness-dependent empirical term in Theorem~\ref{thm:pecl-bound}.


% We next examine how perturbation \emph{type} changes curvature and the loss surface under the same PECL scope.
% Quantitative diagnostics in Table~\ref{table:q1_ci_with_weight_flatness_reorg} show that stronger optimizers reduce sharpness and curvature statistics measured on the trainable LoRA parameters, including $\lambda_{\max}(H)$ and $\operatorname{tr}(H)$, and similarly for Fisher.
% Interestingly, a one-dimensional probe {within the LoRA subspace} (top of Fig.~\ref{fig:perturb_type_landscape}) can yield loss curves that appear comparably flat, which by itself may understate the differences among perturbation types.
% However, when the same trained models are probed under {full-parameter} perturbations (bottom of Fig.~\ref{fig:perturb_type_landscape}), the profiles separate markedly, revealing distinct sensitivity of the effective model $W_0+\Delta W$ in the full space.
% Taken together, these results indicate that, even though learning is restricted to LoRA coordinates, different perturbation geometries reshape the \emph{full-model} landscape of $W_0+\Delta W$ for task, consistent with the sharpness-dependent empirical term in Theorem~\ref{thm:pecl-bound}.

% Under the same LoRA-restricted update scope, perturbation type still induces distinct curvature and loss-surface properties. Table~\ref{table:q1_ci_with_weight_flatness_reorg} shows that stronger objectives reduce LoRA-restricted sharpness and curvature diagnostics (Hessian and Fisher). Together, these results support the bound perspective that, at a fixed scope, type acts primarily through the perturbation geometry inside $\mathcal W_{\Delta,t}$, thereby reshaping the effective landscape.
% Beyond accuracy, perturbation type induces a coherent set of geometric effects.
% % under the same LoRA-restricted update scope.
% As predicted by the bound, changing type modifies the perturbation geometry within $\mathcal W_{\Delta,t}$ and primarily acts by suppressing task-conditioned sharpness on the permissible parameters.
% Empirically, Table~\ref{table:q1_ci_with_weight_flatness_reorg} shows a consistent ordering across objectives: stronger types yield smaller LoRA-restricted sharpness and lower curvature diagnostics (Hessian and Fisher).
% This reduction in subspace sharpness is accompanied by a corresponding change in the \emph{full-parameter} landscape.
% As visualized in Fig.~\ref{fig:scope_type2}, the solutions obtained with increasingly adversarial objectives exhibit flatter loss profiles under the full-parameter probe along the leading eigen-direction of the full curvature.
% , demonstrating reduced global sensitivity of $W_0+\Delta W_t$.
% Taken together, the table and the loss-landscape visualization form a consistent chain of evidence that perturbation type, through lowering sharpness within $\mathcal W_{\Delta,t}$, reshapes both restricted curvature and the resulting full-model geometry, which explains the observed performance gains in continual adaptation.
Beyond accuracy, perturbation type induces distinct curvature and landscape properties under the restricted update scope.
Table~\ref{table:q1_ci_with_weight_flatness_reorg} shows that moving from $\mathrm{RS}^{(0)}$ to $\mathrm{AS}^{(0)}$ and $\mathrm{AS}^{(1)}$ progressively reduces the LoRA-restricted sharpness and curvature diagnostics (Hessian and Fisher), indicating that more concentrated perturbation geometries better suppress task-conditioned sharpness within $\mathcal W_{\Delta,t}$.
Consistently, Fig.~\ref{fig:scope_type2} visualizes the resulting {full-parameter} landscape: solutions trained with stronger objectives exhibit flatter profiles along the leading eigen-direction of the full curvature.

% Together, these results support the bound perspective that, type acts primarily through the perturbation geometry inside $\mathcal W_{\Delta,t}$, which lowers task-conditioned sharpness and translates into flatter full-model solutions, thereby improving continual robustness and accuracy.


% Fig.~\ref{fig:perturb_type_landscape} and Fig.~\ref{fig:evolute} further clarify \emph{how} different perturbation geometries translate into different full-model solutions even when updates are restricted to LoRA.
% Within a LoRA-restricted slice (top of Fig.~\ref{fig:perturb_type_landscape}), the loss profiles can appear similarly flat across perturbation types. However, evaluating the same model under full-parameter perturbations (bottom of Fig.~\ref{fig:perturb_type_landscape}) separates the curves, revealing distinct sensitivity of the effective model $W_0+\Delta W$. 
% Along a one-dimensional direction \emph{within the LoRA subspace} (top of Fig.~\ref{fig:perturb_type_landscape}), the loss curves can look uniformly flat across perturbation types.
% However, the quantitative flatness diagnostics in Table~\ref{table:q1_ci_with_weight_flatness_reorg} reveal clear differences even under the same LoRA restriction: stronger flatness optimizers achieve smaller sharpness and reduced curvature statistics when curvature is evaluated on the trainable LoRA parameters.
% Moreover, when the same trained models are probed under \emph{full-parameter} perturbations (bottom of Fig.~\ref{fig:perturb_type_landscape}), the profiles separate markedly, exposing distinct sensitivity of the effective model $W_0+\Delta W$ in the full parameter space.
% This separation is consistent with the trends of sharpness and projected-curvature statistics reported in Table~\ref{table:q1_ci_with_weight_flatness_reorg}, where stronger flatness optimizers achieve smaller full-model curvature measures and exhibit flatter full-space loss profiles.

% \textbf{why $\mathrm{AS}^{(1)}$ is so better in LoRA based CL?}

The most pronounced effect is the amplified advantage of $\mathrm{AS}^{(1)}$ in LoRA-based PECL. Unlike $\mathrm{AS}^{(0)}$, $\mathrm{AS}^{(1)}$ explicitly penalizes the neighborhood gradient norm and is therefore more sensitive to the dominant modes of the task-conditioned restricted curvature $C_{\Delta,t}$. In a low-dimensional, structured adapter subspace, this curvature-aware control is easier to realize and has larger impact: in SeqLoRA, only $\mathrm{AS}^{(1)}$ induces a persistent increase in the localization ratios $(c_t,\alpha_t)$ (Fig.\ref{fig:evolute}), which coincides with the lowest loss and most stable $\operatorname{Acc}_t$. This coupling suggests that $\mathrm{AS}^{(1)}$ preferentially shapes the few permissible directions, simultaneously concentrating task-relevant curvature into the adapter-induced geometry and steering the effective model toward flatter task-conditioned regions. Consequently, under a fixed PECL scope, perturbation type determines the geometry of perturbations within LoRA coordinates, which in turn determines $\Delta W_t$ and the full-model loss landscape of $W_0+\Delta W_t$, explaining the consistent performance ordering.

% The most striking observation is the amplified advantage of $\mathrm{AS}^{(1)}$ in LoRA-based PECL.
% Compared with SGD and the other perturbation types, $\mathrm{AS}^{(1)}$ achieves the best class-incremental accuracy, the smallest LoRA-restricted sharpness and curvature diagnostics (Table~\ref{table:q1_ci_with_weight_flatness_reorg}), and the flattest full-space loss profiles (Fig.~\ref{fig:perturb_type_landscape}), with a consistently better trajectory across tasks (Fig.~\ref{fig:robustness_imagenet_cp}).
% Notably, such a pronounced separation is less evident under standard full fine-tuning in the original study~\cite{Zhang_GAM}, suggesting that first-order sharpness control becomes particularly effective when adaptation is \emph{subspace-constrained}.
% To understand this amplification, we compare $\mathrm{AS}^{(1)}$ against $\mathrm{AS}^{(0)}$ and SGD in the SeqLoRA setting.
% Among all methods, only $\mathrm{AS}^{(1)}$ induces a persistent increase in the curvature localization ratios $c_t$ and $\alpha_t$, which coincides with the highest and most stable $\operatorname{ACC}t$.
% This coupling supports the following mechanism: by explicitly penalizing the neighborhood gradient norm, $\mathrm{AS}^{(1)}$ is more sensitive to dominant modes of the task-conditioned restricted curvature $C{\Delta,t}$, and thus preferentially shapes the few admissible directions that most strongly affect $W_0+\Delta W$.
% Under LoRA, where optimization is restricted to a low-dimensional structured subspace, this curvature-aware control is easier to realize and has larger consequences: it simultaneously concentrates task-relevant curvature into the adapter-induced subspace and steers the effective model toward flatter task-conditioned regions.


% Consequently, at a fixed PECL scope, perturbation \emph{type} determines the perturbation geometry within LoRA coordinates, which in turn determines the learned updates $\Delta W_t$ and reshapes the \emph{full-model} loss landscape of $W_0+\Delta W_t$, ultimately driving the observed performance ordering.

% In particular, compared with SGD and the other perturbation types, $\mathrm{AS}^{(1)}$ exhibits a markedly stronger advantage in our LoRA-based PECL setting.
% This gap is consistently visible across multiple levels of evidence: $\mathrm{AS}^{(1)}$ achieves the best class-incremental accuracy and the smallest sharpness and curvature diagnostics on the trainable LoRA parameters in Table\ref{table:q1_ci_with_weight_flatness_reorg}, the most flat lossland in Fig\ref{fig:perturb_type_landscape}. and a more favorable accuracy trajectory over the entire task sequence in Fig\ref{fig:robustness_imagenet_c}. Notably, such a pronounced separation is not as evident in the originalstudy under standard full fine-tuning \cite{Zhang_GAM}, which suggests that the benefit of first-order sharpness can be substantially amplified under subspace-constrained adaptation.
% To figure it out, we compare $\mathrm{AS}^{(1)}$ with  $\mathrm{AS}^{(0)}$ and SGD in basic seqLoRA setting. Among all methods, only $\mathrm{AS}^{(1)}$ drives a persistent increase in the curvature localization ratio $c_t$ and $\alpha_t$, coincides with the highest and most stable $\operatorname{Acc}_t$. This indicates that $\mathrm{AS}^{(1)}$ not only concentrates dominant task-conditioned curvature more strongly into the adapter-induced subspace, but also drives optimization toward flatter task-conditioned solutions in the effective model $W_0+\Delta W$. Consequently, $\mathrm{AS}^{(1)}$ learns updates that are both more task-specific and better regularized in the full-model geometry, which is consistent with its consistently improved continual-learning performance. The influence of local curvature is easier to control, and therefore yields larger gains, when adaptation is confined to a small set of trainable coordinates. Under LoRA, where optimization is restricted to a low-dimensional, structured subspace, curvature-aware updates can be aligned with the few admissible high-impact directions, making curvature shaping both feasible and consequentia These results show that, under a fixed PECL scope, the choice of perturbation \emph{type} controls the perturbation geometry within the LoRA coordinates, which in turn determines the resulting $\Delta W_t$ and reshapes the \emph{full-model} landscape of $W_0+\Delta W_t$



% , indicating that the dominant task-conditioned curvature becomes progressively more concentrated in the adapter-induced subspace where $\Delta W_t$ operates.
% At the same time, $\mathrm{AS}^{(1)}$ maintains the max amplification factor $\alpha_t$ throughout training, sugggetings it learen task 


% compared with SGD, $\mathrm{AS}^{(1)}$ yields a larger and steadily increasing curvature localization ratio $c_t$, together with a noticeably flatter loss profile under full-parameter evaluation. This indicates that $\mathrm{AS}^{(1)}$ not only concentrates dominant task-conditioned curvature more strongly into the adapter-induced subspace, but also drives optimization toward flatter task-conditioned solutions in the effective model $W_0+\Delta W$. Consequently, $\mathrm{AS}^{(1)}$ learns updates that are both more task-specific and better regularized in the full-model geometry, which is consistent with its consistently improved continual-learning performance.
% These results show that, under a fixed PECL scope, the choice of perturbation \emph{type} controls the perturbation geometry within the LoRA coordinates, which in turn determines the resulting $\Delta W_t$ and reshapes the \emph{full-model} landscape of $W_0+\Delta W_t$.


% \noindent\textbf{why }\quad

% \noindent\textbf{Interaction with LoRA organization.}\quad
% We further investigate how the effect of perturbation {type} varies with the LoRA organization. 
% As shown in Table~\ref{table:q1_ci_with_weight_flatness_reorg} and  Fig.~\ref{fig:robustness_imagenet_cp}, the ordering of perturbation types is stable across LoRA organizations, but the performance gains vary in magnitude, indicating that perturbation geometry interacts with the organization-defined admissible update directions.
% In particular, SeqLoRA remains more robust on ImageNet-C/P, whereas IncLoRA and OLoRA exhibit larger accuracy drops under shift.
% Since the sharpness term is defined on the task adapter subspace, it is inherently sensitive to how different LoRA organization shapes the admissible scope $\mathcal W_{\Delta,t}$ across tasks. When $\mathcal W_{\Delta,t}$ is larger or expands over time, perturbations can probe a broader set of admissible directions, including more directions with non-negligible projected curvature.


\subsection{Ablation study on task length and adapter rank}
\label{sec:results-ablation}


{%
% \setlength{\intextsep}{4pt}        % [h] 浮动体与正文的上下间距（最关键）
% \setlength{\textfloatsep}{4pt}     % [t]/[b] 时浮动体与正文间距（备用）
% \setlength{\abovecaptionskip}{1pt} % 图与caption之间（上方）间距
% \setlength{\belowcaptionskip}{0pt} % caption下方到正文的间距
% \captionsetup{skip=3pt}            % 你已在用：图与caption的距离（更细粒度）
\begin{figure}[t]
    \centering
\includegraphics[width=1\linewidth]{figures_TRML/ablation_rank_length_imagenetr_NME_T.pdf}
\vspace{-10pt}
    \caption{ Sensitivity to Adapter Capacity and Sequence Difficulty
    % Ablation on adapter rank and task sequence length on ImageNet-R.
% We report the final NME-based accuracy $\operatorname{Acc}_T^{\mathrm{nme}}$ while varying (a) the adapter rank $r$ and (b) the task length $T$.
    }
    \label{fig:ablation}
\end{figure}
}

\vspace{-4pt}

Ablation on adapter rank and task sequence length on ImageNet-R.
We study how adapter capacity and sequence difficulty affect PECL by varying the LoRA rank $r$ and the task sequence length $T$, and report the final NME accuracy $\operatorname{Acc}_T^{\mathrm{nme}}$ in Fig.~\ref{fig:ablation}.
Across all settings, the relative ordering is stable: flatness-aware optimizers consistently outperform SGD.
Accuracy improves from small $r$ but quickly saturates at moderate ranks. This supports our theoretical view that PECL primarily requires a low-dimensional task-permissible update subspace to capture task-specific directions, so enlarging $r$ beyond this intrinsic dimension adds capacity with limited benefit.
%while leaving the bound-relevant geometry largely unchanged




\section{Conclusion}

We develop a sequential hierarchical PAC-Bayes framework for LoRA-based PECL, showing that both the sharpness-dependent empirical term and the KL complexity are supported on the task-permissible adapter subspace. Guided by the bound, we decomposed perturbation design into scope and type and showed empirically that subspace-aligned perturbations, especially those emphasizing curvature-aware directions, reduce perturbation-aware risk along the task sequence and translate into higher continual accuracy.

\section{Ref}

\bibliography{main}
\bibliographystyle{tmlr}

\appendix


% \section{Appendix Proof}

% \subsection{Proof of Theorem~\ref{thm:pathwise-pacbayes}}

% \begin{theorem}[Pathwise PAC--Bayes with time-varying hyper-posterior]
% For any $\delta\in(0,1]$, with probability at least $1-\delta$ over $S_{1:T}$,
% \begin{equation}
% \begin{aligned} 
% & \frac{1}{T} \sum_{t=1}^T \underbrace{\mathbb{E}_{P_t \sim \mathcal{P}_t} \mathbb{E}_{W \sim Q_t(\cdot \mid P_t, S_t)} L_{\mathcal{D}_t}(W)}_{\text{test risk at step } t} 
% \le \ \frac{1}{T} \sum_{t=1}^T \underbrace{\mathbb{E}_{P_t \sim \mathcal{P}_t} \mathbb{E}_{W \sim Q_t(\cdot \mid P_t, S_t)} \hat{L}_{S_t}(W)}_{\text{empirical risk at step } t} \\ 
% & + \sqrt{\frac{ \sum_{t=1}^T \mathbb{E}_{P_t \sim \mathcal{P}_t} \mathrm{KL}\big(Q_t(\cdot \mid P_t, S_t) \,\|\, P_t\big) + \sum_{t=1}^T \mathrm{KL}(\mathcal{P}_t \,\|\, \mathcal{P}_{t-1}) + \log\frac{1}{\delta} }{2 \sum_{t=1}^T (m_t - 1)}}. 
% \end{aligned} 
% \end{equation}
% \end{theorem}


% \begin{proof}
    


% % \subsection*{A.1\quad Two-level variational inequalities}
% Let $B_T:=\sum_{t=1}^T (m_t-1)$ and define $\alpha_t:=(m_t-1)/B_T$ so that $\sum_{t=1}^T\alpha_t=1$.
% % \paragraph{Goal quantity.}
% Define the step-$t$ generalization gap
% \[
% \Delta_t
% :=\mathbb E_{P_t\sim\mathcal P_{t-1}}\,
% \mathbb E_{W\sim Q_t(\cdot\mid P_t,S_t)}
% \big(L_{\mathcal D_t}(W)-\hat L_{S_t}(W)\big),
% \]
% and the size-weighted gap $\bar\Delta:=\sum_{t=1}^T \alpha_t \Delta_t$.
% Fix $\lambda>0$ and set
% \(
% g_t(W):=\lambda\alpha_t\big(L_{\mathcal D_t}(W)-\hat L_{S_t}(W)\big).
% \)

% \paragraph{Step 1 Model level.}
% For any fixed $P$ and $S_t$,
% \begin{equation}
% \label{eq:dv-w}
% \mathbb E_{W\sim Q_t(\cdot\mid P,S_t)} g_t(W)
% \;\le\;
% \mathrm{KL}\!\big(Q_t(\cdot\mid P,S_t)\,\|\,P\big)
% +\log\mathbb E_{W\sim P}\!\big[e^{g_t(W)}\big].
% \end{equation}
% Taking $\mathbb E_{S_t}$ and using the standard Hoeffding--MGF bound for bounded losses,
% \begin{equation}
% \label{eq:mgf-mean-gap}
% \mathbb E_{S_t}\log\mathbb E_{W\sim P} e^{g_t(W)}
% \;\le\; \frac{(\lambda\alpha_t)^2}{8\,(m_t-1)}.
% \end{equation}
% Therefore,
% \begin{equation}
% \label{eq:model-step}
% \mathbb E_{S_t}\mathbb E_{W\sim Q_t(\cdot\mid P,S_t)} g_t(W)
% \;\le\;
% \mathbb E_{S_t}\mathrm{KL}\!\big(Q_t(\cdot\mid P,S_t)\,\|\,P\big)
% +\frac{(\lambda\alpha_t)^2}{8\,(m_t-1)}.
% \end{equation}
%  % (change of measure in $P$-space)
% \paragraph{Step 2 Hyper level.} 
% Define
% \[
% h_t(P):=\mathbb E_{S_t}\mathbb E_{W\sim Q_t(\cdot\mid P,S_t)} g_t(W)
% -\mathbb E_{S_t}\mathrm{KL}\!\big(Q_t(\cdot\mid P,S_t)\,\|\,P\big).
% \]
% By \eqref{eq:model-step}, $h_t(P)\le \frac{(\lambda\alpha_t)^2}{8\,(m_t-1)}$ for all $P$.
% Applying DV in $P$-space to change the measure from $\mathcal P_{t-1}$ to $\mathcal P_t$ gives
% \begin{equation}
% \label{eq:dv-p}
% \log\mathbb E_{P\sim\mathcal P_t}\!\big[e^{h_t(P)}\big]
% \;\le\; \mathrm{KL}(\mathcal P_t\|\mathcal P_{t-1})
% +\log\mathbb E_{P\sim\mathcal P_{t-1}}\!\big[e^{h_t(P)}\big]
% \;\le\; \mathrm{KL}(\mathcal P_t\|\mathcal P_{t-1})
% +\frac{(\lambda\alpha_t)^2}{8\,(m_t-1)}.
% \end{equation}
% Using $\log\mathbb E[e^X]\ge \mathbb E[X]$ yields
% \begin{equation}
% \label{eq:p-expect}
% \mathbb E_{P_t\sim\mathcal P_t} h_t(P_t)
% \;\le\; \mathrm{KL}(\mathcal P_t\|\mathcal P_{t-1})
% +\frac{(\lambda\alpha_t)^2}{8\,(m_t-1)}.
% \end{equation}
% Unfolding $h_t(P)$ and interchanging expectations (Fubini),
% \begin{equation}
% \label{eq:two-term}
% \mathbb E_{S_t}\Big[
% \lambda\alpha_t\Delta_t
% -\alpha_t\,\mathbb E_{P_t\sim\mathcal P_{t-1}}\mathrm{KL}\!\big(Q_t(\cdot\mid P_t,S_t)\,\|\,P_t\big)
% \Big]
% \;\le\;
% \mathrm{KL}(\mathcal P_t\|\mathcal P_{t-1})
% +\frac{(\lambda\alpha_t)^2}{8\,(m_t-1)}.
% \end{equation}

% % \subsection*{A.2\quad Conditional MGF and supermartingale}
% \paragraph{Step 3 Martingale}
% From \eqref{eq:two-term} and the tower property, one obtains the conditional exponential bound
% \begin{equation}
% \label{eq:cond-mgf}
% \mathbb E\!\left[
% \exp\!\Big(
% \lambda\alpha_t\Delta_t
% -\alpha_t\,\mathbb E_{P_t\sim\mathcal P_{t-1}}\mathrm{KL}(Q_t\|P_t)
% -\mathrm{KL}(\mathcal P_t\|\mathcal P_{t-1})
% -\frac{(\lambda\alpha_t)^2}{8\,(m_t-1)}
% \Big)\,\Big|\,\mathcal F_{t-1}\right]\;\le\;1.
% \end{equation}
% Define $X_t$ to be the exponential increment inside \eqref{eq:cond-mgf} and
% \[
% \mathcal{M}_0:=1,\qquad
% \mathcal{M}_t:=\prod_{s=1}^t X_s.
% \]
% Then $(\mathcal{M}_t)_{t\ge0}$ is a nonnegative supermartingale, so
% $\mathbb E[\mathcal{M}_T]\le 1$.
% By Markov's inequality, with probability at least $1-\delta$,
% \begin{equation}
% \label{eq:hpp}
% \sum_{t=1}^T \lambda\alpha_t\Delta_t
% \;\le\;
% \sum_{t=1}^T \alpha_t\,\mathbb E_{P_t\sim\mathcal P_{t-1}}\mathrm{KL}(Q_t\|P_t)
% \;+\; \sum_{t=1}^T \mathrm{KL}(\mathcal P_t\|\mathcal P_{t-1})
% \;+\; \frac{\lambda^2}{8}\sum_{t=1}^T \frac{\alpha_t^2}{(m_t-1)}
% \;+\; \log\tfrac1\delta.
% \end{equation}

% % \subsection*{A.3\quad Size-weighted self-normalization and optimization}

% By $\alpha_t=(m_t-1)/B_T$,
% \[
% \sum_{t=1}^T \frac{\alpha_t^2}{(m_t-1)}
% =\sum_{t=1}^T \frac{(m_t-1)^2/B_T^2}{(m_t-1)}
% =\frac{1}{B_T}.
% \]
% Hence, \eqref{eq:hpp} reads, for all $\lambda>0$ and with probability at least $1-\delta$,
% \begin{equation}
% \label{eq:gap-preopt}
% \bar\Delta
% \;\le\;
% \frac{1}{\lambda}\Big(
% \sum_{t=1}^T \alpha_t\,\mathbb E_{P_t\sim\mathcal P_{t-1}}\mathrm{KL}(Q_t\|P_t)
% + \sum_{t=1}^T \mathrm{KL}(\mathcal P_t\|\mathcal P_{t-1})
% + \log\tfrac1\delta\Big)
% \;+\;
% \frac{\lambda}{8B_T}.
% \end{equation}
% Optimizing the right-hand side over $\lambda>0$ (take $\lambda^\star=\sqrt{8B_T\,A}$ with
% $A:=\sum_{t=1}^T \alpha_t\,\mathbb E_{P_t\sim\mathcal P_{t-1}}\mathrm{KL}(Q_t\|P_t)
% + \sum_{t=1}^T \mathrm{KL}(\mathcal P_t\|\mathcal P_{t-1}) + \log\tfrac1\delta$) yields
% \begin{equation}
% \label{eq:weighted-final}
% \bar\Delta
% \;\le\;
% \sqrt{\frac{A}{2B_T}}
% \;\le\;
% \sqrt{\frac{
% \sum_{t=1}^T \mathbb E_{P_t\sim\mathcal P_{t-1}}\mathrm{KL}(Q_t\|P_t)
% + \sum_{t=1}^T \mathrm{KL}(\mathcal P_t\|\mathcal P_{t-1})
% + \log\tfrac1\delta}{2\,\sum_{t=1}^T (m_t-1)}}.
% \end{equation}
% The last inequality uses $\sum_t \alpha_t x_t \le \sum_t x_t$ for nonnegative $(x_t)$.

% \paragraph{Step 4 From weighted to the statement of Theorem~\ref{thm:pathwise-pacbayes}.}
% The theorem controls the unweighted average $\frac1T\sum_{t=1}^T \Delta_t$ and uses the smoothed empirical risk on the right-hand side.
% When $m_t$ are homogeneous, $\alpha_t\equiv 1/T$ and $\bar\Delta=\frac1T\sum_t\Delta_t$.
% For heterogeneous $m_t$, we keep the same right-hand side and upper-bound the numerator $\sum_t \alpha_t \mathbb E_{P_t}\mathrm{KL}(Q_t\|P_t)$ by $\sum_t \mathbb E_{P_t}\mathrm{KL}(Q_t\|P_t)$, which yields exactly the denominator $2\sum_t(m_t-1)$ in the theorem.

% % \subsection*{A.4\quad Gaussian smoothing (LoRA) corollary}

% If $Q_t(\cdot\mid P_t,S_t)=\mathcal N(\hat W_t,\Sigma_t)$ (e.g., supported on an adapter subspace), then
% \[
% \mathbb E_{W\sim Q_t(\cdot\mid P_t,S_t)}\hat L_{S_t}(W)
% =\mathbb E_{\varepsilon\sim\mathcal N(0,\Sigma_t)}\hat L_{S_t}(\hat W_t+\varepsilon),
% \]
% which converts the empirical term into the smoothed empirical loss (expected sharpness) used in the main text; substituting this identity into \eqref{eq:weighted-final} yields the empirical term in Theorem~\ref{thm:pathwise-pacbayes}.

% \end{proof}


% \subsection{Proof of Lemma~\ref{lem:kl-decomposition}}

% % \begin{lemma}[KL decomposition with a frozen backbone]
% % \label{equation: subspace-full}
% % Let $W\in\mathbb R^d$ be fixed and let $P^{\Delta},Q^{\Delta}$ be probability measures on $(\mathcal W_\Delta,\mathcal B(\mathcal W_\Delta))$ with $Q^{\Delta}\ll P^{\Delta}$. Define their pushforwards onto the affine subspace by
% % \[
% % P^{\mathrm{full}}:=(T_W)_{\text{\#}} P^{\Delta},\qquad
% % Q^{\mathrm{full}}:=(T_W)_{\text{\#}} Q^{\Delta}.
% % \]
% % Then $Q^{\mathrm{full}}\ll P^{\mathrm{full}}$ and
% % \[
% % \mathrm{KL}\left(Q^{\mathrm{full}}\middle|P^{\mathrm{full}}\right)
% % =\mathrm{KL}\left(Q^{\Delta}\middle|P^{\Delta}\right).
% % \]
% % Equivalently, identifying ${W}\times\mathcal W_\Delta$ with $W+\mathcal W_\Delta$ via $\Phi$, one may write
% % $P^{\mathrm{full}}=\Phi_{\text{\#}}(\delta_W\otimes P^{\Delta})$ and
% % $Q^{\mathrm{full}}=\Phi_{\text{\#}}(\delta_W\otimes Q^{\Delta})$, and the same equality holds.
% % \end{lemma}

% \begin{proof}
%   Absolute continuity is preserved by pushforward under measurable maps, hence $Q^{\mathrm{full}}\ll P^{\mathrm{full}}$. Since $T_W$ is a Borel isomorphism, by the Radon-Nikodym chain rule
%  \[
%  \frac{dQ^{\mathrm{full}}}{dP^{\mathrm{full}}}(T_W\Delta)=\frac{dQ^{\Delta}}{dP^{\Delta}}(\Delta)\quad\text{for }P^{\Delta}\text{-a.e.\ }\Delta.
%  \]

% \begin{equation*}
%      \begin{aligned}
%  \mathrm{KL}\left(Q^{\mathrm{full}}\middle|P^{\mathrm{full}}\right)
%  &= \int_{W+\mathcal W_\Delta}
%  \log\Big(\frac{dQ^{\mathrm{full}}}{dP^{\mathrm{full}}}(y)\Big)
%  Q^{\mathrm{full}}(dy)\\
%  &= \int_{\mathcal W_\Delta}
%  \log\Big(\frac{dQ^{\mathrm{full}}}{dP^{\mathrm{full}}}\big(T_W\Delta\big)\Big)
%  Q^{\Delta}(d\Delta) \quad \text{(} Q^{\mathrm{full}}=(T_W)_{\text{\#}}Q^{\Delta}\text{)} 
%  \\
%  &= \int_{\mathcal W_\Delta}
%  \log\Big(\frac{dQ^{\Delta}}{dP^{\Delta}}(\Delta)\Big)
%  Q^{\Delta}(d\Delta) \\
%  &= \mathrm{KL}\left(Q^{\Delta}\middle|P^{\Delta}\right).
%  \end{aligned}
% \end{equation*} 
% \end{proof}

% % \section{ Proof of Lemma~\ref{lemma: subspace-full}}
% % \begin{lemma}[Sharpness decomposition in PECL]
% % \label{lemma: Sharpness in PECL}
% % Let $W_0\in\mathbb R^d$ be fixed and let the trainable update lie in the adapter subspace $\mathcal W_\Delta\subseteq\mathbb R^d$. Consider a Gaussian posterior on $(\mathcal W_\Delta,\mathcal B(\mathcal W_\Delta))$,
% %  $
% %  Q_t^{\Delta}=\mathcal N(\hat{\Delta W}_t,\Sigma_t),
% %  $
% %  with mean $\hat{\Delta W}_t\in\mathcal W_\Delta$ and covariance operator $\Sigma_t\succeq 0$ acting on $\mathcal W_\Delta$ (viewed in $\mathbb R^d$, this covariance is supported in $\mathcal W_\Delta$). Define the full-space posterior on ${W_0}\times\mathcal W_\Delta$ by
% %  $
% %  Q_t^{\mathrm{full}}:=\delta_{W_0}\otimes Q_t^{\Delta},
% %  $
% % and identify ${W_0}\times\mathcal W_\Delta$ with the affine subspace $W_0+\mathcal W_\Delta$ via $(w,\Delta)\mapsto w+\Delta$. Let $\hat L_{S_t}:\mathbb R^d\to\mathbb R$ be measurable and integrable under $Q_t^{\mathrm{full}}$. Then
% %  \[
% %  \mathbb{E}_{W\sim Q_t^{\mathrm{full}}}\big[\hat{L}_{S_t}(W)\big]
% %  = \mathbb{E}_{\varepsilon\sim\mathcal N(0,\Sigma_t)}
% %  \big[\hat{L}_{S_t}(W_0+\hat{\Delta W}_t+\varepsilon)\big]
% %  = \hat{L}_{S_t}(W_0+\hat{\Delta W}_t) +\mathrm{E\text{-}Sh}_t^{\Delta}(\hat{\Delta W}_t,\Sigma_t) \]
% %    where the subspace expected sharpness is
% %    \[
% %    \mathrm{E\text{-}Sh}_t^{\Delta}(\hat{\Delta W}_t,\Sigma_t)
% %    = \mathbb{E}_{\varepsilon\sim\mathcal N(0,\Sigma_t)}
% %    \big[\hat{L}_{S_t}(W_0+\hat{\Delta W}_t+\varepsilon)
% %   - \hat{L}_{S_t}(W_0+\hat{\Delta W}_t)\big].
% %      \]
% % \end{lemma}
% % \begin{proof}
% %     Sampling $W\sim Q_t^{\mathrm{full}}$ is equivalent to sampling $\varepsilon\sim\mathcal N(0,\Sigma_t)$ on $\mathcal W_\Delta$ and setting $W=W_0+\hat{\Delta W}_t+\varepsilon$. Taking expectations and adding-subtracting $\hat L_{S_t}(W_0+\hat{\Delta W}_t)$ inside the expectation gives the decomposition. 
% % \end{proof}


% \subsection{Proof of Theorem \ref{thm:pecl-bound}}

% % \begin{theorem}[Sharpness-aware PAC--Bayes bound for PECL]\label{thm:pecl-bound}
% % For any $\delta\in(0,1]$, with probability at least $1-\delta$ over $S_1,\ldots,S_T$,
% % \begin{equation}\label{eq:pecl-main}
% % \begin{aligned}
% % &\frac{1}{T}\sum_{t=1}^T
% % \mathbb{E}_{P_t\sim\mathcal P_{t-1}}\mathbb{E}_{W\sim Q_t}L_{\mathcal{D}_t}(W)
% % \ \le\
% % \ \frac{1}{T}\sum_{t=1}^T
% % \mathbb{E}_{P_t\sim\mathcal P_{t-1}}
% % \mathbb{E}_{\varepsilon \sim \mathcal{N}(0, \Sigma_t)}[\hat{L}_{S_t}(W_0 + \hat{\Delta W}_t + \varepsilon)
% % \\
% % &\quad +\sqrt{\frac{
% % \displaystyle \sum_{t=1}^T \mathbb{E}_{P_t^\Delta\sim\mathcal P_{t-1}}\!\Big[\mathrm{KL}\!\big(Q_t^\Delta\|P_t^\Delta\big)\Big]
% % \;+\;\displaystyle \sum_{t=1}^T \mathrm{KL}\!\big(\mathcal P_t\|\mathcal P_{t-1}\big)
% % \;+\;\log\!\frac{1}{\delta}
% % }{2\,\displaystyle\sum_{t=1}^T (m_t-1)}}.
% % \end{aligned}
% % \end{equation}
% % \end{theorem}


% \begin{proof}
% The bound follows by specializing Theorem 2 to the PECL setting through two key decompositions:

% By Lemma \ref{lem:kl-decomposition}, the full-space KL terms reduce to their adapter-space counterparts:
% \[
% \mathrm{KL}(Q_t^{\mathrm{full}} \| P_t^{\mathrm{full}}) = \mathrm{KL}(Q_t^\Delta \| P_t^\Delta),
% \quad
% \mathrm{KL}(\mathcal{P}_t^{\mathrm{full}} \| \mathcal{P}_{t-1}^{\mathrm{full}}) = \mathrm{KL}(\mathcal{P}_t^\Delta \| \mathcal{P}_{t-1}^\Delta).
% \]

% % \textbf{Step 2: Sharpness decomposition via Lemma 6.}  
% Under the Gaussian posterior $Q_t^\Delta = \mathcal{N}(\Delta\hat{W}_t, \Sigma_t)$, Lemma \ref{lemma: subspace-full} gives:
% \[
% \mathbb{E}_{W \sim Q_t^{\mathrm{full}}}[\hat{L}_{S_t}(W)] 
% = \mathrm{E}_{\varepsilon \sim \mathcal{N}(0, \Sigma_t)}{L}_{S_t}(W_{t-1}+\Delta\hat{W}_t + \varepsilon)
%  = \hat{L}_{S_t}(W_{t-1} + \Delta\hat{W}_t) + \mathrm{E\text{-}Sh}_t^\Delta(\Delta\hat{W}_t, \Sigma_t),
% \]
% where the perturbations are confined to the adapter subspace.

% % \textbf{Step 3: Substitution into general bound.}
% Substituting these decompositions into Theorem 2 yields the final PECL-specific bound, where all complexity and sharpness terms are now restricted to the trainable LoRA subspace $\mathcal{W}_\Delta$.
% % Substituting these decompositions into Theorem 2 yields the final PECL-specific bound, where all complexity and sharpness terms are now restricted to the trainable LoRA subspace $\mathcal{W}_\Delta$.

% \end{proof}





% \section{Appendix Experiments }


% \subsection{Exp Detailed setting}

%  We set \( \rho=0.05 \) for both \( \mathrm{Sh}^{(0)}(\rho) \)  and \( \mathrm{Sh}^{(1)}(\rho) \) ; \( \mathrm{E\text{-}Sh}(Q) \) is estimated using 100 Gaussian perturbations \( \delta\sim\mathcal N(0,\sigma^2 I) \) with \( \sigma=\rho/\sqrt{d} \), where $d$ is the dimensionality of the parameter vector and table entries report \(\mathrm{E\text{-}Sh}(Q)\times10^{4}\) for readability.ViT-B/16-IN21K; means $\pm$ sample std over 5 seeds with shuffled class orders


% \subsection{EFM analyse}

% \begin{figure}[htbp]
% \centering
% \begin{minipage}{0.9\linewidth}
% \centering
% \includegraphics[width=0.9\linewidth]{figures_TRML/finetune_42_gam-sam-sgd_efm_spectrum_compare.pdf}
% \subcaption{Eigenvalue Spectrum in Sequential Fine-tuning (SeqFT)}
% \label{fig:EFM_seqFT}
% \end{minipage}
% \caption{
% Comparison of EFM eigenvalue spectra in the feature space across incremental learning methods. The horizontal axis represents eigenvalue rank (sorted by magnitude), while the vertical axis shows eigenvalue magnitude. Only the top 200 eigenvalues are non-zero, with remaining dimensions at zero. SAM consistently produces flatter, heavier-tailed spectra compared to SGD across both configurations, indicating reduced energy concentration in the leading principal directions. This spectral profile suggests that SAM learns more isotropic representations with higher effective dimensionality.
% }
% \label{fig:EFM_spectra_comparison}
% \end{figure}


% % \subsection{Ablation Study}


% \subsection{Time consusme analyse}

% \begin{table}[thpb]
% \centering
% \caption{Average per‑task training time (seconds) across methods and optimizers on ImageNet‑R}
% \label{tab:lora_opt_metric}
% \setlength{\tabcolsep}{4pt}
% \renewcommand{\arraystretch}{1}
% \resizebox{\linewidth}{!}{
% \begin{tabular}{l ccccc}
% \toprule
% \textbf{Method} & \textbf{SGD} & \textbf{RWP} & \textbf{SAM} & \textbf{GAM} & \textbf{C-FLAT} \\
% \midrule
% SeqLoRA & 282.72 ± 61.34  & 244.95 ± 6.38  & 472.29 ± 83.70  & 670.20 ± 217.92 & 532.62 ± 35.93   \\
% IncLoRA & 310.54 ± 141.00 & 658.34 ± 83.90 & 777.91 ± 239.29 & 1247.56 ± 5.05  & 1189.40 ± 109.93 \\
% OLoRA   & 412.99 ± 146.84 & 644.19 ± 55.05 & 628.44 ± 181.29 & 836.99 ± 327.65 & 990.88 ± 360.79  \\
% \bottomrule
% \end{tabular}
% }
% \end{table}



\section{Appendix}

\subsection{Scope and Summary of Notation}
\label{app:notation}

This appendix collects the main symbols used in the paper and in the appendices, following (approximately) the order in which they appear.

\begin{itemize}
\item $t\in\{1,\dots,T\}$: task index in the continual learning sequence. % optional
% \item $T$: total number of tasks.
\item $S_t$: training set of task $t$.
\item $m_t := |S_t|$: number of training samples in $S_t$.
\item $D_t$: data distribution of task $t$.
\item $z=(x,y)\sim D_t$: an input label pair drawn from $D_t$.
% \item $f_W$: predictor (model) parameterized by $W$.
\item $\ell(W;z)$: bounded loss evaluated at parameters $W$ on sample $z$.
\item $L_{D_t}(W)=\mathbb E_{z\sim D_t}[\ell(W;z)]$: population risk on task $t$. 
\item $\widehat L_{S_t}(W)=\frac{1}{m_t}\sum_{z\in S_t}\ell(W;z)$: empirical risk on $S_t$.

% \item $W\in\mathbb R^{d}$: full parameter vector (or a vectorized weight block) of the model.

% \item $d$: dimensionality of the ambient parameter (or increment) space.


\item $\rho$: perturbation radius for adversarial sharpness.
% \item $\epsilon$: adversarial perturbation with $\|\epsilon\|\le \rho$.

% \item $\varepsilon\sim\mathcal N(0,\Sigma)$: Gaussian perturbation variable.
\item $\epsilon$ and $\varepsilon$: adversarial and Gaussian perturbations, respectively, where $\|\epsilon\|\le \rho$ and $\varepsilon\sim\mathcal N(0,\Sigma)$.

% \item $\Sigma\succeq 0$: covariance matrix for Gaussian perturbations.
% \item $\varepsilon\sim\mathcal N(0,\Sigma_t)$: Gaussian perturbation used to define expected (random) sharpness. 
% \item $\nabla_W \hat L_S(W)$: gradient of empirical risk with respect to $W$.
% \item $\mathrm{AS}_S^{(0)}(W,\rho)$: adversarial zeroth-order sharpness.
% \item $\mathrm{AS}_S^{(1)}(W,\rho)$: adversarial first-order sharpness.
\item $\mathrm{AS}^{(0)}_S(W,\rho)$: adversarial zeroth-order sharpness from SAM,
$\mathrm{AS}^{(0)}_S(W,\rho)=\max_{\|\epsilon\|\le\rho}\big[\widehat L_S(W+\epsilon)-\widehat L_S(W)\big]$.

\item $\mathrm{RS}_S^{(0)}(W,\Sigma)$: random zeroth-order sharpness, i.e., the expected empirical loss elevation under $\varepsilon\sim\mathcal N(0,\Sigma)$:
$\mathrm{RS}_S^{(0)}(W,\Sigma)=\mathbb E_{\varepsilon\sim\mathcal N(0,\Sigma)}[\hat L_S(W+\varepsilon)-\hat L_S(W)]$.
\item $\mathrm{AS}^{(1)}_S(W,\rho)$: adversarial first-order sharpness  from GAM,
$\mathrm{AS}^{(1)}_S(W,\rho)=\max_{\|\epsilon\|\le\rho}\|\nabla_W\widehat L_S(W+\epsilon)\|$.

\item $P$: prior distribution over parameters (hypotheses) in PAC-Bayes analysis.
\item $Q$: posterior distribution over parameters (hypotheses) in PAC-Bayes analysis.
% \item $\mathrm{KL}(Q\Vert P)$: Kullback--Leibler divergence measuring posterior complexity relative to the prior.
% \item $\delta\in(0,1]$: confidence level in high-probability bounds. 

\item $P_t$: task prior at step $t$ 
% (a distribution over model parameters restricted to the current task). 
\item $Q_t(\cdot \mid P_t,S_t)$: task posterior after learning task $t$, possibly conditioned on $P_t$ and $S_t$.
\item $Q_t^\Delta \ll P_t^\Delta$: absolute continuity assumption
\item $\mathcal P_t$: hyper-posterior (a distribution over task priors), forming a time-varying hyper-process . 
\item $\mathcal F_{t-1}$: filtration ($\sigma$-algebra) capturing the history up to step $t-1$ in the martingale argument.

\item $W_0$: pretrained initialization before continual adaptation.
\item $W_t$: model parameters after training on tasks $1$ to $t$.
\item $W_t^{\mathrm{froz}}$: frozen component when learning task $t$ (the part not trainable during task $t$).
\item $\Delta W_t$: trainable update at task $t$ under parameter-efficient adaptation.
\item $W_t = W_t^{\mathrm{froz}} + \Delta W_t$: decomposition of current parameters into frozen and trainable parts. 





% \item $\mathrm{LoRA}$: low-rank adaptation parameterization of a weight block, typically $\Delta W_t = B_tA_t$ with rank $r$.
% \item $A_t\in\mathbb R^{r\times d_{\text{in}}},\; B_t\in\mathbb R^{d_{\text{out}}\times r}$: LoRA factors for task $t$ (shapes may vary by layer/block).
% \item $r$: LoRA rank (adapter dimension).


\item $\mathcal W_{\Delta,t}\subseteq\mathbb R^d$: the task-admissible linear update subspace at step $t$, containing all feasible increments induced by LoRA while the backbone is frozen.


\item $\widehat W_t$: the full-model parameter vector after learning task $t$ (a point estimate in the ambient parameter space).
\item $\widehat{\Delta W}_t$: the learned trainable increment at task $t$ under PEFT/LoRA, satisfying $\widehat W_t = W_t^{\mathrm{froz}} + \widehat{\Delta W}_t$.

 


 
% \item $\mathcal B(\cdot)$: Borel sigma-algebra (and its trace on subspaces/affine sets). 
% \item $T_t(\Delta)=W_t^{\mathrm{froz}}+\Delta$: translation map from the update space to the affine model class. 
% \item $T_{t\#}\mu$: pushforward measure induced by $T_t$, defined by $(T_{t\#}\mu)(A)=\mu(T_t^{-1}(A))$. 
\item $P_t^\Delta,\;Q_t^\Delta$: prior/posterior supported on $(\mathcal W_{\Delta,t},\mathcal B(\mathcal W_{\Delta,t}))$.
% \item $P_,\; Q_t$: induced full-model prior/posterior supported on $W_t^{\mathrm{froz}}+\mathcal W_{\Delta,t}$.
% \item $Q_t^\Delta \ll P_t^\Delta$: absolute continuity assumption.
% \item $\frac{dQ}{dP}$: Radon--Nikodym derivative appearing in the KL invariance proof. 

% \item $\Sigma_t$: covariance defining a Gaussian perturbation distribution on $\mathcal W_{\Delta,t}$. 



\item $\mathrm{RS}_t^\Delta(\widehat W_t,\Sigma_t)$: task-conditioned random sharpness with perturbations supported on $\mathcal W_{\Delta,t}$, defined as
$\mathrm{RS}_t^\Delta(\widehat W_t,\Sigma_t)=\mathbb E_{\varepsilon\sim\mathcal N(0,\Sigma_t)}[\hat L_{S_t}(\widehat W_t+\varepsilon)-\hat L_{S_t}(\widehat W_t)]$,
where $\varepsilon$ is restricted to $\mathcal W_{\Delta,t}$, $\Sigma_t$ is defined on $\mathcal{W}_{\Delta,t}$
% (equivalently, $\mathrm{supp}(\varepsilon)\subseteq\mathcal W_{\Delta,t}$).


% \item $\widehat W_t$: posterior mean (or a point estimate) used when specializing $Q_t$ to a Gaussian.
% \item $\mathrm{RS}_t^\Delta(\widehat W_t,\Sigma_t)$:  random sharpness(expected sharpness) restricted to $\mathcal W_{\Delta,t}$ .

% \item $a_{tj}$: accuracy on task $j$ after training up to task $t$.
% \item $\mathrm{Acc}_t=\frac{1}{t}\sum_{j=1}^t a_{tj}$: average accuracy across the first $t$ tasks.
% \item $\mathrm{Acc}_t^{\mathrm{cls}}$: class-incremental classification accuracy at step $t$.
% \item $\mathrm{Acc}_t^{\mathrm{ncm}}$: nearest-class-mean accuracy (prototype-based) reflecting representation quality. 

\item $\alpha_t$ : feature amplification factor

% \item $c_t$: curvature projection ratio 

% \item $\lambda>0$: exponential tilting parameter used in the moment generating function (MGF) step of the proof~\ref{proof:theorem all}. 
% \item $B_T:=\sum_{t=1}^T (m_t-1)$: aggregate sample-size factor in Proof~\ref{proof:theorem all}. 
% \item $\omega_t := (m_t-1)/B_T$: sample-size weight used to form the weighted generalization gap.
% \item $\Delta_t$: step-$t$ generalization gap (defined in ~\ref{proof:theorem all}).
% \item $\overline\Delta := \sum_{t=1}^T \alpha_t \Delta_t$: size-weighted average gap. 
% \item $g_t(W):=\lambda\alpha_t\big(L_{D_t}(W)-\widehat L_{S_t}(W)\big)$: MGF test function at step $t$. 
% \item $h_t(P)$: hyper-level functional used in the Donsker--Varadhan (DV) change-of-measure step. 
% \item $\mathcal F_{t-1}$: filtration ($\sigma$-algebra) capturing the history up to step $t-1$ in the martingale argument. 
% \item $X_t$: exponential increment defined by the conditional MGF bound at step $t$.
% \item $\mathcal M_t:=\prod_{s=1}^t X_s$: nonnegative supermartingale used to obtain a high-probability bound.

% \item $\mathcal W_{\Delta,t}\subseteq\mathbb R^d$: permissible update subspace at task $t$ (task-admissible increment space). 
% \item $\Pi_{\Delta,t}$: orthogonal projector onto $\mathcal W_{\Delta,t}$ when $\mathcal W_{\Delta,t}$ is a linear subspace. 
% \item $J_{\mathrm{LoRA}}$: Jacobian of the LoRA map (used to define the local tangent increment space), with $\mathcal W_{\Delta,t}=\mathrm{Im}(J_{\mathrm{LoRA}})$.


% \item $\Delta W_t=U_{\Delta,t}S_{\Delta,t}V_{\Delta,t}^\top$: rank-$r$ truncated SVD of the learned update . 
% \item $\mathcal S_{\Delta,t}=\{U_{\Delta,t}ZV_{\Delta,t}^\top: Z\in\mathbb R^{r\times r}\}$: bilinear diagnostic subspace induced by the truncated SVD of $\Delta W_t$. 

% % \item $\Pi_{\Delta,t}^{\mathrm{SVD}}(M):=U_{\Delta,t}(U_{\Delta,t}^\top M V_{\Delta,t})V_{\Delta,t}^\top$: Frobenius-orthogonal projection onto $\mathcal S_{\Delta,t}$ (used only for curvature localization). 
% \item $\Pi_{\mathcal S,t}(M):=U_{\Delta,t}(U_{\Delta,t}^\top M V_{\Delta,t})V_{\Delta,t}^\top$: Frobenius-orthogonal projection onto $\mathcal S_{\Delta,t}$ .


% % \item $(U_{W,t},V_{W,t})$: rank-$r$ truncated SVD subspace of the corresponding backbone weight block (Appendix A.4.2). 
% % \item $\mathcal S_{W,t}=\{U_{W,t}ZV_{W,t}^\top: Z\in\mathbb R^{r\times r}\}$: backbone reference bilinear subspace. 
% % \item $\Pi_{W,t}^{\mathrm{SVD}}(M):=U_{W,t}(U_{W,t}^\top M V_{W,t})V_{W,t}^\top$: projection onto $\mathcal S_{W,t}$. 

% \item $(U_{W,t},V_{W,t})$: rank-$r$ truncated SVD bases of the corresponding backbone weight block.
% \item $\mathcal S_{W,t}=\{U_{W,t}ZV_{W,t}^\top: Z\in\mathbb R^{r\times r}\}$: backbone reference bilinear subspace.
% \item $\Pi_{\mathcal S_W,t}(M):=U_{W,t}(U_{W,t}^\top M V_{W,t})V_{W,t}^\top$: Frobenius-orthogonal projection onto $\mathcal S_{W,t}$.

% \item $\{(\lambda_{t,i},v_{t,i})\}_{i=1}^k$: top-$k$ eigenpairs of a PSD curvature operator (e.g., GGN/Fisher) at task $t$. 
% \item $(\delta W^{(0)}_{t,i},\delta A_{t,i},\delta B_{t,i})$: decomposition of eigen-direction $v_{t,i}$ into backbone and LoRA-factor components. 
% \item $\delta W_{t,i}=\delta W^{(0)}_{t,i}+(\delta B_{t,i}A_t+B_t\delta A_{t,i})$: effective weight-space perturbation induced by $(\delta W^{(0)}_{t,i},\delta A_{t,i},\delta B_{t,i})$. 
% \item $\rho_{\Delta,t,i}=\frac{\|\Pi^{\mathrm{SVD}}_{\Delta,t}(\delta W_{t,i})\|_F^2}{\|\delta W_{t,i}\|_F^2}$: per-mode projection ratio onto $\mathcal S_{\Delta,t}$.
% \item $\rho_{W,t,i}=\frac{\|\Pi_{W,t}^{\mathrm{SVD}}(\delta W_{t,i})\|_F^2}{\|\delta W_{t,i}\|_F^2}$: per-mode projection ratio onto $\mathcal S_{W,t}$.
\item $R_{\Delta,t}=\frac{\sum_{i=1}^k\lambda_{t,i}\rho_{\Delta,t,i}}{\sum_{i=1}^k\lambda_{t,i}}$: eigenvalue-weighted localization on $\mathcal S_{\Delta,t}$.
\item $R_{W,t}=\frac{\sum_{i=1}^k\lambda_{t,i}\rho_{W,t,i}}{\sum_{i=1}^k\lambda_{t,i}}$: eigenvalue-weighted localization on $\mathcal S_{W,t}$.
\item $c_t=R_{\Delta,t}/R_{W,t}$: curvature projection ratio comparing $\mathcal S_{\Delta,t}$ against the backbone reference. 
\end{itemize}


\subsection{Proof}

\subsubsection{Proof of Theorem \ref{thm:pathwise_pac}}
\label{proof:theorem all}

\paragraph{Preliminary: a supermartingale tail bound.}
Let $(\mathcal M_t)_{t\ge0}$ be a nonnegative supermartingale adapted to $(\mathcal F_t)$ with $\mathbb E[\mathcal M_0]\le 1$.
Then for any $\delta\in(0,1)$,
\[
\mathbb P\!\left(\mathcal M_T \ge \tfrac{1}{\delta}\right)
\le \delta,
\]
by Markov's inequality and $\mathbb E[\mathcal M_T]\le \mathbb E[\mathcal M_0]\le 1$.
Equivalently, with probability at least $1-\delta$, $\log \mathcal M_T \le \log \tfrac{1}{\delta}$.
This is the fixed-time form of the classical supermartingale maximal inequality  \citet{yangshuxue}.


\begin{theorem}[Pathwise PAC--Bayes with time-varying hyper-posterior]
\label{the: full}
For any $\delta\in(0,1]$, with probability at least $1-\delta$ over $S_{1:T}$,
\begin{equation}
\begin{aligned} 
& \frac{1}{T} \sum_{t=1}^T \underbrace{\mathbb{E}_{P_t \sim \mathcal{P}_t} \mathbb{E}_{W \sim Q_t(\cdot \mid P_t, S_t)} L_{\mathcal{D}_t}(W)}_{\text{test risk at step } t} 
\le \ \frac{1}{T} \sum_{t=1}^T \underbrace{\mathbb{E}_{P_t \sim \mathcal{P}_t} \mathbb{E}_{W \sim Q_t(\cdot \mid P_t, S_t)} \hat{L}_{S_t}(W)}_{\text{empirical risk at step } t} \\ 
& + \sqrt{\frac{ \sum_{t=1}^T \mathbb{E}_{P_t \sim \mathcal{P}_t} \mathrm{KL}\big(Q_t(\cdot \mid P_t, S_t) \,\|\, P_t\big) + \sum_{t=1}^T \mathrm{KL}(\mathcal{P}_t \,\|\, \mathcal{P}_{t-1}) + \log\frac{1}{\delta} }{2 \sum_{t=1}^T (m_t - 1)}}. 
\end{aligned} 
\end{equation}
\end{theorem}

\begin{proof}
    


% \subsection*{A.1\quad Two-level variational inequalities}
% \item $B_T:=\sum_{t=1}^T (m_t-1)$: aggregate sample-size factor in Proof~\ref{proof:theorem all}. 
% \item $\omega_t := (m_t-1)/B_T$: sample-size weight used to form the weighted generalization gap.
Let $B_T=\sum_{t=1}^T (m_t-1)$ denote an aggregate sample-size factor and define $\omega_t=(m_t-1)/B_T$ which serves as a sample-size weight satisfying $\sum_{t=1}^T\omega_t=1$.
% $B_T$ is 

% \paragraph{Goal quantity.} 


Define the step-$t$ generalization gap
\[
\Delta_t
:=\mathbb E_{P_t\sim\mathcal P_{t-1}}\,
\mathbb E_{W\sim Q_t(\cdot\mid P_t,S_t)}
\big(L_{\mathcal D_t}(W)-\hat L_{S_t}(W)\big),
\]
and the size-weighted gap $\bar\Delta:=\sum_{t=1}^T \omega_t \Delta_t$.
Fix $\lambda>0$ and set
\(
g_t(W):=\lambda\omega_t\big(L_{\mathcal D_t}(W)-\hat L_{S_t}(W)\big).
\)

\paragraph{Step 1 Model level.}
% For any fixed $P$ and $S_t$,
For any fixed prior distribution over the parameters $P$ and posterior distribution $Q$, and only current task dataset $S_t$,   
\begin{equation}
\label{eq:dv-w}
\mathbb E_{W\sim Q_t(\cdot\mid P,S_t)} g_t(W)
\;\le\;
\mathrm{KL}\!\big(Q_t(\cdot\mid P,S_t)\,\|\,P\big)
+\log\mathbb E_{W\sim P}\!\big[e^{g_t(W)}\big].
\end{equation}
Taking $\mathbb E_{S_t}$ and using the standard Hoeffding--MGF bound for bounded losses,
\begin{equation}
\label{eq:mgf-mean-gap}
\mathbb E_{S_t}\log\mathbb E_{W\sim P} e^{g_t(W)}
\;\le\; \frac{(\lambda\omega_t)^2}{8\,(m_t-1)}.
\end{equation}
Therefore,
\begin{equation}
\label{eq:model-step}
\mathbb E_{S_t}\mathbb E_{W\sim Q_t(\cdot\mid P,S_t)} g_t(W)
\;\le\;
\mathbb E_{S_t}\mathrm{KL}\!\big(Q_t(\cdot\mid P,S_t)\,\|\,P\big)
+\frac{(\lambda\omega_t)^2}{8\,(m_t-1)}.
\end{equation}
 % (change of measure in $P$-space)
\paragraph{Step 2 Hyper level.} 
Define
\[
h_t(P):=\mathbb E_{S_t}\mathbb E_{W\sim Q_t(\cdot\mid P,S_t)} g_t(W)
-\mathbb E_{S_t}\mathrm{KL}\!\big(Q_t(\cdot\mid P,S_t)\,\|\,P\big).
\]
By \eqref{eq:model-step}, $h_t(P)\le \frac{(\lambda\omega_t)^2}{8\,(m_t-1)}$ for all $P$.
Applying Donsker–Varadhan variational representation of the KL divergence~\cite{NEURIPS2021_54a367d6} to change the measure from $\mathcal P_{t-1}$ to $\mathcal P_t$ gives
\begin{equation}
\label{eq:dv-p}
\log\mathbb E_{P\sim\mathcal P_t}\!\big[e^{h_t(P)}\big]
\;\le\; \mathrm{KL}(\mathcal P_t\|\mathcal P_{t-1})
+\log\mathbb E_{P\sim\mathcal P_{t-1}}\!\big[e^{h_t(P)}\big]
\;\le\; \mathrm{KL}(\mathcal P_t\|\mathcal P_{t-1})
+\frac{(\lambda\omega_t)^2}{8\,(m_t-1)}.
\end{equation}
Using $\log\mathbb E[e^X]\ge \mathbb E[X]$ yields
\begin{equation}
\label{eq:p-expect}
\mathbb E_{P_t\sim\mathcal P_t} h_t(P_t)
\;\le\; \mathrm{KL}(\mathcal P_t\|\mathcal P_{t-1})
+\frac{(\lambda\omega_t)^2}{8\,(m_t-1)}.
\end{equation}
Unfolding $h_t(P)$ and interchanging expectations,
\begin{equation}
\label{eq:two-term}
\mathbb E_{S_t}\Big[
\lambda\omega_t\Delta_t
-\omega_t\,\mathbb E_{P_t\sim\mathcal P_{t-1}}\mathrm{KL}\!\big(Q_t(\cdot\mid P_t,S_t)\,\|\,P_t\big)
\Big]
\;\le\;
\mathrm{KL}(\mathcal P_t\|\mathcal P_{t-1})
+\frac{(\lambda\omega_t)^2}{8\,(m_t-1)}.
\end{equation}

% \subsection*{A.2\quad Conditional MGF and supermartingale}
\paragraph{Step 3 Martingale}
% From \eqref{eq:two-term} and the tower property of conditional expectation, 







Using Eq.\eqref{eq:two-term} and the tower property of conditional expectation (iterated expectation) with respect to $\mathcal F_{t-1}$, we obtain the desired conditional exponential bound~\cite{yangshuxue}.
% The conditional exponential bound
\begin{equation}
\label{eq:cond-mgf}
\mathbb E\!\left[
\exp\!\Big(
\lambda\omega_t\Delta_t
-\omega_t\,\mathbb E_{P_t\sim\mathcal P_{t-1}}\mathrm{KL}(Q_t\|P_t)
-\mathrm{KL}(\mathcal P_t\|\mathcal P_{t-1})
-\frac{(\lambda\omega_t)^2}{8\,(m_t-1)}
\Big)\,\Big|\,\mathcal F_{t-1}\right]\;\le\;1.
\end{equation}


From the conditional exponential bound~\eqref{eq:cond-mgf}, define
\[
X_t
:=
\exp\!\Big(
\lambda\omega_t\Delta_t
-\omega_t\,\mathbb E_{P_t\sim\mathcal P_{t-1}}\mathrm{KL}(Q_t\|P_t)
-\mathrm{KL}(\mathcal P_t\|\mathcal P_{t-1})
-\frac{(\lambda\omega_t)^2}{8\,(m_t-1)}
\Big).
\]
Assume $\omega_t$ is deterministic or $\mathcal F_{t-1}$-measurable. Then $X_t\ge 0$ is $\mathcal F_t$-measurable and accoring Eq.~\eqref{eq:cond-mgf}
\[
\mathbb E[X_t\mid \mathcal F_{t-1}] \le 1.
\]
Define the process
\[
\mathcal{M}_0:=1,\qquad
\mathcal{M}_t:=\prod_{s=1}^t X_s.
\]
It follows that $(\mathcal M_t)_{t\ge 0}$ is a nonnegative supermartingale, since
\[
\mathbb E[\mathcal M_t\mid \mathcal F_{t-1}]
=
\mathcal M_{t-1}\,\mathbb E[X_t\mid \mathcal F_{t-1}]
\le \mathcal M_{t-1}.
\]
In particular, $\mathbb E[\mathcal M_T]\le \mathbb E[\mathcal M_0]=1$.
By Markov's inequality,
\[
\mathbb P\!\left(\mathcal M_T \ge \tfrac{1}{\delta}\right)
\le \delta,
\]
so with probability at least $1-\delta$ we have $\log \mathcal M_T \le \log\tfrac{1}{\delta}$.
% Expanding $\log \mathcal M_T = \sum_
Expanding $\log \mathcal M_T = \sum_{t=1}^T \log X_t$ yields~\eqref{eq:hpp}.






% Define $X_t$ to be the exponential increment inside \eqref{eq:cond-mgf} and
% \[
% \mathcal{M}_0:=1,\qquad
% \mathcal{M}_t:=\prod_{s=1}^t X_s.
% \]
% Then $(\mathcal{M}_t)_{t\ge0}$ is a nonnegative supermartingale, so
% $\mathbb E[\mathcal{M}_T]\le 1$.
% By Markov's inequality, with probability at least $1-\delta$,

% so with probability at least $1-\delta$ we have $\log \mathcal M_T \le \log\tfrac{1}{\delta}$

\begin{equation}
\label{eq:hpp}
\sum_{t=1}^T \lambda\omega_t\Delta_t
\;\le\;
\sum_{t=1}^T \omega_t\,\mathbb E_{P_t\sim\mathcal P_{t-1}}\mathrm{KL}(Q_t\|P_t)
\;+\; \sum_{t=1}^T \mathrm{KL}(\mathcal P_t\|\mathcal P_{t-1})
\;+\; \frac{\lambda^2}{8}\sum_{t=1}^T \frac{\omega_t^2}{(m_t-1)}
\;+\; \log\tfrac1\delta.
\end{equation}

% \subsection*{A.3\quad Size-weighted self-normalization and optimization}

By $\omega_t=(m_t-1)/B_T$,
\[
\sum_{t=1}^T \frac{\omega_t^2}{(m_t-1)}
=\sum_{t=1}^T \frac{(m_t-1)^2/B_T^2}{(m_t-1)}
=\frac{1}{B_T}.
\]
Hence, \eqref{eq:hpp} reads, for all $\lambda>0$ and with probability at least $1-\delta$,
\begin{equation}
\label{eq:gap-preopt}
\bar\Delta
\;\le\;
\frac{1}{\lambda}\Big(
\sum_{t=1}^T \omega_t\,\mathbb E_{P_t\sim\mathcal P_{t-1}}\mathrm{KL}(Q_t\|P_t)
+ \sum_{t=1}^T \mathrm{KL}(\mathcal P_t\|\mathcal P_{t-1})
+ \log\tfrac1\delta\Big)
\;+\;
\frac{\lambda}{8B_T}.
\end{equation}
Optimizing the right-hand side over $\lambda>0$ (take $\lambda^\star=\sqrt{8B_T\,A}$ with
$A:=\sum_{t=1}^T \omega_t\,\mathbb E_{P_t\sim\mathcal P_{t-1}}\mathrm{KL}(Q_t\|P_t)
+ \sum_{t=1}^T \mathrm{KL}(\mathcal P_t\|\mathcal P_{t-1}) + \log\tfrac1\delta$) yields
\begin{equation}
\label{eq:weighted-final}
\bar\Delta
\;\le\;
\sqrt{\frac{A}{2B_T}}
\;\le\;
\sqrt{\frac{
\sum_{t=1}^T \mathbb E_{P_t\sim\mathcal P_{t-1}}\mathrm{KL}(Q_t\|P_t)
+ \sum_{t=1}^T \mathrm{KL}(\mathcal P_t\|\mathcal P_{t-1})
+ \log\tfrac1\delta}{2\,\sum_{t=1}^T (m_t-1)}}.
\end{equation}
The last inequality uses $\sum_t \omega_t x_t \le \sum_t x_t$ for nonnegative $(x_t)$.

\paragraph{Step 4 From weighted to the statement of Theorem~\ref{the: full}.}
The theorem controls the unweighted average $\frac1T\sum_{t=1}^T \Delta_t$ and uses the smoothed empirical risk on the right-hand side.
When $m_t$ are homogeneous, $\omega_t\equiv 1/T$ and $\bar\Delta=\frac1T\sum_t\Delta_t$.
For heterogeneous $m_t$, we keep the same right-hand side and upper-bound the numerator $\sum_t \omega_t \mathbb E_{P_t}\mathrm{KL}(Q_t\|P_t)$ by $\sum_t \mathbb E_{P_t}\mathrm{KL}(Q_t\|P_t)$, which yields exactly the denominator $2\sum_t(m_t-1)$ in the theorem.

% \subsection*{A.4\quad Gaussian smoothing (LoRA) corollary}

If $Q_t(\cdot\mid P_t,S_t)=\mathcal N(\hat W_t,\Sigma_t)$ (e.g., supported on an adapter subspace), then
\[
\mathbb E_{W\sim Q_t(\cdot\mid P_t,S_t)}\hat L_{S_t}(W)
=\mathbb E_{\varepsilon\sim\mathcal N(0,\Sigma_t)}\hat L_{S_t}(\hat W_t+\varepsilon),
\]
which converts the empirical term into the smoothed empirical loss (expected sharpness) used in the main text; substituting this identity into \eqref{eq:weighted-final} yields the empirical term in Theorem~\ref{the: full}.

Theorem~A.1 provides a tight complexity penalty of the form.
Using $\sqrt{x+y+z}\le \sqrt{x}+\sqrt{y}+\sqrt{z}$ for $x,y,z\ge 0$ yields the decomposed bound, which is the form stated in Theorem~\ref{thm:pathwise_pac} for interpretability.

\end{proof}

\subsubsection{Proof of Lemma~\ref{lem:kl-decomp-pecl}}
\label{proof:Lemma1}

Let $\mathcal W_{\Delta,t}\subseteq\mathbb R^d$ be a linear subspace and $W_t^\mathrm{froz}+\mathcal W_{\Delta,t}$ be its affine translation. 
We endow both sets with the trace Borel $\sigma$-algebra induced from $\mathbb R^d$. 
Define the translation map
\[
T_t:(\mathcal W_{\Delta,t},\mathcal B(\mathcal W_{\Delta,t}))
\to (W_t^{\mathrm{froz}}+\mathcal W_{\Delta,t},\mathcal B(W_t^{\mathrm{froz}}+\mathcal W_{\Delta,t})),
\qquad
T_t(\Delta)=W_t^{\mathrm{froz}}+\Delta .
\]
For any probability measure $\mu$ on $\mathcal W_{\Delta,t}$, its pushforward
$T_{t\#}\mu$ is the measure on $W_t^{\mathrm{froz}}+\mathcal W_{\Delta,t}$
defined by $(T_{t\#}\mu)(A)=\mu(T_t^{-1}(A))$ for all measurable $A$.
Accordingly, we define the induced full-model prior and posterior supported on
$W_t^{\mathrm{froz}}+\mathcal W_{\Delta,t}$ as
\[
P_t := T_{t\#} P_t^\Delta,
\qquad
Q_t := T_{t\#} Q_t^\Delta,
\]
where $P_t^\Delta$ and $Q_t^\Delta$ are probability measures on
$(\mathcal W_{\Delta,t},\mathcal B(\mathcal W_{\Delta,t}))$.
% $\mathcal B(\mathcal W_{\Delta,t})\ =\{B\cap \mathcal W_{\Delta,t}:\ B\in \mathcal B(\mathbb R^d)\},\quad
% \mathcal B(W^{\mathrm{froz}}_t+\mathcal W_{\Delta,t})\ =\{B\cap (W^{\mathrm{froz}}_t+\mathcal W_{\Delta,t}):\ B\in \mathcal B(\mathbb R^d)\}.$
% These coincide with the Borel $\sigma$-algebras generated by the subspace topologies on $\mathcal W_{\Delta,t}$ and $W^{\mathrm{froz}}_t+\mathcal W_{\Delta,t}$, respectively, so both are Borel spaces.
% To formalize the probabilistic relationship between the LoRA subspace and the full model space, we define the translation map:
% % \begin{align*}
% $T_t: (\mathcal W_{\Delta,t},\mathcal B(\mathcal W_{\Delta,t})) \to (W^{\mathrm{froz}}_t+\mathcal W_{\Delta,t},\mathcal B(W^{\mathrm{froz}}_t+\mathcal W_{\Delta,t})), 
% T_t(\Delta)=W^{\mathrm{froz}}_t+\Delta,$
% % \end{align*}
% which is a Borel isomorphism between the subspace and affine space.  The full-space prior and posterior are defined as pushforwards onto the affine subspace by
% $
% P=(T_t)_{\text{\#}} P^{\Delta}_t,
% Q=(T_t)_{\text{\#}} Q^{\Delta}_t,
% $
% where $P_t^{\Delta}$ and $Q_t^{\Delta}$ are probability measures on $(\mathcal{W}_{\Delta,t}, \mathcal{B}(\mathcal{W}_{\Delta,t}))$.

% \begin{lemma}[KL decomposition with a frozen backbone]
% \label{equation: subspace-full}
% Let $W\in\mathbb R^d$ be fixed and let $P^{\Delta},Q^{\Delta}$ be probability measures on $(\mathcal W_\Delta,\mathcal B(\mathcal W_\Delta))$ with $Q^{\Delta}\ll P^{\Delta}$. Define their pushforwards onto the affine subspace by
% \[
% P:=(T_W)_{\text{\#}} P^{\Delta},\qquad
% Q:=(T_W)_{\text{\#}} Q^{\Delta}.
% \]
% Then $Q\ll P$ and
% \[
% \mathrm{KL}\left(Q\middle|P\right)
% =\mathrm{KL}\left(Q^{\Delta}\middle|P^{\Delta}\right).
% \]
% Equivalently, identifying ${W}\times\mathcal W_\Delta$ with $W+\mathcal W_\Delta$ via $\Phi$, one may write
% $P=\Phi_{\text{\#}}(\delta_W\otimes P^{\Delta})$ and
% $Q=\Phi_{\text{\#}}(\delta_W\otimes Q^{\Delta})$, and the same equality holds.
% \end{lemma}

\begin{lemma}[KL in PECL is adapter-subspace KL]
% Assume $Q_t^\Delta \ll P_t^\Delta$. 
% Assume $Q_t^\Delta \ll P_t^\Delta$, meaning $Q_t^\Delta(A)=0$ for any measurable set $A$ such that $P_t^\Delta(A)=0$.
Assume $Q_t^\Delta \ll P_t^\Delta$ (i.e., $Q_t^\Delta$ is absolutely continuous w.r.t.\ $P_t^\Delta$), meaning $Q_t^\Delta(A)=0$ for any measurable set $A$ such that $P_t^\Delta(A)=0$.
Then $Q_t \ll P_t$ and
$$
\mathrm{KL}\bigl(Q_t\,\Vert\, P_t\bigr)
\;=\;
\mathrm{KL}\bigl(Q_t^\Delta \,\Vert\, P_t^\Delta\bigr).
$$
\end{lemma}

\begin{proof}
  % Absolute continuity is preserved by pushforward under measurable maps, hence $Q_t\ll P_t$. Since $T_t$ is a Borel isomorphism, 
  Since $T_t$ is measurable, absolute continuity is preserved under pushforward, hence
$Q_t=T_{t\#}Q_t^\Delta \ll T_{t\#}P_t^\Delta=P_t$.
Moreover, $T_t$ is a measurable bijection with measurable inverse. 
  By the Radon-Nikodym chain rule
 \[
 \frac{dQ_t}{dP_t}(T_t(\Delta))=\frac{dQ_t^{\Delta}}{dP_t^{\Delta}}(\Delta)\quad\text{for }P^{\Delta}_t\text{-a.e.\ }\Delta.
 \]

\begin{equation*}
     \begin{aligned}
 \mathrm{KL}\left(Q_t\middle|P_t\right)
 &= \int_{W+\mathcal W_{\Delta,t}}
 \log\Big(\frac{dQ_t}{dP_t}(y)\Big)
 Q_t(dy)\\
 &= \int_{\mathcal W_{\Delta,t}}
 \log\Big(\frac{dQ_t}{dP_t}\big(T_t\Delta\big)\Big)
 Q_t^{\Delta}(d\Delta) \quad \text{(} Q_t=(T_t)_{\text{\#}}Q_t^{\Delta}\text{)} 
 \\
 &= \int_{\mathcal W_{\Delta,t}}
 \log\Big(\frac{dQ_t^{\Delta}}{dP_t^{\Delta}}(\Delta)\Big)
 Q_t^{\Delta}(d\Delta) \\
 &= \mathrm{KL}\left(Q_t^{\Delta}\middle|P_t^{\Delta}\right).
 \end{aligned}
\end{equation*} 
\end{proof}

% \section{ Proof of Lemma~\ref{lemma: subspace-full}}
% \begin{lemma}[Sharpness decomposition in PECL]
% \label{lemma: Sharpness in PECL}
% Let $W_0\in\mathbb R^d$ be fixed and let the trainable update lie in the adapter subspace $\mathcal W_\Delta\subseteq\mathbb R^d$. Consider a Gaussian posterior on $(\mathcal W_\Delta,\mathcal B(\mathcal W_\Delta))$,
%  $
%  Q_t^{\Delta}=\mathcal N(\hat{\Delta W}_t,\Sigma_t),
%  $
%  with mean $\hat{\Delta W}_t\in\mathcal W_\Delta$ and covariance operator $\Sigma_t\succeq 0$ acting on $\mathcal W_\Delta$ (viewed in $\mathbb R^d$, this covariance is supported in $\mathcal W_\Delta$). Define the full-space posterior on ${W_0}\times\mathcal W_\Delta$ by
%  $
%  Q_t:=\delta_{W_0}\otimes Q_t^{\Delta},
%  $
% and identify ${W_0}\times\mathcal W_\Delta$ with the affine subspace $W_0+\mathcal W_\Delta$ via $(w,\Delta)\mapsto w+\Delta$. Let $\hat L_{S_t}:\mathbb R^d\to\mathbb R$ be measurable and integrable under $Q_t$. Then
%  \[
%  \mathbb{E}_{W\sim Q_t}\big[\hat{L}_{S_t}(W)\big]
%  = \mathbb{E}_{\varepsilon\sim\mathcal N(0,\Sigma_t)}
%  \big[\hat{L}_{S_t}(W_0+\hat{\Delta W}_t+\varepsilon)\big]
%  = \hat{L}_{S_t}(W_0+\hat{\Delta W}_t) +\mathrm{RS}_t^{\Delta}(\hat{\Delta W}_t,\Sigma_t) \]
%    where the subspace expected sharpness is
%    \[
%    \mathrm{RS}_t^{\Delta}(\hat{\Delta W}_t,\Sigma_t)
%    = \mathbb{E}_{\varepsilon\sim\mathcal N(0,\Sigma_t)}
%    \big[\hat{L}_{S_t}(W_0+\hat{\Delta W}_t+\varepsilon)
%   - \hat{L}_{S_t}(W_0+\hat{\Delta W}_t)\big].
%      \]
% \end{lemma}
% \begin{proof}
%     Sampling $W\sim Q_t$ is equivalent to sampling $\varepsilon\sim\mathcal N(0,\Sigma_t)$ on $\mathcal W_\Delta$ and setting $W=W_0+\hat{\Delta W}_t+\varepsilon$. Taking expectations and adding-subtracting $\hat L_{S_t}(W_0+\hat{\Delta W}_t)$ inside the expectation gives the decomposition. 
% \end{proof}

\subsubsection{Proof of Lemma~\ref{lem:subspace-loss-pecl}}
\label{proof:lemma2}

% \begin{equation*}
% \label{eq:subspace-loss-pecl}
% \begin{aligned}
% &\mathbb E_{W\sim Q_t}\bigl[\hat L_{S_t}(W)\bigr] \\
% &= \mathbb{E}_{\varepsilon\sim\mathcal N(0,\Sigma_t)}
%     \big[\hat{L}_{S_t}(W_\mathrm {frozen }+\hat{\Delta W}_t+\varepsilon)\big],\varepsilon\in\mathcal W_\Delta \\
% & =\hat L_{S_t}(W_{\mathrm{frozen}} + \hat{\Delta W}_t)
% +
% \mathrm{RS}_t^\Delta(\hat{\Delta W}_t,\Sigma_t),\varepsilon\in\mathcal W_\Delta.
% \end{aligned}
% \end{equation*}
We specialize the posterior to a Gaussian supported on $\mathcal W_\Delta$. For task $t$, we consider a Gaussian posterior \( Q_t^{\Delta} = \mathcal{N}(\hat{\Delta W}_t, \Sigma_t) \) on the measurable space \( (\mathcal{W}_\Delta, \mathcal{B}(\mathcal{W}_\Delta)) \), where \( \hat{\Delta W}_t \in \mathcal{W}_\Delta \) is the mean and \( \Sigma_t\) is a covariance on the adapter subspace \( \mathcal{W}_\Delta \). Sampling from \( \Delta W \sim Q_t^{\Delta} \) is equivalent to 
$
\Delta W = \hat{\Delta W}_t + \varepsilon, \quad \varepsilon \sim \mathcal{N}(0, \Sigma_t) \ \text{on } \mathcal{W}_\Delta.
$ 

\begin{lemma}
    The posterior expected empirical loss decomposes as
\begin{equation*}
    \mathbb{E}_{W\sim Q_t}\big[\hat{L}_{S_t}(W)\big]
= \mathbb{E}_{\varepsilon\sim\mathcal N(0,\Sigma_t)}
    \big[\hat{L}_{S_t}(W_\mathrm {frozen }+\hat{\Delta W}_t+\varepsilon)\big]
= \hat{L}_{S_t}(W_\mathrm {frozen }+\hat{\Delta W}_t)
  + \mathrm{RS}_t^{\Delta}(\hat{\Delta W}_t,\Sigma_t),
\end{equation*}
where the random sharpness is
$$
\mathrm{RS}_t^{\Delta}(\hat{\Delta W}_t,\Sigma_t)
:= \mathbb{E}_{\varepsilon\sim\mathcal N(0,\Sigma_t)}
   \big[\hat{L}_{S_t}(W_\mathrm {frozen }+\hat{\Delta W}_t+\varepsilon)
       - \hat{L}_{S_t}(W_\mathrm {frozen }+\hat{\Delta W}_t)\big],\quad
\varepsilon\sim\mathcal N(0,\Sigma_t)\ \text{on } \mathcal W_\Delta$$
\end{lemma}


\begin{proof}
Recall that the adapter posterior is a Gaussian supported on the update subspace,
$
Q_t^\Delta = \mathcal N(\hat{\Delta W}_t,\Sigma_t)
$
on $(\mathcal W_\Delta,\mathcal B(\mathcal W_\Delta))$.
Hence we can represent a draw $\Delta W \sim Q_t^\Delta$ as
$
\Delta W = \hat{\Delta W}_t + \varepsilon
$
with $\varepsilon \sim \mathcal N(0,\Sigma_t)$ on $\mathcal W_\Delta$.

Under the PECL constraint, the full parameter is obtained by translating the adapter update by the frozen part,
$
W = W_{\mathrm{frozen}} + \Delta W.
$
Therefore, for the empirical risk,
\begin{align*}
\mathbb E_{W \sim Q_t}\big[\hat L_{S_t}(W)\big]
&= \mathbb E_{\Delta W \sim Q_t^\Delta}\big[\hat L_{S_t}(W_{\mathrm{frozen}}+\Delta W)\big] \\
&= \mathbb E_{\varepsilon \sim \mathcal N(0,\Sigma_t)}
\big[\hat L_{S_t}(W_{\mathrm{frozen}}+\hat{\Delta W}_t+\varepsilon)\big].
\end{align*}
Finally, add and subtract the mean point loss to obtain
\begin{align*}
\mathbb E_{\varepsilon \sim \mathcal N(0,\Sigma_t)}
\big[\hat L_{S_t}(W_{\mathrm{frozen}}+\hat{\Delta W}_t+\varepsilon)\big]
&=
\hat L_{S_t}(W_{\mathrm{frozen}}+\hat{\Delta W}_t) \\
&\quad +
\mathbb E_{\varepsilon \sim \mathcal N(0,\Sigma_t)}
\Big[
\hat L_{S_t}(W_{\mathrm{frozen}}+\hat{\Delta W}_t+\varepsilon)
-
\hat L_{S_t}(W_{\mathrm{frozen}}+\hat{\Delta W}_t)
\Big].
\end{align*}
The second term is exactly the random Sharpness on the subspace
\[
\mathrm{RS}_t^\Delta(\hat{\Delta W}_t,\Sigma_t)
:=
\mathbb E_{\varepsilon \sim \mathcal N(0,\Sigma_t)}
\Big[
\hat L_{S_t}(W_{\mathrm{frozen}}+\hat{\Delta W}_t+\varepsilon)
-
\hat L_{S_t}(W_{\mathrm{frozen}}+\hat{\Delta W}_t)
\Big],
\]
completes the proof.
\end{proof}



\subsubsection{Proof of Theorem \ref{thm:pecl-bound}}
\label{proof:theorem}

% \subsection{Proof of }





\begin{theorem}[Sharpness-aware PAC--Bayes bound for PECL]
% \label{thm:pecl-bound}

For any $\delta\in(0,1]$, with probability at least $1-\delta$ over $S_1,\ldots,S_T$,
\begin{equation}\label{eq:pecl-main}
\begin{aligned}
&\frac{1}{T}\sum_{t=1}^T
\mathbb{E}_{P_t\sim\mathcal P_{t-1}}\mathbb{E}_{W\sim Q_t}L_{\mathcal{D}_t}(W)
\ \le\
\ \frac{1}{T}\sum_{t=1}^T
\mathbb{E}_{P_t\sim\mathcal P_{t-1}}
\mathbb{E}_{\varepsilon \sim \mathcal{N}(0, \Sigma_t)}[\hat{L}_{S_t}(W_0 + \widehat{\Delta W}_t + \varepsilon)
\\
&\quad +\sqrt{\frac{
\displaystyle \sum_{t=1}^T \mathbb{E}_{P_t^\Delta\sim\mathcal P_{t-1}}\!\Big[\mathrm{KL}\!\big(Q_t^\Delta\|P_t^\Delta\big)\Big]
\;+\;\displaystyle \sum_{t=1}^T \mathrm{KL}\!\big(\mathcal P_t\|\mathcal P_{t-1}\big)
\;+\;\log\!\frac{1}{\delta}
}{2\,\displaystyle\sum_{t=1}^T (m_t-1)}}.
\end{aligned}
\end{equation}
\end{theorem}



\begin{proof}
The bound follows by specializing Theorem \ref{the: full} to the PECL setting through two key decompositions:
By Lemma~\ref{lem:kl-decomp-pecl}, the full-space KL terms reduce to their adapter-space counterparts:
\[
\mathrm{KL}(Q_t \| P_t) = \mathrm{KL}(Q_t^\Delta \| P_t^\Delta).
\]
% \quad
% \mathrm{KL}(\mathcal{P}_t \| \mathcal{P}_{t-1}) = \mathrm{KL}(\mathcal{P}_t^\Delta \| \mathcal{P}_{t-1}^\Delta)
% \textbf{Step 2: Sharpness decomposition via Lemma 6.}  
Under the Gaussian posterior $Q_t^\Delta = \mathcal{N}(\Delta\hat{W}_t, \Sigma_t)$, Lemma \ref{lem:subspace-loss-pecl} gives:
\[
\mathbb{E}_{W \sim Q_t}[\hat{L}_{S_t}(W)] 
= \mathrm{E}_{\varepsilon \sim \mathcal{N}(0, \Sigma_t)}{L}_{S_t}(W^{\mathrm{froz}}_t+\widehat{\Delta W}_t + \varepsilon)
 = \hat{L}_{S_t}(W^{\mathrm{froz}}_t + \widehat{\Delta{W}_t}) + \mathrm{RS}_t^\Delta(\Delta\widehat{W}_t, \Sigma_t),
\]
where the perturbations are confined to the adapter subspace.

% \textbf{Step 3: Substitution into general bound.}
Substituting these decompositions into Lemma \ref{the: full} yields the final PECL-specific bound, where all complexity and sharpness terms are now restricted to the trainable LoRA subspace $\mathcal{W}_\Delta$.
% Substituting these decompositions into Theorem 2 yields the final PECL-specific bound, where all complexity and sharpness terms are now restricted to the trainable LoRA subspace $\mathcal{W}_\Delta$.

\end{proof}



\subsection{Bridging the Permissible Update Subspace and the LoRA Implementation}
\label{app:lora_subspace_bridge}

This appendix clarifies how the abstract permissible update subspace
$\mathcal W_{\Delta,t}$ used in the PAC--Bayesian analysis relates to the concrete LoRA
parameterization implemented in practice.
The key point is that, throughout the paper, $\mathcal W_{\Delta,t}$ denotes a
\emph{linear subspace of the ambient weight-increment space} $\mathbb R^d$ with
$d=\dim(\mathrm{vec}(\Delta W))$, rather than the rank-$r$ factor coordinate space of $(A,B)$.
This distinction is essential because the LoRA mapping $\theta_\Delta=(A,B)\mapsto \Delta W=BA$
is bilinear and therefore does not define a global linear set of reachable increments in weight space.
Our PAC--Bayesian analysis is local and only requires a linear model of permissible \emph{increment directions}
around the learned solution, which is precisely what the Jacobian provides.

% \paragraph{Abstract object in the bound: a linear increment subspace.}
% In the main text, PECL is modeled as restricting task-$t$ updates to a linear subspace
% $\mathcal W_{\Delta,t}\subseteq\mathbb R^d$, leading to the affine model class
% $
% W_{\mathrm{froz}}+\mathcal W_{\Delta,t}.
% $
% All priors, posteriors, and perturbation distributions appearing in
% Theorem~\ref{thm:pathwise_pac} are defined on $(\mathcal W_{\Delta,t},\mathcal B(\mathcal W_{\Delta,t}))$
% and pushed forward to the affine class by translation.
% Concretely, letting $T_{W_{\mathrm{froz}}}(\Delta)=W_{\mathrm{froz}}+\Delta$, we work with
% \[
% P_t := (T_{W_{\mathrm{froz}}})_{\#} P_t^{\Delta},
% \qquad
% Q_t := (T_{W_{\mathrm{froz}}})_{\#} Q_t^{\Delta},
% \]
% where $P_t^\Delta$ and $Q_t^\Delta$ are measures supported on $\mathcal W_{\Delta,t}$.
% Lemma~\ref{lem:kl-decomp-pecl} shows that
% \[
% \mathrm{KL}\!\left(Q_t\Vert P_t\right)
% =
% \mathrm{KL}\!\left(Q_t^{\Delta}\Vert P_t^{\Delta}\right),
% \]
% so under PECL the KL contribution is determined entirely by the distributions supported on
% $\mathcal W_{\Delta,t}$.

% \paragraph{Implemented object: LoRA paramaters.}
For a given adapted weight block (or layer) $\ell$, LoRA introduces trainable coordinates
$\theta_{\Delta,\ell}=(A_\ell,B_\ell)$ and defines the effective increment
\[
\Delta W_\ell(A_\ell,B_\ell)= B_\ell A_\ell,
\qquad
A_\ell\in\mathbb R^{r\times n_\ell},\;
B_\ell\in\mathbb R^{m_\ell\times r}.
\]
Let $\hat\theta_{\Delta,t}=(\hat A_t,\hat B_t)$ denote the learned LoRA parameters after training task $t$.
Consider a small coordinate perturbation
$\delta\theta_\Delta=(\delta A,\delta B)$ around $\hat\theta_{\Delta,t}$, and define
\[
g(\theta_\Delta):=\mathrm{vec}(\Delta W(\theta_\Delta))\in\mathbb R^d,
\qquad d=m_\ell n_\ell .
\]
A first-order expansion yields
\begin{equation}
\label{eq:lora_local_linearization}
g(\hat\theta_{\Delta,t}+\delta\theta_\Delta)
=
g(\hat\theta_{\Delta,t})
+
J_{\mathrm{LoRA}}(\hat\theta_{\Delta,t})\,\delta\theta_\Delta
+
R(\delta\theta_\Delta),
\end{equation}
where $J_{\mathrm{LoRA}}(\hat\theta_{\Delta,t})$ is the Jacobian of $g$ evaluated at $\hat\theta_{\Delta,t}$,
and $\|R(\delta\theta_\Delta)\|_2=O(\|\delta\theta_\Delta\|_2^2)$ as $\|\delta\theta_\Delta\|_2\to 0$.

Crucially, for LoRA the Jacobian image admits an explicit characterization.
Writing $\delta\theta_\Delta=(\delta A,\delta B)$, the induced first-order increment variation in weight space is
\[
\delta(\Delta W_\ell)
=
s_\ell(\delta B\,\hat A_t+\hat B_t\,\delta A),
\]
and hence the Jacobian image at the learned coordinates is the linear set
\begin{equation}
\label{eq:WDelta_Jacobian_image_explicit}
\mathrm{Im}\!\big(J_{\mathrm{LoRA}}(\hat\theta_{\Delta,t})\big)
=
\Big\{
\mathrm{vec}\!\big(s_\ell(\delta B\,\hat A_t+\hat B_t\,\delta A)\big)
:\;
\delta A\in\mathbb R^{r\times n_\ell},\;
\delta B\in\mathbb R^{m_\ell\times r}
\Big\}
\subseteq \mathbb R^d .
\end{equation}
This is a \emph{local subspace in increment space}, not a global reachable set, because it is defined
by the first-order linearization around $\hat\theta_{\Delta,t}$.

% \paragraph{Identification used by the theory: $\mathcal W_{\Delta,t}$ as a local tangent model.}
The abstract subspace $\mathcal W_{\Delta,t}$ in Theorem~\ref{thm:pathwise_pac} is interpreted as a local linear
model of permissible increment \emph{directions} around the task-$t$ solution in the ambient increment space.
Accordingly, we instantiate it via the LoRA-induced tangent space
\begin{equation}
\label{eq:WDelta_t_identification}
\mathcal W_{\Delta,t}
\equiv
\mathrm{Im}\!\big(J_{\mathrm{LoRA}}(\hat\theta_{\Delta,t})\big),
\end{equation}
and define the PAC--Bayesian distributions directly on this linear subspace, as required by
Lemma~\ref{lem:kl-decomp-pecl} and Theorem~\ref{thm:pathwise_pac}.

% \paragraph{How coordinate perturbations induce increment perturbations.}
In implementation, perturbations are applied in the LoRA coordinate space.
Under the local linearization~\eqref{eq:lora_local_linearization}, a coordinate perturbation
$\delta\theta_\Delta=(\delta A,\delta B)$ induces an increment perturbation
\[
\delta\Delta
:=
g(\hat\theta_{\Delta,t}+\delta\theta_\Delta)-g(\hat\theta_{\Delta,t})
\approx
J_{\mathrm{LoRA}}(\hat\theta_{\Delta,t})\,\delta\theta_\Delta,
\]
whose support lies in $\mathrm{Im}(J_{\mathrm{LoRA}}(\hat\theta_{\Delta,t}))=\mathcal W_{\Delta,t}$.
For instance, if $\delta\theta_\Delta\sim\mathcal N(0,\Sigma_{\theta,t})$ and the perturbation radius is chosen
so that the remainder term is negligible with high probability, then the induced increment perturbation is approximately Gaussian in $\mathbb R^d$,
\begin{equation}
\label{eq:Sigma_pushforward}
\delta\Delta \approx \mathcal N(0,\Sigma_t),
\qquad
\Sigma_t := J_{\mathrm{LoRA}}(\hat\theta_{\Delta,t})\,\Sigma_{\theta,t}\,
J_{\mathrm{LoRA}}(\hat\theta_{\Delta,t})^\top,
\end{equation}
and is supported on $\mathcal W_{\Delta,t}$.

\paragraph{Takeaway.}
The bound is formulated on $\mathcal W_{\Delta,t}\subset\mathbb R^d$ because it is a statement about perturbations
and posterior complexity in \emph{weight-increment space} under PECL constraints.
LoRA is a concrete mechanism that realizes such a constraint.
Although the mapping $(A,B)\mapsto BA$ is bilinear and does not define a global linear class of increments,
its local linearization around the learned coordinates induces the tangent increment subspace
$\mathrm{Im}(J_{\mathrm{LoRA}}(\hat\theta_{\Delta,t}))$, which provides an implementation-level realization of
$\mathcal W_{\Delta,t}$ while keeping the KL analysis in Lemma~\ref{lem:kl-decomp-pecl} intact.




\subsection{Feature amplification and curvature projection ratio}
\label{app:weight_projection}

This appendix details two diagnostic measures used to interpret why the perturbation scope matters in LoRA-based continual learning: the feature amplification factor $\alpha_t$ and the curvature projection preference $c_t$. 
This two diagnostics are computed on a single adapted block (layer/module) in experiments rather than the full parameter vector.



\subsubsection{Feature amplification factor}
\label{app:feature_amplification}

Let $\Delta W_t$ denote the learned weight update for a given adapted block at task $t$ (implemented as $\Delta W_t = B_t A_t$).
Let
$
\Delta W_t = U_{\Delta,t} S_{\Delta,t} V_{\Delta,t}^\top
$
be its rank-$r$ truncated SVD. The pair $(U_{\Delta,t},V_{\Delta,t})$ induces the bilinear subspace
$
\mathcal S_{\Delta,t}=\{U_{\Delta,t} Z V_{\Delta,t}^\top : Z\in\mathbb{R}^{r\times r}\}.
$
Note that $\mathcal{S}_{\Delta, t}$ is an empirical diagnostic subspace defined from the learned update $\Delta W_t$ (via its truncated SVD) and is used only for curvature localization. $\mathcal{W}_{\Delta, t}=\operatorname{Im}\left(J_{\text {LoRA }}\right)$ is the local tangent space of LoRA-induced increments in $\operatorname{vec}(\Delta W)$-space, whereas $\mathcal{S}_{\Delta, t}$ captures only the core component that preserves the row and column subspaces of $\Delta W_t$; hence $\mathcal{S}_{\Delta, t}$ is an empirical diagnostic subspace of the current solution and is not assumed to coincide with $\mathcal{W}_{\Delta, t}$.
Under the Frobenius inner product, the orthogonal projection of a matrix $M$ onto $\mathcal S_{\Delta,t}$ is
% $
% \Pi_{\Delta,t}(M)=U_{\Delta,t}(U_{\Delta,t}^\top M V_{\Delta,t})V_{\Delta,t}^\top.
% $
$
\Pi_{\mathcal S,t}(M)=U_{\Delta,t}(U_{\Delta,t}^\top M V_{\Delta,t})V_{\Delta,t}^\top.
$
Equivalently, since $U_{\Delta,t}$ and $V_{\Delta,t}$ have orthonormal columns,
% $
% \|\Pi_{\Delta,t}(M)\|_F = \|U_{\Delta,t}^\top M V_{\Delta,t}\|_F.
% $
$\|\Pi_{\mathcal S,t}(M)\|_F = \|U_{\Delta,t}^\top M V_{\Delta,t}\|_F.
$



Following~\cite{hu2022lora}, we define the feature amplification factor
\[
\alpha_t=\frac{\|\Delta W_t\|_F}{\|U_{\Delta,t}^\top W V_{\Delta,t}\|_F}
=\frac{\|\Delta W_t\|_F}{\|\Pi_{\mathcal S,t}(W)\|_F},
\]
%=\frac{\|\Delta W_t\|_F}{\|\Pi_{\Delta,t}(W)\|_F},
which compares the magnitude of the learned update to the backbone energy restricted to the same bilinear subspace induced by $\Delta W_t$.
Larger $\alpha_t$ indicates that, relative to the backbone component available within $\mathcal S_{\Delta,t}$, the learned update contributes a larger effective change along the same row and column subspaces.

\subsubsection{Curvature projection ratio}
\label{app:curvature_preference}

We next quantify whether dominant task-conditioned curvature modes concentrate more on $\mathcal S_{\Delta,t}$ than on a backbone reference bilinear subspace.

% \paragraph{Curvature eigenpairs in the LoRA-augmented parameter space.}
Let $\{(\lambda_{t,i}, v_{t,i})\}_{i=1}^k$ denote the top-$k$ eigenpairs of a positive semidefinite curvature operator at task $t$ (e.g., GGN/Fisher) computed in the parameter space that includes the backbone block $W$ and the LoRA factors $(A_t,B_t)$ (with $W$ frozen during training).
% \paragraph{Mapping eigenvectors to effective weight-space directions.}
Each eigenvector $v_{t,i}$ is a direction in the LoRA-augmented parameter space.
We decompose it into the corresponding blocks and denote the induced local perturbations by
$(\delta W_{t,i}^{(0)}, \delta A_{t,i}, \delta B_{t,i})$,
where $\delta W_{t,i}^{(0)}$ is reshaped to match the backbone weight block.
Under the first-order pushforward of the LoRA parametrization, this direction induces an effective weight-space perturbation
\[
\delta W_{t,i}
=
\delta W_{t,i}^{(0)} + \big(\delta B_{t,i}A_t + B_t \delta A_{t,i}\big).
\]
% Each eigenvector $v_{t,i}$ is a direction in the full parameter space.
% We map it to an effective direction in the weight-matrix space by aggregating its components on the backbone block and the LoRA factors:
% \[
% dW_{t,i} = dW_{t,i}^{(0)} + s\big(dB_{t,i}A_t + B_t dA_{t,i}\big),
% \]
where $ \delta W_{t,i}^{(0)}$ is the component corresponding to the backbone weight block, $(\delta A_{t,i},\delta B_{t,i})$ are the components corresponding to the LoRA factors.
This mapping matches the implementation, where curvature directions are evaluated after being pushed forward to effective weight-space increments.

% \paragraph{Per-mode projection ratios.}
We define the per-mode projection ratio onto $\mathcal S_{\Delta,t}$ as
% $$
% \rho_{\Delta,t,i}
% =\frac{\|\Pi_{\Delta,t}(\delta W_{t,i})\|_F^2}{\|\delta W_{t,i}\|_F^2}.
% $$
$$
\rho_{\Delta,t,i}
=\frac{\|\Pi_{\mathcal S,t}(\delta W_{t,i})\|_F^2}{\|\delta W_{t,i}\|_F^2}.
$$
For the backbone reference, let $(U_{W,t},V_{W,t})$ be the rank-$r$ truncated SVD subspace of the corresponding backbone weight block (restricted to the same layer/block), inducing
$
\mathcal S_{W,t}=\{U_{W,t} Z V_{W,t}^\top : Z\in\mathbb{R}^{r\times r}\},
\quad
\Pi_{\mathcal S_W,t}(M)=U_{W,t}(U_{W,t}^\top M V_{W,t})V_{W,t}^\top,
$
%\Pi_{W,t}(M)=U_{W,t}(U_{W,t}^\top M V_{W,t})V_{W,t}^\top,
and define
% $$
% \rho_{W,t,i}
% =\frac{\|\Pi_{W,t}(\delta W_{t,i})\|_F^2}{\|\delta W_{t,i}\|_F^2}.
% $$
$$\rho_{W,t,i}
=\frac{\|\Pi_{\mathcal S_W,t}(\delta W_{t,i})\|_F^2}{\|\delta W_{t,i}\|_F^2}.
$$

% \paragraph{Eigenvalue-weighted aggregation and preference ratio.}
We aggregate the localization of the top-$k$ curvature modes by eigenvalue-weighted averages
\[
R_{\Delta,t}=\frac{\sum_{i=1}^k \lambda_{t,i}\rho_{\Delta,t,i}}{\sum_{i=1}^k \lambda_{t,i}},
\qquad
R_{W,t}=\frac{\sum_{i=1}^k \lambda_{t,i}\rho_{W,t,i}}{\sum_{i=1}^k \lambda_{t,i}},
\]
and report the curvature projection ratio
\[
c_t=\frac{R_{\Delta,t}}{R_{W,t}}.
\]
Larger $c_t$ indicates that dominant curvature modes have higher projected Frobenius energy within the bilinear subspace induced by $\Delta W_t$ than within the backbone reference bilinear subspace.








\subsection{Implementation Details}

The code is https://anonymous.4open.science/r/FlatnessICMLNoname9527223/.
% Unless stated otherwise, we use SGD as the base optimizer, train each session for $20$ epochs and SAM radius $\rho = 0.05$. LoRA based variants are trained with learning rate $0.01$ without momentum or weight decay. The same training configuration is used across SeqLoRA, IncLoRA, and OLoRA, so that differences in performance can be attributed to the geometry of the adapter subspace and the choice of flatness optimizer rather than to changes in training budget.

% \subsection{Implementation Details}
\noindent\textbf{Data.}
We evaluate continual learning on ImageNet-R. For robustness under corruptions and perturbations, we use the Tiny ImageNet variants of ImageNet-C and ImageNet-P (200 classes) and construct ImageFolder-style splits that prevent leakage by grouping all corruption/perturbation variants from the same original image (or video) into a single train/test partition. All robustness evaluations use the same $224$-resolution preprocessing pipeline as the main experiments.

\noindent\textbf{Experimental setup.}
Unless stated otherwise, we use SGD as the base optimizer and train each task/session for $20$ epochs with batch size $128$ and a cosine learning-rate schedule. For sharpness-aware optimization, we set the SAM radius to $\rho=0.05$. LoRA-based variants are trained with learning rate $0.01$ and no momentum or weight decay. The same training configuration is used across SeqLoRA, IncLoRA, and OLoRA to isolate the effect of adapter geometry and perturbation design.

\noindent\textbf{LoRA configuration.}
Unless specified otherwise, the LoRA rank is set to $r=16$ for the main ImageNet-R experiments. For the Tiny ImageNet-C/P robustness experiments, we use $r=8$. In rank ablations, we sweep $r\in\{2,8,16\}$.

\noindent\textbf{Evaluation protocol.}
For accuracy reporting, we aggregate results over multiple random seeds. For sharpness and curvature measurements, we fix the random seed to $1993$ to ensure controlled comparisons and reduce estimator variance. Flatness/curvature evaluation is performed on a $10\%$ subset of the evaluation set with batch size $32$. Random sharpness uses $100$ Gaussian samples. Hessian trace and top-eigenvalue estimates are computed with $100$ Lanczos power iterations and $100$ trace samples. Loss-landscape visualizations use radius $1$ with $41$ points along a random basis with norm-filtered directions.

\noindent\textbf{Weight and curvature analysis.}
In the weight/curvature analysis, we include frozen parameters when computing full-parameter Hessian-vector products and use Lanczos with top-$k=16$ eigenvalues. Curvature localization and $\Delta W$ projections are computed on the last transformer block (block 11) and restricted to the $q$ and $v$ subspaces with row ranges $[0,768]$ and $[1536,2304]$, respectively. The localized Rayleigh estimator uses $128$ samples, and the curvature MVP backend is set to GGN. For $\Delta W$ and $\Delta W_{\mathrm{full}}$ projections, we use the QKV weight and its corresponding LoRA factors with rank $16$. The $W$--$\Delta W$ alignment analysis uses rank $16$, seed $42$,, and the seqlora setting.


\subsection{Additional Experimental Results}


\subsubsection{Raw Task-Wise $\operatorname{ACC}^{\mathrm{cls}}_t$ on ImageNet-R and ImageNet-C for Fig.~\ref{fig:robustness_imagenet_cp}}

To complement Fig.~\ref{fig:robustness_imagenet_cp}, we provide the underlying raw task-wise $\operatorname{ACC}^{\mathrm{cls}}_t$ trajectories for ImageNet-C and ImageNet-P in the tables below.
All entries are {means over repeated runs} (averaged across the same set of random seeds used for the main figure).



\begin{table*}[ht]
\centering
\caption{Raw task-wise $\operatorname{ACC}^{\mathrm{cls}}_t$ trajectories ($t=1,\dots,20$) on ImageNet-C under different optimizers (mean over repeated runs).}
\label{tab:taskwise_acc_imagenetr}
\setlength{\tabcolsep}{2.2pt}
\renewcommand{\arraystretch}{0.92}
\scriptsize
\begin{tabular}{@{}l l *{20}{r}@{}}
\toprule
\textbf{Method} & \textbf{Opt.} &
\textbf{t01} & \textbf{t02} & \textbf{t03} & \textbf{t04} & \textbf{t05} &
\textbf{t06} & \textbf{t07} & \textbf{t08} & \textbf{t09} & \textbf{t10} &
\textbf{t11} & \textbf{t12} & \textbf{t13} & \textbf{t14} & \textbf{t15} &
\textbf{t16} & \textbf{t17} & \textbf{t18} & \textbf{t19} & \textbf{t20} \\
\midrule

\multirow{3}{*}{\textbf{SeqLoRA}}
& SGD
& 92.67 & 87.84 & 87.67 & 83.42 & 82.13 & 81.50 & 79.62 & 78.83 & 78.11 & 76.47
& 75.85 & 75.28 & 74.69 & 74.24 & 74.27 & 73.48 & 72.69 & 71.72 & 70.49 & 69.77 \\
& \quad $+\mathrm{AS}^{(0)}_S$
& 92.67 & 87.67 & 87.33 & 84.42 & 83.20 & 82.39 & 80.33 & 79.29 & 79.11 & 78.20
& 77.73 & 77.00 & 76.69 & 77.00 & 77.33 & 77.08 & 76.45 & 75.94 & 75.33 & 73.93 \\
& \quad $+\mathrm{AS}^{(1)}_S$
& 95.67 & 93.67 & 89.78 & 88.84 & 87.67 & 86.61 & 86.67 & 86.33 & 84.85 & 84.14
& 82.61 & 82.17 & 81.87 & 81.60 & 81.42 & 80.44 & 80.10 & 79.96 & 80.21 & 80.03 \\

\addlinespace[2pt]

\multirow{3}{*}{\textbf{IncLoRA}}
& SGD
& 92.33 & 89.84 & 88.22 & 86.58 & 83.40 & 81.22 & 80.38 & 78.46 & 75.89 & 73.77
& 73.21 & 73.61 & 73.02 & 72.95 & 71.62 & 70.17 & 68.82 & 66.30 & 64.86 & 61.77 \\
& \quad $+\mathrm{AS}^{(0)}_S$
& 92.33 & 90.83 & 90.45 & 85.92 & 87.00 & 83.39 & 81.81 & 79.75 & 78.22 & 77.87
& 76.79 & 73.47 & 72.03 & 73.43 & 71.07 & 71.44 & 70.78 & 69.63 & 67.96 & 67.98 \\
& \quad $+\mathrm{AS}^{(1)}_S$
& 96.00 & 92.17 & 88.33 & 86.92 & 85.87 & 85.11 & 86.33 & 85.17 & 84.78 & 84.03
& 82.85 & 82.56 & 81.54 & 81.52 & 80.64 & 79.73 & 79.59 & 80.32 & 79.58 & 79.47 \\

\addlinespace[2pt]

\multirow{3}{*}{\textbf{OLoRA}}
& SGD
& 92.33 & 88.50 & 87.78 & 85.59 & 83.00 & 81.28 & 80.57 & 78.58 & 75.67 & 74.10
& 73.55 & 73.36 & 73.46 & 73.33 & 71.87 & 71.40 & 70.57 & 68.33 & 67.44 & 65.12 \\
& \quad $+\mathrm{AS}^{(0)}_S$
& 92.33 & 88.67 & 90.22 & 85.67 & 86.33 & 83.00 & 81.48 & 79.83 & 78.26 & 78.33
& 77.00 & 73.97 & 71.77 & 72.50 & 70.51 & 70.36 & 69.33 & 68.54 & 66.21 & 66.98 \\
& \quad $+\mathrm{AS}^{(1)}_S$
& 96.00 & 91.17 & 88.00 & 86.42 & 85.47 & 85.83 & 86.33 & 84.96 & 84.19 & 83.07
& 82.30 & 81.53 & 81.46 & 80.69 & 80.27 & 78.75 & 79.14 & 78.98 & 79.28 & 79.35 \\

\bottomrule
\end{tabular}
\end{table*}


\begin{table*}[ht]
\centering
\caption{Raw task-wise $\operatorname{ACC}^{\mathrm{cls}}_t$ trajectories ($t=1,\dots,20$) on ImageNet-P under different optimizers (mean over repeated runs).}
\label{tab:taskwise_acc_new}
\setlength{\tabcolsep}{2.2pt}
\renewcommand{\arraystretch}{0.92}
\scriptsize
\begin{tabular}{@{}l l *{20}{r}@{}}
\toprule
\textbf{Method} & \textbf{Opt.} &
\textbf{t01} & \textbf{t02} & \textbf{t03} & \textbf{t04} & \textbf{t05} &
\textbf{t06} & \textbf{t07} & \textbf{t08} & \textbf{t09} & \textbf{t10} &
\textbf{t11} & \textbf{t12} & \textbf{t13} & \textbf{t14} & \textbf{t15} &
\textbf{t16} & \textbf{t17} & \textbf{t18} & \textbf{t19} & \textbf{t20} \\
\midrule

\multirow{3}{*}{\textbf{SeqLoRA}}
& SGD
& 96.17 & 93.42 & 92.00 & 89.88 & 88.17 & 87.86 & 86.43 & 84.92 & 85.57 & 84.63
& 83.88 & 83.71 & 82.40 & 82.02 & 83.08 & 81.95 & 80.09 & 80.17 & 78.01 & 72.76 \\
& \quad $+\mathrm{AS}^{(0)}_S$
& 95.83 & 92.67 & 92.28 & 89.54 & 88.47 & 88.58 & 87.90 & 87.29 & 87.63 & 86.82
& 85.80 & 85.60 & 84.57 & 84.21 & 84.51 & 83.58 & 83.30 & 83.70 & 82.32 & 80.95 \\
& \quad $+\mathrm{AS}^{(1)}_S$
& 99.17 & 94.92 & 95.00 & 92.71 & 91.60 & 91.92 & 89.55 & 88.69 & 89.29 & 89.10
& 87.70 & 87.71 & 87.68 & 87.62 & 88.19 & 87.51 & 86.89 & 87.02 & 86.93 & 86.08 \\

\addlinespace[2pt]

\multirow{3}{*}{\textbf{IncLoRA}}
& SGD
& 95.67 & 92.50 & 91.78 & 89.71 & 88.80 & 88.08 & 85.69 & 85.29 & 84.59 & 83.25
& 80.94 & 78.58 & 77.72 & 77.79 & 74.13 & 73.24 & 71.70 & 69.29 & 68.85 & 66.92 \\
& \quad $+\mathrm{AS}^{(0)}_S$
& 95.83 & 91.58 & 92.17 & 89.62 & 89.27 & 88.28 & 86.57 & 85.31 & 85.56 & 85.33
& 83.95 & 81.81 & 80.01 & 78.26 & 78.75 & 77.08 & 77.34 & 75.79 & 76.53 & 73.75 \\
& \quad $+\mathrm{AS}^{(1)}_S$
& 98.83 & 96.41 & 95.39 & 93.33 & 91.63 & 90.97 & 89.81 & 88.83 & 88.78 & 88.63
& 87.46 & 86.27 & 86.58 & 86.15 & 86.63 & 86.42 & 86.45 & 85.54 & 85.34 & 85.09 \\

\addlinespace[2pt]

\multirow{3}{*}{\textbf{OLoRA}}
& SGD
& 95.67 & 90.00 & 90.39 & 88.33 & 88.20 & 87.61 & 85.76 & 85.02 & 84.45 & 83.37
& 79.97 & 77.76 & 76.68 & 76.64 & 72.82 & 72.54 & 70.40 & 68.16 & 67.61 & 65.55 \\
& \quad $+\mathrm{AS}^{(0)}_S$
& 95.83 & 90.17 & 91.06 & 88.46 & 88.27 & 87.86 & 86.74 & 84.96 & 85.37 & 84.95
& 83.74 & 82.61 & 81.01 & 79.58 & 78.79 & 77.21 & 76.96 & 75.30 & 75.73 & 73.32 \\
& \quad $+\mathrm{AS}^{(1)}_S$
& 98.83 & 95.33 & 94.39 & 91.46 & 90.17 & 90.14 & 88.60 & 87.48 & 87.02 & 87.23
& 86.23 & 85.47 & 85.99 & 85.79 & 86.19 & 86.23 & 86.01 & 85.29 & 84.61 & 84.47 \\
\bottomrule
\end{tabular}
\end{table*}





\subsubsection{Computational Cost Analysis}


Table~\ref{tab:lora_opt_metric} reports the average per-task training time on ImageNet-R.
Across LoRA variants, adding sharpness regularization increases the computational cost relative to SGD, with $AS^{(1)}_S$ being consistently the most expensive due to the additional gradient-norm maximization, while $AS^{(0)}_S$ shows a moderate overhead.
In contrast, $RS^{(0)}_S$ incurs a smaller and more variable overhead, and can be comparable to SGD in some settings.


\begin{table}[hpb]
\centering
\caption{Average per-task training time (seconds) across methods and optimizers on ImageNet-R.}
\label{tab:lora_opt_metric}
\footnotesize
\setlength{\tabcolsep}{3pt}
\renewcommand{\arraystretch}{0.95}
\begin{tabular}{l cccc}
\toprule
\textbf{Method} & \textbf{SGD} & $+\mathrm{RS}^{(0)}_S$ & $+\mathrm{AS}^{(0)}_S$ & $+\mathrm{AS}^{(1)}_S$ \\
\midrule
SeqLoRA & $282.72 \pm 61.34$  & $244.95 \pm 6.38$  & $472.29 \pm 83.70$  & $670.20 \pm 217.92$ \\
IncLoRA & $310.54 \pm 141.00$ & $658.34 \pm 83.90$ & $777.91 \pm 239.29$ & $1247.56 \pm 5.05$ \\
OLoRA   & $412.99 \pm 146.84$ & $644.19 \pm 55.05$ & $628.44 \pm 181.29$ & $836.99 \pm 327.65$ \\
\bottomrule
\end{tabular}
\end{table}

% \begin{table}[hpb]
% \centering
% \caption{Average per‑task training time (seconds) across methods and optimizers on ImageNet‑R}
% \label{tab:lora_opt_metric}
% \setlength{\tabcolsep}{4pt}
% \renewcommand{\arraystretch}{1}
% \resizebox{\linewidth}{!}{
% \begin{tabular}{l ccccc}
% \toprule
% \textbf{Method} & \textbf{SGD} & $+RS^{(0)}_S$ & $+AS^{(0)}_S$ & $+AS^{(1)}_S$ \\
% \midrule
% SeqLoRA & 282.72 ± 61.34  & 244.95 ± 6.38  & 472.29 ± 83.70  & 670.20 ± 217.92    \\
% IncLoRA & 310.54 ± 141.00 & 658.34 ± 83.90 & 777.91 ± 239.29 & 1247.56 ± 5.05  3 \\
% OLoRA   & 412.99 ± 146.84 & 644.19 ± 55.05 & 628.44 ± 181.29 & 836.99 ± 327.65  \\
% \bottomrule
% \end{tabular}
% }
% \end{table}








\end{document}
